\chapter{Dirichlet 膜与稳定性条件}\label{chap:5}

\SourceParagraph{ch05-intro-001}{303}
现在，借助更好的工具，继续\BookChapterRef{3}的 Dirichlet 膜讨论。此前研究的是拓扑 Dirichlet 膜，即 A、B 拓扑扭曲模型的边界条件，以及这些模型的开弦关联函数。虽然来自物理，这些定义，给出的结构却非常像 A 模型的 Fukaya 范畴，以及 B 模型的凝聚层范畴。

\SourceParagraph{ch05-intro-002}{303}
自然会猜想，这就是正确的识别。不过，如果如此，镜像对称物理就会预言，这两个范畴等价；实际却不成立。最简单的迹象是，凝聚层范畴太小，没有包含足够多的 B 膜。

\SourceParagraph{ch05-intro-003}{303}
Kontsevich 在1994年提出\BookCite{ref309}，正确的 B 膜范畴，应通过进入凝聚层导出范畴得到。\BookSectionRef{sec:5.3}{5.3}将从弦理论内部，给出支持这一主张的详细论证。

\SourceParagraph{ch05-intro-004}{303}
与\BookChapterRef{4}比较，所需的基本形式要素已经具备，即：D 膜自然可以看作 Abel 范畴中的对象，而态射是带有 $\mathbb{Z}$ 分次的 $Q$ 上同调类。不过，还不清楚，为什么按拟同构识别对象，会与物理相关。即使作了这一广泛识别，简单例子也说明，导出范畴中的对象，仍然太多，不能全部是物理 D 膜。

\SourceParagraph{ch05-intro-005}{303}
一个解决问题的线索是，即使在小得多的凝聚层范畴中，也并非每个对象都对应于 D 膜。只有允许 Hermitian--Yang--Mills 方程解的层才是 D 膜。由 Donaldson--Uhlenbeck--Yau 定理，这些层与 $\mu$ 稳定层的子集相对应。

\SourceParagraph{ch05-intro-006}{303}
A 膜图景中也有类似现象。Lagrange 子流形范畴，仍然太一般，不能直接与物理对象对应。更适合考虑 Hamilton 同痕下的等价类；这些类，又与特殊 Lagrange 子流形，有一个猜想中的联系，其中涉及猜想的稳定性概念\BookCite{ref446,ref262}。
\begin{table}[H]
\centering
\caption{镜像对称。前一部分是拓扑膜的对应，后一部分是物理 BPS 膜的对应。}\label{tab:5.1}
\begin{tabular}{@{}p{0.25\linewidth}p{0.30\linewidth}p{0.36\linewidth}@{}}
\toprule
 & A 模型 & B 模型\\
\midrule
几何 & 辛几何，不依赖可积复结构 & 代数几何，不依赖度规\\
范畴 & Fukaya 范畴 & 导出范畴\\
D 膜 & Lagrange 子流形 & 凝聚层复形\\
开弦 & Floer 上同调 & Ext 群\\
参数依赖 & $B+\mathrm{i}J$ & 复结构\\
荷 & $l_i\in H^3$ & $\operatorname{ch}(E)\sqrt{\operatorname{td}(X)}\in H^{\mathrm{even}}(X)$，或 $K(X)$\\
\midrule
BPS A/B 膜 & 特殊 Lagrange 子流形 & $\Pi$ 稳定复形\\
稳定性的参数依赖 & 复结构 & $B+\mathrm{i}J$\\
束缚态 & $A\#B$ & $\operatorname{Cone}(A\to B)$\\
\bottomrule
\end{tabular}
\end{table}
这里表中的 $J$，表示 K\"ahler 形式，也就是前面记作 $\omega$ 的形式。

\SourceParagraph{ch05-intro-007}{304}
这提示，应区别“拓扑”与“物理”D 膜。纯粹用 A、B 模型定义的拓扑膜，分别是 Fukaya 范畴与凝聚层导出范畴的一般对象。Kontsevich 的镜像对称猜想，适用于它们。

\SourceParagraph{ch05-intro-008}{304}
另一方面，物理 SCFT 依赖更多数据，刻画其超对称边界条件，是另一个问题。称这些边界条件为物理 A、B 膜，并尝试将它们，与 Fukaya 范畴或导出范畴中，满足某种稳定性条件的对象相识别。

\SourceParagraph{ch05-intro-009}{304}
虽然能够在 SCFT 中识别稳定性条件的若干要素，但本书使用的技术，尚不足以直接导出其完整定义。不过，已有足够的要素，检验这一一般思想：与大体积极限中已知的稳定性条件比较，并检验镜像对称的预言。首先，预期下面的性质成立。
\BookStatementNumber{1}
\begin{property}\label{property:5.1}\label{pty:5.1}
镜像稳定性条件的模空间同构。
\end{property}

\SourceParagraph{ch05-intro-010}{304}
要正确理解这一点，需要讨论稳定性条件模空间中，单值变换诱导的“量子等价”。随后，还有第二项性质。
\BookStatementNumber{2}
\begin{property}\label{property:5.2}\label{pty:5.2}
给定镜像对 $X,Y$ 上两个对应的稳定性条件，$X$ 上 Fukaya 范畴中的稳定对象范畴，与 $Y$ 上凝聚层导出范畴中的稳定对象范畴等价。
\end{property}

\SourceParagraph{ch05-intro-011}{305}
让这些性质成立的最简单方式，当然是，纯粹根据 Kontsevich 猜想所识别的结构，定义稳定性，也就是使用三角结构，或者还包括 $A_\infty$ 结构。不过，立刻遇到一个问题：基于 GIT 或其简单推广的稳定性定义，只适用于 Abel 范畴。特别地，需要谈论子对象，但一般的导出范畴，没有相同意义的子对象。解决这一问题，是本章其余部分的中心。

\SourceParagraph{ch05-intro-012}{305}
沿着这些思想，可以为三角范畴精确定义稳定性结构\BookCite{ref065,ref063}。定义本身不内在地依赖镜像对称或弦理论，因此，可以纯粹从数学上发展。

\SourceParagraph{ch05-intro-013}{305}
先在\BookSectionRef{sec:5.1}{5.1}简要介绍 D 膜的拓扑分类。\BookSectionRef{sec:5.2}{5.2}回顾 A、B 膜的几何稳定性定义，分别基于特殊 Lagrange 子流形的稳定性，以及 Donaldson--Uhlenbeck--Yau 定理。虽然两者相当不同，共同特征却提示了进一步发展的方向。

\SourceParagraph{ch05-intro-014}{305}
\BookSectionRef{sec:5.3}{5.3}，论证一般拓扑 B 膜，是凝聚层导出范畴的对象。这使我们能够与边界 CFT 的基本定义联系，作出一般论证。特别地，\BookSectionRef{sec:5.3.4}{5.3.4}将说明，整数分次推广为实数分次；它依赖于拓扑理论忽略的“另一组模参数”，即 A 膜的复结构参数，以及 B 膜的 K\"ahler 参数。

\SourceParagraph{ch05-intro-015}{305}
\BookSectionRef{sec:5.4}{5.4}，从互补的“世界体”观点发展物理，定义各种有意思的空间上的膜理论，包括 $\mathbb{C}^3$ 的轨形及其消解。还引入快子凝聚、Seiberg 对偶等物理图景。

\SourceParagraph{ch05-intro-016}{305}
\BookSectionRef{sec:5.5}{5.5}，结合这些物理论证，得到正确稳定性条件必须满足的要求。\BookSectionRef{sec:5.6}{5.6}，通过例子说明，特别是 $\mathbb{C}^3/\mathbb{Z}_3$ 轨形与其消解。

\SourceParagraph{ch05-intro-017}{305}
\BookSectionRef{sec:5.7}{5.7}，给出将这些思想公理化的精确数学定义。公理足够强，能够确定稳定性条件空间的局部结构；它是有限维流形。通过 $\mathbb{C}^3/\mathbb{Z}_3$ 的例子，显式说明稳定性结构空间。

\SourceParagraph{ch05-intro-018}{305}
最后，在\BookSectionRef{sec:5.9}{5.9}讨论一些开放问题。\BookRef{tab:5.1}{表1}，预先概括了相当一部分讨论，展示两个镜像对称对应。

\section{K 理论、相交积与反膜}\label{sec:5.1}

\SourceParagraph{ch05-5.1-intro-001}{306}
先从导致 D 膜“拓扑”分类的物理论证开始。良好的一般介绍，见文献\BookCite{ref363}。

\SourceParagraph{ch05-5.1-intro-002}{306}
有几种可能的方法。一种是，把能够通过“连续物理过程”联系的配置相识别。特别地，所有相对论量子理论都含有“反粒子”，并允许粒子—反粒子湮灭。这启发了 K 理论的 Grothendieck 构造的物理版本\BookCite{ref471}。

\SourceParagraph{ch05-5.1-intro-003}{306}
另一种方法，是使用配置在渐近无穷远处可测量的守恒荷。由定义，它们在物理过程中不变，还满足量子化条件；原型是 Maxwell 理论中，允许的电荷与磁荷的 Dirac 条件。Dirichlet 膜的守恒荷，是 Ramond--Ramond 荷；这一方法与 K 理论的类比，最早见于文献\BookCite{ref358}。

\SourceParagraph{ch05-5.1-intro-004}{306}
原则上，两种方法可能给出不同结果：未必每个不变量都能在无穷远测量，量子化条件允许的荷，也未必全部由物理配置实现。不过，最简单的猜想，是两种方法一致。

\SourceParagraph{ch05-5.1-intro-005}{306}
有些论证需要一定物理背景，教材如\BookCite{ref394}中有说明。这里不作这种要求，只介绍讨论的基本思路与主要结果。虽然主要针对 II 型弦的 Calabi--Yau 紧化，构造其实更一般。扩展包括非零 H 通量\BookCite{ref055,ref344}、带有实丛结构的 I 型弦\BookCite{ref471}，以及定向轨形\BookCite{ref114,ref244}。

\subsection{K 理论的物理实现}\label{sec:5.1.1}

\SourceParagraph{ch05-5.1.1-001}{306}
有一个基本构造，给定物理理论 $T$，产生守恒荷的 Abel 群 $K(T)$。不必考虑微观动力学，只需知道理论所描述的“粒子”集合 $S$，以及粒子结合或分裂成其他粒子的束缚态形成、衰变过程。为简便起见，通常称它们为“束缚过程”。

\SourceParagraph{ch05-5.1.1-002}{306}
如果一串束缚过程，可以把一组粒子转换成另一组，就称它们物理等价。以粒子的有限整系数形式组合群 $\mathbb{Z}^{(S)}$，模去这些过程所生成的关系，定义 $K(T)$。

\SourceParagraph{ch05-5.1.1-003}{307}
\begin{example*}
假定 $T$ 的粒子集合为 $S=\{A,B,C\}$。如果它们完全稳定，且没有转化关系，可以定义三个守恒整数荷，即各自的粒子数，因此，$K(T)\cong\mathbb{Z}^3$。再加入束缚过程
\begin{equation}\label{eq:5.1}
A+B\longleftrightarrow C,
\end{equation}
双向箭头表示过程可以正向或反向发生。此时，显然有 $K(T)\cong\mathbb{Z}^2$。
\end{example*}

\SourceParagraph{ch05-5.1.1-004}{307}
也许会批评，这一构造假定了负粒子数配置的存在。不过，在所有满足狭义相对论约束的物理理论中，带荷粒子都伴有质量相同、荷相反的反粒子。粒子与反粒子，可以湮灭为一组零荷粒子。虽然最初是 Dirac 方程的预言，这也来自量子场论的一般公理，并适用于弦理论。因此，有过程
\[
B+\overline B\longleftrightarrow Z_1+Z_2+\cdots,
\]
其中，$\overline B$ 是反粒子，$Z_i$ 是为保持能量守恒而产生的零荷粒子。定义荷群时，将这样的零荷配置，与群恒等元相识别，得到 $B+\overline B\longleftrightarrow0$。于是，反粒子提供负元素。

\SourceParagraph{ch05-5.1.1-005}{307}
承认反粒子以后，就可以用 Grothendieck 构造更简单地表述。以配置对 $(E,F)$ 表示形式差 $E-F$，并施加关系
\[
(E,F)\sim(E+B,F+B)
\sim(E+L,F+R)\sim(E+R,F+L),
\]
其中，$L,R$ 是某个束缚过程的左右两边。可以从非负配置的交换幺半群作群完备化，也可以直接使用上面的形式整组合。

\subsubsection{Dirichlet 膜的 K 理论}\label{sec:5.1.1.1}

\SourceParagraph{ch05-5.1.1.1-001}{307}
把膜视为粒子，每个膜 $B$ 都应有反膜 $\overline B$。如果 $B$ 包覆子流形 $L$，预期反膜包覆相同子流形，但定向相反。

\SourceParagraph{ch05-5.1.1.1-002}{307}
潜在的问题是，$L$ 的定向是否真的影响物理，并不先验明显；尤其在 $L$ 为一点的退化情形中。在一般理论，例如玻色弦中，它未必成立。所幸，正如\BookSectionRef{sec:3.5.3}{3.5.3}所述，超弦中膜的定义，需要定向选择，以实施 GSO 投影。

\SourceParagraph{ch05-5.1.1.1-003}{308}
继续讨论，还需要定义束缚过程。最直接的情形，是具有几何定义的 sigma 模型边界条件；此时，可以使用连续形变。

\subsection{A 膜}\label{sec:5.1.2}

\SourceParagraph{ch05-5.1.2-001}{308}
为明确起见，取 Calabi--Yau 三维流形；不过，大部分讨论，也适用于一般复维数 $n$。物理 A 膜，由特殊 Lagrange 子流形 $L$ 与平坦丛 $E$ 组成的对 $(L,E)$ 描述。什么时候，两对这样的数据，可以通过束缚过程相联系？

\SourceParagraph{ch05-5.1.2-002}{308}
Feynman 泛函积分给出一个简单的启发式回答：如果配置空间中有连续路径相连，路径本身，或者它的小扰动，就会以某个非零权重出现在泛函积分中。因此，这个问题，在本质上是拓扑的。

\SourceParagraph{ch05-5.1.2-003}{308}
暂时忽略平坦丛，就预期 A 膜的荷群为 $H_3(Y,\mathbb{Z})$，膜的类为 $(\operatorname{rank}E)[L]$。如果 $L$ 上平坦联络的模空间连通，也得到相同图景。不过，如果它不连通，例如 $\pi_1(L)$ 是挠群，情形会怎样？

\SourceParagraph{ch05-5.1.2-004}{308}
\BookStatementNumber{3}
\begin{example}\label{ex:5.3}
设 $Y\subset\mathbb{CP}^4$ 是 Fermat 五次超曲面，由 $x_0^5+x_1^5+x_2^5+x_3^5+x_4^5=0$ 定义。它具有实对合 $\sigma$，取 $L$ 为固定点集。可以说明，$L$ 是特殊 Lagrange 子流形，拓扑为 $\mathbb{RP}^3$，因此，$\pi_1(L)\cong\mathbb{Z}_2$。
\end{example}

\SourceParagraph{ch05-5.1.2-005}{308}
在这一情形中，需要更深的物理论证，才能判断荷群，是 $H_3(Y,\mathbb{Z})$，还是更大的群。不过，基于前面的思想，一个自然猜想是 $K^1(Y)$；通过相应的 K 同调解释，它描述奇维子流形及其丛\footnote{物理上，$K^0$ 与 $K^1$ 的联系，涉及在一个辅助圆周上作 T 对偶。细节见文献\BookCite{ref471}。}。

\SourceParagraph{ch05-5.1.2-006}{308}
对于\BookRef{ex:5.3}{例题5.3}，文献\BookCite{ref069}论证，只在平坦联络选择上不同的两个膜，实际上可以在弦理论中相联系，与 $H_3$ 荷群的图景一致。也可以使用 Atiyah--Hirzebruch 谱序列，说明对单连通的 Calabi--Yau 三维流形 $Y$，$K^1(Y)\cong H_3(Y,\mathbb{Z})$。因此，这个一般猜想，在这里以及其他已经研究的例子中得到支持。不过，目前还没有一般的物理证明。

\subsubsection{相交数}\label{sec:5.1.2.1}

\SourceParagraph{ch05-5.1.2.1-001}{308}
$H_3(Y,\mathbb{Z})$ 上，有由定向相交数给出的自然双线性形式：
\begin{equation}\label{eq:5.2}
I(L_1,L_2)=\#([L_1]\cap[L_2]).
\end{equation}
在中间维数的一般情形中，其对称性为 $(-1)^n$；特别地，三维 Calabi--Yau 的配对反对称。此外，由 Poincar\'e 对偶，模去挠子群后的相交格幺模。所有整同调类，是否都能由特殊 Lagrange 子流形实现，仍是非平凡的问题。

\SourceParagraph{ch05-5.1.2.1-002}{309}
可以将它与膜的 RR 荷物理表达式比较。Dirichlet $p$ 膜带有 RR 荷，意味着，它是十维时空中推广的 $(p+1)$ 形式规范势的源\BookCite{ref393}。这一点，来自世界体作用量中的极小耦合项
\begin{equation}\label{eq:5.3}
\int_{\mathcal W}C^{(p+1)},
\end{equation}
其中，$\mathcal W$ 是 $(p+1)$ 维世界体。对于 $p=0$，$C^{(1)}$ 是1形式，或者“向量”势，如 Maxwell 理论。因此，D0膜，是对于这个十维 Maxwell 理论带电的粒子。

\SourceParagraph{ch05-5.1.2.1-003}{309}
考虑十维时空为 $\mathbb{R}^{3,1}\times X$ 的紧化，让 Dp膜包覆 $p$ 维循环 $L$。其世界体为 $\mathbb{R}\times L$，其中，$\mathbb{R}$ 是四维时空中的类时世界线。用 $L$ 的 Poincar\'e 对偶类 $\omega_L\in H^{2n-p}(X,\mathbb{R})$，将耦合改写为
\begin{equation}\label{eq:5.4}
\int_{\mathbb{R}\times X}C^{(p+1)}\wedge\omega_L.
\end{equation}
对 $X$ 积分以后，变成四维世界线上1形式的积分，正是粒子的极小电耦合。因此，包覆膜携带一个可以在渐近无穷远检测的电荷。

\SourceParagraph{ch05-5.1.2.1-004}{309}
这里将 RR 荷概括为形式表达式
\begin{equation}\label{eq:5.5}
\int_LC=\int_XC\wedge\omega_L,
\end{equation}
其中，$C\in H^*(X,\mathbb{R})$。换句话说，荷由 $H_p(X,\mathbb{R})$ 的类表示。

\SourceParagraph{ch05-5.1.2.1-005}{309}
特别地，三维的 A 膜，具有 $H_3(X,\mathbb{R})$ 中的守恒荷。这弱于整同调陈述。要从物理上看出整数条件，必须识别其中的磁荷，并研究 Dirac--Schwinger--Zwanziger 量子化。它要求电磁场的角动量，满足 $J=\hbar m/2$，$m\in\mathbb{Z}$。在适当的电磁单位中，角动量由 $\int\mathrm{d}^3x\,\vec x\times(\vec E\times\vec B)$ 给出；量子化，正好要求\BookEqRef{eq:5.2}{5.2}取整数。因此，物理与数学的相容条件一致。

\SourceParagraph{ch05-5.1.2.1-006}{309}
余各向同性 A 膜也有类似考察。如果 $X$ 的完整群恰为 $\operatorname{SU}(3)$，非 Lagrange 的余各向同性膜，必须为五维。由于 $H^5(X,\mathbb{Q})=0$，子流形 $L$ 在有理同调中平凡。因此，如果 $E$ 拓扑平凡，所得同调荷与相应 K 类，只可能是挠类。稍后讨论非平凡丛，得到\BookEqRef{eq:5.8}{5.8}。

\SourceParagraph{ch05-5.1.2.1-007}{310}
如果 $X$ 是环面或 K3 曲面，情形更复杂。余各向同性膜的有理荷，也未必只位于中间维数上同调，而是在偶次或奇次总上同调的某个子空间；采用哪种奇偶性，取决于复维数。在半经典层面，该子空间由
\[
(L-\Lambda)\alpha=0,\qquad\alpha\in\bigoplus_pH^p(X,\mathbb{Q})
\]
确定。其中，$L,\Lambda$ 是 Lefschetz $\operatorname{SL}(2,\mathbb{C})$ 作用的生成元：$L$ 是与 K\"ahler 类作杯积，$\Lambda$ 是其对偶\BookCite{ref290,ref284}。

\subsection{作为开弦指标的相交形式}\label{sec:5.1.3}

\SourceParagraph{ch05-5.1.3-001}{310}
几何 D 膜的这个对应，容易直接看出，也并不很深。不过，物理论证还有一个更深的结果：任何 $(2,2)$ SCFT 的边界条件，其荷群上，都必须有这样的整双线性形式。后面将讨论世界面计算；最简单的观察是，\BookEqRef{eq:5.2}{5.2}，也计算从 $L_1$ 到 $L_2$ 的开弦量子理论指标，按指标迹的正规化含义理解，
\begin{equation}\label{eq:5.6}
I(L_1,L_2)=\operatorname{Tr}_{L_1,L_2}(-1)^F.
\end{equation}

\SourceParagraph{ch05-5.1.3-002}{310}
对于 A 膜，如\BookSectionRef{sec:3.6.2.5}{3.6.2.5}所述，开弦与交点相关。可以证明\BookCite{ref150}，如果两个交点的鬼数之差为奇数，它们的定向符号相反。因此，相交数，就是复形\BookEqRef{eq:3.186}{3.186}的 Euler 示性数：
\begin{equation}\label{eq:5.7}
\#([L_1]\cap[L_2])
=\int_Yl_1\wedge l_2
=\sum_i(-1)^i\dim\operatorname{Hom}^i(L_1,L_2),
\end{equation}
其中，$l_a$ 表示相应的 Poincar\'e 对偶类。

\subsection{携带非平凡丛的膜}\label{sec:5.1.4}

\SourceParagraph{ch05-5.1.4-001}{310}
B 膜的类似考察，因可携带非平凡 K 类的丛而更加复杂。最先理解的部分，是耦合\BookEqRef{eq:5.3}{5.3}的推广；它可以从 RR 场与 D 膜的耦合读出。在若干部分世界面计算以后，文献\BookCite{ref190,ref095,ref358}，通过反常流入论证，给出一般表达式。对于嵌入 $f\colon L\hookrightarrow M$ 的膜 $(L,E)$，耦合为
\begin{equation}\label{eq:5.8}
\int_Lf^*C\wedge\operatorname{ch}(E)
\frac{\widehat A(TL)\,\mathrm{e}^{d/2}}
{\sqrt{f^*\widehat A(TM)}}.
\end{equation}
这里，$M$ 是十维时空，$C$ 是偶次或奇次 RR 势的形式和，$\widehat A$ 是 Â 特征类。它与复丛的 Todd 类满足 $\operatorname{td}=\mathrm{e}^{c_1/2}\widehat A$。$d$ 是定义 $L$ 上相应 $\operatorname{spin}^c$ 结构的二次整类，费米子的定义需要它。

\SourceParagraph{ch05-5.1.4-002}{311}
B 膜中，这个公式可以显著简化。不过，完整表达式已经很有启发性，因为可以使用相对 K 理论的推出 $f_!$，改写为整个时空上的积分\BookCite{ref095}：
\begin{equation}\label{eq:5.9}
\int_MC\wedge\operatorname{ch}(f_!E)\sqrt{\widehat A(TM)}.
\end{equation}

\SourceParagraph{ch05-5.1.4-003}{311}
真正重要的是，相交形式的推广，涉及时空中这些荷的配对。令 $Q_a=f_{a!}E_a$ 为两个膜推到时空的 K 类，在相应的相对或紧支撑指标条件下，
\begin{equation}\label{eq:5.10}
\begin{aligned}
I(L_1,L_2)
&=\int_M\operatorname{ch}(Q_1^{\vee})
\operatorname{ch}(Q_2)\widehat A(TM)\\
&=\operatorname{ind}\not D_{Q_1^{\vee}\otimes Q_2}.
\end{aligned}
\end{equation}
右边是耦合相应乘积虚丛的 Dirac 算符指标。它计数从第一个膜到第二个膜的费米 Ramond 开弦，带有符号，因此，弦理论在很一般的情形中实现\BookEqRef{eq:5.6}{5.6}。

\subsubsection{B 膜}\label{sec:5.1.4.1}

\SourceParagraph{ch05-5.1.4.1-001}{311}
现在专门用于 B 型膜。对于 Calabi--Yau $X$ 中复子流形 $L$ 上的全纯丛 $E$，使用 $\operatorname{td}=\mathrm{e}^{c_1/2}\widehat A$ 与 $c_1(X)=0$，将耦合改写为
\begin{equation}\label{eq:5.11}
\int_XC\wedge\operatorname{ch}(f_!E)\sqrt{\operatorname{td}(TX)}.
\end{equation}

\SourceParagraph{ch05-5.1.4.1-002}{311}
对支撑在光滑全纯子簇 $S$ 上的凝聚层，设 $i\colon S\hookrightarrow X$，法丛为 $N$，则也可以写为
\begin{equation}\label{eq:5.12}
\int_XC\wedge Q(i_*E)
=\int_S\operatorname{ch}(E')
\sqrt{\frac{\widehat A(TS)}{\widehat A(N)}}\wedge i^*C,
\end{equation}
其中，$E'=E\otimes K_S^{-1/2}$。半典范扭曲按相应的 $\operatorname{spin}^c$ 结构解释；若平方根存在，就是通常的张量积。\BookSectionRef{sec:5.3.3.2}{5.3.3.2}将使用这一表达式。

\subsection{共形场论中的边界条件}\label{sec:5.1.5}

\SourceParagraph{ch05-5.1.5-001}{312}
最后，与\BookChapterRef{2}的讨论联系。指标公式\BookEqRef{eq:5.6}{5.6}是 Cardy 条件\BookEqRef{eq:2.12}{2.12}的特殊情形：
\[
\theta_b\bigl((\iota_b\circ\iota^a)(1)\bigr)
=\sum_{\psi_\mu\in\mathcal H_{ab}}\psi_\mu\psi^\mu.
\]
这里采用\BookSectionRef{sec:2.2}{2.2}的分次对偶基约定。它表示，具有边界条件 $a,b$ 的圆环，既有闭弦解释，也有开弦解释。闭弦一侧，在完整弦理论中，是无质量 RR 场的交换；如\BookSectionRef{sec:3.3.4}{3.3.4}所述，它们对应于拓扑闭弦态。开弦一侧，则是计数 $\mathbb{Z}_2$ 分次拓扑开弦态的指标。

\SourceParagraph{ch05-5.1.5-002}{312}
将已知的膜与 RR 场的几何耦合代入左边的 $\iota^a(1)$ 及 $\theta_b(\iota_b(\phi))$，就得到前面的表达式。

\SourceParagraph{ch05-5.1.5-003}{312}
原则上，对于任意具体定义的 SCFT，都可以类似地导出膜的拓扑类与相交形式。\BookSectionRef{sec:5.4.5.6}{5.4.5.6}将以轨形为例说明。

\section{稳定性预备知识}\label{sec:5.2}

\subsection{A 膜}\label{sec:5.2.1}

\SourceParagraph{ch05-5.2.1-intro-001}{312}
先回顾特殊 Lagrange 子流形的定义。然后，通过展示只能在复结构模空间的某个“边际稳定性墙”的一侧发生的束缚过程——如\BookEqRef{eq:5.1}{5.1}中的过程——说明：在所有 Lagrange 子流形中，哪些属于“特殊”子流形，以非平凡的方式依赖于 $Y$ 的复结构。

\subsubsection{特殊 Lagrange 子流形}\label{sec:5.2.1.1}

\SourceParagraph{ch05-5.2.1.1-001}{312}
设 $Y$ 是 $m$ 维 Calabi--Yau 流形。回顾\BookSectionRef{sec:1.1.4}{1.1.4}和\BookSectionRef{sec:3.6.1}{3.6.1}，特殊 Lagrange 子流形是一个定向 Lagrange 子流形 $L\subset Y$，对于某个实常数 $\xi(L)$，满足
\begin{equation}\label{eq:5.13}
\operatorname{Im}\bigl(\mathrm{e}^{-\mathrm{i}\pi\xi(L)}\Omega|_L\bigr)=0.
\end{equation}

\SourceParagraph{ch05-5.2.1.1-002}{312}
所谓“定向”，是指在 $L$ 上选择了一个处处非零的最高次形式。可以取这个形式为实体积形式 $\mathrm{d}V_L$。由\BookEqRef{eq:1.2}{1.2}，有
\begin{equation}\label{eq:5.14}
\mathrm{d}V_L=|Z|\,\mathrm{e}^{-\mathrm{i}\pi\xi(L)}\Omega|_L,
\end{equation}
其中，$|Z|$ 是某个正实常数。这确定了 $\xi(L)$，但仍允许加上 $2\mathbb{Z}$ 中的常数；规定 $0\le\xi(L)<2$，便可消除这一歧义。注意，这个参数就是\BookEqRef{eq:3.169}{3.169}中的参数，即
\begin{equation}\label{eq:5.15}
\xi(L)=\frac{1}{\pi}\arg\frac{\Omega|_L}{\mathrm{d}V_L}.
\end{equation}

\SourceParagraph{ch05-5.2.1.1-003}{313}
由于 $\xi(L)$ 是常数，也可以写成
\begin{equation}\label{eq:5.16}
\xi(L)=\frac{1}{\pi}\arg\int_L\Omega.
\end{equation}
除去 $1/\pi$ 因子，这就是与闭链 $L$ 对应的全纯 $m$ 形式周期的辐角。因此，此处的 $\xi(L)$ 只依赖于 $L$ 的同调类。\BookSectionRef{sec:5.3.4}{5.3.4}将把它提升为一个取值于 $\mathbb{R}/2m\mathbb{Z}$ 的量。

\SourceParagraph{ch05-5.2.1.1-004}{313}
当然，$\Omega$ 本身只确定到相差一个非零复常数；为了定义 $\xi(L)$，必须固定这一歧义。事实上，特殊 Lagrange 子流形的一种标准定义，便利用这个自由度令 $\xi=0$，然后要求 $\Omega|_L$ 的虚部为零。不过，我们需要比较不同特殊 Lagrange 子流形的 $\xi(L)$ 值。始终应当记住，只有 $\xi(L)$ 的相对值才有意义，另见\BookSectionRef{sec:5.3.4.2}{5.3.4.2}。

\SourceParagraph{ch05-5.2.1.1-005}{313}
周期本身，
\begin{equation}\label{eq:5.17}
Z(L)=\int_L\Omega,
\end{equation}
也将起重要作用。出于物理上的原因，$Z$ 称为 Dirichlet 膜的 \emph{BPS 中心荷}，或简称中心荷\footnote{这个名称来自该项在时空超对称代数中的作用，本书不使用这一代数。不过，从现在起，说到“中心荷”，指的都是 BPS 中心荷，而非\BookSectionRef{sec:3.2.2}{3.2.2}中的 Virasoro 中心荷 $c$ 与 $\hat c$。}。

\SourceParagraph{ch05-5.2.1.1-006}{313}
\BookStatementNumber{4}
\begin{example}[复环面]\label{ex:5.4}
设 $X\cong\mathbb{C}^n/\{z^i\sim z^i+1,\ z^i\sim z^i+\tau^{ij}\}$，取 $\Omega=\mathrm{d}z^1\wedge\cdots\wedge\mathrm{d}z^n$。选择 $n$ 对互素整数 $(a^i,b^i)$。超平面
\[
z^i=a^i\sigma^i+\sum_j\tau^{ij}b^j\sigma^j,\qquad \sigma^i\in[0,1),
\]
给出特殊 Lagrange 子流形，其中心荷为
\[
Z=\det(a^i\delta^{ij}+b^j\tau^{ij}).
\]
\end{example}

\SourceParagraph{ch05-5.2.1.1-007}{313}
由\BookChapterRef{1}的讨论，特殊 Lagrange 条件蕴含\BookSectionRef{sec:3.5.2.3}{3.5.2.3}所讨论的体积极小条件。物理上，它作为 A 膜满足 BPS 条件的要求而出现\BookCite{ref034}。这个条件通过时空超对称条件定义，与\BookEqRef{eq:1.2}{1.2}一样，
\[
\operatorname{Re}\bigl(\mathrm{e}^{\mathrm{i}\pi\xi}\epsilon^{\mathrm t}\Gamma\epsilon\bigr)|_M
=\operatorname{Vol}|_M.
\]
关键在于：为了让旋量 $\epsilon$ 定义一个超对称，这个关系必须在 $L$ 的每一点，都对同一个协变常旋量 $\epsilon$ 成立\footnote{物理上，参数 $\xi$ 决定从完整的 $\mathcal{N}=2$ 时空超对称中保留下来的 $\mathcal{N}=1$ 子代数。实际上，由于 $d=4$ 中的 $\mathcal{N}$ 个时空超对称带来 $U(\mathcal{N})$ R 对称，似乎应当有一个 $U(2)/U(1)$ 的选择。不过，$\mathcal{N}=2$ 时空超对称的谱流图景，只看到 R 对称的 $U(1)\times U(1)$ 子群，所以，这个参数仅取值于 $U(1)$。}。随后，物理论证表明：膜的“质量”——在当前讨论的几何极限中，就是 $L$ 的体积——必须大于或等于中心荷的绝对值，
\begin{equation}\label{eq:5.18}
M\equiv\int_L\operatorname{Vol}\ge |Z(L)|,
\end{equation}
等号成立当且仅当这个态是 BPS 态。这一物理论证在计入任何和全部修正，包括 $\alpha'$ 与 $g_s$ 修正以后，仍然有效。

\SourceParagraph{ch05-5.2.1.1-008}{314}
来看看 $\sigma$ 模型中的类似条件。首先，把\BookEqRef{eq:5.13}{5.13}改写为
\begin{equation}\label{eq:5.19}
\Omega|_L=\exp(2\mathrm{i}\pi\xi)\overline\Omega|_L.
\end{equation}
$\Omega$ 与 $\overline\Omega$ 在 A 模型中都有对应的算符：
\[
\Omega_{IJK}\psi_+^I\psi_+^J\psi_+^K;\qquad
\overline\Omega_{\bar I\bar J\bar K}\psi_-^{\bar I}\psi_-^{\bar J}\psi_-^{\bar K}.
\]
将它们代入\BookEqRef{eq:5.19}{5.19}，得到场的边界条件
\begin{equation}\label{eq:5.20}
\Omega_{IJK}\psi_+^I\psi_+^J\psi_+^K
=\exp(-2\mathrm{i}\pi\xi)\overline\Omega_{\bar I\bar J\bar K}
\psi_-^{\bar I}\psi_-^{\bar J}\psi_-^{\bar K}.
\end{equation}
这推广了\BookSectionRef{sec:3.6.6}{3.6.6}中的\BookEqRef{eq:3.196}{3.196}。\BookSectionRef{sec:5.3.4}{5.3.4}将利用这一识别，把这些定义和论证推广到任意 $(2,2)$ SCFT，无须假定 A 膜作为特殊 Lagrange 子流形的几何图景。

\SourceParagraph{ch05-5.2.1.1-009}{314}
回顾\BookSectionRef{sec:3.6.2}{3.6.2}中的\BookRef{def:3.3}{定义3.3}，A 膜还必须满足两个条件。显然，对于特殊 Lagrange 子流形，\BookEqRef{eq:3.170}{3.170}中的映射 $\xi_*$ 是平凡的，因此，Maslov 类条件自动满足。另一方面，量子阻碍条件并不自动满足。\BookSectionRef{sec:5.4.3.5}{5.4.3.5}和\BookSectionRef{sec:8.3.3}{8.3.3}将进一步讨论。

\SourceParagraph{ch05-5.2.1.1-010}{314}
注意，没有理由认为，Calabi--Yau 流形中的所有极小三维流形，都是特殊 Lagrange 子流形；这反映了，并非所有稳定 D 膜都一定是 BPS 态。马上就会给出一个简单的非光滑例子。

\subsubsection{几何衰变}\label{sec:5.2.1.2}

\SourceParagraph{ch05-5.2.1.2-001}{314}
在\BookSectionRef{sec:3.6.2}{3.6.2}中已经看到，暂且忽略量子阻碍条件时，一个携带线丛的 A 膜，其形变空间由 $H^1(L,\mathbb{C})$ 给出。\BookSectionRef{sec:6.1.1}{6.1.1}将说明，带平坦联络的特殊 Lagrange 子流形，其无穷小形变空间也由 $H^1(L,\mathbb{C})$ 给出。因此，局部而言，特殊 Lagrange 子流形的模空间，与 Lagrange 子流形模去 Hamilton 形变后的模空间一致。乍看之下，这似乎意味着，在 Lagrange 子流形模去 Hamilton 同痕的每个等价类中，都有唯一的特殊 Lagrange 子流形。

\SourceParagraph{ch05-5.2.1.2-002}{314}
事实并非如此。绝大多数 Lagrange 子流形，都不存在与之 Hamilton 同痕的特殊 Lagrange 子流形。从我们的角度，最适合说明这一点的方法，是考察靶空间 $Y$ 的复结构发生形变时，特殊 Lagrange 子流形如何“消失”或“衰变”。注意，Lagrange 子流形纯粹由 $Y$ 上 K\"ahler 形式诱导的辛结构定义，因而不依赖于复结构。在特殊 Lagrange 子流形中，加入“特殊”二字，才引入了对复结构的依赖。

\SourceParagraph{ch05-5.2.1.2-003}{315}
与\BookSectionRef{sec:5.1}{5.1}一样，衰变过程的逆过程是“束缚”过程，其中，两个或更多膜结合为一个。遵循 Joyce\BookCite{ref262}，现在用相交点附近的局部行为，给出这一过程的具体图景。

\SourceParagraph{ch05-5.2.1.2-004}{315}
考虑一族 Calabi--Yau $m$ 维流形 $Y_z$，相近的复结构由绝对值很小的 $z\in\mathbb{C}$ 参数化。假定 $Y_0$ 包含两个特殊 Lagrange 子流形 $L_1,L_2$，它们在点 $p$ 横截相交，且在 $z=0$ 时，有 $\xi(L_1)=\xi(L_2)$。

\SourceParagraph{ch05-5.2.1.2-005}{315}
可以证明，若一个光滑特殊 Lagrange 子流形 $L$ 在 $z=0$ 存在，那么它也在 $0$ 的某个邻域中存在\BookCite{ref349,ref262}。因此，把注意力限制在 $L_1,L_2$ 都存在的区域，并按照\BookEqRef{eq:5.15}{5.15}定义 $\xi(L_1),\xi(L_2)$。再限制 $z\in\mathbb{R}$，并要求 $|\xi(L_1)-\xi(L_2)|$ 很小。于是，
\BookStatementNumber{5}
\begin{theorem}[Joyce，\BookCite{ref262}，定理9.10]\label{thm:5.5}
存在接近连通和 $L_1\cup L_2$ 的特殊 Lagrange 子流形 $L_1\#L_2\subset Y_z$，当且仅当 $\xi(L_2)\le\xi(L_1)$。
\end{theorem}

\SourceParagraph{ch05-5.2.1.2-006}{315}
根据 $\xi(L_2)-\xi(L_1)$ 的符号，把复结构模空间分成两个区域 $\mathcal{M}^+$ 与 $\mathcal{M}^-$，它们由一个“边际稳定性墙”分隔。可以看到，特殊 Lagrange 子流形 $L_1\#L_2$ 只存在于 $\mathcal{M}^+$ 中。这个记号特意是不对称的，因为 $L_1\#L_2$ 与 $L_2\#L_1$ 很不相同。

\SourceParagraph{ch05-5.2.1.2-007}{315}
相反，在 $\mathcal{M}^-$ 中，没有使 $L_1\cup L_2$ 的体积极小的光滑特殊 Lagrange 子流形；对于很小的 $|z|$，$L_1\cup L_2$ 本身就是体积极小配置。物理上，这是“非 BPS”或“超对称破缺”配置。

\subsubsection{\texorpdfstring{$\mathbb{C}^m$}{ℂᵐ} 中的特殊 Lagrange 平面}\label{sec:5.2.1.3}

\SourceParagraph{ch05-5.2.1.3-001}{315}
\BookRef{thm:5.5}{定理5.5}的证明，考察相交点 $p$ 的一个邻域，用 $\mathbb{C}^m$ 建模，并把膜 $L_1,L_2$ 在其中用实 $m$ 维平面建模。

\SourceParagraph{ch05-5.2.1.3-002}{315}
因此，考虑线性子空间 $\mathbb{R}^m\subset\mathbb{C}^m$。可以通过角 $\phi_j$ 指定其嵌入，
\begin{equation}\label{eq:5.21}
\Pi^{\vec\phi}
=\{(\mathrm{e}^{\mathrm{i}\phi_1}x^1,\mathrm{e}^{\mathrm{i}\phi_2}x^2,
\ldots,\mathrm{e}^{\mathrm{i}\phi_m}x^m):x^j\in\mathbb{R}\}.
\end{equation}
用标准全纯 $m$ 形式 $\Omega=\mathrm{d}z^1\wedge\mathrm{d}z^2\wedge\cdots\wedge\mathrm{d}z^m$，得到
\begin{equation}\label{eq:5.22}
\begin{aligned}
\xi(\Pi^{\vec\phi})
&=\frac{1}{\pi}\arg\int_{\Pi^{\vec\phi}}\Omega\\
&=\frac{1}{\pi}\sum_{j=1}^m\phi_j\pmod 2.
\end{aligned}
\end{equation}
用 $\Pi^{\vec 0}$ 表示定向为 $\mathrm{d}x^1\wedge\cdots\wedge\mathrm{d}x^m$、嵌入角 $\phi_1=\phi_2=\cdots=\phi_m=0$ 的平面，因此，$\xi=0$。

\SourceParagraph{ch05-5.2.1.3-003}{316}
一般而言，角 $\phi_j$ 只确定到相差移位 $\phi_i\mapsto\phi_i+2\pi$，以及保持定向的变换
\[
\begin{aligned}
R_{ij}\colon(x^1,\ldots,x^i,\ldots,x^j,\ldots)
&\longmapsto(x^1,\ldots,-x^i,\ldots,-x^j,\ldots),\\
\phi_i&\longmapsto\phi_i\pm\pi,\qquad
\phi_j\longmapsto\phi_j\pm\pi.
\end{aligned}
\]
可以通过要求
\begin{equation}\label{eq:5.23}
\begin{aligned}
&0\le\phi_1<2\pi,\\
&0\le\phi_i<\pi,\qquad 2\le i\le m,
\end{aligned}
\end{equation}
消除这一歧义。

\SourceParagraph{ch05-5.2.1.3-004}{316}
例如，与 $\Pi^{\vec 0}$ 嵌入相同、定向相反的膜，记作 $\Pi^{(\pi,0,\ldots,0)}$。它的 $\xi\equiv1\pmod2$。由于它是 $\Pi^{\vec 0}$ 的反膜，也可以记作 $\overline\Pi^{\vec 0}$。

\subsubsection{Lawlor 颈}\label{sec:5.2.1.4}

\SourceParagraph{ch05-5.2.1.4-001}{316}
考虑两个 D 膜 $\Pi^{(\pi,0,\ldots,0)}\cup\Pi^{\vec\phi}$，在原点横截相交；这要求所有 $\sin\phi_i\ne0$。根据\BookEqRef{eq:5.22}{5.22}，它们的相位差为
\begin{equation}\label{eq:5.24}
\Delta\xi\equiv-1+\frac{1}{\pi}\sum_{j=1}^m\phi_j\pmod2.
\end{equation}
利用\BookEqRef{eq:5.23}{5.23}，把角取在基本区域中，可得这个相位差的代表值 $\Delta\xi\in(-1,m)$。顺便注意，这大于人们对 $\pi\Delta\xi$ 这样的角变量，直观上会预期的范围；\BookSectionRef{sec:5.3.4}{5.3.4}将进一步说明。

\SourceParagraph{ch05-5.2.1.4-002}{316}
现在假定 $\Delta\xi=0$；显然，只有 $m\ge2$ 时才可能如此。那么，可以把这组膜看作一个奇异特殊 Lagrange 子流形。它是 Lawlor\BookCite{ref321}找到的光滑特殊 Lagrange 子流形的缩放极限，描述如下\BookCite{ref223,ref262}。令
\begin{equation}\label{eq:5.25}
P(x)=\frac{\prod_{j=1}^m(1+a_jx^2)-1}{x^2}.
\end{equation}
还固定一个正实数 $A$，它将确定解的整体尺度。

\SourceParagraph{ch05-5.2.1.4-003}{316}
正实数 $a_1,\ldots,a_m$ 由 $A$ 与 $\phi_1,\phi_2,\ldots,\phi_m$，通过下面的方程隐式且唯一地确定：
\begin{equation}\label{eq:5.26}
\begin{aligned}
\phi_j&=a_j\int_{-\infty}^{\infty}
\frac{\mathrm{d}x}{(1+a_jx^2)\sqrt{P(x)}},\\
A&=\frac{\omega_m}{\sqrt{a_1\cdots a_m}},
\end{aligned}
\end{equation}
其中，$\omega_m$ 是 $\mathbb{R}^m$ 中单位球的体积。这些方程有解，要求 $\sum_j\phi_j=\pi$，恰好就是前面要求的 $\xi=1$ 条件。

\SourceParagraph{ch05-5.2.1.4-004}{316}
现在，定义函数 $\eta_j\colon\mathbb{R}\to\mathbb{C}$，
\begin{equation}\label{eq:5.27}
\eta_j(y)=\exp\!\left(\mathrm{i}a_j\int_{-\infty}^y
\frac{\mathrm{d}x}{(1+a_jx^2)\sqrt{P(x)}}\right)
\sqrt{\frac{1}{a_j}+y^2}.
\end{equation}
由此，可以定义 Lawlor 颈为
\begin{equation}\label{eq:5.28}
\begin{aligned}
L^{\vec\phi,A}=\bigl\{& (\eta_1(y)x^1,\eta_2(y)x^2,\ldots,\eta_m(y)x^m):\\
&y\in\mathbb{R},\ x^j\in\mathbb{R},\ (x^1)^2+\cdots+(x^m)^2=1\bigr\}.
\end{aligned}
\end{equation}

\SourceParagraph{ch05-5.2.1.4-005}{317}
可以证明，这是 $\mathbb{C}^m$ 中的光滑特殊 Lagrange 子流形，拓扑为 $S^{m-1}\times\mathbb{R}$。远离原点的区域由两个极限给出：$y\to-\infty$ 时渐近于 $\overline\Pi^{\vec0}$，$y\to+\infty$ 时渐近于 $\Pi^{\vec\phi}$。取极限 $A\to0$ 时，相交区域缩小到零，极限中得到 $\Pi^{(\pi,0,\ldots,0)}\cup\Pi^{\vec\phi}$。\BookRef{fig:5.1}{图1}给出示意。
\begin{figure}[H]
\centering
\BookImage{01_Numbered_Figures/Chapter_05/Ch05_Fig01_PDFp317_Bookp307.pdf}
\caption{Lawlor 颈。}\label{fig:5.1}
\end{figure}

\subsubsection{角度定理}\label{sec:5.2.1.5}

\SourceParagraph{ch05-5.2.1.5-001}{317}
现在考虑一个相近的平面 $\Pi^{\vec\phi'}$，使 $\Delta\xi'\ne0$。显然，在这种情形中，不存在同时渐近于这两个平面的光滑特殊 Lagrange 子流形。这提示，这样的配置可能不再使体积极小。

\SourceParagraph{ch05-5.2.1.5-002}{317}
要构造一个渐近行为相同、体积却更小的子流形，可以考虑把 $\Pi^{(\pi,0,\ldots,0)}\cup\Pi^{\vec\phi'}$ 与一个 Lawlor 颈 $L^{\vec\phi,A}$ 组合。\BookRef{fig:5.2}{图2}粗略画出 $m=1$ 时的情形，尽管此时 Lawlor 颈实际上并不存在。直观上，当这对平面张开的角更大时，这种方法奏效；这确实正确，条件是
\begin{equation}\label{eq:5.29}
0<\Delta\xi=\xi(\Pi^{\vec\phi'})-\xi(\Pi^{(\pi,0,\ldots,0)}).
\end{equation}
那么，用 Lawlor 颈的内部，替换 $\Pi^{(\pi,0,\ldots,0)}\cup\Pi^{\vec\phi'}$ 的内部，就得到一个渐近行为相同的连续子流形。不过，由于 $L^{\vec\phi,A}$ 是校准子流形，而两个平面的并不是，它在内部具有更小的体积。随后可以把接合处光滑化，得到光滑子流形。
\begin{figure}[H]
\centering
\BookImage{01_Numbered_Figures/Chapter_05/Ch05_Fig02_PDFp318_Bookp308.pdf}
\caption{平面与 Lawlor 颈的组合。}\label{fig:5.2}
\end{figure}

\SourceParagraph{ch05-5.2.1.5-003}{317}
进一步发展这一论证，得到
\BookStatementNumber{6}
\begin{theorem}[Lawlor--Nance，见\BookCite{ref223}，定理7.134]\label{thm:5.6}
$\mathbb{R}^{2p}$ 中两个定向 $p$ 维子空间的并使体积极小，当且仅当，对于复结构的任意选择，它们定向的相对角都不满足\BookEqRef{eq:5.29}{5.29}。
\end{theorem}

\SourceParagraph{ch05-5.2.1.5-004}{318}
如果主要关心的是 $\mathbb{R}^{2p}$，可以通过适当排列平面之间的相对角，重新表述这个条件，使之不再提到复结构。

\subsubsection{嵌入 Calabi--Yau 几何}\label{sec:5.2.1.6}

\SourceParagraph{ch05-5.2.1.6-001}{318}
现在，回到两个 A 膜 $L_1,L_2$ 在 $p\in Y$ 横截相交的情形。可以使用一个 $U(m)$ 变换，把 $L_1$ 的切平面，旋转到上述标准平面 $\Pi^{(\pi,0,\ldots,0)}$。同一个 $U(m)$ 矩阵，把 $L_2$ 的切平面旋转到 $\Pi^{\vec\phi}$。

\SourceParagraph{ch05-5.2.1.6-002}{318}
直观上，把注意力限制在 $p$ 的小邻域中，应当恰好得到刚才描述的情形。文献\BookCite{ref262}证明了这一点，从而得到\BookRef{thm:5.5}{定理5.5}。

\SourceParagraph{ch05-5.2.1.6-003}{318}
回顾一下：在 $\mathcal{M}^+$ 中，有一个 BPS 态 $L_1\#L_2$。也有 BPS 态 $L_1,L_2$，但 $L_1\#L_2$ 的质量，小于 $L_1$ 与 $L_2$ 的质量之和。到达墙时，$L_1\#L_2$ 变成 $L_1\cup L_2$。越过墙，进入 $\mathcal{M}^-$ 后，只剩 BPS 态 $L_1,L_2$，它们放在一起会破坏超对称。刚才描述的是：从 $\mathcal{M}^+$ 进入 $\mathcal{M}^-$ 时，BPS 态 $L_1\#L_2$ 衰变为产物 $L_1,L_2$。在 $\mathcal{M}^+$ 中，把 $L_1\#L_2$ 看作 $L_1,L_2$ 的束缚态。\BookRef{fig:5.3}{图3}描绘这一过程。

\SourceParagraph{ch05-5.2.1.6-004}{318}
远离墙时，$L_1\cup L_2$ 不再是特殊 Lagrange 子流形，因为两个分支的 $\xi$ 值不同。光滑空间 $L_1\#L_2$ 与 $L_1\cup L_2$ 同调，并使体积极小，也就是使缠绕这些闭链的 D 膜能量极小。因此，
\begin{equation}\label{eq:5.30}
\left|\int_{L_1\#L_2}\Omega\right|
<\left|\int_{L_1}\Omega\right|+\left|\int_{L_2}\Omega\right|,
\end{equation}
并且，可以选择
\begin{equation}\label{eq:5.31}
\xi(L_2)<\xi(L_1\#L_2)<\xi(L_1).
\end{equation}
\begin{figure}[H]
\centering
\BookImage{01_Numbered_Figures/Chapter_05/Ch05_Fig03_PDFp319_Bookp309.pdf}
\caption{A 膜的衰变。}\label{fig:5.3}
\end{figure}

\SourceParagraph{ch05-5.2.1.6-005}{318}
物理上，这也与四维 $\mathcal{N}=2$ 理论中 BPS 态的标准性质一致，例如，Seiberg 与 Witten\BookCite{ref415}研究过这些性质。相关介绍见文献\BookCite{ref329}，另见\BookBibText{ref242}{MS1}第19.3节。

\subsection{B 膜}\label{sec:5.2.2}

\SourceParagraph{ch05-5.2.2-intro-001}{319}
刚才讨论的现象，其镜像应当是：随着复化 K\"ahler 结构改变，B 型膜的集合也发生变化。与\BookSectionRef{sec:3.4.1.1}{3.4.1.1}和\BookSectionRef{sec:3.6.2.2}{3.6.2.2}中关于关联函数的讨论一样，预期问题的这一方面会受到世界面瞬子修正，也可能受到其他 $\alpha'$ 修正。不过，可以先考虑 Calabi--Yau 三维流形 $X$ 的大半径极限；在这一极限中，世界体上的联络服从 Yang--Mills 方程。

\SourceParagraph{ch05-5.2.2-intro-002}{319}
考虑缠绕 $X$ 的一个6膜，它对应于全纯向量丛 $E\to X$，曲率 $F$ 为 $(1,1)$ 形式。在大半径极限中，BPS 条件化为 Hermitian--Yang--Mills 条件\BookCite{ref347}。也就是说，曲率张量满足
\begin{equation}\label{eq:5.32}
g^{j\bar k}F^{\ \beta}_{\alpha j\bar k}
=\mu(E)\,\delta_\alpha^\beta,
\end{equation}
其中，$\alpha,\beta$ 是 $E$ 的纤维指标，$\mu(E)$ 是实数，称为 $E$ 的“斜率”。例如，遵循文献\BookCite{ref453}的分析，积分\BookEqRef{eq:5.32}{5.32}，得到\footnote{这里遵循文献\BookCite{ref453}的约定，不过，另一种常见的 $\mu$ 定义，会去掉\BookEqRef{eq:5.33}{5.33}中的 $\operatorname{Vol}(X)$ 因子。}
\begin{equation}\label{eq:5.33}
\mu(E)=\frac{\deg E}{k\,\operatorname{Vol}(X)},
\end{equation}
其中，$k$ 是 $E$ 的秩，而
\begin{equation}\label{eq:5.34}
\deg E=\int_XJ\wedge J\wedge c_1(E)
\end{equation}
是丛 $E$ 的次数。

\SourceParagraph{ch05-5.2.2-intro-003}{320}
与通常的超对称真空条件一样，这是场，即联络，所满足的一阶微分方程，并蕴含场满足运动方程，即作为二阶微分方程的 Yang--Mills 方程。

\SourceParagraph{ch05-5.2.2-intro-004}{320}
Hermitian--Yang--Mills 条件\BookEqRef{eq:5.32}{5.32}显式依赖于度规，因而依赖于 K\"ahler 形式 $J$。所以，解是否存在，可以依赖于 $J$。事实上，它只依赖于下面这个代数几何条件：
\BookStatementNumber{7}
\begin{definition}\label{def:5.7}
局部自由层 $E$ 称为 \emph{$\mu$ 稳定}的，如果其每个满足 $0<\operatorname{rk}F<\operatorname{rk}E$ 的子层 $F$ 都满足
\begin{equation}\label{eq:5.35}
\mu(F)<\mu(E).
\end{equation}
\end{definition}

\SourceParagraph{ch05-5.2.2-intro-005}{320}
于是，
\BookStatementNumber{8}
\begin{theorem}[Donaldson\BookCite{ref128}；Uhlenbeck 与 Yau\BookCite{ref453}]\label{thm:5.8}
一个丛允许不可约 Hermitian--Yang--Mills 联络，当且仅当它是 $\mu$ 稳定的。
\end{theorem}

\SourceParagraph{ch05-5.2.2-intro-006}{320}
因此，在忽略 $\alpha'$ 修正的大体积极限中，不可约 BPS B 膜的稳定性等价于 $\mu$ 稳定性。

\SourceParagraph{ch05-5.2.2-intro-007}{320}
为了开始建立与\BookSectionRef{sec:5.2.1}{5.2.1}的类比，把子层 $F\to E$ 放入短正合列
\begin{equation}\label{eq:5.36}
0\longrightarrow F\xrightarrow{f}E\xrightarrow{g}G\longrightarrow0.
\end{equation}
由此，可以把 $E$ 看作 $F,G$ 的束缚态。这个束缚态是否存在，取决于条件\BookEqRef{eq:5.35}{5.35}。观察\BookEqRef{eq:5.33}{5.33}可知，当 $b^{1,1}>1$ 时，这可以依赖于 K\"ahler 锥中的射线。事实确实如此，人们在研究 Donaldson 不变量时就认识到了这一点\BookCite{ref399}；其物理解释由文献\BookCite{ref424}指出。后面将看到，这个过程正是\BookSectionRef{sec:5.2.1.2}{5.2.1.2}中衰变的镜像。

\SourceParagraph{ch05-5.2.2-intro-008}{320}
低维膜也可以类似分析。可以把注意力集中在子空间 $S\subset X$ 上，在 $S$ 上的丛类中寻找稳定向量丛；如果 $S$ 不是自旋流形，则考虑扭曲丛。这些子空间在如下意义上互不干扰：$X$ 上对应于 $E$ 的丛，不能衰变为只支撑在真子空间 $S$ 上的子层 $F$，因为不存在非零同态 $F\to E$。同样，$S$ 上的丛也不能衰变为在整个 $X$ 上具有正秩的子层，因为那样，\BookEqRef{eq:5.36}{5.36}中的商层 $G$ 就会具有负秩。

\SourceParagraph{ch05-5.2.2-intro-009}{321}
因此，$\mu$ 稳定性在大半径极限中确定了一组稳定 B 膜；它们是凝聚层中的一个子集，依赖于 K\"ahler 锥中的射线。

\subsubsection{作为辛商的 HYM}\label{sec:5.2.2.1}

\SourceParagraph{ch05-5.2.2.1-001}{321}
文献\BookCite{ref129}讨论了一种理解这一结果的好方法：把 HYM 方程\BookEqRef{eq:5.32}{5.32}的左边看作矩映射。首先注意，$E$ 上的联络空间具有自然的辛形式。给定两个切向量 $\beta,\gamma\in\Omega^1(X,\operatorname{ad}(E))$，其值为
\begin{equation}\label{eq:5.37}
\int_X\omega^{n-1}\wedge\operatorname{Tr}(\bar\beta\wedge\gamma).
\end{equation}
那么，以
\[
H=\int_X g^{j\bar k}\operatorname{tr}(F_{j\bar k}\epsilon)\,\operatorname{vol}_X
\]
为“Hamilton 函数”，由\BookEqRef{eq:5.75}{5.75}给出的演化，就生成参数为 $\epsilon$ 的规范变换。

\SourceParagraph{ch05-5.2.2.1-002}{321}
于是，寻找模去规范等价的解，就成为一个无穷维辛商问题。有限维时，根据几何不变量理论，这样的商空间等价于满足某个稳定性条件的复化对称群轨道空间，如\BookSectionRef{sec:4.7.3}{4.7.3}所述。由此类比，可以预期，HYM 解的存在性也受某个稳定性条件控制。

\SourceParagraph{ch05-5.2.2.1-003}{321}
Donaldson--Uhlenbeck--Yau 定理的一侧，即 $\mu$ 稳定性的必要性，可以利用一个消失定理来证明。
\BookStatementNumber{9}
\begin{lemma}\label{lem:5.9}
若 $V$ 允许不可约 Hermitian--Yang--Mills 联络，并且具有非零全纯截面 $s$，则
\[
\deg_J(V)=\int_X\omega^{n-1}\wedge c_1(V)\ge0.
\]
\end{lemma}
\begin{proof}
由
\[
0\le\int_X\omega^{i\bar j}(\nabla_{\bar j}s^*)(\nabla_i s)\,\operatorname{vol}_X
\]
以及分部积分即可得到。
\end{proof}

\SourceParagraph{ch05-5.2.2.1-004}{321}
接着，考虑 $E$ 中一个可能使它不稳定的子层 $F$。\BookEqRef{eq:5.36}{5.36}中的映射 $f$，是 $V\equiv\operatorname{Hom}(F,E)$ 的截面。在 $F$ 也是丛、且 $E,F$ 都允许不可约 Hermitian--Yang--Mills 联络的情形中，$V$ 也带有诱导的 Hermitian--Yang--Mills 联络。因此，由引理，
\[
c_1(V)=\operatorname{rk}(F)c_1(E)-\operatorname{rk}(E)c_1(F)
\]
的次数非负，也就是 $\mu(F)\le\mu(E)$。对于一般子层，可以通过正交投影与第二基本形式得到同样的不等式；不可约性排除真子对象达到等号的情形。

\SourceParagraph{ch05-5.2.2.1-005}{321}
一种更具体的理解方法，是把\BookEqRef{eq:5.36}{5.36}看作 $E$ 上单参数联络族的“Jordan 块”表示。其全纯结构的 $(0,1)$ 部分具体为
\[
A_E^{0,1}=\begin{pmatrix}
A_F^{0,1}&t\,\delta A_{G,F}\\
0&A_G^{0,1}
\end{pmatrix},
\]
其中，$\delta A_{G,F}\in H^{0,1}(X,\operatorname{Hom}(G,F))$，$t$ 是参数。还可以注意到，不同的 $t\in\mathbb{C}^*$ 对应于复规范变换。因此，GIT 稳定性要求：在把问题约化到这一单参数族以后，仍然稳定。

\SourceParagraph{ch05-5.2.2.1-006}{322}
物理上，可以计算 $t=0$ 附近 Yang--Mills 作用量对 $t$ 的依赖。得到
\[
\begin{aligned}
S_{\mathrm{YM}}(E)={}&S_{\mathrm{YM}}(A_F)+S_{\mathrm{YM}}(A_G)\\
&+C\bigl(\mu(F)-\mu(E)\bigr)|t|^2+O(|t|^3),
\end{aligned}
\]
其中，$C$ 是某个正常数。因此，若 $\mu(F)<\mu(E)$，则 $t\ne0$ 的形变使能量降低。相反，若 $\mu(F)>\mu(E)$，取 $t\to0$ 会降低能量，并在 $t=0$ 的极限中得到可约联络。

\SourceParagraph{ch05-5.2.2.1-007}{322}
注意，这与涉及\BookEqRef{eq:5.30}{5.30}的论证相似。事实上，若这些膜是 II 型弦在 Calabi--Yau 三维流形上紧化后得到的粒子，那么，四维中的可观测结果恰好相同：\BookEqRef{eq:5.36}{5.36}描述膜 $F,G$ 形成束缚态 $E$ 的一种可能过程，或相反的 $E$ 的衰变。\BookSectionRef{sec:5.4.4}{5.4.4}将说明，从四维物理，即有效场论的观点看，A 膜与 B 膜的这两个过程无法区分。镜像对称要成为理论之间的物理等价，就必须如此。

\subsubsection{\texorpdfstring{$\mu$}{mu} 稳定性的修正与 MMMSL 方程}\label{sec:5.2.2.2}

\SourceParagraph{ch05-5.2.2.2-001}{322}
文献\BookCite{ref347}计算了 Hermitian--Yang--Mills 条件的首项 $\alpha'$ 修正，得到 \emph{MMMSL 方程}，
\begin{equation}\label{eq:5.38}
\frac{1}{n!}(\omega\,1_E+2\pi\alpha'F)^{\wedge n}
=k\,\frac{\omega^n}{n!}\,1_E,
\end{equation}
其中，$n=\dim_{\mathbb C}X$，$k$ 是依赖于 $E$ 与 $\alpha'$ 的常数，可以通过在 $X$ 上积分\BookEqRef{eq:5.38}{5.38}的两边得到。文献\BookCite{ref347}从6膜上规范场的超对称 Born--Infeld 作用量导出这一方程。也可回顾导出等价条件\BookEqRef{eq:3.198}{3.198}的世界面论证。

\SourceParagraph{ch05-5.2.2.2-002}{322}
严格来说，这两种推导都只适用于 $E$ 为线丛的情形。事实上，\BookEqRef{eq:5.38}{5.38}也适用于一般的 $E$，\BookSectionRef{sec:5.5.5}{5.5.5}将说明这一点。

\SourceParagraph{ch05-5.2.2.2-003}{322}
Leung\BookCite{ref330}出于代数几何的原因，研究了相同的\BookEqRef{eq:5.38}{5.38}。他证明，对于固定的全纯向量丛 $E$，该方程对于所有充分小的 $\alpha'$ 都有解，当且仅当 $E$ 在 Gieseker 意义下稳定。这一稳定性条件仍然采取\BookEqRef{eq:5.35}{5.35}的一般形式，但斜率函数不同；Gieseker 斜率定义为
\begin{equation}\label{eq:5.39}
\begin{aligned}
p(E,r)&=\frac{1}{\operatorname{rk}E}\int_X
\operatorname{ch}(E)\mathrm{e}^{r\omega}\operatorname{td}(X)\\
&=\frac{1}{\operatorname{rk}E}\dim H^0(X,E\otimes\mathcal{L}^r),\qquad r\gg0,
\end{aligned}
\end{equation}
其中，$c_1(\mathcal{L})=\omega$。更明确地说，$E$ 是 Gieseker 稳定的，是指对于其任意非零真子层 $V$，在 $r$ 充分大时，都有 $p(V,r)<p(E,r)$。

\SourceParagraph{ch05-5.2.2.2-004}{322}
与 HYM 方程一样，MMMSL 方程也可以理解为：在 $E$ 的联络空间上，由规范变换群作辛商时的矩映射方程。与前面讨论的唯一区别，是\BookEqRef{eq:5.37}{5.37}被替换为
\[
\int_X\operatorname{Tr}\!\left(
\bar\beta\wedge\gamma\wedge\exp\!\left(\frac{F}{2\pi}+r\omega\right)
\right)\operatorname{td}(X).
\]
在 $r\to\infty$ 的极限中，除去整体归一化因子，这趋近于\BookEqRef{eq:5.37}{5.37}。

\section{B 膜范畴}\label{sec:5.3}

\SourceParagraph{ch05-5.3-intro-001}{323}
首先回顾\BookSectionRef{sec:4.4}{4.4}中导出范畴的基本构造，用更富物理意味的语言，重新表述其各个组成部分。在适当的地方，会说明新的要素如何进入。

\SourceParagraph{ch05-5.3-intro-002}{323}
希望用拓扑弦理论定义一个 B 膜范畴。对象是边界条件，态射则是扭曲超荷 $Q$ 的 $Q$ 上同调元素。

\SourceParagraph{ch05-5.3-intro-003}{323}
回顾\BookSectionRef{sec:3.6.3}{3.6.3}和\BookSectionRef{sec:3.6.4}{3.6.4}的讨论。那里论证了，B 膜范畴的对象包含 $X$ 上的全纯向量丛，态射是 Dolbeault 上同调群 $H_{\bar\partial}^{0,q}(X,\operatorname{Hom}(E,F))$ 中的元素。这一上同调的 $\mathbb{Z}$ 分次，对应于 R 荷，也称鬼数。正如已经讨论过的，这些上同调群也同构于 $\operatorname{Ext}^q(E,F)$。

\SourceParagraph{ch05-5.3-intro-004}{323}
现在希望说明，由这些对象与态射构造的复形，如何从\BookChapterRef{3}的定义中产生。第一步是取对象的直和；显然，Chan--Paton 构造允许这样做。不过，不能直接断言直和
\begin{equation}\label{eq:5.40}
E\equiv\bigoplus_{n\in\mathbb{Z}}E^n
\end{equation}
带有分次；这必须以某种方式蕴含在物理定义之中。

\SourceParagraph{ch05-5.3-intro-005}{323}
归根结底，这来自 R 荷守恒，\BookSectionRef{sec:3.6}{3.6}已经说明过。不过，让我们在这一步多作一点停留，因为它最终将引出许多来自物理的新特征。

\subsection{分次的物理来源}\label{sec:5.3.1}

\SourceParagraph{ch05-5.3.1-001}{323}
要说明的是：原以为完全由全纯丛 $E$ 的选择所确定的边界条件，实际上还携带另一个整数量子数 $n$；因此，\BookEqRef{eq:5.40}{5.40}中的不同项 $E^n$ 是不同的边界条件。此外，这个量子数必须对 R 荷作出贡献，使群 $\operatorname{Ext}^i(E^m,F^n)$ 中的开弦具有 R 荷
\begin{equation}\label{eq:5.41}
J=i+n-m.
\end{equation}
若如此，R 荷就成为\BookSectionRef{sec:4.4}{4.4}所用意义下的 $\mathbb{Z}$ 分次。

\SourceParagraph{ch05-5.3.1-002}{323}
这个量子数从何而来？\BookSectionRef{sec:5.3.4}{5.3.4}将详细解释，这里先给出基本思路。正如\BookSectionRef{sec:3.6}{3.6}和\BookSectionRef{sec:3.3.4.6}{3.3.4.6}讨论的，R 荷是 $U(1)$ 荷，因而由\BookSectionRef{sec:3.2.4}{3.2.4}中分离出来的 $U(1)$ 理论描述。除去\BookEqRef{eq:3.70}{3.70}表达的细节——\BookSectionRef{sec:5.3.4}{5.3.4}会讨论它们——这个 $U(1)$ 理论，基本上就是玻色弦理论中靶空间为 $S^1$ 的理论。因此，可以借用靶空间 $S^1$ 中 Dirichlet 膜的讨论，理解这个扇区中可能的边界条件。

\SourceParagraph{ch05-5.3.1-003}{324}
\BookSectionRef{sec:3.5.4}{3.5.4}讨论过这个问题，结论是，有两种可能：一种是 D0膜，由 $S^1$ 上的坐标参数化；另一种是 D1膜，由平坦联络参数化，等价地，由对偶 $S^1$ 上的坐标参数化。与\BookEqRef{eq:3.167}{3.167}比较可知，B 膜类似于 D1膜，而 A 膜类似于 D0膜。

\SourceParagraph{ch05-5.3.1-004}{324}
无论哪一种，都找到了唯一一个额外量子数，即某个 $S^1$ 上的坐标。虽然将论证，这就是正在寻找的量子数 $n$，但在弄清二者的关系以前，先用不同的符号 $\mu$ 表示它。首先检验能否识别二者：令 $\mu(E)=m$、$\mu(F)=n$，则\BookEqRef{eq:5.41}{5.41}与\BookEqRef{eq:3.189}{3.189}相同，因此，复形内部的分次与 R 荷之间的关系恰好正确。

\SourceParagraph{ch05-5.3.1-005}{324}
当然，这一解释面临一个问题：$\mu$ 取值于 $S^1$，而非 $\mathbb{Z}$。不过，视 $\mu$ 的周期，也就是 $S^1$ 的周长而定，这一区别未必是本质的。\BookSectionRef{sec:4.5}{4.5}已经讨论过，在复维数为 $d$ 的流形上，层的消解可取长度 $d$。因此，只要新量子数模 $N$ 定义（从而得到 $\mathbb{Z}/N\mathbb{Z}$ 分次），且 $N\ge2d$，就能区分复形之间所需的全部映射。在这种情形中，把取值限制为量子化的值，就得到一个 $\mathbb{Z}/N\mathbb{Z}$ 量子数。

\SourceParagraph{ch05-5.3.1-006}{324}
为了使这一解释成立，在验证\BookEqRef{eq:5.41}{5.41}以后，必须检查 $N\ge2\hat c$。物理上，\BookEqRef{eq:5.41}{5.41}中的移位，直接来自定义\BookEqRef{eq:3.144}{3.144}与\BookEqRef{eq:3.167}{3.167}；它们说明，$i$ 与 $n$ 以完全相同的方式进入守恒 R 荷 $\partial_\tau\varphi$。A 模型图景中，R 荷是弦的缠绕数，因此，可以给出更直观的论证。通过 T 对偶，或镜像对称，两种论证同样有效。

\SourceParagraph{ch05-5.3.1-007}{324}
由此，分次的周期 $N$ 应当来自 $S^1$ 的周长 $2\pi R$。仔细追踪\BookSectionRef{sec:3.3.4.6}{3.3.4.6}的定义，就可以计算这一周期，得到 $2\pi R=2\hat c$，与 $N=2\hat c$ 一致。实际上，这种说法略有不精确，因为 R 荷本身是整值守恒荷；使分次 $n$ 成为周期变量的，是\BookEqRef{eq:3.70}{3.70}中求和的性质。不过，最终结论将是正确的。

\SourceParagraph{ch05-5.3.1-008}{324}
还需要理解的一点，是为什么\BookEqRef{eq:5.40}{5.40}中的分次取整数，而量子数 $\mu$ 未必取整数。实际上，此处应当从物理中获得启发，接受取值于 $\mathbb{R}$ 或 $\mathbb{R}/N\mathbb{Z}$ 的分次。当然，复形的定义要求各项的分次之差为整数。在这一限制下，可以检查，作出这个推广以后，后面的定义仍然有意义。称所得结构为 $\mathbb{R}$ 分次范畴。

\SourceParagraph{ch05-5.3.1-009}{325}
\BookSectionRef{sec:5.3.4}{5.3.4}将把 $n$ 与 $\mu$ 都联系到\BookSectionRef{sec:5.2}{5.2}的分次 $\xi$，从而说明\BookEqRef{eq:5.41}{5.41}与\BookEqRef{eq:3.189}{3.189}之间的关系。这将引出分次的流动，以及稳定性方案中的其他重要要素。不过，既然已经解释了\BookEqRef{eq:5.40}{5.40}的来源和意义，就暂且把这些发展留待后文。

\subsection{形变与复形}\label{sec:5.3.2}

\SourceParagraph{ch05-5.3.2-intro-001}{325}
现在已经看到，一旦计入 R 荷，拓扑开弦理论中的边界条件，就是\BookEqRef{eq:5.40}{5.40}形式的分次直和，
\begin{equation}\label{eq:5.42}
E=\bigoplus_{n\in\mathbb{Z}}E^n,
\end{equation}
其中，$E^n$ 是\BookChapterRef{3}讨论的 A 型或 B 型边界条件。现在希望允许这些直和发生形变，并确定这类对象的完整集合。这些对象两两之间的映射，其 $Q$ 上同调类，就构成拓扑开弦的完整集合。

\SourceParagraph{ch05-5.3.2-intro-002}{325}
假定原来的 A 型或 B 型边界条件，是加法范畴 $\mathcal{A}$ 中的对象；为具体起见，可以考虑 $X$ 上的局部自由层。那么，
\BookStatementNumber{10}
\begin{proposition}\label{prop:5.10}
在具有 R 荷的拓扑扭曲理论中，膜的完整集合对应于同伦范畴 $K(\mathcal{A})$。
\end{proposition}

\SourceParagraph{ch05-5.3.2-intro-003}{325}
为说明这一点，先考虑从 $E$ 到自身的开弦，即 $\operatorname{Ext}^*(E,E)$ 中元素的线性组合。按照\BookEqRef{eq:5.41}{5.41}，R 荷为 $q$ 的开弦，是 $\operatorname{Ext}^k(E^n,E^{n-k+q})$ 中的元素，其中，$n,k$ 任意。

\SourceParagraph{ch05-5.3.2-intro-004}{325}
与\BookSectionRef{sec:3.6.4}{3.6.4}一样，R 荷为1的手征算符，可以对应于形变；一般而言，形变可能受阻。那里讨论了 $\operatorname{Ext}^1(E^n,E^n)$ 的情形；对于丛，这对应于改变联络。

\SourceParagraph{ch05-5.3.2-intro-005}{325}
$\operatorname{Hom}\cong\operatorname{Ext}^0$ 的情形，对应于改变复形的微分：
\begin{equation}\label{eq:5.43}
\begin{aligned}
d&=\sum_n d_n,\\
d_n&\in\operatorname{Ext}^0(E^n,E^{n+1})
=\operatorname{Hom}(E^n,E^{n+1}).
\end{aligned}
\end{equation}
令 $W_d^{(1)}$ 是满足
\begin{equation}\label{eq:5.44}
\{Q,W_d^{(1)}\}=\mathrm{d}_\Sigma d
\end{equation}
的算符，其中，$\mathrm{d}_\Sigma$ 是世界面上的 de Rham 算符。把作用量形变为
\begin{equation}\label{eq:5.45}
S=S_0+\oint_{\partial\Sigma}W_d^{(1)}.
\end{equation}

\SourceParagraph{ch05-5.3.2-intro-006}{326}
遵循通常的 Noether 方法，可以说明，这会使 BRST 荷变为
\begin{equation}\label{eq:5.46}
Q=Q_0+d.
\end{equation}
因此，为保持 $Q^2=0$，必须要求
\begin{equation}\label{eq:5.47}
\{Q_0,d\}+d^2=0.
\end{equation}
根据\BookEqRef{eq:5.43}{5.43}，知道 $\{Q_0,d\}=0$，所以，这个条件化为 $d^2=0$。用\BookEqRef{eq:5.43}{5.43}中的相继映射展开，便成为
\begin{equation}\label{eq:5.48}
d_{n+1}d_n=0,\qquad\text{对所有 }n.
\end{equation}
于是，物理自洽条件，即 BRST 荷的幂零性，等价于 $E$ 为复形的数学条件，
\begin{equation}\label{eq:5.49}
\cdots\xrightarrow{d_{n-1}}E^n\xrightarrow{d_n}E^{n+1}
\xrightarrow{d_{n+1}}E^{n+2}\xrightarrow{d_{n+2}}\cdots.
\end{equation}

\SourceParagraph{ch05-5.3.2-intro-007}{326}
因此，更一般地，膜由 $\mathcal{A}$ 中对象构成的复形表示。复形中的映射，表示从最初简单的对象集合出发的形变。注意，对象 $E$ 本身，也以一个相当平凡的方式成为复形：
\begin{equation}\label{eq:5.50}
\cdots\xrightarrow{0}0\xrightarrow{0}E\xrightarrow{0}0\xrightarrow{0}\cdots.
\end{equation}
更一般地，对象在复形中的位置，就是\BookSectionRef{sec:5.3.1}{5.3.1}所述，膜在与 R 荷对应的 $S^1$ 中的位置。

\subsubsection{开弦}\label{sec:5.3.2.1}

\SourceParagraph{ch05-5.3.2.1-001}{326}
上述形变也会影响 B 膜之间的开弦谱。本节计算相应的开弦态 Hilbert 空间。

\SourceParagraph{ch05-5.3.2.1-002}{326}
考虑从由复形 $E$ 表示的 B 膜，到由复形 $F$ 表示的另一个 B 膜的开弦。更准确地说，先取一组对象 $E^0,E^1,\ldots\in\mathcal{A}$，以及另一组局部自由层 $F^0,F^1,\ldots\in\mathcal{A}$，其中，$E^n,F^n$ 的 R 荷为 $n$。现在，对理论作形变，把这组 $E$ 变成\BookEqRef{eq:5.49}{5.49}中的复形，边界映射为 $d_n^E$；对 $F$ 也作同样的处理。

\SourceParagraph{ch05-5.3.2.1-003}{326}
完成以后，总 BRST 算符包含 $Q_0$ 以及两端的复形微分。对于次数为 $|f|$ 的映射，其作用为
\begin{equation}\label{eq:5.51}
Q(f)=Q_0(f)+d_F\circ f-(-1)^{|f|}f\circ d_E.
\end{equation}
它作用于从复形 $E$ 到复形 $F$ 的映射
\[
f^{m,n}\colon E^m\longrightarrow F^n
\]
所构成的直和上。根据定义，拓扑开弦属于这一算符的上同调空间。

\SourceParagraph{ch05-5.3.2.1-004}{326}
假定各个映射都满足 $Q_0f^{m,n}=0$，即它们是 $\operatorname{Hom}(E^m,F^n)$ 中的元素。那么，计入符号以后，条件 $Qf=0$，恰好表示 $f$ 是复形的态射。此外，两组映射若相差一个 BRST 恰当量，即 $f'=f+Qh$，就是\BookSectionRef{sec:4.4.2}{4.4.2}定义的同伦态射。

\SourceParagraph{ch05-5.3.2.1-005}{327}
因此，在这一复形映射模型中，复形之间的拓扑开弦，恰好是同伦范畴 $K(\operatorname{Coh}X)$ 的态射。

\SourceParagraph{ch05-5.3.2.1-006}{327}
既然定义一致，就可以采用\BookSectionRef{sec:4.4.4}{4.4.4}中的记号 $\operatorname{Ext}^n(E,F)$，表示从 $E$ 到 $F$、R 荷为 $n$ 的开弦的上同调空间。也可以采用记号 $[n]$，表示使一组膜的 R 荷发生相应移位的操作。更准确地说，若 $F$ 的第 $q$ 项为 $F^q$，则 $F[n]$ 的第 $q$ 项为 $F^{q+n}$。于是，
\begin{equation}\label{eq:5.52}
\operatorname{Ext}^q(E[m],F[n])
=\operatorname{Ext}^{q-m+n}(E,F),
\end{equation}
所有次数为 $q$ 的态射，都是 $\operatorname{Hom}(E,F[q])$ 中的元素。

\SourceParagraph{ch05-5.3.2.1-007}{327}
\begin{example*}
一个简单例子是：在适当的环境层范畴中，取 $E,F$ 为局部自由层 $\mathcal{E},\mathcal{F}$ 的内射消解。这时，空间 $\operatorname{Ext}^n(E,F)$ 是超 Ext 群，与\BookSectionRef{sec:4.5}{4.5}讨论的 $\operatorname{Ext}^n(\mathcal{E},\mathcal{F})$ 一致。因此，得到的开弦，与按照\BookSectionRef{sec:3.6.4}{3.6.4}在 $\mathcal{E},\mathcal{F}$ 之间找到的开弦相同。
\end{example*}

\subsubsection{其他形变}\label{sec:5.3.2.2}

\SourceParagraph{ch05-5.3.2.2-001}{327}
下面是一类非常广泛的形变。假定有两个由复形 $E,F$ 给出的 D 膜。若打开微分以后，各项的鬼数不变，那么，与 $f\in\operatorname{Hom}(E[-1],F)$ 对应的开弦，其鬼数为1。因此，可以考虑由 $f$ 给出的形变。容易看到，这产生\BookSectionRef{sec:4.4.3}{4.4.3}定义的映射锥 $\operatorname{Cone}(f\colon E[-1]\to F)$，或简称 $\operatorname{Cone}(f)$。所以，它已经包含在原有构造之中。

\SourceParagraph{ch05-5.3.2.2-002}{327}
锥构造涵盖了几乎所有其他可能考虑的形变。例如，复形本身可以看作迭代的锥。把对象看作只有一项的复形，并计入使各项处于正确次数的移位，依次有
\begin{equation}\label{eq:5.53}
\begin{aligned}
K_1&=\operatorname{Cone}(d_0\colon E^0\to E^1)[-1],\\
K_2&=\operatorname{Cone}(d_1\colon K_1\to E^2[-1])[-1],\\
K_3&=\operatorname{Cone}(d_2\colon K_2\to E^3[-2])[-1],\qquad\ldots.
\end{aligned}
\end{equation}
这样继续，就逐项构造出 $E$。

\subsubsection{关联函数与 \texorpdfstring{$A_\infty$}{A∞} 代数}\label{sec:5.3.2.3}

\SourceParagraph{ch05-5.3.2.3-001}{327}
为了完成包含复形的拓扑开弦理论构造，还必须定义开弦算符的关联函数。要定义\BookEqRef{eq:3.176}{3.176}\allowbreak 与\BookEqRef{eq:3.193}{3.193}\allowbreak 中的三点函数，需要能够相乘态射，并取迹。

\SourceParagraph{ch05-5.3.2.3-002}{327}
若把 B 膜取为凝聚层复形，就可以使用\BookChapterRef{4}给出的复形范畴的态射乘法。特别地，在导出范畴中，迹的定义来自 Serre 对偶，即\BookRef{thm:4.68}{定理4.68}，以及 $X$ 为 Calabi--Yau 流形时，对偶化层 $\omega_X$ 平凡这一事实。取完美配对中的第一个态射为恒等态射，就得到迹
\[
\operatorname{Tr}\colon\operatorname{Ext}^n(E,E)\longrightarrow\mathbb{C},
\]
其中，$n=\dim X$。这个迹在相差整体归一化的意义下唯一。

\SourceParagraph{ch05-5.3.2.3-003}{328}
Serre 对偶有一个简单推论，后面将明确使用。
\BookStatementNumber{11}
\begin{corollary}\label{cor:5.11}
若态射 $\alpha\in\operatorname{Ext}^*(E,F)$ 在所有关联函数中都为零，则它本身就是零态射。
\end{corollary}

\SourceParagraph{ch05-5.3.2.3-004}{328}
为了定义高点关联函数，最好发展进一步的结构。特别地，不应直接进入 BRST 上同调，而应把开弦看作微分分次代数的元素。这样，可以直接说明，高点关联函数与对应 $A_\infty$ 代数中的高阶乘积有关。这也更便于把讨论推广到凝聚层以外的情形，因为论证更直接地基于 SCFT 及其拓扑扭曲。关于这些问题，见文献\BookCite{ref022}。

\subsection{导出范畴}\label{sec:5.3.3}

\SourceParagraph{ch05-5.3.3-intro-001}{328}
已经看到，拓扑开弦理论包含同伦范畴 $K(\mathcal{A})$ 的所有对象与态射。沿着\BookSectionRef{sec:4.4.2}{4.4.2}的思路，如果接下来能给出物理论证，说明 $K(\mathcal{A})$ 中拟同构的对象，作为边界条件无法区分，那么，就说明了 B 膜范畴是导出范畴 $D(\mathcal{A})$；此处采用 Abel 范畴的导出构造。

\SourceParagraph{ch05-5.3.3-intro-002}{328}
导出范畴的对象，会比 $K(\mathcal{A})$ 少得多。另一方面，它又远比朴素地对 $\mathcal{A}$ 的对象作分次直和丰富。为说明这一点，假定试图用上同调本身构成的序列替换复形 $E$，也就是下面这个映射全为零的复形：
\begin{equation}\label{eq:5.54}
\cdots\xrightarrow{0}H^0(E)\xrightarrow{0}H^1(E)
\xrightarrow{0}H^2(E)\xrightarrow{0}\cdots.
\end{equation}
一般而言，在 $E$ 与\BookEqRef{eq:5.54}{5.54}之间，无论哪一个方向，都未必存在诱导上同调同构的态射；因此，复形未必拟同构于其上同调。这一重要事实，带来了导出范畴的复杂结构。

\subsubsection{导出范畴与 K 理论}\label{sec:5.3.3.1}

\SourceParagraph{ch05-5.3.3.1-001}{328}
把这里的讨论与\BookSectionRef{sec:5.1}{5.1}比较。\BookSectionRef{sec:5.3.2.2}{5.3.2.2}已经看到，可以把两个 B 膜 $E[-1],F$，形变为由态射 $f\colon E[-1]\to F$ 的锥表示的单个 B 膜。因此，B 膜 $\operatorname{Cone}(f)$ 由 $E[-1]$ 与 $F$ 组成，其中，$E[-1]$ 是 $E$ 的反膜。

\SourceParagraph{ch05-5.3.3.1-002}{328}
因此，可以采用\BookSectionRef{sec:5.1}{5.1}的定义，把 $D^b(X)$ 中的对象看作粒子，把所有形变
\begin{equation}\label{eq:5.55}
\begin{aligned}
[\operatorname{Cone}(f)]&\longleftrightarrow[F]+[E[-1]]\\
&\longleftrightarrow[F]-[E]
\end{aligned}
\end{equation}
看作束缚或衰变过程。所得群是 $X$ 的 Grothendieck 群；更多细节见文献\BookCite{ref295}第77页。它在 D 膜中的应用，可以追溯到文献\BookCite{ref225,ref425}。把全纯丛忘记为拓扑丛以后，它自然映到前面讨论的拓扑 K 群 $K(X)$。

\SourceParagraph{ch05-5.3.3.1-003}{329}
可以如下自然地把导出范畴映到 K 理论。使用局部自由消解，把任意复形替换为与之拟同构的局部自由层复形 $E$。随后，构造 K 理论对象
\begin{equation}\label{eq:5.56}
\cdots\ominus E^{-1}\oplus E^0\ominus E^1\oplus E^2\ominus\cdots,
\end{equation}
其中，$E^i$ 是与相应局部自由层对应的全纯向量丛。可以说明，这给出良好定义的映射 $D(X)\to K(X)$。注意，这个映射未必满。拓扑 $K(X)$ 由所有向量丛生成，而只从 B 膜出发，得到的是全纯向量丛的 K 理论。为了生成所有可能的类，完整的 K 理论，可能需要对应于非 BPS 膜的非全纯丛\BookCite{ref471}。

\SourceParagraph{ch05-5.3.3.1-004}{329}
现在可以把这与\BookChapterRef{2}讨论的半单情形比较。那里出现的范畴是 Abel 范畴，对象对应于复向量空间，态射对应于线性映射。在这种情形中，复形与其上同调之间，总存在拟同构，因而导出范畴具有简单形式。既然每个对象的同构类都由其上同调确定，它的全部内容都包含在\BookSectionRef{sec:5.1}{5.1}的 K 理论中。

\SourceParagraph{ch05-5.3.3.1-005}{329}
不过，在同调维数非零的范畴中，一般不再有这种简化。例如，尽管 $X$ 上的点状0膜具有相同的拓扑 K 理论类，$X$ 中不同的两个点所对应的0膜，却对应于 $D^b(X)$ 中不同构的对象。

\subsubsection{物理同一性与拟同构}\label{sec:5.3.3.2}

\SourceParagraph{ch05-5.3.3.2-001}{329}
从物理观点看，为什么要按拟同构作商？虽然可以根据从这一假设导出的结果彼此一致来论证，但最好找到基于物理公理与可观测量的演绎论证。

\SourceParagraph{ch05-5.3.3.2-002}{329}
\BookRef{def:4.45}{定义4.45}把拟同构定义为诱导上同调同构的态射，这不是一个完全令人满意的出发点，因为复形的上同调并非物理可观测量。

\SourceParagraph{ch05-5.3.3.2-003}{329}
TFT 中，物理可观测量是开弦谱及其关联函数。因此，需要定义一对膜“物理等价”的概念，也就是，两者在所有关联函数中表现相同。由物理动机得到的膜范畴便是：
\BookStatementNumber{12}
\begin{definition}\label{def:5.12}
范畴 $T(\mathcal{C})$ 是 $K(\mathcal{C})$ 按物理等价得到的商，包含对相应态射空间的识别。
\end{definition}
其中，
\BookStatementNumber{13}
\begin{definition}\label{def:5.13}
两个对象 $E,E'\in K(\mathcal{C})$ \emph{物理等价}，当且仅当，对每个 $p$ 及每个 $F\in K(\mathcal{C})$，有
\begin{equation}\label{eq:5.57}
\begin{aligned}
\operatorname{Ext}^p(E,F)&\cong\operatorname{Ext}^p(E',F),\\
\operatorname{Ext}^p(F,E)&\cong\operatorname{Ext}^p(F,E'),
\end{aligned}
\end{equation}
并且，这些识别保持\BookEqRef{eq:3.193}{3.193}中的三点关联函数。
\end{definition}

\SourceParagraph{ch05-5.3.3.2-004}{330}
\begin{example*}
例如，采用\BookEqRef{eq:3.192}{3.192}的记号，取 $E=E_2$、$E'=E'_2$，以及
\begin{equation}\label{eq:5.58}
\begin{aligned}
a'&\in H_{\bar\partial}^{0,p}(X,\operatorname{Hom}(E_1,E'_2)),\\
b'&\in H_{\bar\partial}^{0,q}(X,\operatorname{Hom}(E'_2,E_3)),\\
c&\in H_{\bar\partial}^{0,r}(X,\operatorname{Hom}(E_3,E_1)),
\end{aligned}
\end{equation}
则有
\[
\langle W_aW_bW_c\rangle
=\langle W_{a'}W_{b'}W_c\rangle,\qquad\forall E_1,E_3.
\]
\end{example*}
由 OPE，等价地，由结合律，对三点函数的条件，迫使所有由这些乘积构成的 $n$ 点函数相等，从而使相应物理可观测量相同。这显然是等价关系。

\SourceParagraph{ch05-5.3.3.2-005}{330}
主要结果便是：
\BookStatementNumber{14}
\begin{theorem}[Aspinwall--Lawrence\BookCite{ref023}]\label{thm:5.14}
在上述拓扑开弦实现及非退化关联函数配对的条件下，
\[
T(\mathcal{C})\simeq D(\mathcal{C}).
\]
\end{theorem}

\SourceParagraph{ch05-5.3.3.2-006}{330}
为证明这一点，回顾\BookSectionRef{sec:4.4}{4.4}的讨论：导出范畴 $D(\mathcal{C})$，是使 $C(\mathcal{C})$ 中的拟同构映到同构的泛范畴。因此，具有这一性质的函子 $F\colon C(\mathcal{C})\to T(\mathcal{C})$，会通过 $D(\mathcal{C})$ 唯一地分解为
\[
F=G\circ Q,
\]
其中，$Q$ 是\BookRef{lem:4.46}{引理4.46}中的局部化函子。因此，需要说明，从 $K$ 到 $T$ 的自然函子，把拟同构映到同构，并且，$G$ 是范畴等价。这里的物理开弦上同调，按照前面的消解模型计算；以下 $\operatorname{Hom}^*$ 取相应的导出映射空间。

\SourceParagraph{ch05-5.3.3.2-007}{330}
\BookStatementNumber{15}
\begin{proposition}\label{prop:5.15}
从 $K$ 到 $T$ 的自然函子，把拟同构映到同构。
\end{proposition}
\begin{proof}
给定拟同构 $f\colon E\to E'$，按照\BookSectionRef{sec:4.4.3}{4.4.3}构造其映射锥 $C(f)$。它是无环复形，即上同调为零。因此，在相应消解模型中，对任意 $F$，有
\[
\operatorname{Hom}^*(C(f),F)
\cong\operatorname{Hom}^*(F,C(f))\cong0,
\]
从而得到\BookEqRef{eq:5.57}{5.57}。
\end{proof}

\SourceParagraph{ch05-5.3.3.2-008}{330}
\BookStatementNumber{16}
\begin{proposition}\label{prop:5.16}
$G$ 是范畴等价。
\end{proposition}
\begin{proof}
只有当 $D(\mathcal{C})$ 中两个不等价的对象 $E,E'$ 物理等价时，这才可能失败；换句话说，它们与所有其他对象之间，都具有相同的态射和关联函数。不过，条件\BookEqRef{eq:5.57}{5.57}蕴含态射 $\alpha\in\operatorname{Hom}(E,E')$ 与 $\beta\in\operatorname{Hom}(E',E)$ 的存在，其乘积 $\beta\alpha$，在关联函数中的表现，与恒等态射 $\operatorname{id}_E\in\operatorname{Hom}(E,E)$ 完全相同。因此，态射 $\beta\alpha-\operatorname{id}_E$ 在所有关联函数中都为零；由\BookRef{cor:5.11}{推论5.11}，它本身就是零。交换 $E,E'$，同样得到 $\alpha\beta=\operatorname{id}_{E'}$。于是，$E,E'$ 实际同构，与假设矛盾。
\end{proof}

\SourceParagraph{ch05-5.3.3.2-009}{330}
因此，如果能构造形成 Abel 范畴 $\mathcal{A}$ 的 B 膜，这些论证就定义了对应于 $D(\mathcal{A})$ 的拓扑开弦理论。有些情形中，B 膜显然形成 Abel 范畴，例如，\BookSectionRef{sec:3.6.8}{3.6.8}讨论的 Landau--Ginzburg 模型。

\SourceParagraph{ch05-5.3.3.2-010}{331}
\BookSectionRef{sec:5.4}{5.4}还将说明，Dirichlet 膜的世界体构造，直接把 B 膜描述为箭图表示范畴中的对象；这个范畴是 Abel 范畴。例如，\BookSectionRef{sec:5.4.5}{5.4.5}的轨形构造，就属于这一类。

\subsubsection{消解与凝聚层}\label{sec:5.3.3.3}

\SourceParagraph{ch05-5.3.3.3-001}{331}
现在，把注意力集中在\BookSectionRef{sec:3.6.3}{3.6.3}中，通过扭曲 $X$ 上的 $\sigma$ 模型定义的 B 模型。虽然拥有大量 B 型边界条件，即 $X$ 上所有全纯丛，它们却不构成 Abel 范畴，因为映射的余核未必是丛。这会给刚才的论证带来问题吗？

\SourceParagraph{ch05-5.3.3.3-002}{331}
当然，还有其他 B 型膜，例如 D0膜，以及支撑在全纯子流形上的膜。把它们包括进来，就提示，$\sigma$ 模型的正确出发点，应当是取 $\mathcal{A}=\operatorname{Coh}X$，即凝聚层范畴。虽然它是 Abel 范畴，解决了这个问题，但此时也许会担心：这一步没有明确的物理动机，因而物理可能作出不同的选择。怎么知道，凝聚层就是正确的推广呢？

\SourceParagraph{ch05-5.3.3.3-003}{331}
直接处理这一问题的方法，是把所有 B 膜都定义为 $\sigma$ 模型中的显式边界条件，并检查它们具有与凝聚层相同的性质。下面将讨论一些这类结果，不过，目前尚未完成这一类型的完全一般的分析。

\SourceParagraph{ch05-5.3.3.3-004}{331}
尽管如此，仍有很好的论证表明，这必定可行，且所有 B 型边界条件，都是 $D^b(X)$ 中的对象：在已经说明如何构造的边界条件，即局部自由层复形中，已经有完全合适的候选对象，来表示 D0膜，以及所有其他支撑在子流形上的 B 型膜。这些复形被安排成，使子流形 $Z\subset X$ 以外的全部结构，都从 $Q$ 上同调中抵消。

\SourceParagraph{ch05-5.3.3.3-005}{331}
\begin{example*}
例如，考虑复形 $E_Z$，
\[
\mathcal{I}_Z\hookrightarrow\mathcal{O}_X,
\]
其中，与\BookRef{ex:4.42}{例题4.42}一样，$Z\subset X$ 是局部由单个方程 $f=0$ 定义的有效 Cartier 除子；在局部平凡化中，这个映射就是乘以 $f$。这时，$\mathcal{I}_Z$ 是局部自由层。
\end{example*}

\SourceParagraph{ch05-5.3.3.3-006}{331}
现在考察 $E_Z$ 与 $p\in X$ 处的 D0膜之间的态射。数学上，可以把 D0膜定义为点的结构层 $\mathcal{O}_p$；物理上，可以把它定义为纯 Dirichlet 边界条件。来比较二者。数学上，一个简单练习就能验证：除非 $p\in Z$，这些态射都从 $Q$ 上同调中抵消。另一方面，\BookSectionRef{sec:3.6.5}{3.6.5}已经论证，物理定义也给出相同结果。

\SourceParagraph{ch05-5.3.3.3-007}{332}
回看\BookRef{ex:4.42}{例题4.42}，可知，$E_Z$ 是结构层 $\mathcal{O}_Z$ 的消解；$\mathcal{O}_Z$ 是 $X$ 上的凝聚层，也是用来识别缠绕子流形 $Z$ 的物理 D 膜的自然数学对象。由于凝聚层与其消解拟同构，构造拓扑膜范畴时，无须显式加入凝聚层 $\mathcal{O}_Z$；它已经在那里。这个论证可以推广到支撑在任意子簇 $Z$ 上的任意凝聚层，说明对应的拓扑膜已经存在，并应当对应于某个缠绕 $Z$ 的物理膜。

\SourceParagraph{ch05-5.3.3.3-008}{332}
与 $\mathcal{O}_Z$ 对应的物理 D 膜是什么？起初也许会说，自然的候选是缠绕 $Z$、携带平凡丛的 D4膜。不过，这并非显然正确；为了作出这一识别，可能需要知道更多关于 $Z$ 的几何信息。正如\BookSectionRef{sec:5.1.4}{5.1.4}引述的，另一些基于闭弦理论关系的物理论证，给出 K 理论层面的识别\BookEqRef{eq:5.12}{5.12}，其中出现 $Z$ 的法丛与切丛。对于当前目的，这一关系的精确形式不太重要；重要的是，$\mathcal{O}_Z$ 可以识别为缠绕 $Z$、携带某个丛的 D4膜，这个丛可以从\BookEqRef{eq:5.12}{5.12}读出。

\SourceParagraph{ch05-5.3.3.3-009}{332}
不过，为了论证，假定事情并非如此。也就是说，假定从 D4膜更直接的物理定义，即混合 Dirichlet--Neumann 边界条件出发，所得结果与把它识别为凝聚层所预言的结果不同，例如，某些态射不同。此时，既然刚才已经论证，层确实存在于拓扑开弦理论中，那么，就会有两个不同对象，对应于缠绕 $Z$ 的膜：凝聚层与物理 D4膜。无论从物理观点，还是从数学观点看，这都会非常意外。物理上，在已分析的许多情形中，并未看到这种现象的证据；数学上，也没有已知的候选范畴，既自然地与 $X$ 对应，又真包含 $D^b(X)$。

\SourceParagraph{ch05-5.3.3.3-010}{332}
当然，这个论证还算不上证明。不过，目前还无法得到证明，因为其中一部分，必须说明已经考虑了所有可能的物理自洽边界条件；对于 $\hat c>1$ 的理论，物理中尚做不到这一点。

\subsubsection{Dirichlet 膜与凝聚层的比较}\label{sec:5.3.3.4}

\SourceParagraph{ch05-5.3.3.4-001}{332}
概括上一小节的结论。考虑 $D^b(X)$ 的满子范畴，其对象为局部自由层的复形。把这个子范畴称为 $D_0^b(X)$。在前面采用的光滑情形中，每个凝聚层都有有界局部自由消解，所以，它等价于 $D^b(X)$。另一方面，有一个从 $D_0^b(X)$ 到 B 膜范畴的函子 $F$，把复形赋予对应的 B 膜，并识别相应的态射空间。根据定义，这个函子是满忠实的。

\SourceParagraph{ch05-5.3.3.4-002}{333}
\BookSectionRef{sec:5.3.3.3}{5.3.3.3}的主张，是它也是本质满的，即，任意 B 膜，作为 B 膜范畴的对象，都同构于 $F$ 的像中的某个对象。若如此，范畴 $D_0^b(X)$，因而 $D^b(X)$，就等价于 B 膜范畴。

\SourceParagraph{ch05-5.3.3.4-003}{333}
有些情形中，已经通过计算特定 B 膜之间的态射空间，并与 $D^b(X)$ 中对应对象之间的态射空间比较，检验了这一点。\BookSectionRef{sec:3.6.5}{3.6.5}讨论了 D0膜，结论是，它对应于支撑在 $p$ 的摩天层 $\mathcal{O}_p$。摩天层在 $D^b(X)$ 中的分次自同态代数，是 $T_pX$ 的外代数，与\BookEqRef{eq:3.194}{3.194}所得结果一致。

\SourceParagraph{ch05-5.3.3.4-004}{333}
一个更简单的检验，是计算从对应于全纯向量丛 $E$ 的 B 膜，到 $p$ 处 D0膜的态射空间。在这种情形中，边界条件迫使所有费米场为零，边界算符可以识别为向量空间 $E_p^*$ 的元素。这又与导出范畴中对应对象之间的态射一致。

\SourceParagraph{ch05-5.3.3.4-005}{333}
可以很容易地把这推广到对应于 $X$ 中复子流形的 B 膜。设 $Y$ 是这样的子流形。对应于这种 B 膜的边界条件，迫使玻色场取值于 $Y$，并使 $\eta^{\bar i}$ 与 $\theta_i$，分别取值于 $TY^{0,1}\subset TX^{0,1}|_Y$ 以及 $NY^*\subset T^*X^{1,0}|_Y$。因此，候选 BRST 不变边界算符，是
\[
\bigoplus_{p,q}\Lambda^p NY\otimes\Omega^{0,q}(Y)
\]
的截面。在这一局部模型中，BRST 算符作为 Dolbeault 算符 $\bar\partial$ 作用，因此，自同态的初步计算给出全纯向量丛 $\Lambda NY$ 的 Dolbeault 上同调。若把这个 B 膜对应于层 $i_*\mathcal{O}_Y$，其中，$i\colon Y\hookrightarrow X$ 是嵌入，那么，这些群也出现在计算 $D^b(X)$ 中态射的局部到整体 Ext 谱序列中；谱序列退化时，它们直接给出全部态射，一般情形还须计入其微分。

\SourceParagraph{ch05-5.3.3.4-006}{333}
计算对应于全纯向量丛 $E$ 的 B 膜，与对应于复子流形 $Y$ 的 B 膜之间的态射时，会遇到一点意外\BookCite{ref298}。导出范畴预期的结果，是全纯向量丛 $E^*|_Y$ 的 Dolbeault 上同调；物理计算给出的，却是 $E^*|_Y\otimes K_Y^{1/2}$ 的 Dolbeault 上同调。这提示，与 $Y$ 对应的 B 膜，应当识别为 $i_*K_Y^{1/2}$，而非 $i_*\mathcal{O}_Y$；这里，$K_Y$ 是 $Y$ 的典范线丛。注意，$K_Y$ 的平方根，只有在 $Y$ 为自旋流形时才存在，而且，若 $Y$ 非单连通，它未必唯一。这似乎违背了把 B 膜识别为导出范畴对象的愿望，因为任意复子流形 $Y$，即使不是自旋流形，其结构层也给出 $D^b(X)$ 的合法对象。文献\BookCite{ref156,ref298}指出了解决这一悖论的方法，其中涉及膜上允许有非平凡规范场的可能性。

\SourceParagraph{ch05-5.3.3.4-007}{334}
为理解这一机制，先考虑 $Y$ 为自旋流形、但携带向量丛 $E'$ 的情形。那么，必须给它赋予凝聚层 $i_*(\mathcal{E}'\otimes K_Y^{1/2})$，其中，$\mathcal{E}'$ 是 $Y$ 上与向量丛 $E'$ 对应的局部自由层。若 $Y$ 不是自旋流形，根据文献\BookCite{ref156,ref298}，$\mathcal{E}'$ 与 $K_Y^{1/2}$ 都是扭曲凝聚层，而 $\mathcal{E}'\otimes K_Y^{1/2}$ 是 $Y$ 上的普通凝聚层。与这种 B 膜对应的 $D^b(X)$ 中的对象，是 $i_*(\mathcal{E}'\otimes K_Y^{1/2})$。

\SourceParagraph{ch05-5.3.3.4-008}{334}
文献\BookCite{ref298,ref289}中的进一步计算，证实了 B 膜之间的态射，与 $D^b(X)$ 中对应对象之间的态射一致。虽然这些计算合起来仍不构成证明，但它们强烈表明，B 膜范畴等价于 $D^b(X)$。

\subsubsection{镜像对称恢复了吗？}\label{sec:5.3.3.5}

\SourceParagraph{ch05-5.3.3.5-001}{334}
若 $Y$ 上的 A 模型，与 $X$ 上的 B 模型“相同”，那么，似乎已经为如下主张提供了动机：$Y$ 上的 Fukaya 范畴，等价于 $X$ 上的导出范畴 $D^b(X)$。这正是 Kontsevich 最初的主张。事实上，仍有一点美中不足，\BookSectionRef{sec:5.5.2}{5.5.2}将讨论，不过，这个主张似乎已经十分接近事实。最后，指出这一镜像对称运作良好的几个其他特征。

\SourceParagraph{ch05-5.3.3.5-002}{334}
\BookSectionRef{sec:3.6.2.5}{3.6.2.5}中的蝌蚪项消除条件，带来 A 膜模空间的两个有趣方面：
\begin{enumerate}[label=(\arabic*)]
\item 对应于 $H^1(L)$ 的 Lagrange 子流形的一阶形变，可能受阻，从而不产生真正的 A 膜形变。
\item 有些 A 膜，可能依赖于 $B+\mathrm{i}J$ 的非常特殊的取值，并在一般的 $B+\mathrm{i}J$ 下完全消失。
\end{enumerate}
B 模型中的镜像陈述，都成立：
\begin{enumerate}[label=(\arabic*)]
\item 凝聚层的一阶形变，对应于 $\operatorname{Ext}^1(E,E)$，也可能受阻。例子见文献\BookCite{ref444}。
\item 有些层只在复结构的特殊取值下存在。一个例子是，4膜缠绕 $X$ 中一个以正亏格代数曲线为底的直纹曲面\BookCite{ref462}。
\end{enumerate}
文献\BookCite{ref276,ref277}也分析了这一问题。注意，这是镜像对称的一个典型例子：A 模型中的瞬子效应，即蝌蚪项，映为 B 模型中可以用经典几何理解的效应。

\subsection{分次与边界条件}\label{sec:5.3.4}

\SourceParagraph{ch05-5.3.4-intro-001}{334}
现在更详细地解释，用于定义复形的分次 $\mu$ 的物理来源，以及它与\BookSectionRef{sec:5.2.1}{5.2.1}中分次 $\xi$ 的关系。一开始应当说明，这个问题涉及许多物理内容，包括超对称、BPS 界与玻色化。为了使叙述在合理程度上自成一体，把注意力集中在当前目的所必需的部分。

\SourceParagraph{ch05-5.3.4-intro-002}{335}
显然，正确的定义，应当同时在 A 膜与 B 膜的语言中有意义。事实上，由于它必须镜像对称，应当在产生镜像对称的底层 SCFT 结构中寻找。

\SourceParagraph{ch05-5.3.4-intro-003}{335}
回顾\BookSectionRef{sec:3.2.4}{3.2.4}与\BookSectionRef{sec:3.4.3.2}{3.4.3.2}，镜像对称在 SCFT 中的来源，是 $N=2$ SCA 的 $U(1)$ 子代数的因子分解。特别地，回顾\BookSectionRef{sec:3.3.4.6}{3.3.4.6}，$\mathbb{Z}$ 分次完全在 $U(1)$ 子代数内部定义。因此，因子分解告诉我们，这些分次的理论具有普遍性，可以借助任意 $U(1)$ CFT 的理论来理解，包括\BookChapterRef{3}中详细讨论的、靶空间为 $S^1$ 的简单自由玻色子理论。

\SourceParagraph{ch05-5.3.4-intro-004}{335}
我们知道这个理论中边界条件的完整集合：Dirichlet 与 Neumann 情形，每一种都由单个实参数描述。把\BookEqRef{eq:3.165}{3.165}、\BookEqRef{eq:3.166}{3.166}与\BookEqRef{eq:3.139}{3.139}比较，可知，A 型膜对应于 Dirichlet 情形，而 B 型膜对应于 Neumann 情形。实际上，由于镜像对称的作用\BookEqRef{eq:3.134}{3.134}在这一描述中非常简单，两种情形的讨论完全平行。

\SourceParagraph{ch05-5.3.4-intro-005}{335}
与\BookSectionRef{sec:5.3.1}{5.3.1}所提示的一样，无论哪种情形，分次 $\mu$ 都是边界条件中的实参数，至多相差一个乘法因子。在 A 膜，即 Dirichlet 情形中，可以把这个参数想象为，D0膜在自由玻色子靶空间 $S^1$ 上的位置。Neumann 情形中，可以把它理解为 $U(1)$ 联络绕环路的平行移动，也可以理解为 T 对偶 $S^1$ 上的位置，后一种也许更容易想象。

\SourceParagraph{ch05-5.3.4-intro-006}{335}
这一主张最基本的检验，是开弦的 R 荷，确实由\BookEqRef{eq:3.189}{3.189}与\BookEqRef{eq:5.41}{5.41}给出，其中，$\mu$ 是这里定义的参数。既然 R 荷守恒，由此得到的结构就是分次范畴。

\SourceParagraph{ch05-5.3.4-intro-007}{335}
这一点的论证，在 Dirichlet 图景中稍微简单一些，所以，就在这一图景中说明。使用\BookSectionRef{sec:3.2.3.5}{3.2.3.5}中 $U(1)$ 流代数的玻色化描述，如\BookSectionRef{sec:3.3.4.6}{3.3.4.6}所用的那样。根据\BookEqRef{eq:3.167}{3.167}与\BookEqRef{eq:3.102}{3.102}，采用适于边界参数的玻色子归一化，从膜 $E$ 到膜 $F$ 的开弦，其 R 荷，是算符
\[
J_0=\sqrt{\hat c}\int_0^\pi\partial_\sigma\varphi\,\mathrm{d}\sigma
\]
的期望值。积分横跨开弦，边界条件 $\varphi(0)$ 与 $\varphi(\pi)$，分别由 $E,F$ 边界条件的实参数确定。

\SourceParagraph{ch05-5.3.4-intro-008}{335}
由于这一表达式是局部密度的积分，直观上很清楚：当两条弦在端点接合时，它们的 R 荷相加。此外，由于密度是全导数，可以直接积分。为此，与\BookEqRef{eq:3.54}{3.54}一样，把 $\varphi$ 分解为取值于 $\mathbb{R}$ 的函数，以及量子化的“缠绕”部分 $p=m$（这里，$p$ 的单位由所选坐标的归一化确定）。于是，
\begin{equation}\label{eq:5.59}
J_0=\sqrt{\hat c}\bigl(p+\varphi(\pi)-\varphi(0)\bigr).
\end{equation}

\SourceParagraph{ch05-5.3.4-intro-009}{336}
这具有\BookEqRef{eq:5.41}{5.41}的形式，其中，$q=\sqrt{\hat c}\,m$，而对于 $B=E,F$，
\begin{equation}\label{eq:5.60}
\mu(B)=\sqrt{\hat c}\,\varphi(B).
\end{equation}
换句话说，对于膜 $B$，施加在 $\varphi$ 上的具体 Dirichlet 边界条件，就是它的分次，至多相差一个整体约定因子。

\subsubsection{分次的周期性}\label{sec:5.3.4.1}

\SourceParagraph{ch05-5.3.4.1-001}{336}
先从 $E=F$、$\varphi(0)=\varphi(\pi)$ 的情形开始。这时，重新得到\BookEqRef{eq:3.103}{3.103}中的 $q=\sqrt{\hat c}\,p$。因此，为了与\BookSectionRef{sec:3.3.2}{3.3.2}的 R 荷约定一致，特别是 $Q$ 的荷为1，$p$ 必须以 $1/\sqrt{\hat c}$ 为单位量子化。根据自由玻色子的讨论，这对应于作用量\BookEqRef{eq:3.59}{3.59}中，在所用半径单位下的 $R=1/\sqrt{\hat c}$。

\SourceParagraph{ch05-5.3.4.1-002}{336}
将这一点与\BookEqRef{eq:5.59}{5.59}结合，参数 $\mu(B)=\sqrt{\hat c}\,\varphi(B)$ 的有效周期为 $2\hat c$，也就是说，这样的移位是 SCFT 的对称。换句话说，
\begin{equation}\label{eq:5.61}
\mu(B)\in\mathbb{R}/2\hat c\mathbb{Z}.
\end{equation}
$E\ne F$ 时也一样。它唯一带来的变化，是在\BookEqRef{eq:5.59}{5.59}中产生一个常数的整体移位，对应于\BookEqRef{eq:5.41}{5.41}中的分次之差。

\SourceParagraph{ch05-5.3.4.1-003}{336}
细心的读者，此时也许会有些困惑：前面论证过，镜像对称包含对这个玻色子作 T 对偶；根据\BookSectionRef{sec:3.2.3.6}{3.2.3.6}，其作用为 $R\mapsto1/R$。除去 $R=1$，镜像对称怎么会与其他半径相容？

\SourceParagraph{ch05-5.3.4.1-004}{336}
答案是，$U(1)$ 扇区，并不独立于 $(2,2)$ SCFT 的其余部分。相反，因子分解意味着，$U(1)$ 荷与理论其余部分的算符谱之间存在关联，如\BookEqRef{eq:3.70}{3.70}所表达的。开弦扇区中，只有一个守恒的 $U(1)$ 荷，因此，其形式为
\begin{equation}\label{eq:5.62}
\mathcal{H}_{E,F}=\bigoplus_p\mathcal{H}^{X'}_{E,F,p}\otimes V_p.
\end{equation}
作 T 对偶变换时，必须计入这些关联。虽然\BookSectionRef{sec:3.2.3.7}{3.2.3.7}中的算符关系仍然有效，这却改变了\BookEqRef{eq:3.66}{3.66}以及后续分析步骤。

\SourceParagraph{ch05-5.3.4.1-005}{336}
事实上，由于\BookEqRef{eq:5.62}{5.62}中的关联，把 $\varphi$ 简单地看作取值于半径 $R=1/\sqrt{\hat c}$ 的 $S^1$，并不完全正确。真正起作用的是有效周期参数 $R_{\mathrm{eff}}$；对于它，有同构
\[
\mathcal{H}^{X'}_{E,F,p}\cong\mathcal{H}^{X'}_{E,F,p+R_{\mathrm{eff}}},
\]
因为只有这样的 $\varphi\mapsto\varphi+R_{\mathrm{eff}}$，才给出整个理论的对称。由\BookSectionRef{sec:3.3.5}{3.3.5}末尾的讨论，能被保证为对称的最小移位，是 $\Delta q=2\hat c$ 的谱流，因此断定 $R_{\mathrm{eff}}=2\sqrt{\hat c}$。

\subsubsection{分次的规范对称}\label{sec:5.3.4.2}

\SourceParagraph{ch05-5.3.4.2-001}{337}
虽然详细讨论了周期性 $\mu\mapsto\mu+2\hat c$，人们却认为，它并不那么重要，可以把分次 $\mu$ 看作实值量。原因是，整个 SCFT 具有 $U(1)_R$ 对称，由\BookSectionRef{sec:3.6}{3.6}的边界条件所保持的流 $J$ 生成。它把玻色子 $\varphi$、从而把所有分次作整体移位，例如，
\[
\mu(B)\longmapsto\mu(B)+n,\qquad n\in\mathbb{R}.
\]
与此一致，整体移位基本不改变数学框架，例如，
\begin{equation}\label{eq:5.63}
\operatorname{Hom}(E,F)=\operatorname{Hom}(E[n],F[n]),
\end{equation}
所有乘积也都不受影响。

\SourceParagraph{ch05-5.3.4.2-002}{337}
物理上，称复形的这一整体移位，为理论的规范对称。一旦设定某个膜的分次，规范对称就被“固定”，不再有歧义。

\SourceParagraph{ch05-5.3.4.2-003}{337}
\BookSectionRef{sec:3.3.5}{3.3.5}讨论过，非零开弦关联函数的总 $U(1)_R$ 荷，必须为 $\hat c$。回顾\BookEqRef{eq:3.97}{3.97}与\BookEqRef{eq:3.98}{3.98}，相应手征态的荷满足 $0\le q\le\hat c$，所以，这些关联函数实际上无法探测 $2\hat c$ 的周期性，也尚不知它产生其他什么效应\footnote{文献\BookCite{ref123}中的主张不正确，见文献\BookCite{ref018}。}。因此，从现在起，一般忽略这一周期性；若能证明这样做确实有效，将是理想的。

\subsection{分次与中心荷}\label{sec:5.3.5}

\SourceParagraph{ch05-5.3.5-intro-001}{337}
现在，一个边界条件对应于两个数。一方面，有\BookEqRef{eq:5.60}{5.60}中的参数 $\mu$，即边界条件的 $U(1)_R$ 部分。另一方面，\BookSectionRef{sec:3.6.6}{3.6.6}的讨论，另见\BookSectionRef{sec:5.2.1.1}{5.2.1.1}，涉及依赖于参数 $\xi$ 的匹配条件\BookEqRef{eq:3.196}{3.196}与\BookEqRef{eq:3.197}{3.197}。这个参数在特殊 Lagrange 条件中，作为\BookEqRef{eq:5.19}{5.19}出现；在 SCFT 语言中，A、B 两种情形都可以概括为
\begin{equation}\label{eq:5.64}
\Omega=\exp(-2\mathrm{i}\pi\xi)\overline\Omega.
\end{equation}

\SourceParagraph{ch05-5.3.5-intro-002}{337}
事实上，在相同的归一化约定下，这两个参数相同，
\begin{equation}\label{eq:5.65}
\xi=\mu.
\end{equation}

\SourceParagraph{ch05-5.3.5-intro-003}{337}
这有几个重要后果。首先，虽然 $\xi$ 的定义\BookEqRef{eq:5.16}{5.16}提示，它只模2定义，但根据\BookEqRef{eq:5.61}{5.61}，实际上，必须模 $2\hat c$ 定义。

\SourceParagraph{ch05-5.3.5-intro-004}{337}
形式上，可以很容易地通过连续性做到这一点。回顾\BookSectionRef{sec:5.2.1.2}{5.2.1.2}：考虑复结构的变化时，特殊 Lagrange 子流形在复结构模空间的一个开区域中存在。此外，由\BookEqRef{eq:5.18}{5.18}，$|Z|$ 在这个区域处处非零，因为它就是体积。因此，在区域的一点选择\BookEqRef{eq:5.16}{5.16}的分支，就可以沿路径连续延拓；在所考虑区域的单连通覆盖上，这确定了整个分支。

\SourceParagraph{ch05-5.3.5-intro-005}{338}
对于特殊 Lagrange 子流形，可以用\BookSectionRef{sec:5.2.1.2}{5.2.1.2}中的局部模型，为这一形式论证提供依据；不过，它在 SCFT 中是否正确，并不显然。把\BookEqRef{eq:5.65}{5.65}与\BookSectionRef{sec:5.3.4.1}{5.3.4.1}的论证结合，则在 SCFT 中有效。

\subsubsection{物理论证}\label{sec:5.3.5.1}

\SourceParagraph{ch05-5.3.5.1-001}{338}
\BookEqRef{eq:5.65}{5.65}说明，边界条件除了像\BookSectionRef{sec:3.6}{3.6}那样，把左行模“反射”为右行模，还施加一个由 $\xi$ 参数化的 $U(1)_R$ 变换。最简单的证明，使用玻色化的顶点算符描述。详细解释留给文献\BookCite{ref112}与\BookBibText{ref242}{MS1}第11.3节，因为它需要更多 CFT 背景；不过，论证很短，直观上也应当可信。把玻色子 $\varphi$ 分解为左行模与右行模，就可以写出谱流算符，
\begin{equation}\label{eq:5.66}
\Omega={:}\mathrm{e}^{\mathrm{i}\sqrt{\hat c}\varphi_L}{:},\qquad
\overline\Omega={:}\mathrm{e}^{-\mathrm{i}\sqrt{\hat c}\varphi_R}{:}.
\end{equation}
左行模与右行模 $\varphi_{L,R}$ 由关系
\[
\partial_-\varphi_L=\partial_+\varphi_R=0,\qquad
\varphi_L+\varphi_R=\varphi
\]
定义。这些关系把它们确定到相差一个整体常数，再由用 $\mu$ 表示的 $\varphi$ 边界条件\BookEqRef{eq:5.60}{5.60}固定。利用\BookEqRef{eq:5.66}{5.66}，便确定\BookEqRef{eq:5.64}{5.64}中的 $\xi$。

\SourceParagraph{ch05-5.3.5.1-002}{338}
回顾特殊 Lagrange 的情形，$\xi$ 也是中心荷的复辐角，如\BookEqRef{eq:5.16}{5.16}所示。虽然这个方程，以及 BPS 中心荷的定义\BookEqRef{eq:5.17}{5.17}涉及几何，但若将二者合并为
\begin{equation}\label{eq:5.67}
\xi(L)=\frac{1}{\pi}\arg Z(L)\pmod2,
\end{equation}
这个陈述就可以原样推广到 SCFT。

\SourceParagraph{ch05-5.3.5.1-003}{338}
这个主张直接蕴含前面的主张，即 A 模型中开弦的分次依赖于复结构，以及相应的镜像陈述。

\subsubsection{在 SCFT 中计算中心荷}\label{sec:5.3.5.2}

\SourceParagraph{ch05-5.3.5.2-001}{338}
可以从定义\BookEqRef{eq:5.17}{5.17}出发，把计算翻译为 $\sigma$ 模型语言，然后推广。

\SourceParagraph{ch05-5.3.5.2-002}{338}
为具体起见，考虑 A 型膜。\BookEqRef{eq:5.17}{5.17}中的全纯 $(n,0)$ 形式 $\Omega$，在 $\sigma$ 模型中由\BookEqRef{eq:3.196}{3.196}中的算符 $\Omega$ 表示；这应当是可信的，也很容易验证。

\SourceParagraph{ch05-5.3.5.2-003}{338}
随后，在膜的世界体上对微分形式作积分，就转化为以该膜为边界条件的圆盘图上的期望值。$\sigma$ 模型中，这很清楚，因为圆盘世界面上的泛函积分，包含对零模，即质心的积分，并像\BookSectionRef{sec:3.5}{3.5}那样，受边界条件约束。这一期望值，在一般 SCFT 中也有意义，而且，没有其他候选推广\footnote{可以给出更好的证明，但代价是引入更多物理。}。

\SourceParagraph{ch05-5.3.5.2-004}{339}
因此，$Z(L)$ 是在边界为 $L$ 的圆盘世界面上计算的期望值，算符 $\Omega$ 插入在穿孔处。可以验证，这满足 R 荷反常\BookEqRef{eq:3.110}{3.110}。

\subsubsection{A 与 B 的交换}\label{sec:5.3.5.3}

\SourceParagraph{ch05-5.3.5.3-001}{339}
导出\BookEqRef{eq:5.67}{5.67}的关键是：虽然从 A 型膜出发，实际上，这却是在 B 扭曲模型中作的计算。既然主张 $Z(L)$ 可以在某个拓扑扭曲模型中计算，这就是必然的，因为知道它在大体积极限中只依赖于复结构模参数。不过，像文献\BookCite{ref382}那样，在世界面上明确说明这一点，仍然让人放心。

\SourceParagraph{ch05-5.3.5.3-002}{339}
与前面涉及 A 型膜的计算不同，这实际上是闭弦计算。从通常的物理出发点，即 BPS 中心荷的时空超引力定义，就可以明显看出这一点，因为进入该定义的超引力场，来自闭弦扇区。为了从世界面看出这一点，重新绘制圆盘，把穿孔扩大为第二个边界，如\BookRef{fig:5.4}{图4}所示。所得图自然可看作闭弦图，即圆环。
\begin{figure}[H]
\centering
\BookImage{01_Numbered_Figures/Chapter_05/Ch05_Fig04_PDFp339_Bookp329.pdf}
\caption{具有 D 膜边界与外部闭弦的圆环。}\label{fig:5.4}
\end{figure}

\SourceParagraph{ch05-5.3.5.3-003}{339}
从开弦图转到闭弦图时，需要改变对边界条件的解释：原来，它位于确定的“空间”坐标，开弦端点沿时间运动，因而在空间固定；现在，它成为“时间”方向的边界，即 D 膜发射的一圈闭弦的起点。考虑 SCA 中算符在这种旋转下的变换律，会发现，这个旋转实际上交换了 A、B 边界条件。

\SourceParagraph{ch05-5.3.5.3-004}{339}
最简单的理解方法，是考虑 A 膜对 $U(1)$ 流施加的边界条件\BookEqRef{eq:3.165}{3.165}，即 $J_L=-J_R$。开弦扇区中，这与\BookEqRef{eq:3.100}{3.100}中的 $Q_A=G_{+,0}+\widetilde G_{-,0}$ 相容。

\SourceParagraph{ch05-5.3.5.3-005}{339}
转到闭弦扇区时，反转\BookEqRef{eq:3.138}{3.138}中切向与法向、时间与空间方向之间的对应。由于“左行”与“右行”是通过后者定义的，相对于闭弦约定，这反转了 $J_R$ 的符号，现在有 $J_L=J_R$\footnote{一个更形式化的论证，使用 $J_L\,\mathrm{d}z$ 与 $J_R\,\mathrm{d}\bar z$ 在旋转 $z\mapsto\mathrm{e}^{\mathrm{i}\pi/2}z$ 下的张量变换性质。}。

\SourceParagraph{ch05-5.3.5.3-006}{340}
因此，在计算中心荷时，A 边界条件与拓扑 B 模型配合使用，反之亦然。

\subsubsection{A 膜的中心荷}\label{sec:5.3.5.4}

\SourceParagraph{ch05-5.3.5.4-001}{340}
最重要的后果，来自 B 模型闭弦振幅不受量子修正这一事实。因此，A 膜中心荷的大体积定义\BookEqRef{eq:5.17}{5.17}是精确的，不受 $\alpha'$ 修正。

\SourceParagraph{ch05-5.3.5.4-002}{340}
对于定义为特殊 Lagrange 子流形的 A 膜，\BookEqRef{eq:5.17}{5.17}的意义很清楚；但对于\BookSectionRef{sec:3.6.1}{3.6.1}中更一般的 A 型边界条件，或者可能的“非几何”边界条件，比如在 Gepner 模型中定义的那些，应当如何解释它呢？

\SourceParagraph{ch05-5.3.5.4-003}{340}
物理上，从时空超引力的考虑出发，答案再次更清楚。考虑 IIB 型弦在 Calabi--Yau 流形 $X$ 上紧化，得到 $d=4$、$\mathcal{N}=2$ 的超引力理论。这种理论中，BPS 粒子的中心荷有一个一般公式，推广了\BookEqRef{eq:5.17}{5.17}：
\begin{equation}\label{eq:5.68}
Z(L)=Q_i\Pi^i.
\end{equation}
这里，$Q_i$ 是电荷与磁荷构成的向量，而 $\Pi^i$ 是由超引力理论的“向量多重态”结构确定的复向量。

\SourceParagraph{ch05-5.3.5.4-004}{340}
每个物理 BPS A 型 D 膜，都对应于这个理论中的一个 BPS 粒子，见\BookSectionRef{sec:5.4.2.1}{5.4.2.1}；因此，可以与\BookSectionRef{sec:5.1.2}{5.1.2}的一般讨论联系。那里论证了，A 膜的 RR 荷，或相应 K 理论类的上同调荷，给出 $[L]\in H_3(X,\mathbb{Z})$。选择 $H_3(X,\mathbb{Z})$ 自由部分的显式基 $\Sigma^i$，用它表述\BookEqRef{eq:5.17}{5.17}，写成
\begin{equation}\label{eq:5.69}
[L]=\sum_iQ_i\Sigma^i
\end{equation}
以及
\begin{equation}\label{eq:5.70}
\Pi^i=\int_{\Sigma^i}\Omega,
\end{equation}
就得到\BookEqRef{eq:5.68}{5.68}。对于特殊 Lagrange 子流形 $L$，这与\BookEqRef{eq:5.17}{5.17}一致。不过，这个论证还说明，任意 SCFT 中所有 A 型膜的中心荷，都允许这样的表示；周期不探测可能存在的挠荷。

\SourceParagraph{ch05-5.3.5.4-005}{340}
\BookEqRef{eq:5.68}{5.68}与\BookSectionRef{sec:5.3.5.3}{5.3.5.3}讨论的世界面计算之间的关系是：$Q_i$ 表达圆盘边界条件的 K 理论类，而向量 $\Pi^i$ 给出算符 $\Omega$ 的单点函数的一组完整基。因此，已经把这一世界面计算翻译成几何语言。此外，由于这些单点函数不依赖于 K\"ahler 模参数，若存在大体积极限，就可以在该极限中计算它们。因此，把 $\Pi^i$ 定义为 $\Omega$ 的一组完整周期，也可以用于 SCFT，至少对于具有大体积极限的情形如此。

\SourceParagraph{ch05-5.3.5.4-006}{341}
实际上，有些 $(2,2)$ SCFT，没有严格的大体积极限，因为它们完全没有 K\"ahler 模参数。这些理论是刚性 Calabi--Yau 流形的镜像\footnote{可以设想，有些 $(2,2)$ SCFT 同时具有两类模参数，却没有大体积极限，不过，尚没有已知例子。}。这种 SCFT 不可能是非线性 $\sigma$ 模型，因为那样，其靶空间就必须满足 $b^{1,1}=0$。另一方面，中心荷\BookEqRef{eq:5.70}{5.70}，仍然满足它们作为复簇周期必须满足的关系；已有论证认为，在这种情形中，可以把它们表示为更高维复簇的周期\BookCite{ref083}。

\SourceParagraph{ch05-5.3.5.4-007}{341}
要给出\BookEqRef{eq:5.69}{5.69}的非几何版本，需要讨论具体 SCFT。已经完成这种工作的例子，包括 Gepner 模型\BookCite{ref069}与轨形\BookCite{ref115,ref137}。

\subsubsection{B 膜的中心荷}\label{sec:5.3.5.5}

\SourceParagraph{ch05-5.3.5.5-001}{341}
B 膜的中心荷，在拓扑 A 模型中计算，因而可以受到弦性的 $\alpha'$ 修正。不过，先忽略这些修正；在大体积极限中，这应当是很好的近似。这时，有近似表达式
\begin{equation}\label{eq:5.71}
Z(E)=\int_X\mathrm{e}^{-(B+\mathrm{i}J)}
\operatorname{ch}(E)\sqrt{\operatorname{td}(X)}+O(\alpha'),
\end{equation}
其中，$O(\alpha')$ 表示弦性修正。

\SourceParagraph{ch05-5.3.5.5-002}{341}
对膜的 K 理论类的依赖，由\BookSectionRef{sec:5.1.4.1}{5.1.4.1}与公式\BookEqRef{eq:5.69}{5.69}便可看出。此外，还主张
\begin{equation}\label{eq:5.72}
\Pi_i=\int_{\Sigma_i}\mathrm{e}^{-(B+\mathrm{i}J)}+O(\alpha'),
\end{equation}
其中，$\Sigma_i$ 是偶数维同调 $\bigoplus_pH_{2p}(X,\mathbb{Z})$ 自由部分的一组完整基。

\SourceParagraph{ch05-5.3.5.5-003}{341}
承认弦理论中的若干事实，就有一个简单论证。首先，\BookSectionRef{sec:3.5.2.7}{3.5.2.7}与\BookSectionRef{sec:3.6.1}{3.6.1}已经看到：根据规范对称\BookEqRef{eq:3.150}{3.150}，D 膜所携带的丛的曲率 $F$，必须始终与二形式势 $B$，以 $B-F$ 的组合出现。因此，采用把曲率归一化吸入 $F$ 的记号，$\operatorname{ch}(E)=\operatorname{Tr}\mathrm{e}^F$，\BookEqRef{eq:5.71}{5.71}必须包含组合 $\mathrm{e}^{-B}\operatorname{ch}(E)$。此外，根据 $\mathcal{N}=2$ 超引力的一般性质，本节使用的中心荷 $Z(E)$，必须是复化 K\"ahler 模参数 $B+\mathrm{i}J$ 的全纯函数。这迫使表达式具有\BookEqRef{eq:5.71}{5.71}及其推论\BookEqRef{eq:5.72}{5.72}的形式。

\SourceParagraph{ch05-5.3.5.5-004}{341}
后面将把\BookEqRef{eq:5.71}{5.71}代入\BookEqRef{eq:5.67}{5.67}，并按 $1/J$ 展开，导出大体积极限中 $\xi(E)$ 的公式\BookEqRef{eq:5.142}{5.142}。只要可以忽略 $\alpha'$ 修正，通过连续性，就能把 $\xi(E)$ 确定到模 $2\hat c$。

\SourceParagraph{ch05-5.3.5.5-005}{342}
事实上，物理论证排除了\BookEqRef{eq:5.71}{5.71}，等价地，\BookEqRef{eq:5.72}{5.72}中几乎所有幂次形式的 $\alpha'$ 修正。不过，由镜像对称，或\BookSectionRef{sec:3.4}{3.4}的论证，很清楚仍然必须有弦性修正。文献\BookCite{ref085}最早讨论了这些修正，它们以两种方式出现：
\begin{enumerate}[label=(\arabic*)]
\item 微扰修正，来自文献\BookCite{ref198}分析的非线性 $\sigma$ 模型中的四圈修正。这些修正的 $B+\mathrm{i}J$ 次数，比\BookEqRef{eq:5.71}{5.71}中的首项低三次。因此，对于 Calabi--Yau 三维流形，它至多产生常数项。
\item 非微扰修正产生关于 $q_i$ 的无常数项幂级数，其中，取 $H_2(X)$ 的某组基 $\{C_i\}$，令
\[
q_i=\exp\!\left(2\pi\mathrm{i}\int_{C_i}(B+\mathrm{i}J)\right).
\]
\end{enumerate}

\SourceParagraph{ch05-5.3.5.5-006}{342}
不必深入这些修正的具体计算，最简单的处理方法，是直接承认镜像对称，用它通过刚才给出的 A 膜定义，来定义 B 膜中心荷，如\BookSectionRef{sec:3.4.3}{3.4.3}所概述。乍看似乎是在取巧，但从物理观点看，正如\BookSectionRef{sec:3.4.3.2}{3.4.3.2}讨论的，镜像对称在 SCFT 中明确可见；因此，一旦开始把讨论建立在 SCFT 上，这就完全合理。当然，在把某个 SCFT 定义，与更常规的几何定义相识别时，仍须逐步给出依据。

\SourceParagraph{ch05-5.3.5.5-007}{342}
因此，采用公式\BookEqRef{eq:5.68}{5.68}，作为 B 膜中心荷的基本定义。除了镜像对称，这也来自把 B 膜看作 IIA 型弦在 Calabi--Yau 流形 $X$ 上紧化后得到的 BPS 粒子。许多物理论证使用这种语言，不过，这里不依赖于它。

\SourceParagraph{ch05-5.3.5.5-008}{342}
由闭弦镜像对称，推知周期 $\Pi^i$ 可以定义为：镜像 Calabi--Yau 流形 $Y$ 上全纯 $n$ 形式 $\Omega$ 的周期\BookEqRef{eq:5.70}{5.70}的一组完整基。一旦建立 SCFT 计算，这就很明显，因为镜像对称只把一组顶点算符变成另一组。

\subsubsection{复形的中心荷}\label{sec:5.3.5.6}

\SourceParagraph{ch05-5.3.5.6-001}{342}
定义了单个 A、B 膜的中心荷以后，可以由线性性推出复形的中心荷：
\[
Z(E)=\sum_n(-1)^nZ(E^n).
\]
这直接来自\BookEqRef{eq:5.68}{5.68}，以及 $Q_i$ 是可加守恒荷这一事实；分次只通过它对膜定向的作用，进入 $Q_i$ 的计算。

\SourceParagraph{ch05-5.3.5.6-002}{342}
这只把 $\xi(E)$ 确定到模2。为了把它提升到 $[0,2\hat c)$，再次先在一个点进行选择，然后使用连续性。

\subsection{谱流与分次}\label{sec:5.3.6}

\SourceParagraph{ch05-5.3.6-intro-001}{342}
以上讨论，都在单个 SCFT，实际上，是单个开弦扇区中进行。现在论证，SCFT 的某些变化，会改变边界条件的分次，却不影响拓扑理论内部开弦的性质。

\SourceParagraph{ch05-5.3.6-intro-002}{343}
例如，由于知道拓扑 B 模型只依赖于复结构数据，就可以在非线性 $\sigma$ 模型中，改变 Calabi--Yau 流形的 K\"ahler 类，而不改变开弦谱或关联函数。新的主张是，这仍会影响它们的分次。镜像陈述是：在 A 模型中改变复结构，也有相同效果。

\SourceParagraph{ch05-5.3.6-intro-003}{343}
让 SCFT 作这样的变化，并考虑它对开弦的影响。例如，固定 $A$ 的分次 $a_0$，让 $B$ 的分次 $a_1$ 缓慢增加到 $a_1+x$。这时，可以把 Hilbert 空间 $\mathcal{H}_{AB}$，以及其中任意特定开弦 $\psi\in\mathcal{H}$，看作在过程中基本不变，只有一个例外：随着弦的端点移动，其总缠绕量增加 $x$。

\SourceParagraph{ch05-5.3.6-intro-004}{343}
特别地，没有任何先验理由，要求 $x$ 必须为整数。在物理 SCFT 中，它可以是任意实数。甚至从几何观点看，A 膜最自然的分次定义，也具有这一性质；马上就会回顾。

\SourceParagraph{ch05-5.3.6-intro-005}{343}
继续把刚才在 SCFT 中引入的记号，与\BookChapterRef{4}基于数学引入的记号比较。边界条件 $B$ 的任何显式 CFT 定义，例如，携带丛的 D 膜、分数膜或其他定义，都会给出计算所有以它为起点或终点的开弦分次的某种规定。随后，给 $B$ 赋予任意“分次” $x$，与这一规定一起使用。换句话说，从 $A$ 到 $B$ 的开弦，若在这一 $B$ 的定义中算得费米数为 $n$，就记作
\[
\psi\in\operatorname{Ext}^n(A,B[x]).
\]

\SourceParagraph{ch05-5.3.6-intro-006}{343}
若作某个操作，把 $B$ 的分次变为 $x'$，那么，开弦 $\psi$ 的 $U(1)_R$ 荷“流动”到 $n+x'-x$，记作
\[
\psi\in\operatorname{Ext}^n(A,B[x']).
\]
若 $A$ 的分次从 $y$ 移到 $y'$，也发生类似的流动，这时变为 $n+y-y'$。

\SourceParagraph{ch05-5.3.6-intro-007}{343}
此外，由于 $a_0,a_1$ 与缠绕贡献 $w$，在\BookEqRef{eq:5.41}{5.41}中出现的方式没有区别，也可以选择遵循\BookChapterRef{4}的约定，写成
\[
\psi\in\operatorname{Hom}(A,B[n+x]),
\]
把所有分次的变化，放入开弦终点所在的 $B$ 的“像”的选择中；对于 $A$，则是起点所在的像。这并不表示，开弦谱在分次的任意整数移位下都不变。由于关联\BookEqRef{eq:5.62}{5.62}，只能保证在 $2\hat c$ 的移位下不变。

\subsubsection{特殊 Lagrange 子流形}\label{sec:5.3.6.1}

\SourceParagraph{ch05-5.3.6.1-001}{344}
这时，分次的流动，也可以从\BookSectionRef{sec:3.6.2.5}{3.6.2.5}引述的 Floer 上同调性质导出。回顾\BookSectionRef{sec:5.2.1.2}{5.2.1.2}关于一对横截相交的特殊 Lagrange 子流形 $L_1,L_2$ 的讨论。那里论证了，复结构模空间中有一个区域，两个膜都在其中存在，$\xi(L_i)$ 由\BookEqRef{eq:5.22}{5.22}给出。

\SourceParagraph{ch05-5.3.6.1-002}{344}
若 $L_1,L_2$ 在点 $p$ 相交，Floer 指标（即鬼数）为 $\mu(p)$，那么，
\begin{equation}\label{eq:5.73}
\xi(L_2)-\xi(L_1)+\mu(p)=\frac{1}{\pi}\sum_{j=1}^m\phi_j.
\end{equation}
这来自连续性，以及下面的事实。假定 $L_1,L_2$ 在两点 $p_1,p_2$ 相交，且每个 Lagrange 子流形都具有\BookSectionRef{sec:3.6.1}{3.6.1}中的平凡 Maslov 类。使用文献\BookCite{ref150}的论证，可以说明，$\mu(p_1)$ 与 $\mu(p_2)$ 的差，等于两点处 $\sum_j\phi_j$ 的差。

\SourceParagraph{ch05-5.3.6.1-003}{344}
方程\BookEqRef{eq:5.73}{5.73}，把\BookSectionRef{sec:3.6.2.5}{3.6.2.5}中定义鬼数 $\mu(p)$ 的歧义，与 $\xi$ 定义中的歧义联系起来。给每个 A 膜 $L$ 标记一个整数鬼数 $\mu(L)$，就固定了 $\mu(p)$ 的歧义。现在，令 $L[n]$ 为与 $L$ 完全相同的 A 膜，只把鬼数增加 $n$。由\BookEqRef{eq:3.189}{3.189}与\BookEqRef{eq:5.73}{5.73}，得到
\begin{equation}\label{eq:5.74}
\xi(L[n])=\xi(L)+n.
\end{equation}

\SourceParagraph{ch05-5.3.6.1-004}{344}
当然，\BookSectionRef{sec:5.2.1.2}{5.2.1.2}的主要结论，是连通和，即束缚态 $L_1\#L_2$ 的形成，由差 $\xi(L_1)-\xi(L_2)$ 控制。根据\BookSectionRef{sec:5.3.6}{5.3.6}，这将转化为：负责形成这一束缚态的开弦，其 $U(1)_R$ 荷必须满足某个条件。\BookSectionRef{sec:5.4.4}{5.4.4}将讨论这一点。

\section{有效世界体理论}\label{sec:5.4}

\SourceParagraph{ch05-5.4-intro-001}{344}
到目前为止，讨论集中在弦理论与 Dirichlet 膜的世界面方法，例如 $\sigma$ 模型与 SCFT。在这些方法中，从弦的二维世界面出发，导出其余结构。

\SourceParagraph{ch05-5.4-intro-002}{344}
另一种物理方法很不相同：先从四维时空物理的自洽条件出发，分类它们可能的实现，再尝试找出满足其他约束的实现。这些约束可能来自实验数据，也可能是理论约束，例如理论已知的对称性。这种方法通常称为“有效场论”，或 EFT。

\SourceParagraph{ch05-5.4-intro-003}{345}
这是一个庞大的主题，其中与这里直接相关的，是四维 $\mathcal{N}=1$ 超对称箭图规范理论\footnote{用 $\mathcal{N}$ 表示超对称的数目，把 $N$ 留给其他用途，例如规范群的秩。}。根据十分一般的物理理由，可以论证，一组 A 型或 B 型 Dirichlet 膜，必须对应于这种理论。因此，EFT 方法能很快触及主要问题之一，即对各种可能性给出简单的一般描述。

\SourceParagraph{ch05-5.4-intro-004}{345}
主要讨论量子化以前的经典理论。在很大程度上，可以把它看作\BookSectionRef{sec:4.2}{4.2}所讨论的带关系箭图的另一种语言。数学与物理最初研究这些概念的动机很不相同，这种联系因而相当引人注目。

\SourceParagraph{ch05-5.4-intro-005}{345}
$\mathcal{N}=1$ 超对称箭图 EFT 特有的一点，是阻碍理论由超势决定，也就是\BookSectionRef{sec:4.2.5}{4.2.5}引入的超势。这一点的物理论证，并不直接依赖弦理论。不过，这在数学上意味着，箭图的路代数是同调维数为 3 的 Calabi--Yau 代数。这个代数表示从紧化维度得到的物理，因此，这类理论已经显露出某种“十维”性质。

\SourceParagraph{ch05-5.4-intro-006}{345}
量子化这些理论以后，还会出现额外的物理内容，它们在数学上尚未得到同样清楚的表述。大部分物理研究关注这一点；\BookSectionRef{sec:5.4.6}{5.4.6}将简要介绍其中几个方面。

\subsection{\texorpdfstring{$\mathcal{N}=1$}{N=1} 超对称规范理论}\label{sec:5.4.1}

\SourceParagraph{ch05-5.4.1-intro-001}{345}
标准的物理参考文献是\BookCite{ref460}，面向数学读者的介绍见文献\BookCite{ref153}。这里给出适合当前目的的提纲式描述。

\SourceParagraph{ch05-5.4.1-intro-002}{345}
四维有效场论，是从四维时空 $M$，例如 $M\cong\mathbb{R}^{3,1}$，到某个配置空间 $\mathcal{C}$ 的映射所构成的量子场论。$\mathcal{N}=1$ 超对称 EFT 中，场必须允许 $\mathcal{N}=1$ 超对称代数的作用，与\BookEqRef{eq:3.87}{3.87}有些类似。这对场及作用量施加了约束，其形式与一般解，见文献\BookCite{ref460}。

\SourceParagraph{ch05-5.4.1-intro-003}{345}
结果是，$d=4$、$\mathcal{N}=1$ 的 EFT，本质上是把\BookSectionRef{sec:3.3.2}{3.3.2}定义的超对称非线性 $\sigma$ 模型推广到包含规范不变性\footnote{这种说法有所简化，不过，其核心依据是，两类理论都具有4个实超对称：\BookSectionRef{sec:3.3.2}{3.3.2}讨论的二维理论中，有左右行的 $(2,2)$ 超对称；四维理论中，则对应于 Lorentz 群，或 Euclid 群，不可约旋量表示的4个实分量。此外，总可以把 $d=4$、$\mathcal{N}=1$ 理论作维数约化，得到 $d=2$、$\mathcal{N}=(2,2)$ 理论。不过，二维理论还有不能从四维得到的额外可能性。}。针对当前目的，把它简化为下面的定义。
\BookStatementNumber{17}
\begin{definition}\label{def:5.17}
$\mathcal{N}=1$ EFT 是四元组 $(\mathcal{C},H,\mu,W)$，其中：
\begin{itemize}
\item $\mathcal{C}$，即配置空间，是具有 K\"ahler 度规的复流形。度规记为 $g_{i\bar j}$，K\"ahler 形式记为 $\omega$。
\item $H$，即规范群，是紧 Lie 群，以等距变换作用于 $\mathcal{C}$，其复化 $H_{\mathbb{C}}$ 全纯地作用。
\item $\mu$ 是 $H$ 作用的矩映射。
\item $W$ 是 $\mathcal{C}$ 上的全纯 $H$ 不变函数，即超势。
\end{itemize}
\end{definition}

\subsubsection{矩映射}\label{sec:5.4.1.1}

\SourceParagraph{ch05-5.4.1.1-001}{346}
关于超对称理论中的群作用与矩映射，文献\BookCite{ref239}是很好的数学取向参考。回顾一下，辛流形上的群作用的矩映射，是函数 $\mu\colon\mathcal{C}\to(\operatorname{Lie}H)^*$，如下确定无穷小群作用 $h\colon\operatorname{Lie}H\to\Gamma(T\mathcal{C})$。任意 $t\in\operatorname{Lie}H$，确定函数 $F_t$，
\[
F_t(x)=\langle\mu(x),t\rangle,
\]
且这个函数满足
\begin{equation}\label{eq:5.75}
\iota(h(t))\omega=\mathrm{d}F_t,
\end{equation}
其中，$\iota(h(t))\omega$ 表示向量场 $h(t)$ 与二形式 $\omega$ 的缩并。若 $H$ 半单，则 $\omega$ 与 $h$，连同附加条件 $F_{[t_1,t_2]}=\{F_{t_1},F_{t_2}\}$，确定矩映射 $\mu$；这里，$\{\ ,\ \}$ 表示 Poisson 括号。若 $H$ 包含 Abel 因子 $H_{\mathrm{abelian}}$，则仍允许加入 $(\operatorname{Lie}H_{\mathrm{abelian}})^*$ 中的元素。物理上，称这一自由度为加入“Fayet--Iliopoulos 项”的可能性。

\SourceParagraph{ch05-5.4.1.1-002}{346}
经常，$\mathcal{C}$ 是 Euclid 空间，$H$ 在线性地作用。这时，可以等价地指定 $\mathcal{C}\cong R$，其中，$R$ 是 $H_{\mathbb{C}}$ 的线性表示。矩映射 $\mu$ 就是 $\mathcal{C}$ 自然坐标的二次函数。这也很好地近似描述 $\mathcal{C}$ 中的一个小邻域。

\SourceParagraph{ch05-5.4.1.1-003}{346}
\BookSectionRef{sec:3.3.2}{3.3.2}中的 $\sigma$ 模型，是 $H$ 平凡、$W=0$，并约化到 $d=2$ 的特殊情形。由于在经典层面工作，这里不施加 $\mathcal{C}$ 为 Ricci 平坦的约束。

\SourceParagraph{ch05-5.4.1.1-004}{346}
虽然不给出类似于\BookEqRef{eq:3.86}{3.86}的一般 Lagrange 量，额外数据 $(H,\mu,W)$ 的用途如下。令 $M$ 为时空，例如 $M\cong\mathbb{R}^{3,1}$。与\BookChapterRef{3}一样，引入作为映射 $M\to\mathcal{C}$ 的“场”，以及费米型“超伴”。随后，在 $M$ 上引入 $H$ 联络，用它把 Lagrange 量中所有时空导数协变化。再寻找允许类似\BookEqRef{eq:3.87}{3.87}的超对称变换的 Lagrange 量。使用超空间形式体系，这会相当简单：通过对函数 $W$ 与 $\mu$，以及 $\mathcal{C}$ 上的 K\"ahler 势作超空间积分，确定 Lagrange 量；细节见文献\BookCite{ref460}。

\subsubsection{标量势}\label{sec:5.4.1.2}

\SourceParagraph{ch05-5.4.1.2-001}{347}
超对称的一项关键约束，是标量势必须具有如下形式：
\begin{equation}\label{eq:5.76}
V=\|\partial W\|^2+\frac{1}{2}\|\mu\|^2.
\end{equation}
最一般的 EFT 中，标量势原本可以是 $\mathcal{C}$ 上的实值函数；这里，它必须是全纯超势 $W$ 的梯度范数平方，与矩映射范数平方的和。
\BookStatementNumber{18}
\begin{definition}\label{def:5.18}
\emph{真空}是 $V$ 的局部极小点 $p\in\mathcal{C}$ 的 $H$ 轨道。这些点满足
\begin{equation}\label{eq:5.77}
\partial V|_p=0.
\end{equation}
\end{definition}

\SourceParagraph{ch05-5.4.1.2-002}{347}
注意，使势极小的 $H$ 轨道，未必孤立。若不孤立，物理学家把 $\mathcal{C}/H$ 中 $V$ 保持常数的路径方向，称为“平坦方向”。下面是一个特殊情形。
\BookStatementNumber{19}
\begin{definition}\label{def:5.19}
\emph{超对称真空}是满足
\begin{equation}\label{eq:5.78}
\partial W|_p=0\qquad\text{F 项条件}
\end{equation}
及
\begin{equation}\label{eq:5.79}
\mu(p)=0\qquad\text{D 项条件}
\end{equation}
的点 $p$ 的 $H$ 轨道。
\end{definition}

\SourceParagraph{ch05-5.4.1.2-003}{347}
F 与 D 的名称，在物理文献中普遍使用。一般而言，还可以有其他局部极小点，称为非超对称真空。

\SourceParagraph{ch05-5.4.1.2-004}{347}
把关注范围限制为超对称真空。它们的模空间，是 $W$ 的临界点簇的辛 $H$ 商，
\begin{equation}\label{eq:5.80}
\mathcal{M}=\{p\in\mathcal{C}:\partial W|_p=0\}\mathbin{//}(H,\mu).
\end{equation}

\SourceParagraph{ch05-5.4.1.2-005}{347}
当然，这是经典定义。\BookSectionRef{sec:5.4.6.1}{5.4.6.1}将讨论量子化的若干方面，以及相应自洽条件。

\subsubsection{超对称箭图规范理论}\label{sec:5.4.1.3}

\SourceParagraph{ch05-5.4.1.3-001}{347}
这是一个特殊情形：$H_{\mathbb C}$ 是与箭图对应的一组向量空间 $V_i$ 的自同构群，如\BookSectionRef{sec:4.2}{4.2}所述，
\begin{equation}\label{eq:5.81}
H=\prod_{i=1}^r U(\alpha_i),
\end{equation}
而 $\mathcal{C}$ 是那里引入的线性映射空间，
\[
\mathcal{C}=\bigoplus_a\operatorname{Hom}(V_{t(a)},V_{h(a)}),
\]
带 Euclid 度规。给定这个 $H$，矩映射由权向量 $\vec\theta$ 确定，每个 $U(n)$ 因子对应一个权。

\SourceParagraph{ch05-5.4.1.3-002}{347}
并非所有代数都对应于箭图规范理论；箭图的关系，必须能够表述为单个函数的梯度\BookEqRef{eq:5.78}{5.78}。这样的代数，称为“带超势的代数”。

\SourceParagraph{ch05-5.4.1.3-003}{348}
反过来，并非所有超势 $W$ 都导出箭图关系。若 $W$ 具有“单迹”形式，即，它是与箭图中闭环对应的单项式的加权和的迹，就会如此。

\SourceParagraph{ch05-5.4.1.3-004}{348}
\BookSectionRef{sec:5.4.5.5}{5.4.5.5}将说明，D 项条件，正是\BookRef{thm:4.88}{定理4.88}讨论的条件。因此，
\BookStatementNumber{20}
\begin{proposition}\label{prop:5.20}
超对称箭图规范理论，对应于一个带关系的箭图，其关系来自超势。超对称真空是箭图的 $\theta$ 半稳定表示的 S 等价类；等价地，是 $\theta$ 多稳定表示的同构类。
\end{proposition}

\SourceParagraph{ch05-5.4.1.3-005}{348}
\BookStatementNumber{21}
\begin{example}\label{ex:5.21}
一个很基本的例子，是 $\mathcal{N}=4$ 超 Yang--Mills 理论，简称 SYM。这时，$H=U(N)$，$\mathcal{C}\cong\mathbb{C}^3\otimes\operatorname{Mat}_N$，带 Euclid 度规，且
\begin{equation}\label{eq:5.82}
W=\operatorname{Tr}\bigl([Z^1,Z^2]Z^3\bigr),
\end{equation}
其中，$Z^i$ 是 $\mathcal{C}$ 的一个显式 $N\times N$ 复矩阵参数化\footnotemark。F 项方程为
\[
[Z^i,Z^j]=0.
\]
因此，很容易验证，当 $\theta=0$ 时，超对称真空的模空间是
\[
\mathcal{M}\cong(\mathbb{C}^3)^N/S_N
=\operatorname{Sym}^N\mathbb{C}^3,
\]
即 $\mathbb{C}^3$ 的第 $N$ 个对称积；当 $\theta\ne0$ 时，模空间为空。比较\BookRef{ex:4.33}{例题4.33}。
\end{example}
\footnotetext{虽然这里是与箭图相关的情形，实际上，$H$ 可以是任意紧 Lie 群。这时，$Z^i$ 取在复化伴随表示中，超势\BookEqRef{eq:5.82}{5.82}是三次反对称不变量。}

\SourceParagraph{ch05-5.4.1.3-006}{348}
\BookStatementNumber{22}
\begin{example}\label{ex:5.22}
考虑与\BookRef{ex:4.31}{例题4.31}的箭图对应的理论，它有两个顶点，由 $n$ 条箭连接，维数向量为 $(1,1)$。因此，$H=U(1)\times U(1)$，$\mathcal{C}\cong\mathbb{C}^n$。取 $\mathcal{C}$ 上的 Euclid 度规，并引入显式坐标 $\psi^i$。

\SourceParagraph{ch05-5.4.1.3-007}{348}
回看\BookRef{def:5.17}{定义5.17}，已经几乎唯一地指定了 EFT。剩余数据，是 $H$ 的矩映射，由一对实数 $\theta_1,\theta_2$ 指定。整体的对角 $U(1)$ 不作用于 $\psi^i$，其矩映射水平必须为零，即 $\theta_1+\theta_2=0$；否则没有超对称真空。假定这些数满足该条件，那么，超对称真空就是矩映射条件\BookEqRef{eq:5.79}{5.79}的解模去规范等价。对实际作用的相对 $U(1)$，采用相应归一化，具体条件为
\begin{equation}\label{eq:5.83}
0=\mu=\theta_1-\theta_2+\sum_i|\psi^i|^2.
\end{equation}
因此，在上述整体水平条件下，真空存在，当且仅当 $\theta_2\ge\theta_1$。若 $\theta_2=\theta_1$，有唯一真空；若 $\theta_2>\theta_1$，模空间为 $\mathbb{CP}^{n-1}$。

\SourceParagraph{ch05-5.4.1.3-008}{348}
可以看到，用\BookChapterRef{4}的语言说，这些超对称真空，是箭图多稳定表示的同构类。表示的不可分解性，在物理语言中对应于规范对称的破缺。更准确地说，给定表示 $R$，其自同构群的极大紧子群 $H_R\subset\operatorname{Aut}R$，是未破缺的规范群。因此，$\theta_2=\theta_1$ 时的唯一真空，具有未破缺的 $U(1)^2$；而当 $\theta_2>\theta_1$ 时，说“$\psi$ 的期望值把规范对称破缺到 $U(1)$”。
\end{example}

\SourceParagraph{ch05-5.4.1.3-009}{349}
\BookStatementNumber{23}
\begin{example}[\texorpdfstring{$\mathcal{N}=2$}{N=2} 超对称 Yang--Mills 理论]\label{ex:5.23}
这是 $\mathcal{N}=1$ 的特殊情形，满足
\[
\mathcal{C}\cong(\operatorname{ad}H\otimes\mathbb{C})\oplus R\oplus R^*,
\]
其中，$R\cong\mathbb{C}^N$ 是 $H$ 的线性表示，$R^*$ 是其对偶；由于 $H$ 幺正，这等价于复共轭表示。取 $\mathcal{C}$ 为 Euclid 空间，把三个因子中的自然坐标记为 $A,Z,\widetilde Z$，则有
\[
W=\langle\widetilde Z,\iota(A)Z\rangle,
\]
这里的配对，就是相应指标的迹缩并。例如，$\mathcal{N}=4$ SYM 是 $R\cong\operatorname{ad}H\otimes\mathbb{C}$ 的特殊情形。
\end{example}

\SourceParagraph{ch05-5.4.1.3-010}{349}
\BookStatementNumber{24}
\begin{example}\label{ex:5.24}
超对称 QCD 对应于\BookRef{fig:5.5}{图5}的箭图。\BookSectionRef{sec:5.4.4.5}{5.4.4.5}将更详细地讨论。
\end{example}
\begin{figure}[H]
\centering
\BookImage{01_Numbered_Figures/Chapter_05/Ch05_Fig05_PDFp349_Bookp339.pdf}
\caption{$\mathcal{N}=1$ 超对称 QCD。}\label{fig:5.5}
\end{figure}

\subsubsection{质量}\label{sec:5.4.1.4}

\SourceParagraph{ch05-5.4.1.4-001}{349}
为了方便更偏数学的读者，陈述一些基本物理定义与论证。
\BookStatementNumber{25}
\begin{definition}\label{def:5.25}
真空 $p$ 的\emph{质量矩阵} $M^2$，是在正交归一标架中表示的 Hessian $\partial_i\partial_jV|_p$。
\end{definition}

\SourceParagraph{ch05-5.4.1.4-002}{349}
由于 $\partial V|_p=0$，它是张量。而且，它是实对称矩阵，所以，可以对角化。它的本征值称为质量平方，其核给出“无质量场”的方向；物理上还须去除纯规范方向。

\SourceParagraph{ch05-5.4.1.4-003}{349}
因此，真空模空间的切方向，包含在无质量场的子空间中。当然，若 $V$ 含有三次或更高次非线性项，真正的模参数方向，可以只构成这个空间的一部分。

\subsubsection{积去}\label{sec:5.4.1.5}

\SourceParagraph{ch05-5.4.1.5-001}{349}
假定在 $\mathcal{C}$ 的某个区域，某个场 $\phi$ 始终有质量；那么，可以求解
\[
0=\frac{\partial V}{\partial\phi}=M^2\phi+O(\phi^2),
\]
解出 $\phi$，并消去它。这称为\emph{积去}场 $\phi$，因为它是对一部分场作泛函积分这一量子过程的经典近似。

\SourceParagraph{ch05-5.4.1.5-002}{349}
$\mathcal{N}=1$ 理论中，若 $\partial^2W/\partial\phi^2\ne0$，同样的过程也可用于超势。当然，当 $\phi$ 在 $W$ 中至多以二次形式出现时，这最有用，因为解唯一。

\SourceParagraph{ch05-5.4.1.5-003}{350}
若 $\phi$ 在 $W$，或 $V$，中线性出现，则称它为约束 $0=\partial W/\partial\phi$ 的 Lagrange 乘子。若该约束定义的子簇 $S$ 是子流形，那么，在施加相应规范条件以后，超对称真空由 $W|_S$ 的临界点给出。

\subsection{EFT 与 Dirichlet 膜的一般说明}\label{sec:5.4.2}

\SourceParagraph{ch05-5.4.2-intro-001}{350}
考虑 II 型弦在 Minkowski 时空 $\mathbb{R}^{3,1}$ 与 Calabi--Yau 流形 $X$ 的乘积上紧化；更一般地，内部理论可以是 $\hat c=3$ 的 $(2,2)$ SCFT。现在加入一组边界条件，嵌入于 $\mathbb{R}^{3,1}\times L$，其中，$L\subset X$ 是某个 $n$ 维子流形。由于这些膜充满 Minkowski 方向，常称为“充满空间的膜”。

\SourceParagraph{ch05-5.4.2-intro-002}{350}
本书几乎总是按照子流形 $L\subset X$ 的维数，称呼 Dirichlet 膜。不过，现在讨论的膜，其世界体有 $3+n$ 个空间维度，所以，按照标准物理术语，它们是 Dirichlet $(3+n)$ 膜。

\SourceParagraph{ch05-5.4.2-intro-003}{350}
若所有体积相对于弦尺度都很大，可以用十维超引力描述闭弦，并把它与一个 $(4+n)$ 维超对称 Born--Infeld 作用量耦合，后者具有\BookEqRef{eq:3.142}{3.142}的一般形式。若 $L$ 是极小子流形，Yang--Mills 场强也很小，还可以进一步近似为膜世界体上的超 Yang--Mills 理论，
\begin{equation}\label{eq:5.84}
S=\int_{\mathbb{R}^{3,1}\times L}\mathrm{d}^{4+n}x\,
\operatorname{tr}\bigl(F^2+(DX_N)^2+\text{费米项}\bigr).
\end{equation}
这里，$X_N$ 是取值于 $L$ 法丛的标量场，同样采用小场近似。

\SourceParagraph{ch05-5.4.2-intro-004}{350}
当然，弦理论预言的底层时空十维性质，仍是一个假设。要直接看到这一点，需要在很高的能量 $E\sim\hbar c/R$ 下做实验，其中，$R$ 是 $X$ 或 $L$ 的直径。在技术进步使这种实验成为可能以前，弦理论研究者主要关心的，不会是完整的 $(4+n)$ 维场论，而是它到 $\mathbb{R}^{3,1}$ 的“Kaluza--Klein 约化”，或简称 KK 约化。

\SourceParagraph{ch05-5.4.2-intro-005}{350}
对于当前目的，$(d+4)$ 维场论的 KK 约化，是一个四维场论；保留膜的轻模参数以后，其配置空间 $\mathcal{C}$，由膜配置的模空间描述。一个基本例子，是 $L$ 为 $X$ 上一点，即 D0膜。这时，单个膜的模空间就是 $X$ 本身。另一类一般例子，是给定作用量为\BookEqRef{eq:5.84}{5.84}的 $(4+n)$ 维理论，其 KK 约化得到的理论，以携带 Yang--Mills 方程解的体积极小子流形的模参数为变量。

\SourceParagraph{ch05-5.4.2-intro-006}{350}
假定取其中一种膜的 $N$ 个相同副本。由于这些膜是完全相同的对象（物理上它们是玻色子），直观上会预期，它们的模空间是对称化的轨形乘积 $\operatorname{Sym}^N\mathcal{C}$。例如，$N$ 个相同 D0膜的模空间，是 $\operatorname{Sym}^NX$。

\SourceParagraph{ch05-5.4.2-intro-007}{351}
不过，知道这种乘积不是流形，就应当问，重合产生的奇点如何处理。物理答案是：当其中 $n\le N$ 个膜重合时，四维场论具有更大的规范对称，从 $U(1)^N$ 增强为 $U(n)\times U(1)^{N-n}$。若多个子集分别重合，规范群就是相应点层直和的自同构群的幺正子群。这来自涉及 Higgs 机制的一个简单物理论证，留作练习。

\SourceParagraph{ch05-5.4.2-intro-008}{351}
因此，尽管模空间可能奇异，完整规范理论却不因此奇异。至少在弦树级层面，这一“机制”是处理 D 膜模空间奇点的一般方式。物理上，其关键方面之一，是在模空间的奇点处，额外的场变为无质量；这里，它们是增强规范对称的“非对角规范玻色子”。这一观察可以推广到其他场变为无质量的情形。例如，还有另一种奇点“消解”机制，其中，缠绕的 D 膜变为无质量\BookCite{ref432}；不过，要看到它如何起作用，必须超越弦微扰理论。

\SourceParagraph{ch05-5.4.2-intro-009}{351}
由于这一机制在数学上自然，发现它在例子中与前面的奇点消解数学讨论一致，并不应当令人意外。

\SourceParagraph{ch05-5.4.2-intro-010}{351}
离开大体积极限以后，\BookEqRef{eq:5.84}{5.84}会受到 $\alpha'$ 修正；若考虑更高亏格曲面，还可能受到弦耦合 $g_s$ 修正。不过，只要仅考虑在 $\mathbb{R}^{3,1}$ 中缓慢变化的映射，即 $\alpha'|\partial\phi|^2\ll1$，其中，$\partial$ 是 $\mathbb{R}^{3,1}$ 中的导数，就仍然能导出类似的四维场论，也就是有效场论。因此，无论 $X$ 中膜的完整物理如何，都有一部分可以在本节建立的 EFT 框架中描述。这部分包括决定可能真空的标量势，以及轻粒子的谱与耦合，所以，具有很大的物理意义。

\subsubsection{Dirichlet 膜与箭图}\label{sec:5.4.2.1}

\SourceParagraph{ch05-5.4.2.1-001}{351}
若限制在弦树级，即球面与圆盘图，一组 Dirichlet 膜就由箭图规范理论描述。为了从\BookSectionRef{sec:3.5.2.8}{3.5.2.8}的 Chan--Paton 规定看出这一点，引入向量空间 $V=\bigoplus_{i\in\mathcal{I}}\mathbb{C}_i$，其中，$\mathcal{I}$ 是 Chan--Paton 标记集，每个标记给出一条复直线。再把\BookSectionRef{sec:4.2}{4.2}的向量空间 $V_i$，识别为 $V$ 中对应于同一种类型 $B_i$ 的所有膜的子空间。按照 Chan--Paton 规定，用接近边缘的边界算符空间，可以得到 $\mathcal{C}$ 的一个小的线性坐标片。

\SourceParagraph{ch05-5.4.2.1-002}{351}
规范理论 Lagrange 量中的项，对应于这些边界算符的关联函数；其 Chan--Paton 因子沿边界依次缩并。对于只有一个边界的 Riemann 曲面，例如圆盘，所得相互作用满足单迹性质。

\SourceParagraph{ch05-5.4.2.1-003}{352}
与\BookChapterRef{4}的定义比较，就得到一直在使用的对应：边界条件成为箭图表示范畴的对象，开弦成为态射。这个论证不要求超对称，也不要求特定维数。

\SourceParagraph{ch05-5.4.2.1-004}{352}
还需要说明，规范场只来自保持边界条件的算符，即按 $\operatorname{End}(V_i)$ 变换的算符，而非改变边界条件的算符。这里使用 CFT 分解为 $\mathbb{R}^{3,1}$ 因子与“内部”CFT 的乘积，以及规范场边界算符\BookEqRef{eq:3.144}{3.144}已经是边缘算符这一事实。因此，为了得到边缘算符，边界算符的内部部分必须具有维数零。在幺正理论中，这意味着，它是恒等算符，因而保持边界条件。

\SourceParagraph{ch05-5.4.2.1-005}{352}
内部 CFT 的边缘算符，对应于真空模空间的切向量；如上所述，在线性化理论中，这是无质量标量场。更一般地\BookCite{ref394}，EFT 中的质量平方，与 CFT 中对应边界算符的维数 $h$ 的关系为\footnote{对于闭弦，$h$ 是左维数或右维数，比例常数也不同。超弦中也有类似公式，把 $-1$ 换成 $-1/2$，对应 NS 扇区，或 $-5/8$，对应 R 扇区。}
\begin{equation}\label{eq:5.85}
m^2=\frac{h-1}{\alpha'}.
\end{equation}
$m^2=0$、$m^2>0$ 与 $m^2<0$ 三种情形，直接对应于\BookSectionRef{sec:3.2.5}{3.2.5}讨论的三种情形。改变 $m^2>0$ 的有质量场，会使配置离开真空集合。\BookSectionRef{sec:5.4.4}{5.4.4}将讨论 $m^2<0$ 的情形。

\SourceParagraph{ch05-5.4.2.1-006}{352}
这些论证都不要求超对称，所以，也不要求 $X$ 为 Calabi--Yau 流形。另一方面，其他已知的 Ricci 平坦流形，也与超对称有关，下面将讨论。打开 H 通量，并允许伸缩子场不为常数，可以定义更一般的一类共形场论。关于这些背景中的 Dirichlet 膜，一些结果见文献\BookCite{ref363,ref440}。

\subsubsection{D0 量子力学与束缚态}\label{sec:5.4.2.2}

\SourceParagraph{ch05-5.4.2.2-001}{352}
若不取充满空间的膜，而在 Minkowski 时空中采用其他嵌入，有些论证会更简单。例如，\BookSectionRef{sec:5.3.4}{5.3.4}使用的“BPS 粒子”语言，来自把膜嵌入 $\mathbb{R}^{0,1}\times L$。虽然这两种膜不同，物理行为也很不相同，但为了在弦树级导出 EFT，它们本质上相同。BPS 粒子语言有利于讨论中心荷；另一方面，由于四维中超对称约束更强，充满空间的膜的语言，有利于导出 EFT。

\SourceParagraph{ch05-5.4.2.2-002}{353}
所谓“本质上相同”，具体指如下内容。假定有一个 D3膜的世界体理论，具有前面使用的显式形式。特别地，它包含规范联络 $A$，即 $\mathbb{R}^{3,1}$ 上取值于 $H$ 的 Lie 代数的1形式。用 Lie 代数的一组显式基 $t_a$，以及1形式 $\mathrm{d}x^i$，把它写为
\[
A=\sum_{\substack{1\le a\le\dim H\\0\le i\le3}}A_i^at_a\,\mathrm{d}x^i.
\]
若直接去掉 $\mathbb{R}^3$ 中所有空间导数，并把 $i=1,2,3$ 的类空分量 $A_i^a$ 改记为 $X_i^a$，就得到 D0膜的量子力学。

\SourceParagraph{ch05-5.4.2.2-003}{353}
把这一规定用于规范理论作用量，计算标量势，会找到三组项。用非负范数表示，它们为
\begin{equation}\label{eq:5.86}
V=V_1(\phi)+\sum_i\|X_i\cdot\phi\|^2
+\sum_{i<j}\|[X_i,X_j]\|^2.
\end{equation}
这里，$V_1(\phi)$ 是\BookEqRef{eq:5.76}{5.76}，$X_i\cdot\phi$ 表示规范 Lie 代数对标量的作用；在箭图的每条箭上，它为 $X_i^{h(a)}\phi_a-\phi_aX_i^{t(a)}$。第二项来自协变导数 $D\phi$，第三项来自 Yang--Mills 作用量\footnote{符号选择使 $V\ge0$。采用相应矩阵约定时，这些非负范数项，也可写成负的迹平方。}。

\SourceParagraph{ch05-5.4.2.2-004}{353}
计算这个 D0 量子力学中可能的时间无关解。首先，取 $\phi$ 对应于一个可能可约的箭图表示 $R$，使 $V_1(\phi)$ 极小。随后，若 $X_i$ 取值于 $\mathbb{R}^3\otimes\mathfrak{c}_R$，其余项就达到极小；这里，$\mathfrak{c}_R$ 是未破缺规范群 $H_R$ 的 Cartan 子代数。最后，模去 $H_R$ 的 Weyl 群。

\SourceParagraph{ch05-5.4.2.2-005}{353}
因此，若表示 $R$ 满足
\begin{equation}\label{eq:5.87}
H_R=\prod_{i=1}^rU(N_i),
\end{equation}
则得到模空间
\[
\mathcal{M}_R=\prod_i\bigl((\mathbb{R}^3)^{N_i}/S_{N_i}\bigr).
\]
自然地把它解释为：一组配置，包含某一类型的 $N_1$ 个相同粒子、第二类型的 $N_2$ 个相同粒子，依此类推。

\SourceParagraph{ch05-5.4.2.2-006}{353}
特别地，可以考虑这些位置 $\vec X^{\,a}=(X_i^a)$，在 $\mathbb{R}^3$ 中彼此相距很远的区域。直观上会预期，在该区域中，粒子之间的相互作用可以忽略；这也可以证明\footnote{虽然 $\operatorname{Hom}(B_i,B_j)$ 中的开弦仍然存在，\BookEqRef{eq:5.86}{5.86}的第二项，使它们获得很大的质量平方 $m^2>0$，所以，在寻找经典真空时，可以忽略它们。结合\BookSectionRef{sec:5.4.4}{5.4.4}，这一论证很有用，因为它给出定义“快子场”的精确方法：否则，这些场会导致不稳定性。有质量场也可以影响量子理论。}。因此，EFT 的物理变成一组独立 EFT 的物理，每个 $H_R$ 中的 $U(1)$ 因子，对应一个 EFT。这就是满足\BookEqRef{eq:5.87}{5.87}的表示 $R$，描述 $N=\sum_iN_i$ 个粒子的精确含义。

\SourceParagraph{ch05-5.4.2.2-007}{354}
这一结构，至少在量子力学的经典极限中，给出“束缚态”与“束缚过程”的物理定义。满足真空条件的不可分解表示 $R$，对应于单个粒子，因为它位于 $\mathbb{R}^3$ 的一点。假定可以把它连续改变为另一个表示 $R'$，满足 $H_{R'}\cong U(1)^2$。这样的表示对应于两个粒子，所以，这就是一个显式的粒子衰变模型。

\subsection{超对称 Dirichlet 膜}\label{sec:5.4.3}

\SourceParagraph{ch05-5.4.3-intro-001}{354}
现在假定，边界条件是\BookSectionRef{sec:3.6.6}{3.6.6}定义的 A 型或 B 型。这些定义，特别是\BookEqRef{eq:3.196}{3.196}与\BookEqRef{eq:3.197}{3.197}，在物理上来自边界条件保持四维 $\mathcal{N}=1$ 超对称的要求。因此，刚才的一般论证说明，这些膜由 $\mathcal{N}=1$ 超对称 EFT 描述，其中，$H$ 如\BookEqRef{eq:5.81}{5.81}所示。

\SourceParagraph{ch05-5.4.3-intro-002}{354}
超对称 EFT 的一般约束，为这些问题的许多一般方面提供简短论证。例如，它们说明，A 膜模空间的光滑部分是 K\"ahler 流形；\BookChapterRef{6}将直接验证这一点。

\subsubsection{BPS 态}\label{sec:5.4.3.1}

\SourceParagraph{ch05-5.4.3.1-001}{354}
一般而言，BPS 态是超对称理论中保持一部分超对称的局域配置。

\SourceParagraph{ch05-5.4.3.1-002}{354}
对于本书的目的，BPS 态是某组 D 膜的超对称箭图规范理论中的超对称真空，如\BookRef{prop:5.20}{命题5.20}所概括。BPS 粒子，是把规范群破缺到 $U(1)$ 的真空，换句话说，是满足这些真空条件的不可分解表示。

\SourceParagraph{ch05-5.4.3.1-003}{354}
必须指出，这是“经典”定义；许多弦理论文献中的“BPS 态”，则指其量子版本\footnote{这里的“量子”，是亏格展开，即弦耦合展开意义下的量子。}。特别地，所谓“BPS Dirichlet 膜态”，通常指 D0膜量子力学的 Hilbert 空间 $\mathcal{H}$ 中的一个向量。按照刚才讨论的维数约化，这是 $\mathcal{N}=(2,2)$ 超对称量子力学。根据\BookSectionRef{sec:3.1.1}{3.1.1}的论证，在剥离质心等自由方向，并满足适当的紧性与谱条件时，若超对称真空模空间 $\mathcal{M}$ 是非奇异流形，那么，$\mathcal{H}$ 由其 de Rham 上同调给出。

\SourceParagraph{ch05-5.4.3.1-004}{354}
可以像\BookSectionRef{sec:3.1.3}{3.1.3}那样，定义计数 BPS 态的指标；它隐含地是箭图维数向量的函数。在上述光滑情形及相应条件下，这个指标就是 $\mathcal{M}$ 的 Euler 示性数。

\SourceParagraph{ch05-5.4.3.1-005}{354}
通常，$\mathcal{M}$ 会奇异。不过，由于它由显式商\BookEqRef{eq:5.80}{5.80}得到，量子力学为 $\mathcal{H}$ 提供了显式定义。可以预期，这对应于某种已知的上同调理论\BookCite{ref225,ref110}，但据我们所知，这一点尚未得到良好理解。相反，目前使用“虚基本类”的思想，把 Euler 示性数推广到带有奇点与阻碍的模空间，以定义这个指标；见\BookBibText{ref242}{MS1}第27节。

\SourceParagraph{ch05-5.4.3.1-006}{355}
不过，除去\BookSectionRef{sec:5.9.2}{5.9.2}的简短评论，本书其余部分仍采用 BPS 态的“经典”定义。

\subsubsection{从世界面导出 EFT}\label{sec:5.4.3.2}

\SourceParagraph{ch05-5.4.3.2-001}{355}
假定有一个 $N=2$ SCFT $X$，其闭弦 $(c,c)$ 与 $(a,c)$ 代数具有整数 $U(1)_R$ 荷，以及一组不同的边界态 $B_i$，全部为 A 型，或全部为 B 型。那么，$\mathcal{C}$ 是手征算符空间经过 GSO 投影的部分，如\BookSectionRef{sec:3.5.3.2}{3.5.3.2}所述；$W$ 是\BookSectionRef{sec:3.6.4}{3.6.4}定义的圆盘关联函数的生成函数；$\mu$ 则由膜的分次\BookEqRef{eq:5.67}{5.67}确定。下面解释并论证这些陈述。

\SourceParagraph{ch05-5.4.3.2-002}{355}
先从 $B_i$ 保持相同超对称的特殊情形开始，即，它们的相位\BookEqRef{eq:5.67}{5.67}都相等，模2而言。此时，所有开弦手征算符的 $U(1)_R$ 荷都是整数，合并的世界体理论具有超对称真空 $\phi=0$。

\SourceParagraph{ch05-5.4.3.2-003}{355}
于是，在规范代数的复化意义下，有
\begin{equation}\label{eq:5.88}
\begin{aligned}
\operatorname{Lie}(H_{\mathbb C})&=\bigoplus_{i,j}\operatorname{Hom}(B_i,B_j),\\
\mathcal{C}&=\bigoplus_{i,j}\operatorname{Ext}^1(B_i,B_j),
\end{aligned}
\end{equation}
超势为
\begin{equation}\label{eq:5.89}
W[\phi]=\left\langle
\mathrm{e}^{\sum_i\phi^i\int O_i}\right\rangle,
\end{equation}
按\BookEqRef{eq:3.180}{3.180}解释。Fayet--Iliopoulos 项为零，即，在把 $\mathcal{C}$ 的度规近似取为 Euclid 度规时，$\mu$ 是二次齐次函数。

\SourceParagraph{ch05-5.4.3.2-004}{355}
GSO 投影把 $U(1)$ 荷与时空场的类型联系起来：偶荷对应向量场，因而与规范对称有关；奇荷对应标量场。根据\BookEqRef{eq:5.85}{5.85}的超对称版本，$q\gg1$ 的手征算符，不给出无质量场。

\SourceParagraph{ch05-5.4.3.2-005}{355}
\BookStatementNumber{26}
\begin{example}[\texorpdfstring{$\mathbb{C}^3$}{C^3} 中的 D0膜]\label{ex:5.26}
开弦手征环，由\BookSectionRef{sec:3.3.2}{3.3.2}中的自由费米子 $\psi^i$，$1\le i\le3$，生成，它们满足边界条件\BookEqRef{eq:3.156}{3.156}。因此，这是 Grassmann 代数 $\Lambda\mathbb{C}^3$，带有自然的迹。取 $O_i=\psi^i$、$\phi^i=Z^i$，计算 $W$，就重新得到\BookEqRef{eq:5.82}{5.82}。
\end{example}

\subsubsection{超势与阻碍理论}\label{sec:5.4.3.3}

\SourceParagraph{ch05-5.4.3.3-001}{355}
并非 $\operatorname{Ext}^1(B,B)$ 中的每个元素，都对应于对象 $B$ 的真正形变；一般而言，形变可能受阻。

\SourceParagraph{ch05-5.4.3.3-002}{356}
用 $\mathcal{N}=1$ EFT 的语言说，$\mathcal{C}$ 中的非零点，并非每个都是超对称真空；还必须满足方程 $\partial W=0$。因此，三维 Calabi--Yau 上 A、B 膜的阻碍理论，由超势描述。

\SourceParagraph{ch05-5.4.3.3-003}{356}
这一点相当非平凡，一般只在 Calabi--Yau 三维流形中成立。为了理解原因，回顾一下，阻碍对应于第二上同调群，或 Ext 群的元素。例如，$X$ 上全纯丛 $V$ 的两个形变 $\alpha,\beta\in H^1(X,\operatorname{End}V)$，在组合时出现由乘积给出的阻碍类，
\[
\alpha\wedge\beta\in H^2(X,\operatorname{End}V),
\]
并在完整的二阶形变方程中计入相应交叉项。若所有阻碍，都来自单个全纯函数的梯度 $\partial W/\partial\phi$，相关阻碍方程的数目，就不超过 $\dim H^1(X,\operatorname{End}V)$。此外，阻碍与形变之间会有对应。一般而言，两者都不必成立；但在 Calabi--Yau 三维流形上，Serre 对偶给出 $H^2$ 与 $H^1$ 的对偶关系，再结合循环乘积结构，得到上述超势描述。

\SourceParagraph{ch05-5.4.3.3-004}{356}
对于全纯曲线，文献\BookCite{ref098,ref296}给出了这一事实的一般证明。更一般的论证见文献\BookCite{ref022}。

\subsubsection{B 膜的超势}\label{sec:5.4.3.4}

\SourceParagraph{ch05-5.4.3.4-001}{356}
对于三维 Calabi--Yau 流形 $X$ 上的全纯丛，超势有一个用全纯 Chern--Simons 作用量表示的“普遍表达式”。这是联络的泛函，其临界点条件，恰好是 $F^{0,2}=0$。把联络的 $(0,1)$ 部分，写成一个参考全纯联络加上 $(0,1)$ 形式，
\[
\bar\nabla=\bar\nabla_0+A^{0,1},
\]
则有
\begin{equation}\label{eq:5.90}
\begin{aligned}
S_{\mathrm{holo\,CS}}=\int_X\Omega\wedge\operatorname{Tr}\bigl(
&A^{0,1}\wedge\bar\nabla_0A^{0,1}\\
&+\tfrac23 A^{0,1}\wedge A^{0,1}\wedge A^{0,1}\bigr).
\end{aligned}
\end{equation}
可以验证，其变分为 $2\int_X\Omega\wedge\operatorname{Tr}(\delta A^{0,1}\wedge F^{0,2})$，因而临界条件就是 $F^{0,2}=0$。

\SourceParagraph{ch05-5.4.3.4-002}{356}
它本身还不是超势，因为超势应当是线性化形变空间上的函数，使用与 $H^1(X,\operatorname{End}V)$ 对偶的坐标。不过，在一个参数化联络族上计算它，就能得到超势的三次项，
\[
W(t)=S_{\mathrm{holo\,CS}}\!\left[\sum_\alpha t^\alpha A_\alpha^{0,1}\right]+O(t^4).
\]
更高次项，可以使用\BookSectionRef{sec:8.2.1}{8.2.1}讨论的与 $A_\infty$ 代数的关系来获得。

\subsubsection{A 膜的超势}\label{sec:5.4.3.5}

\SourceParagraph{ch05-5.4.3.5-001}{356}
这里，超势是\BookSectionRef{sec:3.6.2.5}{3.6.2.5}中圆盘关联函数的生成函数。预期这些函数再现\BookSectionRef{sec:8.3.3}{8.3.3}所讨论的 Fukaya 范畴的 $A_\infty$ 结构。由于超势是单迹，这要求 $A_\infty$ 结构是循环的。蝌蚪项条件，就是\BookEqRef{eq:5.78}{5.78}。

\subsubsection{膜与反膜}\label{sec:5.4.3.6}

\SourceParagraph{ch05-5.4.3.6-001}{357}
接下来，考虑两组膜的中心荷符号相反的情形，例如，
\[
\xi(B_i)\equiv0,\qquad \xi(\overline B_a)\equiv1\pmod2.
\]
在这个意义上，后一组是“反膜”。

\SourceParagraph{ch05-5.4.3.6-002}{357}
正如\BookSectionRef{sec:3.5.3.2}{3.5.3.2}解释的，在世界面定义中，膜与反膜的唯一区别是 GSO 投影。因此，可以先取膜 $B_i$ 与反膜 $\overline B_a$ 各自 EFT 的直和，再加入新的场，即“膜—反膜开弦”，得到这种情形中的 EFT：
\[
\mathcal{C}'=\bigoplus_{i,a}\bigl(
\operatorname{Hom}(B_i,\overline B_a)
\oplus\operatorname{Hom}(\overline B_a,B_i)\bigr).
\]
换句话说，合并理论的 $\mathcal{C}$，作为向量空间，是各自理论的 $\mathcal{C}$ 的直和，再加上新因子 $\mathcal{C}'$。注意，根据\BookSectionRef{sec:5.4.2.1}{5.4.2.1}的论证，$\operatorname{Ext}^1(B_i,\overline B_a)$ 不给 $H$ 增添新因子，因为这些对应有质量规范玻色子。

\SourceParagraph{ch05-5.4.3.6-003}{357}
新自由度的 $W$ 与 $\mu$ 呢？回顾\BookSectionRef{sec:3.6.2.2}{3.6.2.2}与\BookSectionRef{sec:3.6.4}{3.6.4}中的拓扑关联函数讨论，完全不依赖 GSO 投影，可知，同一个超势表达式\BookEqRef{eq:5.89}{5.89}，仍适用于这种情形，只须把它理解为，保留总 $U(1)_R$ 荷为3的所有项。

\SourceParagraph{ch05-5.4.3.6-004}{357}
特别地，$W$ 没有二次项。这也来自如下物理论证。根据\BookSectionRef{sec:3.3.4.1}{3.3.4.1}，所有 Ramond 基态，都具有相同维数 $h=\hat c/8$。使用\BookEqRef{eq:5.85}{5.85}的相应版本，这意味着，Ramond 基态总是对应于四维中的无质量费米子。另一方面，$\mathcal{N}=1$ EFT 中的费米子质量，恰好来自 $W$ 的二次项。

\SourceParagraph{ch05-5.4.3.6-005}{357}
$\mathcal{C}'$ 中这些场的玻色分量，却不是无质量的。实际上，它们具有质量平方 $m^2<0$，因而是\BookSectionRef{sec:5.4.4}{5.4.4}将解释的“快子”。讨论以后，就会清楚，所有这些内容，也适用于具有一般分次的一般膜集合。

\subsubsection{其他背景中的超对称膜}\label{sec:5.4.3.7}

\SourceParagraph{ch05-5.4.3.7-001}{357}
这些背景包括超 K\"ahler 流形，以及完整群为 $G_2$ 或 $\operatorname{Spin}(7)$ 的流形。一个很好的介绍是文献\BookCite{ref172}。

\SourceParagraph{ch05-5.4.3.7-002}{357}
II 型弦在 $G_2$ 流形上紧化时，Dirichlet 膜由具有两个超荷的三维 EFT 描述；通常称为 $d=3$ 的 $\mathcal{N}=1$ 理论。这与 $d=2$ 的 $(1,1)$ 超对称相当。靶空间是实流形，势来自类似于\BookEqRef{eq:5.76}{5.76}的实超势。已有一些研究，例如，文献\BookCite{ref148,ref228,ref214}在 $G\subset G_2$ 的情形，研究了\BookSectionRef{sec:5.4.5}{5.4.5}的轨形构造。

\SourceParagraph{ch05-5.4.3.7-003}{357}
II 型弦在 $\operatorname{Spin}(7)$ 流形上紧化时，Dirichlet 膜由具有 $(0,1)$ 超对称的二维 EFT 描述。到目前为止，这个主题仍是未知领域，不过，可以计算轨形例子。

\subsection{快子凝聚}\label{sec:5.4.4}

\SourceParagraph{ch05-5.4.4-intro-001}{358}
\BookSectionRef{sec:5.2.1.2}{5.2.1.2}给出了 A 膜束缚态形成与衰变的相当具体的几何图景。从 $\mathcal{N}=1$ 有效场论的观点，也可以使用快子凝聚这一物理思想，为同一个过程提供动机\BookCite{ref423}；另见文献\BookCite{ref363}及其中参考文献。

\subsubsection{膜的组合}\label{sec:5.4.4.1}

\SourceParagraph{ch05-5.4.4.1-001}{358}
先回顾\BookRef{ex:5.22}{例题5.22}，即以两个顶点之间有 $n$ 条箭的箭图为基础的理论。EFT 语言中，这个理论具有规范群 $U(1)^2$、$n$ 个带荷物质多重态，以及按\BookEqRef{eq:5.83}{5.83}进入的 Fayet--Iliopoulos 项，即矩映射参数 $\theta_1,\theta_2$。

\SourceParagraph{ch05-5.4.4.1-002}{358}
按照\BookSectionRef{sec:5.4.2}{5.4.2}与\BookSectionRef{sec:5.4.3}{5.4.3}的思路，这个 EFT 可以作为一对 Dirichlet 膜 $B_1,B_2$ 的世界体理论出现，两个膜由 $n$ 条开弦连接，即 $\dim\operatorname{Ext}^1(B_1,B_2)=n$。因此，可以把它用作超对称 Dirichlet 膜形成束缚态的一般模型。

\SourceParagraph{ch05-5.4.4.1-003}{358}
根据定义，束缚态是单个 Dirichlet 膜，因而具有未破缺规范群 $U(1)$。只要任一物质场 $\psi^i\ne0$，就会如此。由\BookEqRef{eq:5.83}{5.83}，这种真空配置存在，当且仅当 $\theta_2>\theta_1$。

\SourceParagraph{ch05-5.4.4.1-004}{358}
若 $\theta_1=\theta_2$，矩映射条件\BookEqRef{eq:5.83}{5.83}由 $\psi=0$ 满足。这个配置保持 $U(1)^2$ 规范对称，因此，得到两个膜的超对称组合，但它不是束缚态。

\SourceParagraph{ch05-5.4.4.1-005}{358}
若 $\theta_2<\theta_1$，会怎样？为理解这一点，考虑标量势\BookEqRef{eq:5.76}{5.76}。令 $|\psi|^2=\sum_i|\psi^i|^2$，则
\begin{align}
V&=\frac12\bigl(\theta_1-\theta_2+|\psi|^2\bigr)^2,\label{eq:5.91}\\
&=V_0+\frac{M^2}{2}|\psi|^2+\frac12|\psi|^4,\label{eq:5.92}
\end{align}
其中，
\begin{equation}\label{eq:5.93}
M^2=2(\theta_1-\theta_2),
\end{equation}
而 $V_0=\tfrac12(\theta_1-\theta_2)^2$。

\SourceParagraph{ch05-5.4.4.1-006}{358}
由此可知，对于任意 $\theta_2\le\theta_1$，$\psi=0$ 都是真空。区别在于，当 $\theta_2<\theta_1$ 时，真空不是超对称的；例如，真空能量不为零。因此，虽然 $B_1,B_2$ 各自保持超对称，组合在一起却破坏超对称。

\SourceParagraph{ch05-5.4.4.1-007}{358}
通常，破坏超对称的真空，是势的局部极小点，却不是全局极小点。虽然容易验证，当 $\theta_2<\theta_1$ 时，$\psi=0$ 是\BookEqRef{eq:5.91}{5.91}的全局极小点，但在更现实的问题中，不会拥有势的精确表达式。因此，希望更好地理解这种情形。

\SourceParagraph{ch05-5.4.4.1-008}{359}
还应注意，即使 $\theta_2>\theta_1$，配置 $\psi=0$ 仍然是\BookEqRef{eq:5.91}{5.91}的临界点，同样具有非零能量。它仍表示 $B_1,B_2$ 的组合，因此，对于任意 $\theta_1\ne\theta_2$，这种组合都会破坏超对称。需要额外的物理论证，才能更好地理解两种情形的区别。

\subsubsection{快子凝聚}\label{sec:5.4.4.2}

\SourceParagraph{ch05-5.4.4.2-001}{359}
遵循\BookSectionRef{sec:5.4.2}{5.4.2}的物理论证，从 D3 膜 $B_1,B_2$ 出发，不过让它们在 Minkowski 空间方向上相距很大的距离 $\Delta X$。根据局域性，膜之间的相互作用很小；这里，相互作用首次出现在单圈阶，暂且忽略。于是，预期这个组合是稳定配置。在势\BookEqRef{eq:5.91}{5.91}中，这体现为很大的附加项 $(\Delta X)^2|\psi|^2/2$，迫使 $\psi=0$。

\SourceParagraph{ch05-5.4.4.2-002}{359}
然后，把两个膜移到一起。接下来发生什么，取决于\BookEqRef{eq:5.93}{5.93}中 $M^2$ 的符号。若为正，配置 $\psi=0$ 仍是唯一临界点；尽管破坏超对称，这个膜组合在物理上仍然稳定。

\SourceParagraph{ch05-5.4.4.2-003}{359}
若 $M^2<0$，当进入 $|\Delta X|<\sqrt{-M^2}$ 的范围，配置 $\psi=0$ 就从极小点变成极大点或鞍点。线性化运动方程预言，场 $\psi$ 的任意小变化，都会随时间指数增长。现实中，这种涨落无处不在，既来自量子力学，也来自与环境的微弱相互作用。因此，称这种配置为“变得不稳定”。

\SourceParagraph{ch05-5.4.4.2-004}{359}
十分一般地，物理与工程意义上的线性化稳定性分析，就是寻找所有玻色场的质量平方。只要任一玻色场的 $M^2<0$，临界点就在物理意义上不稳定：任意小扰动，都会把系统驱离临界点。

\SourceParagraph{ch05-5.4.4.2-005}{359}
粒子物理中，$M^2<0$ 的场常称为快子\footnote{与\BookSectionRef{sec:3.3.5.4}{3.3.5.4}一样，这一带有科幻意味的名称，有其历史解释。这里，它来自对粒子能量、动量与质量的标准关系 $E^2=\vec P^{\,2}+M^2$ 的应用。这使人朴素地把 $M^2<0$ 解释为，比光速更快的粒子。实际上，合理的物理理论，并不这样描述这类粒子；正如刚才解释的，$M^2<0$ 导致不稳定性。因此，借用“快子”一词，称呼驱动不稳定性的场。}；导致出现 $\psi\ne0$ 配置的动力学，称为\emph{快子凝聚}。

\SourceParagraph{ch05-5.4.4.2-006}{359}
几乎总是，线性化分析只是一种近似；要确定不稳定系统的归宿，需要精确的非线性分析。不过，在大多数现实物理系统中，人们通常不是通过分析运动方程来做到这一点，而是使用下面这个简单论证。假定像通常那样，系统与其他自由度，即“环境”，有微弱相互作用，允许交换能量。那么，熵的极大化倾向于这样的配置：所有可用能量都“辐射”到环境，原来的系统则处于基态。因此，系统最终会达到原来配置附近的某个真空。

\SourceParagraph{ch05-5.4.4.2-007}{360}
可以用由势能函数控制的下降梯度流，有效近似这种动力学，
\[
\frac{\partial\psi^i}{\partial t}
=-G^{ij}\frac{\partial V}{\partial\psi^j},
\]
其中，$G^{ij}$ 是配置空间的逆度规。其他各种物理论证，例如\BookSectionRef{sec:3.2.6.2}{3.2.6.2}的重整化群论证，也导出基本相同的结果。

\SourceParagraph{ch05-5.4.4.2-008}{360}
当前情形的结论是：若 $M^2<0$，快子凝聚导致稳定配置，也就是 $V(\psi)$ 的极小点，由\BookEqRef{eq:5.83}{5.83}给出；它是膜 $B_1,B_2$ 的超对称束缚态。

\subsubsection{特殊 Lagrange 子流形}\label{sec:5.4.4.3}

\SourceParagraph{ch05-5.4.4.3-001}{360}
回顾\BookSectionRef{sec:5.2.1.2}{5.2.1.2}中一对相交 A 膜 $L_1,L_2$ 的讨论。在有效场论中，这由 $H\cong U(1)^2$ 的超对称规范理论描述，由两个膜各自的世界体理论合并而来。稍作推广，取两个膜的相交数为 $+n$，于是\BookEqRef{eq:5.7}{5.7}预言的从 $L_1$ 到 $L_2$ 的开弦，属于
\[
\psi\in\operatorname{Hom}(E_1,E_2)\cong\mathbb{C}^n.
\]
那么，EFT 正是刚才讨论的理论。

\SourceParagraph{ch05-5.4.4.3-002}{360}
因此，\BookSectionRef{sec:5.4.4.1}{5.4.4.1}给出的，是\BookSectionRef{sec:5.2.1.1}{5.2.1.1}中已经在 A 膜上看到的现象的 EFT 描述。回顾，特殊 Lagrange 条件\BookEqRef{eq:5.13}{5.13}依赖于参数 $\xi$。假定有两个特殊 Lagrange 子流形 $L_1,L_2$；若 $\xi(L_1)\ne\xi(L_2)$，它们的并就不是特殊 Lagrange 子流形。这对应于\BookSectionRef{sec:5.4.4.1}{5.4.4.1}的讨论。

\SourceParagraph{ch05-5.4.4.3-003}{360}
此外，把相交点附近的膜，近似为 $\mathbb{C}^m$ 中的特殊 Lagrange 平面 $\Pi^{(\pi,0,\ldots,0)}$ 与 $\Pi^{\vec\phi}$，已经说明，当且仅当
\[
\xi(L_2)<\xi(L_1)
\]
时，才存在接近并集 $L_1\cup L_2$ 的特殊 Lagrange 子流形 $L_1\#L_2$。若对某个实常数 $c>0$，作识别
\begin{equation}\label{eq:5.94}
\xi(L_2)-\xi(L_1)=c(\theta_1-\theta_2),
\end{equation}
这就对应于\BookSectionRef{sec:5.4.4.2}{5.4.4.2}的讨论。因此，可以暂时把矩映射分量 $\theta_i$，识别为分次 $\xi(L_i)$ 的线性函数。

\SourceParagraph{ch05-5.4.4.3-004}{360}
当然，这个论证，并未证明 $\xi$ 与 $\theta$ 之间的关系，确实具有\BookEqRef{eq:5.94}{5.94}中的简单线性形式。虽然计算完整的势很困难，质量公式\BookEqRef{eq:5.93}{5.93}却可以通过世界面弦理论中的直接计算来检验\footnote{这就是把该过程称为快子凝聚，而非仅仅寻找势的极小点的原因。}。可以遵循文献\BookCite{ref042}，或文献\BookCite{ref394}第13.4节，另见\BookBibText{ref242}{MS1}，分析从 $\Pi^{\vec0}$ 出发、在 $\Pi^{\vec\phi}$ 结束的开弦的质量。结果是，R 扇区有始终无质量的弦，而 NS 扇区的相关标量场，具有质量平方，以 $1/\alpha'$ 为单位，
\begin{equation}\label{eq:5.95}
M^2=\frac{1}{2\pi}\left(\sum_{j=1}^m\phi_j-\pi\right).
\end{equation}
若其中一个 D 膜被看作反 D 膜，这些标量场不会被 GSO 过程投影掉。

\SourceParagraph{ch05-5.4.4.3-005}{361}
若考虑鬼数为 $n$ 的开弦，即 Fukaya 范畴中的 $\operatorname{Ext}^n(L_1,L_2)$，则
\begin{equation}\label{eq:5.96}
M^2=\frac12\bigl(\xi(L_2)-\xi(L_1)-n\bigr).
\end{equation}
因此，与上一节比较，在 $\mathcal{M}^+$ 中有 $M^2<0$，因而 $\operatorname{Hom}(L_1,L_2)$ 中的开弦是快子。这完全符合存在基态 $L_1\#L_2$，其能量低于 $L_1\cup L_2$ 这一事实。这个快子凝聚，形成 $L_1\#L_2$。在 $\mathcal{M}^-$ 中，开弦不是快子，不发生凝聚。

\SourceParagraph{ch05-5.4.4.3-006}{361}
因此，快子凝聚图景，为 Joyce 完成、并在\BookSectionRef{sec:5.2.1.2}{5.2.1.2}回顾的困难分析，提供了非常简单的描述。自然会猜想，这个快子图景，给出在复结构模空间中移动时，A 膜如何衰变的完整判据。物理上，这确实有充分动机；不过，要使这样的陈述严格，所需的微分几何很困难。文献\BookCite{ref446,ref442}在这一方向已有进展。

\SourceParagraph{ch05-5.4.4.3-007}{361}
注意，在相应 Floer 指标的相对角约定下，$0\le\phi_j<\pi$。因此，由\BookEqRef{eq:5.73}{5.73}，若 $Y$ 是 Calabi--Yau $m$ 维流形，则
\begin{equation}\label{eq:5.97}
0\le\xi(L_2)-\xi(L_1)+\mu(p)<m.
\end{equation}
这个关系，很好地符合\BookSectionRef{sec:3.3.4}{3.3.4}讨论的，幺正 SCFT 中手征初级场的维数界\BookEqRef{eq:3.97}{3.97}。共形权为 $h$ 的手征初级场的顶点算符，对应于时空玻色场，其质量 $M^2=h-\tfrac12$。与\BookEqRef{eq:5.95}{5.95}及\BookEqRef{eq:5.73}{5.73}比较，可见结果一致。

\SourceParagraph{ch05-5.4.4.3-008}{361}
最后，回顾标准弦理论分析：NS 扇区中，也有对应于未紧化维度内向量粒子的开弦态。其质量为
\begin{equation}\label{eq:5.98}
M^2=\frac12\bigl(\xi(L_2)-\xi(L_1)\bigr).
\end{equation}
因此，当 $L_1=L_2$ 时，它们始终无质量。换句话说，与 $L_1$ 对应的向量，由 $\operatorname{Hom}(L_1,L_1)$ 给出。这些向量对应于 D 膜中的规范群；准确地说，$\operatorname{Hom}(L_1,L_1)$ 是规范代数的复化。单个不可约 D 膜的情形中，预期规范群为 $U(1)$，因而 $\operatorname{Hom}(L_1,L_1)=\mathbb{C}$。若规范群被增大——或因为有两个不同的 D 膜，或因为有重合的 D 膜——规范群、从而 $\operatorname{Hom}(L_1,L_1)$ 也会变大。D 膜的不可约性，等价于 $\operatorname{Hom}(L_1,L_1)=\mathbb{C}$；这一事实可以看作表示论中 Schur 引理的一个版本。

\subsubsection{膜—反膜湮灭}\label{sec:5.4.4.4}

\SourceParagraph{ch05-5.4.4.4-001}{362}
也许会问，为什么先从特殊 Lagrange 子流形开始？最简单的衰变过程，不是膜与自身反膜的湮灭吗？事实上，不能完全用 EFT 理解这一过程；必须引入弦理论的要素。不过，发展这种描述，仍然有用，只须记住它的局限。

\SourceParagraph{ch05-5.4.4.4-002}{362}
基本图景，使用刚才描述的同一个 EFT。通过弦世界面计算，发现膜与其反膜之间的相关玻色开弦，总有
\[
M^2=-\frac12,
\]
因此，总会发生束缚态形成。

\SourceParagraph{ch05-5.4.4.4-003}{362}
另一种检验这一结果的方法，是注意到，对于 $X\cong\mathbb{C}^3$，所考虑的平坦 A 膜 $\Pi_0$ 与其反膜之间的区别，只是空间中的一个整体旋转。考虑一族 $\operatorname{SU}(n)$ 变换，
\[
z_i\longmapsto\mathrm{e}^{\mathrm{i}\pi\alpha}z_i.
\]
它们对分次的作用为
\[
\xi\longmapsto\xi+n\alpha.
\]
若 $\alpha=1$，且 $n$ 为奇数，结果就是反转定向。

\SourceParagraph{ch05-5.4.4.4-004}{362}
随后，追踪这一旋转对开弦计算的影响，可以看到，质量平方为 $M^2=1$ 的 NS 标量开弦场，连续地减小到 $M^2=-\tfrac12$。

\SourceParagraph{ch05-5.4.4.4-005}{362}
虽然这说明，膜—反膜湮灭确实是快子凝聚，但试图用\BookRef{ex:5.22}{例题5.22}中 $n=1$ 的 EFT 描述它，并不真正奏效。问题是，快子凝聚以后，剩下的不是无膜真空，而是描述非平凡 D 膜的 $U(1)$ EFT。这在箭图规范理论中不可避免，因为所有双基本表示场，对对角 $U(1)$ 都不带荷，所以，在这里讨论的 EFT 物理中，无法修复这一点。

\SourceParagraph{ch05-5.4.4.4-006}{362}
这一过程的完整描述，必须涉及弦世界面的考虑，某种奇异 EFT，或弦场论。\BookSectionRef{sec:5.9}{5.9}将简要讨论其中一些思想。不过，通过在 EFT 中再加入一个要素，可以得到这部分物理的一些内容。

\subsubsection{Seiberg 对偶与倾斜}\label{sec:5.4.4.5}

\SourceParagraph{ch05-5.4.4.5-001}{363}
考虑超对称 QCD，即\BookRef{ex:5.24}{例题5.24}，其中，$H\cong U(N_1)\times U(N_2)\times U(N_3)$，且
\begin{equation}\label{eq:5.99}
R=(V_1^\vee\otimes V_2)\oplus(V_2^\vee\otimes V_3),
\end{equation}
$V_i\cong\mathbb{C}^{N_i}$ 是 $U(N_i)$ 的定义表示。分别把这些场，即 $\mathcal{C}$ 上的坐标，记为 $\widetilde Q^i,Q^j$。

\SourceParagraph{ch05-5.4.4.5-002}{363}
根据前面的论证，这是描述三个 Dirichlet 膜 $B_1,B_2,B_3$ 的普遍理论，其相交数满足 $B_1\cap B_2=B_2\cap B_3=+1$，而 $B_1\cap B_3=0$。

\SourceParagraph{ch05-5.4.4.5-003}{363}
超对称真空的模空间依赖于矩映射，后面将系统讨论这种依赖。不过，像\BookRef{ex:5.22}{例题5.22}那样，容易看出，在相应稳定区域 $\theta_1<\theta_2<\theta_3$，可以用一般的复化规范群轨道描述表示。对于除去整体标量以外没有自同构的表示，即规范群破缺到 $U(1)$ 的表示，模空间的预期复维数为\footnote{类似公式适用于任意无关系箭图，不过，稳定性可能更复杂。}
\begin{equation}\label{eq:5.100}
\begin{aligned}
\dim\mathcal{M}&=1+\dim\mathcal{C}-\dim H_{\mathbb C}\\
&=1+N_1N_2+N_2N_3-N_1^2-N_2^2-N_3^2.
\end{aligned}
\end{equation}
多出的1，来自整体 $\mathbb{C}^*\subset H_{\mathbb C}$，即“对角 $U(1)$”的复化，作用平凡。若它非负，则可以证明，它就是 $\mathcal{M}$ 的实际维数。

\SourceParagraph{ch05-5.4.4.5-004}{363}
文献\BookCite{ref414}的物理论证说明，当这一理论作为量子理论有意义时\footnote{如\BookSectionRef{sec:5.4.6.1}{5.4.6.1}将讨论的，这要求 $N_1=N_3\le3N_2$，并且只量子化 $U(N_2)$ 因子。这里的考虑是经典的，不要求这一条件。}，它在红外等价于另一个\emph{对偶}量子理论，箭图如\BookRef{fig:5.6}{图6}所示，规范群为 $H\cong U(N_1)\times U(N_1-N_2)\times U(N_3)$，超势为
\begin{equation}\label{eq:5.101}
W_{\mathrm{dual}}=\operatorname{Tr}(M\widetilde q q).
\end{equation}
\begin{figure}[H]
\centering
\BookImage{01_Numbered_Figures/Chapter_05/Ch05_Fig06_PDFp363_Bookp353.pdf}
\caption{超对称 QCD 的 Seiberg 对偶。}\label{fig:5.6}
\end{figure}

\SourceParagraph{ch05-5.4.4.5-005}{363}
更一般地，对于具有相当一般超势的箭图理论，Seiberg 对偶可以应用于其中的子群。因此，选择一个顶点 $V_i$，把它识别为\BookRef{fig:5.5}{图5}中的中心顶点 $V_2$；该顶点的所有入箭与出箭，分别识别为上述 $\widetilde Q$ 与 $Q$\footnote{这里讨论没有伴随物质的情形，即对应顶点的 $\operatorname{Ext}^1(V_i,V_i)$ 为零。这可以推广，见文献\BookCite{ref317}。}。随后，用\BookRef{fig:5.6}{图6}的子箭图替换它，并取超势
\begin{equation}\label{eq:5.102}
W_{\mathrm{new}}=W_{\mathrm{original}}/\{Q\widetilde Q\mapsto M\}
+W_{\mathrm{dual}}(M,q,\widetilde q),
\end{equation}
其中，操作 $/$ 表示代入替换，$W_{\mathrm{dual}}$ 如\BookEqRef{eq:5.101}{5.101}所示，由此得到对偶理论。

\SourceParagraph{ch05-5.4.4.5-006}{364}
连续应用这些对偶，就得到一族对偶理论。朴素地看，在同一顶点上作两次对偶，会得到与出发点不同的理论。例如，从超对称 QCD 出发，会得到具有额外场 $M_1,M_2$，以及超势 $W=\operatorname{tr}(M_1M_2)$ 的理论。不过，按照\BookSectionRef{sec:5.4.1.5}{5.4.1.5}的定义，把 $M_1,M_2$ 积去以后，它就与原理论等价。

\SourceParagraph{ch05-5.4.4.5-007}{364}
遵循文献\BookCite{ref040}，可以这样理解 Seiberg 对偶：对偶理论描述相同的 Dirichlet 膜配置，但把它们描述为，另一组基本膜“基”的束缚态。特别地，这种经典膜描述的等价，不需要量子场论，在当前讨论的经典层面也成立。

\SourceParagraph{ch05-5.4.4.5-008}{364}
假定把\BookSectionRef{sec:5.4.4.3}{5.4.4.3}的膜组合过程，应用于 $B_1,B_2$，得到束缚态 $B_4$。K 理论关系
\[
[B_1]=[B_4]-[B_2]
\]
提示，可以用 $(B_4,\overline B_2,B_3)$，替换原来的基 $(B_1,B_2,B_3)$；至少对于一部分箭图表示如此，即，那些在两组基下，秩 $N_i$ 都非负的表示。

\SourceParagraph{ch05-5.4.4.5-009}{364}
容易看到，这个变换对应于 K 理论等式
\[
\begin{aligned}
N_1[B_1]+N_2[B_2]+N_3[B_3]
={}&N_1[B_4]+(N_1-N_2)[\overline B_2]\\
&+N_3[B_3],
\end{aligned}
\]
因而重新得到对偶理论的规范群。再多做一点工作，并承认超势\BookEqRef{eq:5.101}{5.101}，就可以验证相应区域中模空间的等价。

\SourceParagraph{ch05-5.4.4.5-010}{364}
在这个意义上，可以在 EFT 中同时处理膜 $B_2$ 与反膜 $\overline B_2$，不过，使用的是描述同一物理的两个不同作用量。

\SourceParagraph{ch05-5.4.4.5-011}{364}
其底层数学，是\BookSectionRef{sec:4.4.7}{4.4.7}的倾斜形式体系。这一物理给出的一般猜想是：倾斜一个带超势的代数，会产生新的带超势代数，超势由显式公式\BookEqRef{eq:5.102}{5.102}给出。文献\BookCite{ref457}最近证明了这一点。

\subsection{轨形奇点处的膜}\label{sec:5.4.5}

\SourceParagraph{ch05-5.4.5-intro-001}{365}
在 D 膜理论中看出箭图如何出现，最初且最直接的物理方法，是考虑轨形背景\BookCite{ref141}。

\SourceParagraph{ch05-5.4.5-intro-002}{365}
设 $X$ 是 $\mathbb{C}^3$ 对有限子群 $G\subset\operatorname{SU}(3)$ 作商得到的轨形。把覆盖空间 $\mathbb{C}^3$ 也记为 $W\cong\mathbb{C}^3$，把 $G$ 的作用记为 $\rho_W$。

\SourceParagraph{ch05-5.4.5-intro-003}{365}
希望导出一组 $N$ 个 D0膜的 EFT，换句话说，每个膜的 $L\subset X$ 是一点。物理上的一种自然方法，是在 $\mathbb{C}^3$ 中分析对应于商映射原像的一组 D 膜，再施加 $G$ 不变性。

\SourceParagraph{ch05-5.4.5-intro-004}{365}
已经知道，$\mathbb{C}^3$ 中的 D 膜，由\BookRef{ex:5.21}{例题5.21}中的四维 $\mathcal{N}=4$ 超对称规范理论描述。它具有各种场：标量 $Z^i$、费米子，以及四维规范联络 $A$。最后可以用 $\mathcal{N}=1$ 超对称推知费米子上的作用。规范场相对于 $\mathbb{C}^3$ 的内部旋转是标量，所以，最适合从它们开始。

\SourceParagraph{ch05-5.4.5-intro-005}{365}
朴素地说，这样的标量在 $G$ 作用下不变。不过，由于 D 膜理论具有规范不变性，应当允许 $G$ 通过规范变换作用。$\mathbb{C}^3$ 可缩，因此，可以把注意力限制在整体规范变换上；它们对联络 $A$ 的作用为
\begin{equation}\label{eq:5.103}
g(A)=\rho_V(g)A\rho_V(g)^{-1},
\end{equation}
其中，$\rho_V$ 是 $G$ 的某个表示。

\SourceParagraph{ch05-5.4.5-intro-006}{365}
换句话说，$A$ 按表示 $\operatorname{End}(V)=V\otimes V^*$ 变换，其 $G$ 不变部分可写为 $\operatorname{Hom}_G(V,V)$。所得规范群，是 $\rho_V$ 的幺正交换子群。把它分解为不可约表示 $\rho_k$，
\[
V=\bigoplus_k\rho_k\otimes\mathbb{C}^{m_k},
\]
由 Schur 引理，规范群就是 $U(m_1)\times U(m_2)\times\cdots$。以下记 $d_k=\dim\rho_k$。

\SourceParagraph{ch05-5.4.5-intro-007}{365}
三个标量场 $Z^i$，参数化法丛；这里，法丛典范地同构于 $\mathbb{C}^3$。因此，$G$ 通过表示 $\rho_W$ 作用于它们，再与规范作用结合。换句话说，标量场的 $G$ 不变子空间，是 $\operatorname{Hom}_G(V,W\otimes V)$。

\SourceParagraph{ch05-5.4.5-intro-008}{365}
回看\BookSectionRef{sec:4.2.6}{4.2.6}，只是重新得到了 $G$ 的 McKay 箭图的定义。规范群的各个因子，是按表示标记的顶点；$Z$ 则是箭图表示中，对应于箭的矩阵。

\SourceParagraph{ch05-5.4.5-intro-009}{365}
最后，箭图关系由超势得到；这个超势，是 $\mathcal{N}=4$ 超势\BookEqRef{eq:5.82}{5.82}在 $G$ 不变子空间上的限制。其临界点，恰好满足\BookSectionRef{sec:4.2.6}{4.2.6}给出的关系。

\SourceParagraph{ch05-5.4.5-intro-010}{365}
因此，得到一族箭图规范理论，由不可约表示 $\rho_k$ 的重数 $m_k$ 参数化；等价地，由各顶点对应的规范群秩参数化。直接对应于 $X$ 上一个 D0膜的情形，是取 $\rho_V$ 为正则表示；$N$ 个普通 D0膜则对应于它的 $N$ 倍。不过，也可以取任意表示；这些理论在物理上都有意义。

\SourceParagraph{ch05-5.4.5-intro-011}{366}
这种情形下，一般的箭图表示 $R$ 的未破缺规范群为 $H_R\cong U(1)^{\sum_km_k}$。按照\BookSectionRef{sec:5.4.2.1}{5.4.2.1}的讨论，可以把每个 $U(1)$ 因子，解释为一个粒子；其他粒子可以由它们形成束缚态。因此，轨形上有一组特定的 D 膜 $F_k$，与 $G$ 的不可约表示 $\rho_k$ 对应。文献\BookCite{ref116}把这些膜称为\emph{分数膜}。

\SourceParagraph{ch05-5.4.5-intro-012}{366}
由有限群表示论，对于一个 D0膜，即正则表示，箭图表示中每个顶点所附的整数 $m_i$，等于对应不可约表示的维数。0膜的位置，由箭图各条箭对应的矩阵决定。研究这些 B 膜的稳定性以后，将看到，这些稳定箭图表示的模空间，就是预期的 $X$。

\SourceParagraph{ch05-5.4.5-intro-013}{366}
还可知，0膜总是由分数膜的非平凡和组成，这也是“分数”之名的来源。后面将看到，在轨形点处，0膜对于衰变为这些底层分数膜，总是边际稳定的。

\subsubsection{其他解释}\label{sec:5.4.5.1}

\SourceParagraph{ch05-5.4.5.1-001}{366}
数学上，可以把同样的箭图，作为 $Q=\mathbb{C}^d$ 上平移不变的 $G$ 等变全纯丛的模空间得到\BookCite{ref409}。设向量丛的纤维由 $G$ 的表示 $V$ 给出。那么，为了满足 $G$ 等变性，这个丛上的联络系数，再次属于 $\operatorname{Hom}_G(V,Q\otimes V)$。把联络1形式写为
\begin{equation}\label{eq:5.104}
\sum_{\mu=1}^d\bigl(Z_\mu\,\mathrm{d}z_\mu
-Z_\mu^\dagger\,\mathrm{d}\bar z_\mu\bigr).
\end{equation}
曲率的 $(2,0)$ 部分为
\begin{equation}\label{eq:5.105}
\sum_{\mu,\nu}\left(-\frac{\partial Z_\mu}{\partial z_\nu}
+\frac12[Z_\mu,Z_\nu]\right)\mathrm{d}z_\mu\wedge\mathrm{d}z_\nu.
\end{equation}
若施加平移不变性，便要求 $Z_\mu$ 的导数为零。对于这一幺正联络，丛全纯的条件，于是等价于 $[Z_\mu,Z_\nu]=0$，再次施加箭图表示的关系。

\SourceParagraph{ch05-5.4.5.1-002}{366}
因此，箭图表示有三种解释：
\begin{enumerate}[label=(\arabic*)]
\item 一个 $G$ 等变层。
\item 轨形上的 D 膜世界体上的标量场。
\item 平移不变的全纯 $G$ 等变向量丛上的联络。
\end{enumerate}

\subsubsection{分数膜的性质}\label{sec:5.4.5.2}

\SourceParagraph{ch05-5.4.5.2-001}{366}
当然，很想立刻像\BookSectionRef{sec:4.7}{4.7}那样，把分数膜识别为凝聚层。不过，应当用物理论证来检验这一点。建立基本图景以后，\BookSectionRef{sec:5.4.5.6}{5.4.5.6}将这样做。

\SourceParagraph{ch05-5.4.5.2-002}{367}
世界体理论中的标量场 $Z_\mu$，必须在世界面描述中作为开弦态出现。也就是说，它们属于 D 膜之间的某些 Ext 群。Ext 箭图的语言立刻告诉我们，这些标量实际上对应于分数膜之间的 $\operatorname{Ext}^1$。因此，可以用顶点 $F_j,F_k$ 之间带符号的箭数，计算相交积\BookEqRef{eq:5.6}{5.6}：
\begin{equation}\label{eq:5.106}
\begin{aligned}
I(F_j,F_k)={}&\dim_{\mathbb C}\operatorname{Hom}_G(\rho_j,W\otimes\rho_k)\\
&-\dim_{\mathbb C}\operatorname{Hom}_G(\rho_k,W\otimes\rho_j).
\end{aligned}
\end{equation}

\SourceParagraph{ch05-5.4.5.2-003}{367}
此外，围绕\BookEqRef{eq:5.96}{5.96}的讨论，以及\BookSectionRef{sec:5.5.4}{5.5.4}，说明，对应于 $\operatorname{Ext}^1$ 的标量态，无质量，当且仅当开弦两端的 $\xi$ 分次相等。因此，已经证明：
\BookStatementNumber{27}
\begin{claim}\label{claim:5.27}
轨形点处，所有分数膜的分次相等。
\end{claim}

\SourceParagraph{ch05-5.4.5.2-004}{367}
这一断言的后果是，轨形点处，分数 D 膜的中心荷必须对齐，具有相同辐角。\BookSectionRef{sec:5.6.3}{5.6.3}将通过例子明确看到这一点。

\SourceParagraph{ch05-5.4.5.2-005}{367}
更一般地，分次对齐，保证一组 D 膜可以纯粹用 EFT 描述，而不遇到\BookSectionRef{sec:5.4.4.4}{5.4.4.4}提到的细节。这可以推广到近似对齐 $|\xi-\xi'|\ll1$，如\BookSectionRef{sec:5.4.4}{5.4.4}解释的那样。

\subsubsection{奇点消解}\label{sec:5.4.5.3}

\SourceParagraph{ch05-5.4.5.3-001}{367}
已经建立，轨形 $\mathbb{C}^d/G$ 的拓扑 B 膜范畴，对应于带关系箭图表示的导出范畴。此外，这种情形中，还知道物理膜的稳定性条件。根据\BookSectionRef{sec:5.4.1}{5.4.1}的基本定义，它们由给定矩映射下，对规范群 $H$ 作辛约化得到。正如\BookSectionRef{sec:5.4.1.3}{5.4.1.3}所提到的，所得模空间，就是\BookSectionRef{sec:4.7.3}{4.7.3}定义的 $\theta$ 稳定配置的集合。因此，这一情形可以得到相当具体的图景，文献\BookCite{ref138,ref137,ref149}研究了它；比较\BookSectionRef{sec:4.7.3}{4.7.3}。

\SourceParagraph{ch05-5.4.5.3-002}{367}
自然地，先从 $\theta=0$ 开始，这时，矩映射条件\BookEqRef{eq:5.79}{5.79}为
\begin{equation}\label{eq:5.107}
\sum_i[Z^i,(Z^i)^\dagger]=0.
\end{equation}
\BookSectionRef{sec:4.7}{4.7}已经看到，$G$ 的正则表示，应当对应于摩天层 $\mathcal{O}_x$。容易说明，与正则表示对应、并满足\BookEqRef{eq:5.107}{5.107}的 $Z$ 的模空间，包含簇 $\mathbb{C}^d/G$\footnote{一般而言，还可以有其他分支，物理上对应于分数膜配置。更完整的分析见文献\BookCite{ref103}。}。因此，这是模参数空间中的“轨形点”，也可以通过显式世界面计算验证\BookCite{ref141}。

\SourceParagraph{ch05-5.4.5.3-003}{368}
现在希望稍微消解轨形。这意味着，修改理论，使与正则表示对应的 $Z$ 的模空间，变为轨形的某个消解 $X$。由于保持复结构，自然猜想，这是通过改变矩映射的分量，即在作用量中加入 Fayet--Iliopoulos 项，来实现的\BookCite{ref141,ref139,ref194}。规范群的每个未破缺 $U(1)$，也就是 $G$ 的每个不可约表示，都允许这样的项。把这些项的系数记为 $\zeta_i$，$i$ 遍历 $G$ 的不可约表示。

\SourceParagraph{ch05-5.4.5.3-004}{368}
结果，\BookEqRef{eq:5.107}{5.107}变为
\begin{equation}\label{eq:5.108}
\sum_\mu[Z_\mu,Z_\mu^\dagger]
=\bigoplus_i\zeta_i\,1_{\mathbb{C}^{d_i}}.
\end{equation}
这里采用正则表示 $m_i=d_i$；因此，每个标量 $\zeta_i$ 在对应对角块中重复 $d_i$ 次。

\SourceParagraph{ch05-5.4.5.3-005}{368}
物理上，预期所得0膜模空间的相应分支，是 $\mathbb{C}^d/G$ 的消解\BookCite{ref141,ref139}；在 $d\le3$、适当的一般参数条件下，这来自文献\BookCite{ref067}。对于 $d=2$ 及任意 $G$，文献\BookCite{ref316}最早证明这一点。对于 $d=3$ 的 Abel 群 $G$，文献\BookCite{ref409}处理了一些例子，文献\BookCite{ref256}给出一般证明。

\SourceParagraph{ch05-5.4.5.3-006}{368}
轨形奇点的爆破，由 $B+\mathrm{i}J$ 的形变得到，并应当局域在奇点附近。因此，消解应当通过加入扭曲闭弦边缘算符产生。这些算符由 $G$ 的非平凡共轭类标记\BookCite{ref125}，记这些类为 $C$。令 $\phi_C$ 为共轭类 $C$ 对应、出现在闭弦拓扑 A 模型中的扭曲算符，考虑在作用量中加入
\begin{equation}\label{eq:5.109}
a_C\int_\Sigma\mathrm{d}^2z\,\phi_C
\end{equation}
所产生的 $B+\mathrm{i}J$ 形变。打开 $\zeta_i$，应当对应于打开某些 $a_C$；这意味着，$\phi_C$ 在 D 膜背景中，应当获得非零单点函数。文献\BookCite{ref141,ref139}计算了这个单点函数。其结果是，至少对于很小的爆破，在线性近似下，
\begin{equation}\label{eq:5.110}
a_C=\sum_i\chi_i(C)\zeta_i,
\end{equation}
其中，$\chi_i$ 是 $G$ 的不可约特征标。这个变换可以逆转，以原来的 $B+\mathrm{i}J$ 参数确定 $\zeta_i$。

\SourceParagraph{ch05-5.4.5.3-007}{368}
注意，恒等共轭类不属于扭曲扇区，因此，不列入这些扭曲形变。于是
\begin{equation}\label{eq:5.111}
0=\sum_i\zeta_i.
\end{equation}

\subsubsection{部分消解}\label{sec:5.4.5.4}

\SourceParagraph{ch05-5.4.5.4-001}{369}
顺便指出，一旦具体定义了三维商奇点附近 D 膜的世界体理论，原则上，就能直接导出一大类三维奇点的世界体理论；这些奇点由轨形奇点的部分消解得到。文献\BookCite{ref368}最早指出这一点，其中及后续文献计算了许多例子。

\SourceParagraph{ch05-5.4.5.4-002}{369}
\BookStatementNumber{28}
\begin{example}[普通双点（ODP），或锥形奇点（conifold）]\label{ex:5.28}
考虑 $\mathbb{C}^4$ 中由
\begin{equation}\label{eq:5.112}
z_1z_2=z_3z_4
\end{equation}
定义的超曲面 $X$。这种奇点，可以用两种方式改为光滑三维流形。显然的一种，是形变方程\BookEqRef{eq:5.112}{5.112}。所得空间将在\BookRef{ex:7.3}{例题7.3}中描述；其拓扑为 $\mathbb{R}^3\times S^3$，所以，$b_2=0$，没有相应 K\"ahler 参数。

\SourceParagraph{ch05-5.4.5.4-003}{369}
还有一种不那么显然的可能性，即所谓小消解。一个简单的构造方法，是引入新变量 $a_1,a_2,b_1,b_2\in\mathbb{C}$，把\BookEqRef{eq:5.112}{5.112}解为
\[
z_1=a_1b_1,\quad z_2=a_2b_2,\quad
z_3=a_2b_1,\quad z_4=a_1b_2.
\]
这给出 $\mathbb{C}^4$ 到 $X$ 的映射；非零点的原像，是 $\mathbb{C}^*$ 作用
\[
a_i\longmapsto\lambda a_i,\qquad
b_i\longmapsto\lambda^{-1}b_i
\]
的轨道，原点上则须取闭轨道商。因此，
\begin{equation}\label{eq:5.113}
X\cong\mathbb{C}^4\mathbin{//}H,
\end{equation}
其中，$H=U(1)$，使用齐次矩映射。现在，通过在这个理论中打开“Fayet--Iliopoulos 项” $\zeta>0$，也就是改变矩映射，就得到小消解。这个空间的拓扑为 $\mathbb{R}^4\times S^2$；容易说明，$\zeta=C\int_{S^2}\omega$，其中，$C$ 是某个实归一化常数。
\begin{figure}[H]
\centering
\BookImage{01_Numbered_Figures/Chapter_05/Ch05_Fig07_PDFp369_Bookp359.pdf}
\caption{锥形奇点的箭图规范理论。}\label{fig:5.7}
\end{figure}

\SourceParagraph{ch05-5.4.5.4-004}{369}
顺便注意，取 $\zeta<0$ 也得到相同的局部拓扑，但这时 $\zeta=-C\int_{S^2}\omega$。不过，若把这个非紧几何识别为紧三维流形的一个子区域，两侧的整体几何就可以不同，并由保典翻转相联系。\BookSectionRef{sec:5.5.7}{5.5.7}将再讨论这一点。

\SourceParagraph{ch05-5.4.5.4-005}{369}
为了把 $U(1)$ 商\BookEqRef{eq:5.113}{5.113}解释为箭图规范理论，引入第二个作用平凡的 $U(1)$，得到\BookRef{fig:5.7}{图7}的箭图。这很好地描述消解几何中的单个 D0膜。不过，还没有结束，因为还希望用具有 $U(N)$ 规范群的类似理论，描述一组 $N$ 个 D0膜。

\SourceParagraph{ch05-5.4.5.4-006}{370}
为此，还必须指定关系，它们来自超势\BookCite{ref305}，
\begin{equation}\label{eq:5.114}
W=\sum_{1\le i,j,k,l\le2}\epsilon^{ij}\epsilon^{kl}
\operatorname{Tr}(A_iB_kA_jB_l),
\end{equation}
其中，$\epsilon^{ij}$ 是反对称不变张量。

\SourceParagraph{ch05-5.4.5.4-007}{370}
为了系统地给出依据，可以通过部分消解得到这个奇点，例如，从轨形 $\mathbb{C}^3/(\mathbb{Z}_2\times\mathbb{Z}_2)$ 出发。计算消解后的模空间，会发现某些场处处非零；把它们代入轨形超势，即\BookEqRef{eq:5.82}{5.82}的 $G$ 不变投影以后，会使其他场具有质量，类似于\BookSectionRef{sec:5.4.4.5}{5.4.4.5}的讨论。解出这些有质量场，就重新得到\BookEqRef{eq:5.114}{5.114}；更多细节见文献\BookCite{ref368}。
\end{example}

\subsubsection{\texorpdfstring{$\theta$}{theta} 稳定性}\label{sec:5.4.5.5}

\SourceParagraph{ch05-5.4.5.5-001}{370}
回到 $\mathbb{C}^3/G$，考虑一般表示 $V=\bigoplus_i\rho_i\otimes\mathbb{C}^{m_i}$。\BookEqRef{eq:5.108}{5.108}右边，变成每个 $\zeta_i$ 在对应等型块中重复 $m_i$ 次的对角矩阵。取方程的迹，可知，超对称真空存在，要求
\begin{equation}\label{eq:5.115}
\sum_i d_i\zeta_i=0.
\end{equation}
一般而言，这与正则表示的条件\BookEqRef{eq:5.111}{5.111}不相容；此时，无法解出 D 项条件\BookEqRef{eq:5.79}{5.79}。

\SourceParagraph{ch05-5.4.5.5-002}{370}
乍看似乎破坏了超对称。不过，更仔细地观察世界体理论，会发现还存在一个与对角 $U(1)$ 对应的第二超对称，即其规范费米子的常数移位。这允许 $\mu$ 非零、但与整体常向量 $e=(1,1,\ldots,1)$ 成比例的解，保持两个超对称的一个线性组合，因此，这些仍是 BPS 态\BookCite{ref116,ref138}。

\SourceParagraph{ch05-5.4.5.5-003}{370}
考虑对于某些实数 $\theta_i$，满足
\begin{equation}\label{eq:5.116}
\sum_\mu[Z_\mu,Z_\mu^\dagger]
=\bigoplus_i\theta_i\,1_{\mathbb{C}^{m_i}}
\end{equation}
的配置。采用相应迹范数，标量势 $\|\mu\|^2$ 的这部分为
\begin{equation}\label{eq:5.117}
\sum_i(\zeta_i-\theta_i)^2.
\end{equation}
在迹条件 $\sum_i d_i\theta_i=0$ 下，它由
\begin{equation}\label{eq:5.118}
\theta_i=\zeta_i-
\frac{\sum_jm_j\zeta_j}{\sum_jm_j}
\end{equation}
达到极小。这把\BookEqRef{eq:5.108}{5.108}解到相差一个整体常数移位的程度，因而给出 BPS 态。

\SourceParagraph{ch05-5.4.5.5-004}{371}
方程\BookEqRef{eq:5.116}{5.116}可以写成更适合箭图的形式。令 $a$ 是箭图的一条箭，箭头为 $h(a)$，箭尾为 $t(a)$；在给定箭图表示中，对应矩阵 $Z_a$ 的大小为 $m_{h(a)}\times m_{t(a)}$。则方程成为
\begin{equation}\label{eq:5.119}
\sum_{h(a)=i}Z_aZ_a^\dagger
-\sum_{t(a)=i}Z_a^\dagger Z_a=\theta_i\,\operatorname{id}.
\end{equation}

\SourceParagraph{ch05-5.4.5.5-005}{371}
这恰好是 King\BookCite{ref304}研究、并在\BookRef{thm:4.88}{定理4.88}中讨论的方程。回顾那里的 $\theta$ 稳定性定义。固定与 $G$ 表示 $V=\bigoplus_i\rho_i\otimes\mathbb{C}^{m_i}$ 对应的箭图表示。对 $G$ 的任意表示 $W=\bigoplus_i\rho_i\otimes\mathbb{C}^{n_i}$，定义稳定性权
\begin{equation}\label{eq:5.120}
\theta(W)=\sum_i\theta_i n_i.
\end{equation}
于是，由\BookEqRef{eq:5.116}{5.116}的迹为零，可知 $\theta(V)=0$。若对于与 $G$ 表示 $W$ 对应的任意非零真箭图子表示，都有 $\theta(W)>0$，就称箭图表示为 $\theta$ 稳定的。

\SourceParagraph{ch05-5.4.5.5-006}{371}
在\BookRef{thm:4.88}{定理4.88}的条件下，箭图表示满足\BookEqRef{eq:5.119}{5.119}，其内积除去显然的自同构以外唯一，当且仅当，它是 $\theta$ 稳定表示的直和。

\SourceParagraph{ch05-5.4.5.5-007}{371}
因此，在十分接近轨形点的地方，得到一个纯粹用箭图表达的稳定性条件。随着 K\"ahler 模参数远离这个点，最终预期会看到\BookSectionRef{sec:5.4.4.4}{5.4.4.4}与\BookSectionRef{sec:5.4.4.5}{5.4.4.5}讨论的更非平凡的膜—反膜现象。

\subsubsection{与 McKay 对应的关系}\label{sec:5.4.5.6}

\SourceParagraph{ch05-5.4.5.6-001}{371}
最后，为把分数膜识别为 McKay 对应的基本对象，给出依据。为此，在理论中加入 D6膜。需要选择 $G$ 作用；令 $V_k$ 为对应不可约表示 $\rho_k$ 的“分数”D6膜。合并理论具有 D0--D6 开弦；可以通过世界面论证验证，它们包括在 $\mathbb{C}^3$ 的内部旋转下作为标量变换的无质量手征费米子，以及有质量标量，这也来自\BookEqRef{eq:5.96}{5.96}与\BookEqRef{eq:5.142}{5.142}。费米子上类似于\BookEqRef{eq:5.103}{5.103}的投影，为
\begin{equation}\label{eq:5.121}
\psi=\rho_k(g)\psi\rho_V(g)^{-1}.
\end{equation}
$V$ 中每出现一次不可约表示 $\rho_k$，就给出一个独立解。

\SourceParagraph{ch05-5.4.5.6-002}{371}
由\BookEqRef{eq:5.6}{5.6}，这些无质量费米子的指标，是 D0膜与 D6膜之间的 K 理论相交数。因此，分数 D0膜与分数 D6膜，构成 K 理论的对偶基；准确地说，分数 D0膜提供紧支撑 K 理论的基。事实上，还有更强的陈述，
\begin{equation}\label{eq:5.122}
\dim\operatorname{Ext}^*(V_k,F_i)=\delta_{i,k}.
\end{equation}

\SourceParagraph{ch05-5.4.5.6-003}{372}
\BookEqRef{eq:5.121}{5.121}也给出 $V_k$ 的几何解释：它是 $G$ 等变向量丛 $\mathcal{O}_{\mathbb{C}^3}\otimes\rho_k$，秩为 $d_k$，在一维不可约表示的情形中是线丛。物理术语中，它具有离散 Wilson 线。其来源是，渐近区域，例如 $|z|>R$，在轨形基本群的意义下保留群 $G$；若群在该区域自由作用，普通基本群也为 $G$。表示 $\rho_k(g)$ 可以解释为平行移动表示 $\pi_1\to H$。

\SourceParagraph{ch05-5.4.5.6-004}{372}
若形变 K\"ahler 模参数并消解轨形，每个 $G$ 等变丛 $V_k$，就对应于消解上的某个层 $\widetilde V_k$。具体对应，依赖于所取的消解；一旦确定，对偶关系\BookEqRef{eq:5.122}{5.122}与\BookSectionRef{sec:4.7}{4.7}的数学，就唯一确定分数膜的识别。

\SourceParagraph{ch05-5.4.5.6-005}{372}
为了找到这一对应，可以考虑一个正则 D0膜，因为其非零 FI 项下的模空间，现在给出具体消解。D0--D6 开弦，是从 $V_k$ 的截面到某一点，即 D0膜位置的取值映射。追踪定义，就重新得到\BookEqRef{eq:4.12}{4.12}中的重言丛 $\mathcal{R}_\rho$ 的纤维定义，从而完全与数学联系起来。采用第4章的协变 Fourier--Mukai 约定，取值映射空间是 $\widetilde V_k^\vee$ 的纤维，因此，对应丛为 $\widetilde V_k=\mathcal{R}_{\rho_k}^\vee$。

\SourceParagraph{ch05-5.4.5.6-006}{372}
\BookSectionRef{sec:4.7.3}{4.7.3}的结果，现在给出 FI 项空间每个腔室中的分数膜识别；$\mathbb{C}^3/\mathbb{Z}_3$ 的例子如\BookRef{tab:4.3}{表3}所示。这还不是完整答案，因为还需要把这一识别延拓到轨形描述的有效范围以外，例如，大体积极限。可以通过考虑单值化做到这一点，\BookSectionRef{sec:5.6}{5.6}将说明。

\subsection{其他物理结果简述}\label{sec:5.4.6}

\SourceParagraph{ch05-5.4.6-intro-001}{372}
关于 Dirichlet 膜的 $\mathcal{N}=1$ EFT，已有大量研究，因为它们提供一条相当直接的路径，从超弦紧化通向具有直接物理意义的理论，例如量子色动力学与最小超对称标准模型。这里无法深入讨论，下面只尝试给出一个简短而不完整的指南，并指出一些文献。

\subsubsection{关于量子化的评论}\label{sec:5.4.6.1}

\SourceParagraph{ch05-5.4.6.1-001}{372}
从超弦紧化得到 $n+1$ 维时空中的量子场论，一种标准方法，当然绝非唯一方法，是考虑一组 D$(n+p)$ 膜，使其世界体的 $n+1$ 个维度，嵌入同一个超平面，
\[
\mathbb{R}^{n,1}\times L_i\subset\mathbb{R}^{3,1}\times X,\qquad n\le3.
\]
若进一步假定，闭链 $L_i$ 的度规尺度，相对于 $\mathbb{R}^{n,1}$ 中感兴趣的长度尺度很小，就可以先按\BookSectionRef{sec:5.4.2.1}{5.4.2.1}的方式，约化为 $(n+1)$ 维场论，再进行量子化，得到很好的描述。

\SourceParagraph{ch05-5.4.6.1-002}{372}
通常，所得量子场论只依赖于这些闭链附近的局部几何。若它只覆盖小区域 $U\subset X$，对于低能问题，就可以把弦理论看作有效地“紧化”在非紧三维流形上，即 $\mathbb{R}^{3,1}\times U$。这正是物理文献如此关注非紧情形，以及奇点消解与形变的主要原因。

\SourceParagraph{ch05-5.4.6.1-003}{373}
量子化的讨论，很依赖于 $n$。\BookSectionRef{sec:3.1.1}{3.1.1}与\BookSectionRef{sec:5.4.2.1}{5.4.2.1}已经讨论了 $n=0$，\BookSectionRef{sec:3.2}{3.2}讨论了 $n=1$。这里集中在 $n=3$，不过，$n=2$ 也有很多研究，文献\BookCite{ref254,ref005}可以作为很好的起点。

\SourceParagraph{ch05-5.4.6.1-004}{373}
也许还会问 $n>3$ 的情形。考虑 $\dim_{\mathbb R}X<6$ 的弦紧化中的膜，例如，取 $X\cong T^4$ 或 K3，就可以得到 $n>3$ 的量子场论。长期以来，人们认为，这种量子理论只能是无相互作用的理论，因为它们在微扰理论中不可重整化。不过，近些年来，基于弦理论，已经给出相当有说服力的论证，表明在 $4+1$ 与 $5+1$ 维中，存在少数相互作用超对称理论。它们没有经典极限，相关研究仍处于初期。当前背景下，也许最有趣的例子，来自把 M 理论紧化在 $\mathbb{R}^{4,1}\times X$ 上，其中，$X$ 是 Calabi--Yau 三维流形，而非来自 Dirichlet 膜。它们与奇点消解的关系，见文献\BookCite{ref252}。

\subsubsection{\texorpdfstring{$3+1$}{3+1} 维}\label{sec:5.4.6.2}

\SourceParagraph{ch05-5.4.6.2-001}{373}
现在固定 $n=3$，概述一些一般结果。一个很好的介绍，是文献\BookCite{ref253}。

\SourceParagraph{ch05-5.4.6.2-002}{373}
首先，偶数时空维度中，带费米子的规范理论，一般会有妨碍量子化的反常\BookCite{ref224}。在 $3+1$ 维中，这些反常由 $H$ 的对称不变量 $d\colon(\operatorname{ad}^*H)^3\to\mathbb{R}$ 控制；$\operatorname{SU}(n)$（$n\ge3$）具有这种不变量。用显式矩阵表示 $t_R\colon\operatorname{ad}H\to\operatorname{Mat}(\mathbb{C})$ 表述，反常消除要求
\[
\operatorname{Tr}t_R=0,\qquad\operatorname{Tr}t_R^3=0.
\]
对于实表示 $R$，这自动满足；其他情形中，则非常严格。

\SourceParagraph{ch05-5.4.6.2-003}{373}
考虑\BookSectionRef{sec:5.4.1.3}{5.4.1.3}的例子，\BookRef{ex:5.22}{例题5.22}具有手征规范反常；超对称 QCD，即\BookRef{ex:5.24}{例题5.24}，在 $N_1=N_3$，并且只量子化规范群 $U(N_2)$ 时，是自洽的量子场论。其余群不纳入路径积分，而作为整体对称处理。对于 $\mathbb{C}^3$ 的 McKay 箭图，通常，无反常理论对应于正则表示的倍数。

\SourceParagraph{ch05-5.4.6.2-004}{373}
其次，可以要求理论可重整化。虽然这不是严格必需的，但若不满足，理论的物理，就会依赖于\BookRef{def:5.17}{定义5.17}中指定的数据 $(H,R,\mu,W)$ 以外的更多信息。可重整化性要求，$\mathcal{C}$ 是复 Euclid 空间，$H$ 线性地作用，而 $W$ 是 $\mathcal{C}$ 的 Euclid 坐标中，次数至多为3的非齐次多项式。

\SourceParagraph{ch05-5.4.6.2-005}{374}
接下来，决定理论物理的最重要因素之一，是 RG 与 beta 函数，沿着\BookSectionRef{sec:3.2.5}{3.2.5}讨论的思路。回顾，对于不相关算符的相应耦合，在所用能标约定下，若 beta 函数为正，那么，它在长距离会流向零。

\subsubsection{超势的不重整化}\label{sec:5.4.6.3}

\SourceParagraph{ch05-5.4.6.3-001}{374}
关于超对称规范理论 beta 函数，最重要的结果是：在全纯 Wilson 有效作用量的约定下，超势 $W$ 的直接微扰重整化为零\footnote{这里指，在该全纯约定下，$W$ 中每个单项式的系数，都没有直接的微扰顶点重整化。}，在微扰理论的所有阶都如此。这就是文献\BookCite{ref170}著名的“不重整化定理”。

\SourceParagraph{ch05-5.4.6.3-002}{374}
不重整化定理，仍然允许超势受到非微扰量子修正；这些修正在规范耦合的小耦合极限中，呈指数微小。指数微小的修正，物理上极其有趣，因为它提供一种自然方法，产生非常小的比值，例如，观测中粒子间引力强度，与其他作用强度之间约 $10^{-40}$ 的比值。

\SourceParagraph{ch05-5.4.6.3-003}{374}
对于这里的形式讨论，不重整化定理最主要的意义，是允许许多这类量子理论，拥有非平凡的真空模空间；条件是，指数微小的修正也消失，或受到适当约束。可以用“非微扰不重整化定理”来说明这种情形，见文献\BookCite{ref253}。

\subsubsection{规范耦合的 beta 函数}\label{sec:5.4.6.4}

\SourceParagraph{ch05-5.4.6.4-001}{374}
这些 beta 函数（近似地，且相差一个整体乘法因子）为
\begin{equation}\label{eq:5.123}
\beta=\Lambda\frac{\partial}{\partial\Lambda}\frac{1}{g_{\mathrm{YM}}^2}
=C_2(R)-3C_2(\mathrm{adjoint}),
\end{equation}
这里，$C_2(R)$ 是表示的“第二 Casimir”，即 $\operatorname{Lie}H$ 泛包络代数中的二次不变量。这个结果对 $\mathcal{N}=2$ 与 $\mathcal{N}=4$ 是精确的；$\mathcal{N}=1$ 的精确结果，见文献\BookCite{ref253}及其中参考文献。

\SourceParagraph{ch05-5.4.6.4-002}{374}
例如，对于 $\mathcal{N}=4$ 理论，使用 $R\cong(\operatorname{ad}H\otimes\mathbb{C})^{\oplus3}$，就得到 $C_2(R)=3C_2(\mathrm{adj})$，因而 $\beta=0$。

\subsubsection{超共形场论}\label{sec:5.4.6.5}

\SourceParagraph{ch05-5.4.6.5-001}{374}
与\BookChapterRef{3}一样，所有 beta 函数为零的情形，特别有趣，因为这时，量子场论具有尺度不变性与共形不变性。刚才看到，$\mathcal{N}=4$ 理论如此；$\mathcal{N}=1,2$ 中，某些 $H,R$ 的选择，也如此。这包括具有适当耦合的 $\mathbb{C}^3/\Gamma$ McKay 箭图，取正则表示的倍数。

\SourceParagraph{ch05-5.4.6.5-002}{374}
一般而言，证明超共形不变性时，也会用到超势。一项必要条件，来自四维 SCA 包含 $U(1)_R$ 子代数这一事实；它类似于\BookSectionRef{sec:3.3.3}{3.3.3}的 $U(1)_R$，后者在\BookChapterRef{3}中起了核心作用。这个 $U(1)_R$ 在超共形点附近作用于 $\mathcal{C}$，因而作用于超势；可以说明，$W$ 必须在这一作用下具有齐次权2。$\mathcal{C}$ 的坐标可看作手征场，因此，得到\BookSectionRef{sec:3.3.4}{3.3.4}中手征代数与 $U(1)_R$ 荷界的直接类似物。另一方面，虽然可以使用 $U(1)_R$ 扭曲这些理论，在这里考虑的完整群约化下，只得到二维中已经见过的一部分拓扑理论，因为四维中的 $U(1)$ 完整群，要求 $M^4$ 局部为 $\Sigma\times M^2$，其中，$M^2$ 平坦。

\SourceParagraph{ch05-5.4.6.5-003}{375}
前面讨论的其他一些 $d=2$、$N=2$ 概念，在这种情形中也有类似物。例如，可以通过 $U(1)$ 流与能动张量的二点函数，定义某种类似于 Virasoro 中心荷的量。它们与极值原则有关，特别是确定超共形 R 对称的 $a$ 极大化原则\BookCite{ref251}。这在一定程度上，使这些理论的构造与 RG 流系统化。

\subsubsection{Seiberg 对偶与级联}\label{sec:5.4.6.6}

\SourceParagraph{ch05-5.4.6.6-001}{375}
\BookEqRef{eq:5.123}{5.123}的解，有一个值得注意的特征：在近似流中，$1/g_{\mathrm{YM}}^2$ 可以到达零。这对应于“无穷耦合”，起初似乎很成问题。

\SourceParagraph{ch05-5.4.6.6-002}{375}
不过，可以定义穿过这一表观奇点的“RG 流”或“跃迁”。结果是不同的 EFT，但其 $(H,R,\mu,W)$，与原理论的数据有简单关系。在新理论中继续流动，就能找到一串，或一个“级联”的跃迁。

\SourceParagraph{ch05-5.4.6.6-003}{375}
由于新理论对应于同一组 Dirichlet 膜，箭图表示的模空间，至少对于一大部分维数向量，必须相同。因此，可以预期，事实上也确实如此，二者的关系是\BookSectionRef{sec:5.4.4.5}{5.4.4.5}讨论的 Seiberg 对偶。

\SourceParagraph{ch05-5.4.6.6-004}{375}
对偶级联的概述，见文献\BookCite{ref217}，其中有来自 $\mathbb{C}^3/\mathbb{Z}_3$ 模型的显式例子。

\subsubsection{规范—引力对偶}\label{sec:5.4.6.7}

\SourceParagraph{ch05-5.4.6.7-001}{375}
物理提供的最深远联系，是前面讨论的场论，与其“引力”或弦理论对偶之间的联系。它们来自比较同一组 D 膜的场论描述，与其作为弦理论解的原始描述。若能独立描述后者，就可以提出“对偶”，即某些场论问题与对偶弦理论问题之间的精确对应。

\SourceParagraph{ch05-5.4.6.7-002}{375}
一般而言，对偶是关于完整量子场论的陈述。因此，需要满足上述自洽条件，即反常与可重整性。它把理论中的强耦合问题，联系到对偶描述中的弱耦合问题，即经典或半经典问题。

\SourceParagraph{ch05-5.4.6.7-003}{375}
这一主题有许多变体。一种是，把 Dirichlet 膜重新解释为对偶弦理论中的膜。例如，可以用这种方式“导出” $\mathcal{N}=2$ SYM 理论的 Seiberg--Witten 解\BookCite{ref470}。另一种是，建立不同维数的膜的相交配置，具有多个世界体解释；这导出单极子模空间的 Nahm 构造的重新推导\BookCite{ref113}、更多 $\mathcal{N}=2$ 理论的解\BookCite{ref218}，等等。

\SourceParagraph{ch05-5.4.6.7-004}{376}
也许研究得最充分的对偶，是 AdS/CFT 对应\BookCite{ref343,ref004}。考虑规范群为 $U(N)$ 的 $\mathcal{N}=4$ SYM，即\BookRef{ex:5.21}{例题5.21}。如\BookSectionRef{sec:5.4.2}{5.4.2}所论证，它是 $\mathbb{C}^3$ 中 $N$ 个点处 D3膜的世界体理论。规范理论的真空模空间，必须是这些点配置的空间，即 $(\mathbb{C}^3)^N/S_N$；这个简单观察，就是一个基本例子。

\SourceParagraph{ch05-5.4.6.7-005}{376}
更仔细地考察与这些 Dirichlet 膜对应的超引力解。它包含十维度规 $g_{\mu\nu}$，以及超引力的其他场。可以显式找到它；度规为
\[
\begin{aligned}
\mathrm{d}s^2={}&H(r)^{-1/2}(-\mathrm{d}t^2+\mathrm{d}\vec x^{\,2})\\
&+H(r)^{1/2}(\mathrm{d}r^2+r^2\mathrm{d}\Omega_{S^5}^2),\\
H(r)={}&1+\frac{L^4}{r^4},\qquad L^4=4\pi g_sN\alpha'^2.
\end{aligned}
\]
这里，$(t,\vec x)$ 是 $\mathbb{R}^{3,1}$ 上的坐标，把 $\mathbb{C}^3$ 写为 $\{0\}\cup\mathbb{R}_{>0}\times S^5$，$r\in\mathbb{R}_{>0}$，$\mathrm{d}\Omega_{S^5}^2$ 是 $S^5$ 上的圆度规。常数因子与弦尺度，也可以像通常的弦单位写法那样吸收到 $r$ 中。

\SourceParagraph{ch05-5.4.6.7-006}{376}
现在考虑极限 $r\to0$，这时，系数中的“1”可以略去。极限度规，分解为两个齐性度规的直积：一个在 $S^5$ 上，另一个在所谓反 de Sitter 空间上，后者是双曲空间的解析延拓。

\SourceParagraph{ch05-5.4.6.7-007}{376}
虽然这个具体例子，与镜像对称没有直接联系，也可以对轨形奇点处的 D3膜作同样处理，得到\BookSectionRef{sec:5.4.5}{5.4.5}讨论的理论的对偶描述，即 $\operatorname{AdS}_5\times S^5/G$ 上的超引力。文献\BookCite{ref278,ref368}最早研究了这一点，并启发了大量物理工作。数学上，这种方法，与空间 $S^5/G$ 上的度规几何联系起来。例如，人们用这种方式，得到并研究了一大族新的显式 Sasaki--Einstein 度规\BookCite{ref350,ref351}。

\SourceParagraph{ch05-5.4.6.7-008}{376}
也提一下文献\BookCite{ref151}，它是一项代表性工作，把前面提到的许多发展，与其他内容，例如与二聚体模型的关系，联系在一起。

\section{稳定性}\label{sec:5.5}

\SourceParagraph{ch05-5.5-intro-001}{376}
现在开始讨论，如何把“物理”D 膜，即 $(2,2)$ 超共形场论中的相应边界条件，识别为扭曲所得 A、B 模型边界条件中的一个子集。\BookSectionRef{sec:5.2}{5.2}已经看到，在大体积极限中，这个子集随“另一组”模参数改变：A 模型中是复结构参数，B 模型中是 K\"ahler 参数。因此，对于完整 SCFT，显然也应当预期如此。在 B 模型的大体积极限中，这个子集是 $\mu$ 稳定凝聚层；而特殊 Lagrange 子流形，似乎是在某种尚待精确定义的意义下，A 模型中 Lagrange 子流形的同痕类的稳定子集。

\SourceParagraph{ch05-5.5-intro-002}{377}
此外，承认镜像对称——知道它在 $(2,2)$ SCFT 中成立——就可以把一个 $\sigma$ 模型中高度“弦性”的区域，联系到镜像 $\sigma$ 模型的大体积极限。因此，物理膜与稳定性的关系，必须在大体积极限以外也有意义。这使下面的看法很可信：一般 $(2,2)$ SCFT 中的物理 D 膜，不必假定存在大体积极限，也由某种稳定性条件确定。

\SourceParagraph{ch05-5.5-intro-003}{377}
原则上，这一稳定性条件，可能依赖于各种 SCFT 数据。不过，最简单的假设，是它只依赖于拓扑扭曲模型的数据，但现在同时依赖 A、B 两个模型，回顾本章引言的\BookRef{tab:5.1}{表1}。这些数据包括范畴结构，它概括拓扑膜与开弦关联函数；使用 B 模型语言，称它为导出范畴 $D(\operatorname{Coh}X)$ 的选择。数据还包括 $X$ 上的 A 模型，不过，只需要\BookSectionRef{sec:5.3.4}{5.3.4}讨论中出现的那部分，即决定圆盘单点函数的拓扑闭弦理论。

\SourceParagraph{ch05-5.5-intro-004}{377}
这一思想的第一个检验，是大体积极限：必须看到，$\mu$ 稳定性，即\BookRef{def:5.7}{定义5.7}，的各个要素，如何从 SCFT 数据中出现。对于斜率 $\mu$，这比较容易，\BookSectionRef{sec:5.5.5}{5.5.5}将说明。更难理解的是，如何决定一个对象何时是另一个对象的子对象，因为这一概念，并不直接推广到导出范畴。因此，先从这一点开始。

\subsection{三角形}\label{sec:5.5.1}

\SourceParagraph{ch05-5.5.1-001}{377}
与 A 模型一样，B 模型本身并不知道稳定性。不过，确实要求 B 模型给出某种判据，判断导出范畴中的给定对象，是否有可能衰变为另外两个对象。

\SourceParagraph{ch05-5.5.1-002}{377}
\BookSectionRef{sec:5.4.4}{5.4.4}通过快子凝聚讨论 A 膜衰变时，已经说明，在边际稳定性墙上，开弦无质量。因此，它在共形场论中表现为边缘算符，但未必是严格边缘算符。在这一意义上，束缚过程看起来像形变。因此，\BookSectionRef{sec:4.4.3}{4.4.3}的映射锥，定义了两个 D 膜的一个可能束缚态。

\SourceParagraph{ch05-5.5.1-003}{377}
导出范畴中的映射锥构造，赋予范畴三角结构。这个数学结构，对 D 膜稳定性的概念起核心作用。

\SourceParagraph{ch05-5.5.1-004}{377}
回顾\BookSectionRef{sec:4.4.3}{4.4.3}中三角范畴的定义。对于任意 Abel 范畴 $\mathcal{C}$，导出范畴 $D(\mathcal{C})$ 是三角范畴。平移函子就是移位函子，正合三角形则由映射锥给出。起初，映射锥的定义，看起来可能不如三角形对称，但实际上并非如此：只要把“$[1]$”移到适当的边，并计入所需移位，正合三角形的任意顶点，都可表示为对边的映射锥。

\SourceParagraph{ch05-5.5.1-005}{378}
\begin{example*}
从\BookRef{ex:4.52}{例题4.52}开始。回顾，若 $B$ 是 $\mathcal{C}$ 中 $C$ 被 $A$ 的扩张，那么，短正合列
\begin{equation}\label{eq:5.124}
0\longrightarrow A\xrightarrow{a}B\xrightarrow{b}C\longrightarrow0
\end{equation}
在 $D(\mathcal{C})$ 中诱导正合三角形，
\begin{equation}\label{eq:5.125}
\BookImage{02_Unnumbered_Diagrams/Chapter_05/Ch05_Diagram_PDFp378_01_Bookp368.pdf}
\end{equation}
其新边，是元素 $c\in\operatorname{Ext}^1(C,A)\cong\operatorname{Hom}(C,A[1])$。
\end{example*}

\SourceParagraph{ch05-5.5.1-006}{378}
由于 $D(\mathcal{C})$ 一般不是 Abel 范畴，没有在 $\mu$ 稳定性与 $\theta$ 稳定性定义中使用的、同样意义的短正合列与子对象。因此，在使用过子对象的地方，都需要找到只使用三角结构的替代物。

\SourceParagraph{ch05-5.5.1-007}{378}
物理上，希望把三角形\BookEqRef{eq:5.125}{5.125}读成：D 膜 $A,C$，可以通过可能是快子的开弦 $c$ 结合，形成 $B$。逐一考察各条公理，并解释其含义。

\SourceParagraph{ch05-5.5.1-008}{378}
\begin{letteraxiom}{TR1}\label{interpretation:5.TR1}
\begin{enumerate}[label=\alph*)]
\item $A$ 与 $0$，即空膜，结合，产生 $A$。
\item 把 $D(X)$ 中两个同构对象，看作同一个 D 膜。因此，自洽性要求这一规则。
\item 从 $A$ 到 $B$ 存在开弦，意味着，$B$ 有可能衰变为 $A$ 与某个其他衰变产物 $C$。这并不显然，不过，在下一条公理以后，可以重新表述。
\end{enumerate}
\end{letteraxiom}

\SourceParagraph{ch05-5.5.1-009}{378}
\begin{letteraxiom}{TR2}\label{interpretation:5.TR2}
若 $B$ 有可能衰变为 $A,C$，则 $C$ 有可能衰变为 $A[1],B$。这与\BookSectionRef{sec:5.3.3.1}{5.3.3.1}的观察一致，即，可以把 $A[1]$ 解释为 $A$ 的反膜。

注意，使用这一公理，可以把\BookRef{interpretation:5.TR1}{公理TR1}的 c) 重新表述为：给定从 $A$ 到 $B$ 的开弦，有可能形成这两个 D 膜的相应束缚态；移位决定哪个对象作为反膜参与。
\end{letteraxiom}

\SourceParagraph{ch05-5.5.1-010}{378}
\begin{letteraxiom}{TR3}\label{interpretation:5.TR3}
给定 D 膜 $A,A'$ 之间，以及 $B,B'$ 之间的相容开弦，使相应方块交换，就可以构造对应束缚态之间的开弦。
\end{letteraxiom}

\SourceParagraph{ch05-5.5.1-011}{378}
\begin{letteraxiom}{TR4}\label{interpretation:5.TR4}
这条看起来令人畏惧的公理，基本上只是 D 膜结合规则中的结合律。若粗略地用加法表示结合规则，使用\BookRef{interpretation:5.TR2}{公理TR2}，\BookEqRef{eq:4.7}{4.7}中的正合三角形可以读为
\begin{equation}\label{eq:5.126}
\begin{aligned}
C&=A[1]+B\\
&=A[1]+(E+D[-1])\\
&=(A[1]+E)+D[-1]\\
&=F+D[-1].
\end{aligned}
\end{equation}
\end{letteraxiom}

\SourceParagraph{ch05-5.5.1-012}{379}
可以选择，把这些 D 膜衰变规则看作自明的；也可以认为，它们已由前面的范畴识别得到证明，因为 B 膜范畴是导出范畴，因而具有三角结构。

\subsection{终于得到范畴镜像对称}\label{sec:5.5.2}

\SourceParagraph{ch05-5.5.2-001}{379}
当然，\BookSectionRef{sec:5.5.1}{5.5.1}每次提到束缚或衰变时，“有可能”一词都有些模糊。已经明确说过，若从 $A$ 到 $B$ 有开弦，就把 $A+B$ 看作相应分次下的可能束缚态。要使它真正发生，模空间中必须有某个区域，其中，$A,B$ 本身都稳定，且从 $A$ 到 $B$ 的相关开弦是快子。这并不保证。因此，三角结构体现了一种关于哪些对象能相互结合的乐观看法，而这也是拓扑场论所能知道的全部。

\SourceParagraph{ch05-5.5.2-002}{379}
\BookSectionRef{sec:5.2.1}{5.2.1}直接讨论了 A 膜稳定性，而没有使用拓扑场论语言。因此，几何 Fukaya 范畴未必具有三角结构：它确实知道真正稳定的 A 膜，但未必包含拓扑场论中那些有可能稳定、却从未真正成为稳定膜的对象。特别地，没有理由假定，原始 Fukaya 范畴本身是三角范畴；它可能不满足\BookRef{interpretation:5.TR1}{公理TR1}的 c)。

\SourceParagraph{ch05-5.5.2-003}{379}
若 Fukaya 范畴不是三角范畴，那么，\BookSectionRef{sec:5.3.3.5}{5.3.3.5}的镜像对称主张，就不可能按原样成立。导出范畴 $D(X)$ 是三角范畴，因而不可能等价于一个非三角范畴。解决方法，是把额外的“可能稳定”的 A 膜，形式地加入 Fukaya 范畴，使结果成为三角范畴。可以遵循 Bondal 与 Kapranov\BookCite{ref050}的方法做到这一点，\BookSectionRef{sec:8.3.4}{8.3.4}将讨论。

\SourceParagraph{ch05-5.5.2-004}{379}
由拓扑场论构造得到的、目前最完整的镜像对称猜想，便是：
\BookStatementNumber{29}
\begin{conjecture}\label{conj:5.29}
若 $X,Y$ 是互为镜像的 Calabi--Yau 三维流形，则范畴 $D(X)$ 等价于范畴 $\operatorname{Tw}\mathcal{F}(Y)$。
\end{conjecture}

\SourceParagraph{ch05-5.5.2-005}{379}
\BookChapterRef{8}将详细讨论这一猜想的现状，以及 $\operatorname{Tw}\mathcal{F}(Y)$ 的定义。

\subsection{单值化}\label{sec:5.5.3}

\SourceParagraph{ch05-5.5.3-intro-001}{379}
在更深入地讨论稳定性以前，应当认识到：若想在 $B+\mathrm{i}J$ 模空间的某一点，识别唯一的一组稳定 B 膜，会遇到根本障碍。真正的弦理论模空间，具有不可缩闭环；它们对应于 $D(X)$ 中非平凡的单值化。沿模空间中的闭环行走时，这种单值化改变稳定对象的标记集合。

\SourceParagraph{ch05-5.5.3-intro-002}{380}
在镜像 A 模型的模空间中，这更容易看到；它是镜像流形 $Y$ 的复结构模空间。考虑复结构退化的极限。例如，射影簇中的超曲面，由截面 $f$ 的零点定义；在族中，非退化条件 $\partial f\ne0$，通常沿复余维1的轨迹失效。挖去这些退化点，便可以产生非单连通的模空间。

\SourceParagraph{ch05-5.5.3-intro-003}{380}
这些退化极限，对应于“奇异 CFT”，它们以尚未完全理解的方式，违反\BookChapterRef{3}的公理。一种很常见的情形，是某个物理 A 膜变为无质量，即 $Z(B)\to0$。根据文献\BookCite{ref432,ref195}的论证，这在 CFT 中诱导奇点。

\SourceParagraph{ch05-5.5.3-intro-004}{380}
由\BookEqRef{eq:5.128}{5.128}，在这种奇异 CFT 附近，$\xi(B)$ 多值。通常，$Z(B)$ 是很好的局部坐标；因此，绕奇点一周的闭环，诱导单值化 $\xi(B)\mapsto\xi(B)+2$。很快就会更深入地讨论。

\SourceParagraph{ch05-5.5.3-intro-005}{380}
对于某些目的，消除这一歧义最简单的方法，是采用下面的定义。
\BookStatementNumber{30}
\begin{definition}\label{def:5.30}
本节所称的 \emph{Teichm\"uller 空间} $\mathcal{T}(Y)$，是非奇异复流形 $Y$ 的模空间的泛覆盖。
\end{definition}

\subsubsection{单值化在导出范畴上的作用}\label{sec:5.5.3.1}

\SourceParagraph{ch05-5.5.3.1-001}{380}
另一种理解单值化的方法，是与 A 膜图景比较；在那里，它纯粹是经典的。在复结构模空间中绕不可缩闭环时，全纯3形式在整3闭链上的周期，经历非平凡单值化。可以把它解释为 $H_3(Y,\mathbb{Z})$ 的自同构 $T$，保持3闭链之间的相交形式。

\SourceParagraph{ch05-5.5.3.1-002}{380}
更准确地说，考虑复模空间上的 Calabi--Yau 三维流形族，以及族中适当固定的辛结构。沿模空间中的路径，产生两个纤维之间的辛同胚。这个辛同胚，把 A 膜映到 A 膜，不过，单个 A 膜的稳定性会经历复杂变化。闭合路径，恢复原来的度规几何时，原来的稳定 A 膜集合必须恢复；这一点，可以通过在一组显式基中表示稳定膜的同调类来描述。此外，初始 Calabi--Yau 流形上得到一个辛自同构，诱导单值化算符 $T$。它也诱导 Fukaya 范畴的自等价，因为它把 A 膜映到 A 膜。

\SourceParagraph{ch05-5.5.3.1-003}{381}
通过\BookSectionRef{sec:3.4.3}{3.4.3}的镜像映射，可以把同调上的这一单值化作用，从 A 模型转移到镜像 B 模型。再通过例子，可以看出 B 模型单值化作用的一些图景。事实上，不难找到这样的例子：从凝聚层出发，得到 $H^{\mathrm{even}}(X,\mathbb{Z})$ 中的一个元素，它不可能对应任何层的 Chern 特征标，例如，丛的秩可能为负。这再次说明，必须把导出范畴纳入讨论。

\SourceParagraph{ch05-5.5.3.1-004}{381}
由于 B 模型在这些模参数变化下不变，可知，每种单值化，都必须对应于 $D(X)$ 的自等价，如\BookSectionRef{sec:4.6.6}{4.6.6}所述。回顾\BookRef{thm:4.82}{定理4.82}，在相应的光滑射影情形中，$D(X)$ 的这类自等价，由 Fourier--Mukai 变换诱导。因此，有同态
\begin{equation}\label{eq:5.127}
\pi_1(\mathcal{M}(Y))\longrightarrow\operatorname{Auteq}(D(X)),
\end{equation}
从镜像 $Y$ 的复结构模空间的基本群，映到 $X$ 上凝聚层导出范畴的自等价群，等价按自然同构识别。$\mathcal{T}(Y)$ 中由单值化联系的两点，其稳定对象子范畴，必须由对应自等价联系。\BookSectionRef{sec:5.6.2}{5.6.2}与\BookSectionRef{sec:5.8.3.1}{5.8.3.1}将看到例子。

\subsection{\texorpdfstring{$\Pi$}{Π} 稳定性}\label{sec:5.5.4}

\SourceParagraph{ch05-5.5.4-intro-001}{381}
假定镜像对称成立，就可以把\BookSectionRef{sec:5.2.1}{5.2.1}中 A 膜稳定性的描述，转移到 B 膜。回顾\BookSectionRef{sec:5.3.5.5}{5.3.5.5}对中心荷 $Z(E)$ 的定义。给定 $E$，可以选择 $\xi(E)$，使
\begin{equation}\label{eq:5.128}
\xi(E)=\frac1\pi\arg Z(E)\pmod2,
\end{equation}
并要求，只要 $E$ 稳定，$\xi(E)$ 就随 $B+\mathrm{i}J$ 连续变化。按照\BookEqRef{eq:5.74}{5.74}，有
\begin{equation}\label{eq:5.129}
\xi(E[n])=\xi(E)+n.
\end{equation}

\SourceParagraph{ch05-5.5.4-intro-002}{381}
最后，借用\BookRef{fig:5.3}{图3}的图景，断言：若 $D(X)$ 中有如下正合三角形，
\begin{equation}\label{eq:5.130}
\BookImage{02_Unnumbered_Diagrams/Chapter_05/Ch05_Diagram_PDFp381_01_Bookp371.pdf}
\end{equation}
其中，$A,B$ 稳定，那么，相对于这个三角形代表的衰变，$C$ 稳定，当且仅当 $\xi(B)<\xi(A)+1$。此外，若 $\xi(B)=\xi(A)+1$，则 $C$ 边际稳定，可以写为
\begin{equation}\label{eq:5.131}
\xi(C)=\xi(B)=\xi(A)+1.
\end{equation}
使用\BookSectionRef{sec:5.5.1}{5.5.1}中的\BookRef{interpretation:5.TR2}{公理TR2}，可以重新表述为：
\BookStatementNumber{31}
\begin{definition}[\BookCite{ref133}]\label{def:5.31}
若对所有形如\BookEqRef{eq:5.130}{5.130}、且 $A,C$ 均为 $\Pi$ 稳定的三角形，都有
\begin{equation}\label{eq:5.132}
\xi(A)<\xi(B)<\xi(C),
\end{equation}
则称膜 $B$ 为 \emph{$\Pi$ 稳定}的。相反，若任一这种三角形违反\BookEqRef{eq:5.132}{5.132}，则 $B$ 是 $\Pi$ 不稳定的。
\end{definition}

\SourceParagraph{ch05-5.5.4-intro-003}{382}
这一条件的必要性，也可以直接在 SCFT 中看出。对于任意两个物理边界条件 $A,B$，可以定义开弦 Hilbert 空间 $\mathcal{H}_{AB}$，带有 $N=2$ SCA 的作用。它们必须满足\BookSectionRef{sec:3.3.3}{3.3.3}导出的约束，特别地，$U(1)$ 荷必须满足\BookEqRef{eq:3.97}{3.97}。利用\BookEqRef{eq:5.41}{5.41}中 $U(1)$ 荷与分次的关系，得到
\begin{equation}\label{eq:5.133}
\xi(A)>\xi(B)\quad\Longrightarrow\quad\operatorname{Hom}(A,B)=0.
\end{equation}
这直接类似于 $\mu$ 稳定性讨论中的\BookRef{lem:5.9}{引理5.9}，并在大体积极限中约化为它。

\subsubsection{\texorpdfstring{$\Pi$}{Π} 稳定性的分析}\label{sec:5.5.4.1}

\SourceParagraph{ch05-5.5.4.1-001}{382}
像\BookSectionRef{sec:5.3.6}{5.3.6}那样改变分次时，除非 $\Pi$ 稳定膜的集合以某种方式改变，一般就会违反\BookEqRef{eq:5.133}{5.133}。也许希望类比 $\mu$ 稳定性的讨论，把所有涉及稳定膜 $A,C$ 的三角形的不等式\BookEqRef{eq:5.132}{5.132}，看作 $B$ 存在的必要条件，再证明它们充分。不过，这个定义内在的循环性，使人不清楚如何进行。

\SourceParagraph{ch05-5.5.4.1-002}{382}
遵循文献\BookCite{ref138,ref137,ref149,ref133,ref018}，可以探索，\BookRef{def:5.31}{定义5.31}在多大程度上，唯一确定一组自洽的稳定 B 膜。虽然直接完成这一点，并不显然，但若知道模空间某个基点处的稳定 B 膜集合，包括它们的 $\xi$，就可以应用下列规则，确定沿模空间路径移动时，稳定谱如何改变。
\begin{itemize}
\item 从一组稳定 B 膜出发，并为每个膜指定 $\xi$ 值。这个集合必须符合 $\Pi$ 稳定性规则，即，不应有正合三角形，允许其中某个稳定 B 膜，衰变为另外两个稳定 B 膜。
\item 沿模空间路径移动时，$\xi$ 连续变化。
\item 两个稳定 B 膜，可以结合，形成新的稳定态。
\item 一个稳定 B 膜，可以衰变为其他边际稳定态或稳定态。
\end{itemize}
注意，最后一种情形中，膜可以衰变为另一个态，而后者恰好在模空间的同一点变得不稳定。还应注意，这些规则没有使用不稳定对象的 $\xi$ 值；这个值可能并没有定义。

\SourceParagraph{ch05-5.5.4.1-003}{383}
于是，猜想是：这些规则，在路径终点确定唯一的 $\Pi$ 稳定子集。按路径表述，避开了单值化引起的困难。另一方面，必须证明，结果在路径的连续形变下不变。文献\BookCite{ref018}给出了这方面的论证，但它们已被接下来讨论的表述所涵盖。承认这一点，就可以预期，对于 $\mathcal{T}(Y)$ 中任意一点，$X$ 上的稳定 B 膜集合唯一确定。

\subsubsection{多重衰变}\label{sec:5.5.4.2}

\SourceParagraph{ch05-5.5.4.2-001}{383}
对于 $\mathcal{T}(Y)$ 中给定一点，$D(X)$ 的每个对象，要么稳定，要么不稳定。可以预期，不稳定粒子必须衰变为一组稳定粒子。因此，稳定对象的集合必须足够大，能够描述这种衰变。这比上一小节对稳定性的约束更强。例如，截至目前的考虑，甚至与完全没有稳定对象的情形相容。

\SourceParagraph{ch05-5.5.4.2-002}{383}
若不稳定对象衰变为两个稳定对象，已经知道，可以用正合三角形描述。现在希望描述，一个对象衰变为三个稳定对象。使用下面的八面体，描述这一过程\BookCite{ref018}：
\begin{equation}\label{eq:5.134}
\BookImage{02_Unnumbered_Diagrams/Chapter_05/Ch05_Diagram_PDFp383_01_Bookp373.pdf}
\end{equation}

\SourceParagraph{ch05-5.5.4.2-003}{383}
假定从 Teichm\"uller 空间中的点 $p_0$ 出发，此处 $\xi(A_1)<\xi(A_2)<\xi(A_3)$，到达点 $p_1$，此处 $\xi(A_1)>\xi(A_2)>\xi(A_3)$。因此，从 $p_0$ 到 $p_1$ 的过程中，\BookEqRef{eq:5.134}{5.134}中与 $f_1,f_2$ 对应的开弦，从快子变成有质量的场。

\SourceParagraph{ch05-5.5.4.2-004}{383}
在 $p_0$，相对于八面体中的这些三角形，$E_2,F$ 稳定。也可以规定 $E_3$ 稳定，不过，这并非真正必需。假定在 $p_0,p_1$ 之间有两道墙 $W_1,W_2$，使 $\xi(A_i)-\xi(A_{i+1})$，在 $W_i$ 的 $p_0$ 一侧为负，在 $p_1$ 一侧为正。那么，从 $p_0$ 到 $p_1$，要考虑两种可能。
\begin{enumerate}[label=(\arabic*)]
\item 先穿过 $W_1$，再穿过 $W_2$。穿过 $W_2$ 时，对象 $F$ 衰变为 $A_2,A_3$。这时，$\xi(F)=\xi(A_2)<\xi(A_1)$，所以，知道 $E_3$ 一定已经衰变为 $F,A_1$。因此，到达 $p_1$ 时，$E_3$ 已衰变为 $A_1,A_2,A_3$。
\item 先穿过 $W_2$，再穿过 $W_1$。穿过 $W_1$ 时，对象 $E_2$ 衰变为 $A_1,A_2$。这时，$\xi(E_2)=\xi(A_2)>\xi(A_3)$，所以，知道 $E_3$ 一定已经衰变为 $E_2,A_3$。因此，到达 $p_1$ 时，$E_3$ 已衰变为 $A_1,A_2,A_3$。
\end{enumerate}
无论哪种顺序，$E_3$ 衰变为 $A_1,A_2,A_3$ 的条件，都是
\begin{equation}\label{eq:5.135}
\xi(A_1)>\xi(A_2)>\xi(A_3).
\end{equation}

\SourceParagraph{ch05-5.5.4.2-005}{384}
可以把这推广到衰变为任意多个对象的情形。对于任意对象 $E$，定义下面这一组正合三角形，
\begin{equation}\label{eq:5.136}
\BookImage{02_Unnumbered_Diagrams/Chapter_05/Ch05_Diagram_PDFp384_01_Bookp374.pdf}
\end{equation}
只要
\begin{equation}\label{eq:5.137}
\xi(A_1)>\xi(A_2)>\cdots>\xi(A_n),
\end{equation}
$E$ 就可以衰变为 $A_1,\dots,A_n$。由此得到下面的猜想。
\BookStatementNumber{32}
\begin{conjecture}[\BookCite{ref065}]\label{conj:5.32}
对于 $\mathcal{T}(Y)$ 中任意一点，$D(X)$ 中的稳定对象子集，满足：任意对象 $E$ 都可以（如\BookEqRef{eq:5.136}{5.136}）分解为唯一一组满足\BookEqRef{eq:5.137}{5.137}的 $n$ 个稳定对象 $A_k$。
\end{conjecture}

\SourceParagraph{ch05-5.5.4.2-006}{384}
虽然物理上合理，这并不完全显然，因为在 SCFT 中，不清楚正在讨论的不稳定粒子究竟指什么。类比 Donaldson--Uhlenbeck--Yau 定理，一个自然定义是：使用相关对象的正确全纯结构定义边界条件，但还不知道 beta 函数条件的实际解。随后，可以尝试通过 RG 流，构造共形边界条件；可以预期，不稳定出发点会流向这些稳定因子的直和。\BookSectionRef{sec:5.9}{5.9}将进一步讨论这一思想。

\SourceParagraph{ch05-5.5.4.2-007}{384}
无论如何，承认\BookRef{conj:5.32}{猜想5.32}，就得到一组确定稳定对象集合的公理，\BookSectionRef{sec:5.7}{5.7}将说明。

\subsection{与 \texorpdfstring{$\mu$}{mu} 稳定性的比较}\label{sec:5.5.5}

\SourceParagraph{ch05-5.5.5-001}{384}
为了确定 $\Pi$ 稳定对象的集合，最好在 $B+\mathrm{i}J$ 模空间中，选择一个事先知道稳定对象的基点，再像\BookSectionRef{sec:5.5.4}{5.5.4}那样，沿路径行走。显然，最适合把基点放在 Calabi--Yau 三维流形 $X$ 的大半径极限附近，因为预期，那里的 B 膜，对应于支撑在，或“缠绕”，子流形 $L\subset X$ 的向量丛。

\SourceParagraph{ch05-5.5.5-002}{385}
考虑缠绕 $L$ 的膜，把其全纯向量丛记为 $F\to L$，曲率记为 $(1,1)$ 形式 $\mathcal{F}$。在大半径极限中，BPS 条件约化为 Hermitian--Yang--Mills 条件，因此，若 $F$ 是 $\mu$ 稳定的，就存在对应的不可约解，见\BookSectionRef{sec:5.2.2}{5.2.2}。这意味着，对任意满足适当秩条件的非零真子层 $E\subset F$，\BookEqRef{eq:5.33}{5.33}的斜率满足
\begin{equation}\label{eq:5.138}
\mu(E)<\mu(F).
\end{equation}

\SourceParagraph{ch05-5.5.5-003}{385}
当然，若 $E$ 是 $F$ 的子层，就有短正合列，参见\BookEqRef{eq:5.36}{5.36}，
\begin{equation}\label{eq:5.139}
0\longrightarrow E\longrightarrow F\longrightarrow G\longrightarrow0,
\end{equation}
其中，$G$ 是某个层。因此，用 $\Pi$ 稳定性描述这一过程时，应当涉及对应的三角形。

\SourceParagraph{ch05-5.5.5-004}{385}
也许会问，不同维数的膜组成的三角形，是否也会影响这个判据。在大体积极限中，它们并不相关：$X$ 上对应于 $F$ 的丛，不能衰变为只支撑在真子空间 $S$ 上的子层 $E$，因为没有非零同态 $E\to F$。同样，$S$ 上的丛，也不能衰变为在整个 $X$ 上具有正秩的子层，因为\BookEqRef{eq:5.139}{5.139}中的商层 $G$，那样会具有负秩。

\SourceParagraph{ch05-5.5.5-005}{385}
因此，若条件\BookEqRef{eq:5.132}{5.132}在大体积极限中约化为\BookEqRef{eq:5.138}{5.138}，我们就恢复出作为 $\Pi$ 稳定性极限的 $\mu$ 稳定性。

\SourceParagraph{ch05-5.5.5-006}{385}
由\BookEqRef{eq:5.71}{5.71}可知，$J$ 很大时，中心荷 $Z(A)$ 的首项贡献，来自 $\operatorname{ch}(A)$ 中次数最低的微分形式。\BookSectionRef{sec:5.1.4.1}{5.1.4.1}说明，$\operatorname{ch}(i_*F)$ 的最低次项，是与 $S$ Poincar\'e 对偶的 $(6-2\dim S)$ 形式 $s$，乘以丛的秩。这里，$\dim S$ 是 $S$ 的复维数。因此，对很大的 $J$，省略正实系数，有
\begin{equation}\label{eq:5.140}
\begin{aligned}
Z&\sim\int_X(-\mathrm{i}J)^{\dim S}\wedge s\\
&\sim\int_S(-\mathrm{i}J|_S)^{\dim S}\\
&\sim(-\mathrm{i})^{\dim S}\operatorname{Vol}(S),
\end{aligned}
\end{equation}
从而，在极限中，
\begin{equation}\label{eq:5.141}
\xi(i_*F)\equiv-\frac12\dim S\pmod2.
\end{equation}

\SourceParagraph{ch05-5.5.5-007}{385}
若任意选择 $\xi$ 的值，以消除模2歧义，就会违反幺正性条件\BookEqRef{eq:5.97}{5.97}。例如，令 $\mathcal{O}_X$ 为缠绕 $X$ 的6膜，$\mathcal{O}_p$ 为 $p\in X$ 处的0膜，即摩天层。于是，$\xi(\mathcal{O}_X)=-\tfrac32\pmod2$，$\xi(\mathcal{O}_p)=0\pmod2$。把 $X$ 上的函数限制为它在 $p$ 的值，可以看到，$\operatorname{Hom}(\mathcal{O}_X,\mathcal{O}_p)=\mathbb{C}$，所以，若这些 B 膜稳定，必须有 $\xi(\mathcal{O}_X)<\xi(\mathcal{O}_p)$。此外，由 Serre 对偶，$\operatorname{Hom}(\mathcal{O}_p,\mathcal{O}_X[3])=\mathbb{C}$，因而 $\xi(\mathcal{O}_p)<\xi(\mathcal{O}_X)+3$\footnote{$\mathcal{O}_X[3]$ 是把 $\mathcal{O}_X$ 放在次数 $-3$，其余项均为零的复形。}。所以，大半径处唯一可能的极限，是 $\xi(\mathcal{O}_X)=\xi(\mathcal{O}_p)-\tfrac32$。

\SourceParagraph{ch05-5.5.5-008}{386}
一个一致的分支选择，是令
\begin{equation}\label{eq:5.142}
\lim_{J\to\infty}\xi(i_*F)=-\frac12\dim S.
\end{equation}
相应的开弦 $U(1)_R$ 荷移位，也可以直接在非线性 $\sigma$ 模型中看到，作为 Dirichlet--Neumann 开弦量子化的后果。

\SourceParagraph{ch05-5.5.5-009}{386}
看看 $J$ 的次首阶发生什么。限制在6膜，即局部自由层的情形。现在，$\operatorname{ch}(E)=k+c_1+\cdots$，其中，$k$ 是对应向量丛的秩。应用\BookEqRef{eq:5.71}{5.71}，得到
\begin{equation}\label{eq:5.143}
\xi(E)=-\frac32+\frac1\pi\tan^{-1}\frac{\mu(E)}{2}+\cdots,
\end{equation}
$\mu(E)$ 是\BookEqRef{eq:5.33}{5.33}定义的斜率。注意，在大体积极限中，$|\mu(E)|\ll1$。

\SourceParagraph{ch05-5.5.5-010}{386}
短正合列\BookEqRef{eq:5.139}{5.139}，在 $D(X)$ 中诱导三角形
\begin{equation}\label{eq:5.144}
\BookImage{02_Unnumbered_Diagrams/Chapter_05/Ch05_Diagram_PDFp386_01_Bookp376.pdf}
\end{equation}
如\BookSectionRef{sec:5.5.1}{5.5.1}所解释。在大半径附近，若 $E,F$ 都局部自由，$\xi(E),\xi(F)$ 都很接近 $-\tfrac32$。$G$ 要么局部自由，要么支撑在复余维1的子空间上。因此，$\xi(G)$ 分别接近 $-\tfrac32$ 或 $-1$。根据\BookSectionRef{sec:5.1}{5.1}，D 膜的荷相加，故 $Z(F)=Z(E)+Z(G)$，这意味着，$\xi(F)$ 位于 $\xi(E),\xi(G)$ 之间。$F$ 的 $\Pi$ 稳定性条件为 $\xi(E)<\xi(G)$，因而等价于 $\xi(E)<\xi(F)$。由\BookEqRef{eq:5.143}{5.143}，这又等价于 $\mu(E)<\mu(F)$。因此，$\Pi$ 稳定性，在大体积极限中约化为 $\mu$ 稳定性；文献\BookCite{ref138}最早指出这一点。

\subsection{五次三维流形上的例子}\label{sec:5.5.6}

\SourceParagraph{ch05-5.5.6-intro-001}{386}
给出一些稳定性变化的具体例子。也许最重要的出发点，是给出一个不是凝聚层的 B 膜，从而至少从物理观点，为已经引入的相当复杂的形式体系提供依据。最初的这种“奇特”膜的例子，来自文献\BookCite{ref134}对五次三维流形的研究。

\SourceParagraph{ch05-5.5.6-intro-002}{386}
设 $X$ 是 $\mathbb{P}^4$ 中的五次超曲面。它具有 $b^{1,1}=1$，其偶 K 群 $K^0(X)\cong\mathbb{Z}^4$。把 $H^2(X,\mathbb{Z})$ 的生成元记为 $e$。

\SourceParagraph{ch05-5.5.6-intro-003}{387}
考虑从 $\mathbb{P}^4$ 限制到 $X$ 得到的凝聚层。虽然这并非所有凝聚层，但它们张成 $K^0(X)\otimes\mathbb{R}$，对于当前目的，已经足够。令 $\mathcal{O}_X(m)\cong\mathcal{O}_{\mathbb{P}^4}(m)|_X$；它们是秩1局部自由层。

\SourceParagraph{ch05-5.5.6-intro-004}{387}
无须引入全部形式体系，使用 B 膜中心荷的大体积近似，并考虑相关微扰 $\alpha'$ 修正，就能看到关键点。使用精确周期的更完整讨论，见文献\BookCite{ref017}。回顾\BookEqRef{eq:5.71}{5.71}中使用的几何近似式，
\begin{equation}\label{eq:5.145}
Z(E)\simeq\int_X\mathrm{e}^{-(B+\mathrm{i}J)}
\operatorname{ch}(E)\sqrt{\operatorname{td}(X)}.
\end{equation}
这里暂且略去指数微小的瞬子修正；下面的示意讨论，也不保留微扰常数项。对于当前讨论，还可以略去因子
\[
\sqrt{\operatorname{td}(X)}=\sqrt{1+\frac56e^2}
=1+\frac5{12}e^2.
\]

\SourceParagraph{ch05-5.5.6-intro-005}{387}
于是，采用这一级近似，
\begin{equation}\label{eq:5.146}
\begin{aligned}
\xi(\mathcal{O}_X(m))
&\simeq\frac1\pi\arg\int_X\exp\bigl((m-B-\mathrm{i}J)e\bigr)\\
&=\frac1\pi\arg\!\left(\frac56(m-B-\mathrm{i}J)^3\right)\\
&=\frac3\pi\theta_m-3,
\end{aligned}
\end{equation}
其中，$\theta_m$ 是复 $B+\mathrm{i}J$ 平面中的角，即从 $m+\mathrm{i}0$ 到 $B+\mathrm{i}J$ 的直线，与正实轴之间的夹角。

\subsubsection{五次三维流形上4膜的稳定性}\label{sec:5.5.6.1}

\SourceParagraph{ch05-5.5.6.1-001}{387}
考虑下面的层短正合列：
\begin{equation}\label{eq:5.147}
0\longrightarrow\mathcal{O}_X(a)\xrightarrow{f}\mathcal{O}_X(b)
\longrightarrow\mathcal{O}_S(b)\longrightarrow0,
\end{equation}
其中，$b>a$，$S$ 是 $f$ 的零点除子。这里考虑的五次三维流形上的4膜族，对应于不同 $f,b$ 给出的 $\mathcal{O}_S(b)$。

\SourceParagraph{ch05-5.5.6.1-002}{387}
把 $\Pi$ 稳定性判据用于\BookEqRef{eq:5.147}{5.147}对应的正合三角形，得到 $\xi(\mathcal{O}_X(b))-\xi(\mathcal{O}_X(a))<1$，也就是
\begin{equation}\label{eq:5.148}
\theta_b-\theta_a<\frac\pi3.
\end{equation}
满足这一条件时，\BookEqRef{eq:5.147}{5.147}中与 $f$ 对应的开弦，是快子。简单几何表明，这对应于上半 $B+\mathrm{i}J$ 平面中，某条圆弧以上的点；圆心为 $\tfrac12(a+b)+\tfrac{b-a}{2\sqrt3}\mathrm{i}$，半径为 $\tfrac{b-a}{\sqrt3}$。

\SourceParagraph{ch05-5.5.6.1-003}{387}
因此，相对于这一衰变道，这些4膜在大半径极限稳定，符合预期。不过，在这一边际稳定性圆弧以下，4膜衰变为 $\mathcal{O}_X(b)$ 与 $\mathcal{O}_X(a)[-1]$。物理上，这是6膜与反6膜，都带有一些4膜荷。另一方面，在不可约支撑 $S$ 上秩1丛的意义下，这种4膜自动 $\mu$ 稳定，因为没有需要斜率判据检验的、更小正秩子层。因此，这是两个稳定性条件给出不同预言的简单例子。
\begin{figure}[H]
\centering
\BookImage{01_Numbered_Figures/Chapter_05/Ch05_Fig08_PDFp388_Bookp378.pdf}
\caption{各种 D4膜在 $B+\mathrm{i}J$ 平面中的稳定性。如\BookSectionRef{sec:5.5.6.1}{5.5.6.1}所解释，第 $N$ 条圆弧以下，D4膜 $\mathcal{O}_S(N)$ 按正合列\BookEqRef{eq:5.147}{5.147}衰变。取自文献\BookCite{ref018}。}\label{fig:5.8}
\end{figure}

\SourceParagraph{ch05-5.5.6.1-004}{388}
可以验证，使用\BookEqRef{eq:5.146}{5.146}的精确、经瞬子修正的表达式，会预言同样的衰变，只是在稍有不同的 $B+\mathrm{i}J$ 处发生。因此，$\mu$ 稳定性的问题，并非没有计入足够多的 $\alpha'$ 修正，而是它只适合使用凝聚层子对象关系的情形。特别地，它永远无法看见由反膜引起的衰变。

\SourceParagraph{ch05-5.5.6.1-005}{388}
还应注意，有时，即使在大半径处，$\Pi$ 稳定性修正也重要。对于非常大的 $b-a$，这条边际稳定性曲线，可以延伸到很大的 $J$。$\mu$ 稳定性近似成立的条件，应当是 $J\gg|b-a|$；更一般地，是 K\"ahler 尺度相对于相关层的 Chern 类所给出的尺度足够大。

\SourceParagraph{ch05-5.5.6.1-006}{388}
当然，还没有得到边际稳定性曲线的精确形式，因为\BookEqRef{eq:5.146}{5.146}使用大半径近似。取 $a=0$、$b=N$，使用文献\BookCite{ref085}的精确周期得到的精确曲线，如\BookRef{fig:5.8}{图8}所示。该图画出 Teichm\"uller 空间的一部分，每个围起的区域表示一个基本区域。坐标 $B+\mathrm{i}J$，通过文献\BookCite{ref085}的镜像映射定义。可以看到，精确边际稳定性曲线，与圆弧相差不大。

\SourceParagraph{ch05-5.5.6.1-007}{389}
必须强调，截至目前的考虑，既未证明4膜在\BookRef{fig:5.8}{图8}的曲线以上稳定，也未证明它们在曲线以下不稳定。4膜可能在到达这些曲线以前，就通过其他衰变道衰变；衰变产物也可能在到达边际稳定性曲线时，已经不稳定。很快会再讨论这一点。

\subsubsection{锥形点}\label{sec:5.5.6.2}

\SourceParagraph{ch05-5.5.6.2-001}{389}
注意，\BookRef{fig:5.8}{图8}的边际稳定性曲线有端点。这些是“锥形点”，其中，某个衰变产物 $\mathcal{O}_X(a)$ 或 $\mathcal{O}_X(b)$，变为无质量，即 $Z=0$，所以，分次不再良好定义。

\SourceParagraph{ch05-5.5.6.2-002}{389}
这些点，超出\BookEqRef{eq:5.71}{5.71}的有效范围，因此，需要精确分析镜像五次三维流形的周期，才能恰当讨论。由于很快会对 $\mathbb{C}^3/\mathbb{Z}_3$ 轨形作这种分析，得到十分相似的结果，这里略去，参见文献\BookCite{ref017}。

\SourceParagraph{ch05-5.5.6.2-003}{389}
这种无质量膜存在的底层原因，在镜像 A 模型中很清楚：它们是消灭闭链上的特殊 Lagrange 子流形，与复结构退化对应。数学上，可以尝试使用退化附近有效的 Ricci 平坦度规近似，证明它们存在\footnote{截至2008年，还没有完成这种证明。}；物理上，则诉诸 SCFT 奇点必须与无质量粒子相联系的原则，这里的奇点，就是复结构退化。

\SourceParagraph{ch05-5.5.6.2-004}{389}
因此，物理图景表明，B 膜 $\mathcal{O}_X(m)$，在所提出的边际稳定性曲线端点附近，以及大体积附近，稳定。这为所提出的曲线提供了一定依据。也可以假设，4膜还有更曲折的衰变方式，并检查这些过程在更小半径发生。因此，很难想象\BookRef{fig:5.8}{图8}整体不正确。

\SourceParagraph{ch05-5.5.6.2-005}{389}
更一般地，可以变为无质量的膜，起非常核心的作用：它们决定单值化，并把弦性 K\"ahler 模空间，分解为由不同 Abel 子范畴上的稳定性条件控制的区域。很快会发展这一思想。

\subsubsection{一个奇特 B 膜}\label{sec:5.5.6.3}

\SourceParagraph{ch05-5.5.6.3-001}{389}
对刚才讨论的4膜衰变，应用 Serre 对偶\BookCite{ref134}。其中可能的快子，属于 $\operatorname{Hom}(\mathcal{O}_X(a),\mathcal{O}_X(b))$，当 $b>a$ 时，这个空间非零。由 Serre 对偶，这个开弦 Hilbert 空间同构于 $\operatorname{Hom}(\mathcal{O}_X(b),\mathcal{O}_X(a)[3])$。因此，可以考虑正合三角形
\begin{equation}\label{eq:5.149}
\BookImage{02_Unnumbered_Diagrams/Chapter_05/Ch05_Diagram_PDFp389_01_Bookp379.pdf}
\end{equation}
其中，$\mathcal{X}_{a,b}$ 定义为映射 $g$ 的锥。
\begin{figure}[H]
\centering
\BookImage{01_Numbered_Figures/Chapter_05/Ch05_Fig09_PDFp390_Bookp380.pdf}
\caption{奇特对象 $\mathcal{X}_{0,N}$ 在 $B+\mathrm{i}J$ 平面中的稳定性。第 $N$ 条圆弧，对应于\BookEqRef{eq:5.149}{5.149}描述的衰变。}\label{fig:5.9}
\end{figure}

\SourceParagraph{ch05-5.5.6.3-002}{390}
现在，结合上述分析与关系\BookEqRef{eq:5.129}{5.129}，可以确定 $\mathcal{X}_{a,b}$ 的稳定性，至少相对于三角形\BookEqRef{eq:5.149}{5.149}如此。稳定性要求
\[
\xi(\mathcal{O}_X(a)[3])-\xi(\mathcal{O}_X(b))
=3+\xi(\mathcal{O}_X(a))-\xi(\mathcal{O}_X(b))<1.
\]
使用近似\BookEqRef{eq:5.146}{5.146}，可知，$\mathcal{X}_{a,b}$ 在上半平面中的某条圆弧以下稳定；圆心为 $\tfrac12(a+b)-\tfrac{b-a}{2\sqrt3}\mathrm{i}$，半径为 $\tfrac{b-a}{\sqrt3}$。特别地，对于固定 $a,b$，$\mathcal{X}_{a,b}$ 在大半径极限总是不稳定。不过，对任意给定、任意大的半径，都可以选择足够大的 $b-a$，使它稳定。

\SourceParagraph{ch05-5.5.6.3-003}{390}
再次，可以使用数值方法，画出更精确的边际稳定性曲线。\BookRef{fig:5.9}{图9}，同样取自文献\BookCite{ref018}，画出一些例子。应当注意，$\mathcal{X}_{0,2}$ 发生了很有趣的现象：某个衰变产物本身还能衰变，迫使边际稳定性曲线，终止在另一条边际稳定性曲线上，而非锥形点。完整讨论见文献\BookCite{ref018}。

\SourceParagraph{ch05-5.5.6.3-004}{390}
那么，$\mathcal{X}_{a,b}$ 究竟是什么？它是在大半径极限中不存在、且没有直接 $\sigma$ 模型解释的 B 膜；不过，可以从一组 $\sigma$ 模型边界态出发，改变 K\"ahler 模参数，构造出它。更仔细地考察这一点。假定在适当的环境层范畴中，有 $\mathcal{O}_X$ 的单射消解，
\begin{equation}\label{eq:5.150}
0\longrightarrow\mathcal{O}_X\longrightarrow\mathcal{I}^0
\xrightarrow{i_0}\mathcal{I}^1\xrightarrow{i_1}\mathcal{I}^2
\xrightarrow{i_2}\mathcal{I}^3\xrightarrow{i_3}\cdots.
\end{equation}
根据\BookChapterRef{4}，$\operatorname{Ext}^3(\mathcal{O}_X(1),\mathcal{O}_X)$ 的元素，可以由满足 $i_3g=0$ 的映射 $g\colon\mathcal{O}_X(1)\to\mathcal{I}^3$ 表示。于是，$\mathcal{X}_{0,1}=\operatorname{Cone}(g)$，对应于复形
\begin{equation}\label{eq:5.151}
\cdots\to0\to\mathcal{I}^0\xrightarrow{i_0}\mathcal{I}^1
\xrightarrow{\binom{0}{i_1}}
\begin{matrix}\mathcal{O}_X(1)\\\oplus\\\mathcal{I}^2\end{matrix}
\xrightarrow{(g\ \ i_2)}\underset{\scriptscriptstyle\cdots}{\mathcal{I}^3}
\xrightarrow{i_3}\cdots,
\end{equation}
下方的点线，标示次数零的位置。

\SourceParagraph{ch05-5.5.6.3-005}{391}
这个复形的上同调层，为 $\mathcal{H}^{-3}=\mathcal{O}_X$ 与 $\mathcal{H}^{-1}=\mathcal{O}_X(1)$。因此，朴素地看，这个 B 膜，像一个反6膜，加上另一个带有4膜荷的反6膜。当 $g$ 作为 Ext 类为零时，确实如此。更准确地说，\BookEqRef{eq:5.151}{5.151}拟同构于
\begin{equation}\label{eq:5.152}
\cdots\longrightarrow\mathcal{O}_X\xrightarrow{0}0
\xrightarrow{0}\mathcal{O}_X(1)\xrightarrow{0}0\longrightarrow\cdots,
\end{equation}
当且仅当 $g=0$。使用长正合列，显式计算 $\operatorname{Hom}(\mathcal{X}_{0,1},\mathcal{X}_{0,1})$，可以确认这一点，得到
\begin{equation}\label{eq:5.153}
\operatorname{Hom}(\mathcal{X}_{0,1},\mathcal{X}_{0,1})
=\begin{cases}
\mathbb{C}^2,&g=0,\\
\mathbb{C},&g\ne0.
\end{cases}
\end{equation}

\SourceParagraph{ch05-5.5.6.3-006}{391}
这很好地呼应\BookSectionRef{sec:5.4.4}{5.4.4}末尾的评论。若 $g\ne0$，得到单个不可约 D 膜。若 $g=0$，得到两个 D 膜的直和，因而规范群为 $U(1)\times U(1)$。因此，打开快子以后，$\mathcal{X}_{0,1}$ 成为不同的对象。

\SourceParagraph{ch05-5.5.6.3-007}{391}
这一切说明，从经典几何观点看，$\mathcal{X}_{0,1}$ 是真正奇特的对象。它不拟同构于只有单个凝聚层的复形，所以，不能看作支撑在 $X$ 某个子空间上的向量丛。

\SourceParagraph{ch05-5.5.6.3-008}{391}
注意，能够构造这些奇特 D 膜，是因为能够使用 $n>1$ 的 $\operatorname{Ext}^n$。这把复形的各项，分隔得足够远，使它们不至于全部缩并为单项复形。因此，得到一个明确依赖\BookSectionRef{sec:5.3.4}{5.3.4}那些考虑的显式例子。

\SourceParagraph{ch05-5.5.6.3-009}{391}
事实上，Gepner 点处的某些 D 膜态，已经作为边界共形场论显式构造出来\BookCite{ref402,ref069,ref163}。其中，包含与这一奇特态具有相应性质的边界态，包括相同的荷，以及没有带荷手征物质。可以使用 Gepner 模型的 B 扭曲，与从\BookSectionRef{sec:3.6.8}{3.6.8}讨论的矩阵分解构造的导出范畴之间的识别，确认这一对应。

\subsection{保典翻转}\label{sec:5.5.7}

\SourceParagraph{ch05-5.5.7-001}{392}
非紧 Calabi--Yau 三维流形，提供更简单的稳定性例子。很有启发性的一个情形，是保典翻转。它在三维镜像对称中的作用，物理上已有很多研究，概述见文献\BookCite{ref192}。

\SourceParagraph{ch05-5.5.7-002}{392}
基本例子，来自一个包含\BookRef{ex:5.28}{例题5.28}的锥形点的奇异代数簇 $X_0$。如那里所见，可以用两种不同方式，以一条 $\mathbb{P}^1$ 替换锥形点，消解奇异空间，得到光滑 Calabi--Yau 流形 $X$ 或 $X'$。一般而言，二者的整体拓扑可以不同，并由保典翻转相联系。

\SourceParagraph{ch05-5.5.7-003}{392}
几何上，把 $X$ 或 $X'$ 收缩回 $X_0$ 的过程，可以看作 K\"ahler 形式 $J$ 的形变。实际上，在 K\"ahler 形式空间中，可以把 $X,X'$ 看作位于对应于 $X_0$ 的一道墙的两侧。令 $C$ 为 $X$ 内的那条 $\mathbb{P}^1$。从 $X$ 一侧接近墙时，$C$ 的面积缩小到零。继续越过墙，形式上会给 $C$ 负面积；不过，用 $X'$ 重新解释几何，就给新的 $C'\subset X'$ 正面积。

\SourceParagraph{ch05-5.5.7-004}{392}
非线性 $\sigma$ 模型中，K\"ahler 类被提升为复模参数 $B+\mathrm{i}J$。这带来深刻影响，如文献\BookCite{ref020}所述。只要与 $C$，或 $C'$，对应的 $B$ 分量不为零，模去其整数周期，奇异靶空间 $X_0$ 对应的共形场论，就可以非奇异。因此，$J$ 模空间中的实余维1奇异空间之墙，变成复余维1的奇异共形场论子空间，符合镜像对称与\BookSectionRef{sec:5.5.3}{5.5.3}的讨论。于是，可以绕过这个奇异子集，从 $X$ 光滑地过渡到 $X'$，得到连通的 SCFT 模空间。

\SourceParagraph{ch05-5.5.7-005}{392}
不过，一旦把 D 膜纳入图景，就会看到，某种意义上必须有跳跃。至少在一定程度的近似下，Calabi--Yau 靶空间，就是0膜模空间。因此，即使避开 $B=0$ 处的奇异共形场论，0膜模空间也必须在保典翻转过程中经历某种跃迁。文献\BookCite{ref061}，在文献\BookCite{ref016}的构造基础上，说明，0膜稳定性带来了这种不连续性。

\SourceParagraph{ch05-5.5.7-006}{392}
设处于 Calabi--Yau 三维流形 $X$ 中，除参与翻转的 $\mathbb{P}^1$ 以外，所有 K\"ahler 尺度都满足 $J\gg1$。这个极限中，镜像上的周期可以简单分析，如文献\BookCite{ref019}所解释。基本上，唯一感兴趣的 $B+\mathrm{i}J$ 分量，是
\begin{equation}\label{eq:5.154}
t=\int_C(B+\mathrm{i}J).
\end{equation}

\SourceParagraph{ch05-5.5.7-007}{392}
这个参数的紧化模空间，由\BookRef{fig:5.10}{图10}中的 $\mathbb{P}^1$ 给出。保典翻转发生在赤道，奇异共形场论位于标记为 $O$ 的点，在相应分支中 $t=0$。令 $z$ 为这个 $\mathbb{P}^1$ 上的仿射坐标，使 $z=0$ 为大 $C$ 极限，$z=1$ 对应于 $O$，而 $z=\infty$ 为保典翻转以后，大 $C'$ 的极限。
\begin{figure}[H]
\centering
\BookImage{01_Numbered_Figures/Chapter_05/Ch05_Fig10_PDFp393_Bookp383.pdf}
\caption{保典翻转的 $B+\mathrm{i}J$ 模空间。}\label{fig:5.10}
\end{figure}

\SourceParagraph{ch05-5.5.7-008}{393}
镜像上的周期，具有一般形式\BookCite{ref019}，
\begin{equation}\label{eq:5.155}
\Phi=A_1+A_2\log z.
\end{equation}
因此，镜像映射给出的精确关系为
\begin{equation}\label{eq:5.156}
t=\frac{1}{2\pi\mathrm{i}}\log z.
\end{equation}

\SourceParagraph{ch05-5.5.7-009}{393}
令 $i\colon C\to X$ 为包含映射，用 $\mathcal{O}_C(m)$ 表示 $i_*\mathcal{O}(m)$；用 $\mathcal{O}_x$ 表示与点 $x\in X$ 对应的摩天层，即0膜。由 Grothendieck--Riemann--Roch 定理，参见\BookSectionRef{sec:5.1}{5.1}，
\begin{equation}\label{eq:5.157}
\int_X\mathrm{e}^{-(B+\mathrm{i}J)}
\operatorname{ch}(\mathcal{O}_C(m))\sqrt{\operatorname{td}(TX)}
=-t+m+1.
\end{equation}
因此，$Z(\mathcal{O}_C(m))=-t+m+1$ 是精确的，因为 $t$ 与1都是周期。于是，最自然的说法似乎是，与 $\mathcal{O}_C(-1)$ 对应的膜，在 $O$ 处变为无质量。不过，可以自由选择对数分支；为了简单，也可以改为，让 $\mathcal{O}_C$ 在那里变为无质量。这等价于，在大半径极限附近选择基点，但在走向 $O$ 以前，先绕大半径极限一周。采用这个选择，令 $t=1+\epsilon\mathrm{e}^{\mathrm{i}\theta}$，其中，$\epsilon$ 为很小的正实数，把注意力集中在 $O$ 的这一邻域。\BookRef{fig:5.11}{图11}给出示意。
\begin{figure}[H]
\centering
\BookImage{01_Numbered_Figures/Chapter_05/Ch05_Fig11_PDFp394_Bookp384.pdf}
\caption{$O$ 的邻域。}\label{fig:5.11}
\end{figure}

\SourceParagraph{ch05-5.5.7-010}{394}
设 $x\in C$。有短正合列
\begin{equation}\label{eq:5.158}
0\longrightarrow\mathcal{O}_C(-1)\longrightarrow\mathcal{O}_C
\longrightarrow\mathcal{O}_x\longrightarrow0.
\end{equation}
它给出正合三角形，写为
\begin{equation}\label{eq:5.159}
\BookImage{02_Unnumbered_Diagrams/Chapter_05/Ch05_Diagram_PDFp394_01_Bookp384.pdf}
\end{equation}

\SourceParagraph{ch05-5.5.7-011}{394}
$O$ 附近，$Z(\mathcal{O}_C)$ 很小，而 $Z(\mathcal{O}_x)=1$。因此，在 $\epsilon\to0$ 的近似中，$\xi(\mathcal{O}_x)=\xi(\mathcal{O}_C(-1)[1])=0$，$\xi(\mathcal{O}_C)=\theta/\pi-1$。\BookEqRef{eq:5.159}{5.159}中的态射，按箭头与箭尾的 $\xi$ 差标记，并计入 $[1]$。于是，相对于这个三角形，一个顶点 $\Pi$ 稳定，当且仅当对边的标记小于1。因此，当 $\theta$ 增大并越过 $\pi$，0膜 $\mathcal{O}_x$ 衰变为 $\mathcal{O}_C$ 与 $\mathcal{O}_C(-1)[1]$。

\SourceParagraph{ch05-5.5.7-012}{394}
还有正合三角形
\begin{equation}\label{eq:5.160}
\BookImage{02_Unnumbered_Diagrams/Chapter_05/Ch05_Diagram_PDFp394_02_Bookp384.pdf}
\end{equation}
它说明，当 $\theta<0$，$\mathcal{O}_x$ 衰变为 $\mathcal{O}_C(1)$ 与 $\mathcal{O}_C[1]$。无论哪种方向，都看到，从\BookRef{fig:5.11}{图11}的 $X$ 相，移入 $X'$ 相时，$\mathcal{O}_x$ 衰变。

\SourceParagraph{ch05-5.5.7-013}{394}
设 $y\notin C$。虽然 $\mathcal{O}_y$ 与 $\mathcal{O}_x$ 具有完全相同的 D 膜荷，它却不通过\BookEqRef{eq:5.159}{5.159}或\BookEqRef{eq:5.160}{5.160}衰变，因为没有从 $\mathcal{O}_C$ 或 $\mathcal{O}_C(1)$ 到 $\mathcal{O}_y$ 的态射。确实，也不预期远离 $C$ 的0膜，受到保典翻转跃迁的影响。

\SourceParagraph{ch05-5.5.7-014}{394}
还应当能够看到，$X$ 的导出范畴中，哪些对象，在保典翻转以后，充当 $C'$ 上的点。在这以前，需要引入计算相关 Ext 群的方法。设 $E,F$ 是 $C$ 上的层。给定嵌入 $i\colon C\to X$，法丛为 $C$ 上的 $N$，则有谱序列，
\begin{equation}\label{eq:5.161}
E_2^{p,q}=\operatorname{Ext}_C^p(E,F\otimes\Lambda^qN),
\end{equation}
汇聚到 $\operatorname{Ext}_X^n(i_*E,i_*F)$ 的滤过，其伴随分次为 $\bigoplus_{p+q=n}E_\infty^{p,q}$\footnote{一种快速证明方法，是使用 $i_*$ 的右伴随函子 $i^!$。}。

\SourceParagraph{ch05-5.5.7-015}{395}
\begin{example*}
考虑 $\operatorname{Ext}_X^n(\mathcal{O}_C,\mathcal{O}_C(-1))$。由于 $N=\mathcal{O}(-1)\oplus\mathcal{O}(-1)$，谱序列的 $E_2$ 页，由
\[
E_2^{p,q}=H^p(C,\mathcal{O}_C(-1)\otimes\Lambda^qN)
\]
给出，即
\begin{equation}\label{eq:5.162}
\begin{array}{c|cc}
q\backslash p&0&1\\\hline
2&0&\mathbb{C}^2\\
1&0&\mathbb{C}^2\\
0&0&0
\end{array}
\end{equation}
因此，
\begin{equation}\label{eq:5.163}
\operatorname{Ext}_X^n(\mathcal{O}_C,\mathcal{O}_C(-1))
=\begin{cases}
\mathbb{C}^2,&n=2\text{ 或 }3,\\
0,&\text{其他情形}.
\end{cases}
\end{equation}
\end{example*}

\SourceParagraph{ch05-5.5.7-016}{395}
于是，可以使用态射 $f\in\operatorname{Hom}(\mathcal{O}_C[-1],\mathcal{O}_C(-1)[1])\cong\mathbb{C}^2$，构造新对象 $D_f$：
\begin{equation}\label{eq:5.164}
\BookImage{02_Unnumbered_Diagrams/Chapter_05/Ch05_Diagram_PDFp395_01_Bookp385.pdf}
\end{equation}
当 $\theta$ 增大并越过 $\pi$，这些对象变得稳定。如\BookSectionRef{sec:5.3.2.2}{5.3.2.2}所指出，用非零复数缩放 $f$，不改变 $D_f$ 的同构类，因此，有一族由 $\mathbb{P}^1$ 参数化的 $D_f$。它们确实表示 $C'$ 上的点，在 $\theta$ 增大并越过 $\pi$、翻转进入 $X'$ 时，变得稳定，正如文献\BookCite{ref016}所论证。

\SourceParagraph{ch05-5.5.7-017}{395}
对象 $D_f$，与\BookSectionRef{sec:5.5.6.3}{5.5.6.3}的对象，在同一意义上是奇特的。文献\BookCite{ref016}把它们描述为“反常凝聚层”。

\SourceParagraph{ch05-5.5.7-018}{395}
可以与\BookRef{conj:4.81}{猜想4.81}比较；它已在三维情形证明，蕴含 $D(X)$ 与 $D(X')$ 等价。物理论证说明，只要 $X,X'$ 上的非线性 $\sigma$ 模型，可以通过改变 $B+\mathrm{i}J$ 连续联系，就有这一等价，因为它们给出同一个 B 模型。

\section{轨形 \texorpdfstring{$\mathbb{C}^3/\mathbb{Z}_3$}{C³/Z₃} 及其消解}\label{sec:5.6}

\SourceParagraph{ch05-5.6-intro-001}{396}
现在详细分析这一情形。三维轨形上 D 膜最初的大部分分析，都是针对这个最简单的例子进行的，例如，文献\BookCite{ref139,ref116,ref115,ref137}。

\SourceParagraph{ch05-5.6-intro-002}{396}
回顾\BookRef{ex:4.35}{例题4.35}。取 $g$ 为 $\mathbb{Z}_3$ 的生成元，作用为 $g\colon(z_1,z_2,z_3)\mapsto(\omega z_1,\omega z_2,\omega z_3)$，其中，$\omega=\exp(2\pi\mathrm{i}/3)$。令 $V_i$，$i=0,1,2$，为由 $\rho_i(g)=\omega^i$ 给出的一维不可约表示。那么，表示 $V=\bigoplus_i V_i\otimes\mathbb{C}^{m_i}$，对应于箭图
\begin{equation}\label{eq:5.165}
\BookImage{02_Unnumbered_Diagrams/Chapter_05/Ch05_Diagram_PDFp396_01_Bookp386.pdf}
\end{equation}

\SourceParagraph{ch05-5.6-intro-003}{396}
用 $\Delta_{m_0m_1m_2}$，表示具有\BookEqRef{eq:5.165}{5.165}形式及该维数向量的一个箭图表示。例如，分数膜为 $F_0=\Delta_{100}$、$F_1=\Delta_{010}$、$F_2=\Delta_{001}$。

\SourceParagraph{ch05-5.6-intro-004}{396}
众所周知，这个轨形的消解，具有例外除子 $E\cong\mathbb{P}^2$。这里，$X$ 可以看作对应于层 $\mathcal{O}_E(-3)$ 的线丛的总空间。这个空间具有 $h^{1,1}=1$，因此，弦性模空间为复一维。根据很快要讨论的论证，在分析周期与单值化时，可以把它看作删去三个特殊点的 $\mathbb{P}^1$。其中一个，对应于例外除子 $E$ 无限大的大半径极限。另一个极端，是没有爆破的轨形点。$\mathbb{P}^1$ 上第三个点，记为 $P_0$，人们认为 SCFT 在那里变得奇异，原因将在后面讨论。

\subsection{周期}\label{sec:5.6.1}

\SourceParagraph{ch05-5.6.1-intro-001}{396}
再次，可以用镜像对称，分析这一维模空间。$\mathbb{C}^3/\mathbb{Z}_3$ 及其消解的镜像，可以描述为 $\mathbb{C}^2\times(\mathbb{C}^*)^2$ 中的非紧三维流形
\[
xy=1+s+t+\frac{z}{st},
\]
$x,y$ 是 $\mathbb{C}^2$ 的坐标，$s,t$ 是 $(\mathbb{C}^*)^2$ 的坐标。此外，$z$ 给出一维复结构模参数，与消解中例外除子 $E$ 的体积互为镜像，参见文献\BookCite{ref243}。现在分析这个复结构模空间，特别是其中的 Gauss--Manin 联络，以找到周期。

\SourceParagraph{ch05-5.6.1-intro-002}{396}
相关的 Picard--Fuchs 方程为
\begin{equation}\label{eq:5.166}
\left(z\frac{\mathrm{d}}{\mathrm{d}z}\right)^3\Phi
+27z\left(z\frac{\mathrm{d}}{\mathrm{d}z}\right)
\left(z\frac{\mathrm{d}}{\mathrm{d}z}+\frac13\right)
\left(z\frac{\mathrm{d}}{\mathrm{d}z}+\frac23\right)\Phi=0.
\end{equation}
显然，任何常数都解这个微分方程。令
\[
z=(3\mathrm{e}^{-\pi\mathrm{i}}\psi)^{-3},
\]
则在 $\psi=0$ 附近，其余解的一组基，可写为
\begin{equation}\label{eq:5.167}
\varpi_j=\frac{1}{2\pi\mathrm{i}}\sum_{n=1}^{\infty}
\frac{\Gamma(n/3)\omega^{nj}}
{\Gamma(n+1)\Gamma(1-n/3)^2}(3\psi)^n,
\end{equation}
其中，$\omega=\exp(2\pi\mathrm{i}/3)$。因此，$1,\varpi_0,\varpi_1$，构成 Picard--Fuchs 方程解空间的一组基。

\SourceParagraph{ch05-5.6.1-intro-003}{397}
随后，作解析延拓，得到小 $z$ 时的一组基，适用于 $|\arg z|<\pi$。以下 Mellin--Barnes 积分，沿向上的竖直围道 $\mathcal{C}$，$-1/3<\operatorname{Re}s<0$，进行：
\begin{equation}\label{eq:5.168}
\begin{aligned}
\Phi_0&=1,\\
\Phi_1&=\frac{1}{2\pi\mathrm{i}}\frac{3}{2\pi\mathrm{i}}
\int_{\mathcal{C}}\frac{\Gamma(3s)\Gamma(-s)}{\Gamma(1+s)^2}z^s\,\mathrm{d}s\\
&=\frac{1}{2\pi\mathrm{i}}\log z+O(z)=t=\varpi_0,\\[4pt]
\Phi_2&=-\frac{1}{4\pi^2}\frac{-6}{2\pi\mathrm{i}}
\int_{\mathcal{C}}\frac{\Gamma(3s)\Gamma(-s)^2}{\Gamma(1+s)}
(\mathrm{e}^{-\pi\mathrm{i}}z)^s\,\mathrm{d}s\\
&=-\frac{1}{4\pi^2}(\log z-\mathrm{i}\pi)^2-\frac5{12}+O(z)\\
&=t^2-t-\frac16+O(\mathrm{e}^{2\pi\mathrm{i}t})\\
&=-\frac23(\varpi_0-\varpi_1).
\end{aligned}
\end{equation}
镜像映射为 $t=\int_C(B+\mathrm{i}J)=\tfrac{1}{2\pi\mathrm{i}}\log z+O(z)$，其中，$C$ 是 $E$ 中的一条超平面 $\mathbb{P}^1$。这里的 K\"ahler 坐标 $t$，与前面镜像方程中的坐标 $t$，按上下文区别。

\SourceParagraph{ch05-5.6.1-intro-004}{397}
$\mathbb{C}^3/\mathbb{Z}_3$ 轨形点附近的弦性 K\"ahler 模空间，如\BookRef{fig:5.12}{图12}所示。注意，轨形点恰好位于 $B=J=0$\footnote{文献\BookCite{ref019}中的一个错误假设，使 $B$ 移位了 $1/2$。正确论证见文献\BookCite{ref367}。}。
\begin{figure}[H]
\centering
\BookImage{01_Numbered_Figures/Chapter_05/Ch05_Fig12_PDFp398_Bookp388.pdf}
\caption{$\mathbb{Z}_3$ 轨形模空间的基本区域。每个由实线围起的阴影区域，是\BookEqRef{eq:5.168}{5.168}的映射 $\psi\mapsto t$ 下，平面 $\psi^3\in\mathbb{C}$ 的一个像。轨形单值化置换这三个阴影区域；无阴影部分，则分为无穷多个额外基本区域，由其他单值化联系。}\label{fig:5.12}
\end{figure}

\SourceParagraph{ch05-5.6.1-intro-005}{397}
使用\BookEqRef{eq:5.12}{5.12}与\BookEqRef{eq:5.71}{5.71}，匹配精确周期，就可以计算精确中心荷。例如，考虑缠绕例外除子的4膜 $\mathcal{O}_E(m)$：
\begin{equation}\label{eq:5.169}
Z(\mathcal{O}_E(m))
=-\left(m+\frac43\right)\varpi_0+\frac13\varpi_1
+\frac12m^2+\frac32m+\frac43.
\end{equation}

\subsubsection{无质量膜}\label{sec:5.6.1.1}

\SourceParagraph{ch05-5.6.1.1-001}{397}
考虑模空间中 $\psi=\mathrm{e}^{2\pi\mathrm{i}/3}$ 的点 $P_0$，此处有奇异共形场论。$P_0$ 处，$\varpi_0=t=\tfrac12+\mathrm{i}J_0$，其中，$J_0\approx0.4628$。由\BookEqRef{eq:5.167}{5.167}，$\varpi_1$ 在 $P_0$ 的值，显然等于 $\varpi_0$ 在 $\psi=\mathrm{e}^{4\pi\mathrm{i}/3}$ 的值，即 $-\tfrac12+\mathrm{i}J_0$。因此，$P_0$ 处有 $\varpi_0-\varpi_1=1$。由\BookEqRef{eq:5.169}{5.169}，$\mathcal{O}_E(-1)$ 在 $P_0$ 变为无质量。类似地，$\mathcal{O}_E(-2)$ 在 $\psi=\mathrm{e}^{4\pi\mathrm{i}/3}$ 处变为无质量。

\SourceParagraph{ch05-5.6.1.1-002}{398}
暂且假定，$P_0$ 处的共形场论奇点，是稳定 B 膜 $\mathcal{O}_E(-1)$ 变为无质量造成的。很快会检验这一点。

\subsection{单值化}\label{sec:5.6.2}

\SourceParagraph{ch05-5.6.2-intro-001}{398}
回顾\BookSectionRef{sec:5.5.3.1}{5.5.3.1}，沿模空间中的非平凡闭路径行走，产生导出范畴及稳定对象集合上的自等价。现在有一个多连通模空间的显式例子，可以更深入地探索这一点。

\SourceParagraph{ch05-5.6.2-intro-002}{398}
先讨论绕大半径极限的单值化。它对应于 $B\mapsto B+\omega$，其中，$\omega\in H^2(X,\mathbb{Z})$ 与某个除子 $D$ 对偶。根据\BookSectionRef{sec:3.5.2.7}{3.5.2.7}，这等价于丛曲率的 $F\mapsto F+\omega$，可由张量积
\[
\mathcal{E}\longmapsto\mathcal{E}\otimes\mathcal{O}_X(D)
\]
实现。用 Fourier--Mukai 变换的语言，这由核 $\mathcal{L}=\Delta_*\mathcal{O}_X(D)=\mathcal{O}_{\Delta_X}(D)$ 实现。也用 $\mathcal{L}$ 表示这个变换。

\SourceParagraph{ch05-5.6.2-intro-003}{399}
在这个非紧例子中，取 $D$ 与正在变得很大的 K\"ahler 形式分量 Poincar\'e 对偶。因此，要求 $D$ 与 $C$，即 $E$ 的超平面 $\mathbb{P}^1$，相交于一点。把 $\mathcal{O}_X(D)$ 也记为 $\mathcal{O}_X(1)$。这个记号，也与 $\mathcal{O}_X(1)|_E=\mathcal{O}_E(1)$ 一致。还应注意，由于 $E$ 的法丛对应于 $\mathcal{O}_E(-3)$，有短正合列
\begin{equation}\label{eq:5.170}
0\longrightarrow\mathcal{O}_X(3)\longrightarrow\mathcal{O}_X
\longrightarrow\mathcal{O}_E\longrightarrow0.
\end{equation}

\subsubsection{绕轨形点 \texorpdfstring{$\psi=0$}{ψ=0} 的单值化}\label{sec:5.6.2.1}

\SourceParagraph{ch05-5.6.2.1-001}{399}
一般而言，弦理论中的轨形，可以具有“量子对称”；它在由 $g\in G$ 扭曲的态上，通过乘以 $q(g)$ 作用，其中，
\begin{equation}\label{eq:5.171}
q\in\operatorname{Hom}(G,\mathbb{C}^*).
\end{equation}
因此，$\operatorname{Hom}(G,\mathbb{C}^*)$ 是量子对称群。它自然是 $G$ 的 Abel 化 $G/[G,G]$ 的特征标群；作为有限 Abel 群，它与 $G/[G,G]$ 非典范同构。另一个解释，是 $G$ 的一维表示群，群运算为张量积。给定一维表示 $U$，可以用 $U\otimes-$ 作用于表示集合。这给出 McKay 箭图的对称，因而明确说明，轨形量子对称如何作用于 D 膜范畴。

\SourceParagraph{ch05-5.6.2.1-002}{399}
若 $G$ 为 Abel 群，量子对称群与 $G$ 同构，并传递地作用于 McKay 箭图的顶点。这里，量子对称群 $\mathbb{Z}_3$，通过把 $V$ 与表示 $V_1$ 作张量积作用，换句话说，把 McKay 箭图旋转 $2\pi/3$。把这个变换记为 $\mathcal{G}$，则有
\[
\mathcal{G}\colon(F_0,F_1,F_2)\longmapsto(F_1,F_2,F_0).
\]

\SourceParagraph{ch05-5.6.2.1-003}{399}
回顾\BookSectionRef{sec:5.4.5.6}{5.4.5.6}对分数膜的识别。它依赖于腔室的选择，这些腔室也被 $\mathbb{Z}_3$ 作用置换。该选择，对应于把周期\BookEqRef{eq:5.168}{5.168}从大体积，解析延拓到轨形点的某个具体方式。为了检验这些选择是否正确匹配，可以先提出 $F_i$ 对应于某些大体积层的识别，由 K 理论类计算其大体积中心荷，再把表达式延拓回轨形点。正确识别，应当保持量子对称。

\SourceParagraph{ch05-5.6.2.1-004}{400}
\begin{example*}
采用\BookChapterRef{4}的\BookRef{tab:4.3}{表3}中，$\theta\in C_2$ 腔室的对应：
\begin{equation}\label{eq:5.172}
\BookImage{02_Unnumbered_Diagrams/Chapter_05/Ch05_Diagram_PDFp400_01_Bookp390.pdf}
\end{equation}
把它代入\BookEqRef{eq:5.168}{5.168}，沿 $\psi$ 平面的负实轴延拓，在轨形点附近，得到
\begin{equation}\label{eq:5.173}
\begin{aligned}
Z(F_0)&=\frac13(1-\varpi_0+\varpi_1),\\
Z(F_1)&=\frac13(1-\varpi_0-2\varpi_1),\\
Z(F_2)&=\frac13(1+2\varpi_0+\varpi_1).
\end{aligned}
\end{equation}
因此，轨形点处，所有 $F_i$ 的中心荷都为 $1/3$，与量子对称一致。这也确认\BookSectionRef{sec:5.6.1.1}{5.6.1.1}中的无质量膜识别。
\end{example*}

\SourceParagraph{ch05-5.6.2.1-005}{400}
注意，把 $\mathcal{G}$ 指定为置换 $F_j$ 的变换，已经很接近确定导出范畴的自等价；还须指定生成元之间态射上的相容作用，而量子对称的张量积构造，提供了这一数据。由 $\{F_0,F_1,F_2\}$ 三角生成的充满子范畴，通过第4章的等价，对应于支撑在轨形原点的有限维幂零箭图表示的导出范畴，也可用相应的矩阵分解语言描述。它包含所有支撑在例外纤维的凝聚层\footnote{这里不作进一步展开。不过，若考虑有限生成、而不要求有限维的箭图表示，可以得到完整导出范畴。一般紧支撑还可以包含轨形原点以外的有限支撑，因此，应与这里只由 $F_i$ 生成的子范畴区别。}。因此，$\mathcal{G}$ 至少在这个子范畴上，由上述完整的量子对称作用确定。

\subsubsection{绕锥形点的单值化}\label{sec:5.6.2.2}

\SourceParagraph{ch05-5.6.2.2-001}{400}
现在希望找到这一情形的 Fourier--Mukai 变换，记为 $\mathcal{K}$。它在 D 膜荷上的作用，可以从周期\BookEqRef{eq:5.168}{5.168}的单值化得到。固定使下面扭转函子对应于该单值化的环路方向；在所选 $C_2$ 腔室与 $P_0$ 的分支中，无质量对象为 $\mathcal{E}=\mathcal{O}_E(-1)$。用数学 Euler 配对 $\chi(\mathcal{E},F)$ 表示，荷的变换为
\begin{equation}\label{eq:5.174}
\begin{aligned}
\operatorname{ch}(\mathcal{K}(F))
&=\operatorname{ch}(F)-\chi(\mathcal{E},F)\operatorname{ch}(\mathcal{E})\\
&=\operatorname{ch}(F)+I(\mathcal{E},F)\operatorname{ch}(\mathcal{E}),
\end{aligned}
\end{equation}
其中，$I$ 是\BookEqRef{eq:5.106}{5.106}的相交形式；按该式的三维箭图约定，它与数学 Euler 配对相差一个负号。其他分支中的无质量对象，随相应的张量积或量子对称一起变换。

\SourceParagraph{ch05-5.6.2.2-002}{400}
与前面情形不同，这个点附近没有直接的 SCFT 物理定义，因此，没有直接给出具体 $\mathcal{K}$ 的物理论证。最直接的办法，也许是使用已有结果，以及三穿孔球面的基本群关系。在上述分数膜识别与环路约定下，关系写为
\begin{equation}\label{eq:5.175}
\mathcal{K}=\mathcal{L}\circ\mathcal{G}.
\end{equation}

\SourceParagraph{ch05-5.6.2.2-003}{400}
不过，遵循历史发展，先基于突变理论，对这一单值化提出猜想\BookCite{ref310}；突变理论参见文献\BookCite{ref406}。随后，文献\BookCite{ref422}广泛研究了它，文献\BookCite{ref245,ref021}作了重要推广。再用\BookEqRef{eq:5.175}{5.175}检验。

\SourceParagraph{ch05-5.6.2.2-004}{401}
这一思想，适用于任意由无质量膜 $\mathcal{E}$ 诱导的单值化，因此，在这个一般性下定义它，而不限于当前的例子。对于 $A,B\in D(X)$，定义外张量积 $A\boxtimes B=\mathrm{L}p_1^*A\otimes^{\mathrm{L}}\mathrm{L}p_2^*B$，其中，$p_1,p_2\colon X\times X\to X$，是直积的第一与第二投影。
\BookStatementNumber{33}
\begin{definition}\label{def:5.33}
扭转函子 $\Phi_{\mathcal{E}}$，定义为
\begin{equation}\label{eq:5.176}
\Phi_{\mathcal{E}}(F)=\operatorname{Cone}\!\left(
\mathrm{R}\operatorname{Hom}_X(\mathcal{E},F)
\otimes^{\mathrm{L}}\mathcal{E}\xrightarrow{r}F\right),
\end{equation}
其中，$r$ 是取值映射，或“重言映射”，由映射复形的元素在 $\mathcal{E}$ 上的作用给出。
\end{definition}

\SourceParagraph{ch05-5.6.2.2-005}{401}
它对应于核为
\begin{equation}\label{eq:5.177}
\mathcal{K}=\operatorname{Cone}\!\left(
\mathcal{E}^{\vee}\boxtimes\mathcal{E}\xrightarrow{r}
\mathcal{O}_{\Delta_X}\right)
\end{equation}
的 Fourier--Mukai 变换。这里，$\mathcal{E}^{\vee}$ 是在导出范畴中定义的对偶 $\mathrm{R}\mathcal{H}\!om(\mathcal{E},\mathcal{O}_X)$，$r$ 是自然的限制与取值映射。也用 $\mathcal{K}$ 表示该核诱导的 Fourier--Mukai 变换，并说明，它与扭转函子 $\Phi_{\mathcal{E}}$ 相同。有
\[
\begin{aligned}
\mathcal{K}(F)
&=\mathrm{R}p_{2*}\operatorname{Cone}\!\left(
(F\otimes^{\mathrm{L}}\mathcal{E}^{\vee})\boxtimes\mathcal{E}
\to\mathcal{O}_{\Delta_X}\otimes^{\mathrm{L}}\mathrm{L}p_1^*F\right)\\
&\cong\operatorname{Cone}\!\left(
\mathrm{R}\Gamma(X,F\otimes^{\mathrm{L}}\mathcal{E}^{\vee})
\otimes^{\mathrm{L}}\mathcal{E}\to F\right)\\
&\cong\operatorname{Cone}\!\left(
\mathrm{R}\operatorname{Hom}_X(\mathcal{E},F)
\otimes^{\mathrm{L}}\mathcal{E}\to F\right).
\end{aligned}
\]
因此，
\begin{equation}\label{eq:5.178}
\mathcal{K}=\Phi_{\mathcal{E}}.
\end{equation}

\SourceParagraph{ch05-5.6.2.2-006}{401}
注意，根据构造，这个 Fourier--Mukai 变换给出
\[
\operatorname{ch}(\mathcal{K}(F))
=\operatorname{ch}(F)-\chi(\mathcal{E},F)\operatorname{ch}(\mathcal{E}),
\]
推广了\BookEqRef{eq:5.174}{5.174}。在把无质量对象识别为例外除子结构层的适当分支中，取 $\mathcal{E}=\mathcal{O}_E$；在当前 $P_0$ 分支中，则取 $\mathcal{O}_E(-1)$。

\SourceParagraph{ch05-5.6.2.2-007}{401}
另一个简单检验，是 $\mathcal{K}$ 的作用，应使 $\xi(\mathcal{E})$ 改变两个单位。按当前扭转约定，确实有
\[
\mathcal{K}(\mathcal{E})=\mathcal{E}[-2],
\]
因而分次改变 $-2$；反向环路使用逆函子，给出 $\mathcal{E}[2]$。这要求 $\mathcal{E}$ 满足进一步条件，即
\begin{equation}\label{eq:5.179}
\operatorname{Ext}_X^k(\mathcal{E},\mathcal{E})
=\begin{cases}
\mathbb{C},&k=0\text{ 或 }3,\\
0,&\text{其他情形}.
\end{cases}
\end{equation}

\SourceParagraph{ch05-5.6.2.2-008}{402}
这种对象，称为“球面”对象，因为其自 Ext 群，与 $S^3$ 的上同调相同。这意味着，$\mathrm{R}\operatorname{Hom}_X(\mathcal{E},\mathcal{E})$，可以用如下向量空间复形表示：
\begin{equation}\label{eq:5.180}
\cdots\to0\to\underset{\scriptscriptstyle0}{\mathbb{C}}
\xrightarrow{0}0\xrightarrow{0}0\xrightarrow{0}
\underset{\scriptscriptstyle3}{\mathbb{C}}\to0\to\cdots.
\end{equation}
与 $\mathcal{E}$ 作张量积，只是把各个 $\mathbb{C}$ 换成 $\mathcal{E}$。锥构造把源复形左移一次；由于取值映射在次数零的 $\mathcal{E}$ 上是恒等，第一个 $\mathcal{E}$ 与锥右边的 $\mathcal{E}$ 抵消，剩下 $\mathcal{E}[-2]$。

\SourceParagraph{ch05-5.6.2.2-009}{402}
在上述球面条件、适当的完美性与支撑条件，以及 Calabi--Yau 对偶配对下，可以证明，$\Phi_{\mathcal{E}}$ 是自等价\BookCite{ref422}，因此，\BookEqRef{eq:5.178}{5.178}是所需单值化的候选。为确认这一点，可以计算 $\mathcal{L}^{-1}\circ\mathcal{K}$ 在 $\{F_0,F_1,F_2\}$，以及相应生成元态射上的作用。它按与 $\mathcal{G}$ 相同的方式置换这些对象，从而确认\BookEqRef{eq:5.175}{5.175}，至少在相应的有限维箭图表示子范畴上。

\SourceParagraph{ch05-5.6.2.2-010}{402}
值得提一下，\BookSectionRef{sec:5.5.6}{5.5.6}中五次三维流形的类似分析，有一个微妙区别。仍可分别定义绕大半径极限与锥形点的单值化 $\mathcal{L},\mathcal{K}$，并在相应约定下令 $\mathcal{G}=\mathcal{L}^{-1}\circ\mathcal{K}$。由于 Gepner 点的量子 $\mathbb{Z}_5$ 对称，似乎会预期 $\mathcal{G}^5=1$。实际却发现，$\mathcal{G}^5$ 对应于导出范畴中的移位“$[2]$”。可以用矩阵分解理解这一点\BookCite{ref459}，不过，这里不再展开。

\subsection{稳定性的例子}\label{sec:5.6.3}

\SourceParagraph{ch05-5.6.3-intro-001}{402}
五次三维流形的情形中，用大半径极限作为确定稳定性的基点。轨形的情形中，可以利用量子对称，把轨形点用作基点。

\SourceParagraph{ch05-5.6.3-intro-002}{402}
首先，由\BookEqRef{eq:5.173}{5.173}，可以计算轨形点附近的 $\xi$。如\BookSectionRef{sec:5.4.5}{5.4.5}末尾所论证，轨形点处，所有分数膜具有相同 $\xi$；规定这个值为零。附近，在 $\psi$ 的线性阶，有
\begin{equation}\label{eq:5.181}
\begin{aligned}
\xi(\Delta_{m_0m_1m_2})
&=-c\,\frac{(-m_0-m_1+2m_2)\operatorname{Re}\psi
+(m_0-2m_1+m_2)\operatorname{Re}(\omega\psi)}
{m_0+m_1+m_2}\\
&=-c\,\frac{\sum_i m_i\zeta_i}{\sum_i m_i},
\end{aligned}
\end{equation}
其中，$c>0$。按\BookSectionRef{sec:5.4.5.5}{5.4.5.5}的矩映射与 King 权的符号约定，定义
\begin{equation}\label{eq:5.182}
\zeta_k=-\sqrt3\operatorname{Re}\!\left(
\mathrm{e}^{\frac{\pi\mathrm{i}}6(4k-1)}\psi\right),
\end{equation}
于是，
\begin{equation}\label{eq:5.183}
\begin{aligned}
\zeta_0+\zeta_1+\zeta_2&=0,\\
\zeta_0+\omega\zeta_1+\omega^2\zeta_2
&=-\frac{3\sqrt3}{2}\mathrm{e}^{\pi\mathrm{i}/6}\bar\psi,\\
\zeta_0+\omega^2\zeta_1+\omega\zeta_2
&=-\frac{3\sqrt3}{2}\mathrm{e}^{-\pi\mathrm{i}/6}\psi.
\end{aligned}
\end{equation}

\SourceParagraph{ch05-5.6.3-intro-003}{403}
显然，这类似于\BookEqRef{eq:5.110}{5.110}，刚才引入的 $\zeta$，对应于\BookSectionRef{sec:5.4.5.5}{5.4.5.5}的参数。现在可以明确说明，在轨形点附近，$\theta$ 稳定性是 $\Pi$ 稳定性的极限形式。设有箭图表示的短正合列，
\begin{equation}\label{eq:5.184}
\BookImage{02_Unnumbered_Diagrams/Chapter_05/Ch05_Diagram_PDFp403_01_Bookp393.pdf}
\end{equation}
轨形点附近，这三个 D 膜的 $\xi$ 都很接近零。由于中心荷相加，\BookEqRef{eq:5.184}{5.184}中间对象的 $\xi$，必须位于另外两个的 $\xi$ 之间。为判断中间对象的 $\Pi$ 稳定性，画出正合三角形，
\begin{equation}\label{eq:5.185}
\BookImage{02_Unnumbered_Diagrams/Chapter_05/Ch05_Diagram_PDFp403_02_Bookp393.pdf}
\end{equation}
并要求 $f$ 为快子。由\BookEqRef{eq:5.181}{5.181}，这恰好是
\begin{equation}\label{eq:5.186}
\frac{\sum_i n_i\zeta_i}{\sum_i n_i}
>\frac{\sum_i m_i\zeta_i}{\sum_i m_i}.
\end{equation}
结合\BookEqRef{eq:5.118}{5.118}，这等价于\BookSectionRef{sec:5.4.5.5}{5.4.5.5}中，King 的 $\theta$ 稳定性陈述。

\SourceParagraph{ch05-5.6.3-intro-004}{403}
这种 $\theta$ 稳定性表述，允许完整分类轨形点附近稳定的不可约 B 膜。如\BookSectionRef{sec:5.4.4}{5.4.4}所提到的，在稳定对象的范围内，不可约对象 $A$ 满足 $\operatorname{Hom}(A,A)=\mathbb{C}$。满足这个自同态条件的箭图表示，称为“Schur 表示”。文献\BookCite{ref137}讨论了寻找这类表示的问题。

\SourceParagraph{ch05-5.6.3-intro-005}{403}
判断一个带关系的箭图表示，是否为 Schur 表示，纯粹是代数问题，但可能很棘手。在许多情形中，用\BookSectionRef{sec:4.7}{4.7}的等价，把问题改写为凝聚层问题，反而更方便。

\SourceParagraph{ch05-5.6.3-intro-006}{403}
\begin{example*}
作为非 Schur 箭图表示的例子，考虑各条箭上的映射一般的 $\Delta_{211}$。稍作计算，可以说明短正合列
\begin{equation}\label{eq:5.187}
\BookImage{02_Unnumbered_Diagrams/Chapter_05/Ch05_Diagram_PDFp403_03_Bookp393.pdf}
\end{equation}
分裂。这立刻说明，$\operatorname{Hom}(\Delta_{211},\Delta_{211})\supset\mathbb{C}^2$。用层来表述，这更明显：$\Delta_{111}$ 是0膜，一般位于远离例外除子 $E$ 的地方；$\Delta_{100}$ 则是缠绕 $E$ 的4膜 $\mathcal{O}_E(-1)$。因此，$\Delta_{211}$ 是两个互不相交的 D 膜的和，显然可约。注意，箭上的映射不一般时，这个表示实际上可能是 Schur 表示。这对应于0膜位于 $E$ 上，使它有可能与4膜形成不可约束缚态。
\end{example*}

\SourceParagraph{ch05-5.6.3-intro-007}{404}
这里不尝试明确给出完整分类，不过，一些有意义的 Schur 表示，包括适当一般映射下的 $\Delta_{111}$，以及 $\Delta_{abc}$，其中，$\{a,b,c\}$ 是 $\{0,1,n\}$ 的任意排列，$0\le n\le3$。

\SourceParagraph{ch05-5.6.3-intro-008}{404}
分数膜 $F_k$，在轨形点附近显然始终稳定，因为在箭图表示范畴中，它们没有非零真子对象。接下来集中于 $\Delta_{111}$，其中一些对应于0膜。箭图
\begin{equation}\label{eq:5.188}
\BookImage{02_Unnumbered_Diagrams/Chapter_05/Ch05_Diagram_PDFp404_01_Bookp394.pdf}
\end{equation}
若每对相邻顶点之间，至少有一个非零映射，就没有非零真子表示，因此稳定。它表示远离例外除子 $E$ 的0膜；若尚未爆破，则远离轨形原点。确实，预期这种0膜的稳定性，不受轨形原点附近现象的影响。

\SourceParagraph{ch05-5.6.3-intro-009}{404}
现在考虑短正合列：
\begin{equation}\label{eq:5.189}
\BookImage{02_Unnumbered_Diagrams/Chapter_05/Ch05_Diagram_PDFp404_02_Bookp394.pdf}
\end{equation}
按当前的矩映射参数约定，这个 $\Delta_{111}$ 对衰变为 $\Delta_{100}$ 稳定，要求 $\zeta_0>0$。还有短正合列
\begin{equation}\label{eq:5.190}
\BookImage{02_Unnumbered_Diagrams/Chapter_05/Ch05_Diagram_PDFp404_03_Bookp394.pdf}
\end{equation}
给出进一步约束 $\zeta_2<0$。因此，这个0膜，在从轨形点伸出的一个 $2\pi/3$ 扇区中稳定。向这个扇区内稍微爆破以后，它对应于例外除子上的一点。显然，把零映射循环地移到箭图的另一条边，会得到参数循环置换后的类似陈述。因此，另外两类在其他边上切断环路的 $\Delta_{111}$，在扇区 $\zeta_0>0,\zeta_2<0$ 中不稳定，不能在那里作为这个消解的0膜；\BookEqRef{eq:5.188}{5.188}中远离例外除子的表示，仍然稳定。

\SourceParagraph{ch05-5.6.3-intro-010}{404}
左边上的映射为零的箭图，与 $E$ 上的层有紧密联系。对于这种一般表示 $\Delta_{abc}$，有短正合列
\begin{equation}\label{eq:5.191}
\BookImage{02_Unnumbered_Diagrams/Chapter_05/Ch05_Diagram_PDFp404_04_Bookp394.pdf}
\end{equation}
于是，迭代得到
\begin{equation}\label{eq:5.192}
\BookImage{02_Unnumbered_Diagrams/Chapter_05/Ch05_Diagram_PDFp405_01_Bookp395.pdf}
\end{equation}

\SourceParagraph{ch05-5.6.3-intro-011}{405}
继续这一过程，得到
\begin{equation}\label{eq:5.193}
\BookImage{02_Unnumbered_Diagrams/Chapter_05/Ch05_Diagram_PDFp405_02_Bookp395.pdf}
\end{equation}
这显式地把箭图表示，映到 $X$，或 $E$，上的凝聚层导出范畴中。它恰好是\BookSectionRef{sec:4.6.2}{4.6.2}中的 Beilinson 构造，即 $\mathbb{P}^2$ 上层的构造；文献\BookCite{ref137}识别了这一对应。因此，左边上的映射为零的箭图表示，与 $E$ 上的 D 膜对应。当然，只有当\BookEqRef{eq:5.193}{5.193}中复形的上同调集中于一项时，才得到单个层，允许相应移位。

\SourceParagraph{ch05-5.6.3-intro-012}{405}
\BookRef{fig:5.13}{图13}画出一些 $\Pi$ 稳定性的边际稳定性曲线。每种情形中，箭头表示，沿哪个方向越过曲线，会使相应对象衰变。自然地，在原点附近，它与 $\theta$ 稳定性一致。图用复 $-\mathrm{i}\psi$ 平面表示模空间，这样选择，是为了与\BookRef{fig:5.12}{图12}对齐，即，大半径极限在上方。注意，对应于 $p\in E$ 的 $\mathcal{O}_p$，也就是前面相应的 $\Delta_{111}$，其边际稳定性曲线，即使在计入 $\Pi$ 稳定性的非微扰效应以后，也沿着 $\arg\psi$ 的等值线。
\begin{figure}[H]
\centering
\BookImage{01_Numbered_Figures/Chapter_05/Ch05_Fig13_PDFp406_Bookp396.pdf}
\caption{$\mathbb{Z}_3$ 轨形的一些边际稳定性曲线。}\label{fig:5.13}
\end{figure}

\SourceParagraph{ch05-5.6.3-intro-013}{405}
一些值得注意的衰变如下。
\begin{enumerate}[label=(\arabic*)]
\item $\mathcal{O}_C(-1)$：这个层，放入正合列
\begin{equation}\label{eq:5.194}
0\longrightarrow\mathcal{O}_E(-2)\longrightarrow\mathcal{O}_E(-1)
\longrightarrow\mathcal{O}_C(-1)\longrightarrow0.
\end{equation}
因此，它按 $\Pi$ 稳定性衰变，方式基本与\BookSectionRef{sec:5.5.6.1}{5.5.6.1}中五次三维流形上的4膜相同。这些2膜，在大半径处稳定，却在到达轨形点以前衰变。注意，\BookEqRef{eq:5.194}{5.194}说明，在箭图表示导出范畴中，有
\begin{equation}\label{eq:5.195}
\mathcal{O}_C(-1)=\operatorname{Cone}(F_2[-2]\to F_0).
\end{equation}
也就是说，这个 D 膜，本质上是箭图的复形，不能表示为单个箭图表示。换句话说，它不在箭图表示的 Abel 范畴中。因此，它在接近轨形点以前衰变，与整个图景一致。

\SourceParagraph{ch05-5.6.3-intro-014}{406}
\item $\Delta_{101}$：遵循\BookSectionRef{sec:5.5.6.3}{5.5.6.3}的思路，可以取\BookEqRef{eq:5.195}{5.195}的“Serre 对偶”，寻找奇特 D 膜。这给出
\[
\operatorname{Cone}(F_0[-1]\to F_2)\text{，}
\]
即 $F_0$ 被 $F_2$ 的扩张。这恰好是 $\Delta_{101}$。如\BookSectionRef{sec:5.5.6.3}{5.5.6.3}所预期，这些对象，在大半径处不稳定；不过，缩小例外除子时，可以变得稳定。边际稳定性曲线如\BookRef{fig:5.13}{图13}所示。注意，实际上，必须缩小到轨形点，或越过轨形点，它们才变得稳定。这些对象，在三角箭图左边上，一般有非零映射，所以，不由 Beilinson 构造\BookEqRef{eq:5.193}{5.193}分类。

D 膜 $\Delta_{101}$ 与 $\mathcal{O}_C(-1)$ 之间，有很好的互补性：$\mathcal{O}_C(-1)$ 是凝聚层范畴的对象，却是箭图意义下的复形；$\Delta_{101}$ 是箭图表示范畴的对象，却在层的导出范畴中，成为奇特复形 $\operatorname{Cone}(\mathcal{O}_E(-1)[-1]\to\mathcal{O}_E(-2)[2])$。

\SourceParagraph{ch05-5.6.3-intro-015}{407}
\item $\Delta_{n10}$：它放入短正合列
\begin{equation}\label{eq:5.196}
\BookImage{02_Unnumbered_Diagrams/Chapter_05/Ch05_Diagram_PDFp407_01_Bookp397.pdf}
\end{equation}
产生\BookRef{fig:5.13}{图13}所示的衰变。识别\BookEqRef{eq:5.193}{5.193}，连同短正合列
\begin{equation}\label{eq:5.197}
0\longrightarrow\Omega_E\longrightarrow\mathcal{O}_E(-1)^{\oplus3}
\longrightarrow\mathcal{O}_E\longrightarrow0,
\end{equation}
可以用来说明，
\begin{equation}\label{eq:5.198}
\Delta_{210}\cong\mathcal{O}_C,\qquad
\Delta_{310}\cong\mathcal{O}_E.
\end{equation}
因此，这些 D 膜，在箭图与层两种语言中，都是单个对象。所以，它们在大半径极限与轨形点附近，都可以稳定，并不意外。

另一方面，$\Delta_{110}$ 不对应于单个层，因为\BookEqRef{eq:5.193}{5.193}的复形，在不止一个位置具有上同调。这时，有短正合列
\begin{equation}\label{eq:5.199}
\BookImage{02_Unnumbered_Diagrams/Chapter_05/Ch05_Diagram_PDFp407_02_Bookp397.pdf}
\end{equation}
它使 $\Delta_{110}$ 在大半径极限中，只是边际稳定的。

\SourceParagraph{ch05-5.6.3-intro-016}{407}
\item $\Delta_{0n1}$：这类似于 $\Delta_{n10}$ 的情形，\BookRef{fig:5.13}{图13}也画出相应边际稳定性曲线。这次，在\BookEqRef{eq:5.172}{5.172}的识别下，$\Delta_{011}$ 对应于点理想层的移位 $\mathcal{I}_{E,p}[1]$，同样，在大半径极限中只边际稳定。
\end{enumerate}

\section{稳定性结构}\label{sec:5.7}

\SourceParagraph{ch05-5.7-intro-001}{407}
现在已经得到各种要求，物理 D 膜作为拓扑边界条件中的子集，必须满足它们。Bridgeland 在文献\BookCite{ref065}中，提出稳定性条件的精确数学定义。虽然受到物理讨论启发，它本身自成一体，因而可以纯粹从数学上分析。特别地，在某些合理假设下，可以说明，一个合理三角范畴的稳定性条件空间，是有限维复流形，从而为这种范畴提供有趣的几何不变量。

\SourceParagraph{ch05-5.7-intro-002}{407}
截至目前的证据，提示这一稳定性概念，确实描述前面讨论的“物理 D 膜”或 SCFT 边界态，但尚未证明。事实上，它与弦理论讨论，有一些有趣差别：定义不要求 Calabi--Yau 条件，稳定性条件空间的维数，也一般高于 SCFT 的 K\"ahler 模空间。\BookSectionRef{sec:5.9}{5.9}将再讨论这些问题。另一方面，这一概念，在数学上已非常富有成果，广泛出现，并把看起来不同的领域联系起来。

\SourceParagraph{ch05-5.7-intro-003}{408}
来看 Bridgeland 的定义。
\BookStatementNumber{34}
\begin{definition}\label{def:5.34}
三角范畴 $D$ 上的稳定性条件 $\sigma=(Z,\mathcal{P})$，由群同态 $Z\colon K(D)\to\mathbb{C}$，称为中心荷，以及对每个 $\phi\in\mathbb{R}$ 指定的充满加法子范畴 $\mathcal{P}(\phi)\subset D$ 组成，满足下列公理：
\begin{enumerate}[label=(\alph*)]
\item 若 $0\ne E\in\mathcal{P}(\phi)$，则 $Z(E)=m(E)\exp(\mathrm{i}\pi\phi)$，其中，$m(E)\in\mathbb{R}_{>0}$。
\item 对每个 $\phi\in\mathbb{R}$，$\mathcal{P}(\phi+1)=\mathcal{P}(\phi)[1]$。
\item 若 $\phi_1>\phi_2$，且 $A_j\in\mathcal{P}(\phi_j)$，则 $\operatorname{Hom}_D(A_1,A_2)=0$。
\item 对每个非零对象 $E\in D$，存在有限实数序列
\[
\phi_1>\phi_2>\cdots>\phi_n
\]
以及一组正合三角形，
\[
\BookImage{02_Unnumbered_Diagrams/Chapter_05/Ch05_Diagram_PDFp408_01_Bookp398.pdf}
\]
其中，对每个 $j$，都有 $0\ne A_j\in\mathcal{P}(\phi_j)$。
\end{enumerate}
\end{definition}

\SourceParagraph{ch05-5.7-intro-004}{408}
与\BookSectionRef{sec:5.5}{5.5}比较，主要的新要素，也许是公理(d)。它断言，存在 Harder--Narasimhan 滤过，类似于稳定凝聚层的讨论\BookCite{ref220,ref327}。虽然这一主张，在物理上合理，如\BookRef{conj:5.32}{猜想5.32}，却不完全显然，\BookSectionRef{sec:5.9.3}{5.9.3}将讨论。

\SourceParagraph{ch05-5.7-intro-005}{408}
给定定义中的稳定性条件 $\sigma=(Z,\mathcal{P})$，每个子范畴 $\mathcal{P}(\phi)$ 都是 Abel 范畴，下面将说明。$\mathcal{P}(\phi)$ 的非零对象，称为 $\sigma$ 中相位为 $\phi$ 的半稳定对象；它的单对象，称为稳定对象。由其他公理，公理(d)给出的非零对象 $E$ 的分解，在同构意义下唯一。记 $\phi_\sigma^+(E)=\phi_1$、$\phi_\sigma^-(E)=\phi_n$。$E$ 的质量，定义为正实数
\[
m_\sigma(E)=\sum_i|Z(A_i)|.
\]
由三角不等式，对每个非零 $E$，有
\[
m_\sigma(E)\ge|Z(E)|.
\]
等号成立，当且仅当所有因子的中心荷，位于同一条正实射线上。在相位取于同一适当分支，例如同一个心的相位区间时，这等价于 $E$ 半稳定；这与 BPS 条件\BookEqRef{eq:5.18}{5.18}对应。

\SourceParagraph{ch05-5.7-intro-006}{408}
对任意区间 $I\subset\mathbb{R}$，把 $\mathcal{P}(I)$ 定义为，由 $\phi\in I$ 的各个 $\mathcal{P}(\phi)$ 生成的、对扩张封闭的子范畴。因此，例如，充满子范畴 $\mathcal{P}((a,b))$，由零对象，以及满足 $a<\phi_\sigma^-(E)\le\phi_\sigma^+(E)<b$ 的非零对象组成。

\subsection{与 Abel 范畴中稳定性的关系}\label{sec:5.7.1}

\SourceParagraph{ch05-5.7.1-001}{409}
物理讨论，大部分基于这样一种思想：在弦性 K\"ahler 模空间的特定区域，导出范畴的某个 Abel 子范畴，会成为优先选择。例如，大体积极限中的凝聚层，或场论极限中的箭图表示。当前框架中，这由下面的命题精确表述。
\BookStatementNumber{35}
\begin{proposition}[\BookCite{ref065}，命题5.3]\label{prop:5.35}
在 $D$ 上给出稳定性条件，等价于给出 $D$ 上的有界 $t$ 结构，以及其心上具有 Harder--Narasimhan 性质的稳定性函数。
\end{proposition}

\SourceParagraph{ch05-5.7.1-002}{409}
如\BookSectionRef{sec:4.4.6}{4.4.6}所述，$t$ 结构确定 $D$ 的一个 Abel 子范畴 $\mathcal{A}$，即它的心。
\begin{definition*}[稳定性函数]
Abel 范畴 $\mathcal{A}$ 上的稳定性函数，是群同态 $Z\colon K(\mathcal{A})\to\mathbb{C}$，满足
\[
0\ne E\in\mathcal{A}\quad\Longrightarrow\quad
Z(E)\in\mathbb{R}_{>0}\exp(\mathrm{i}\pi\phi(E)),
\qquad 0<\phi(E)\le1.
\]
实数 $\phi(E)\in(0,1]$，称为对象 $E$ 的相位。
\end{definition*}
相位允许排序对象，从而定义稳定性：若非零对象 $E\in\mathcal{A}$ 的每个非零子对象 $A\subset E$，都满足 $\phi(A)\le\phi(E)$，就称 $E$ 对 $Z$ 半稳定。

\SourceParagraph{ch05-5.7.1-003}{409}
现在，非零对象 $E\in\mathcal{A}$ 的 Harder--Narasimhan 滤过，是有限子对象链
\[
0=E_0\subset E_1\subset\cdots\subset E_{n-1}\subset E_n=E,
\]
其因子 $F_j=E_j/E_{j-1}$，是 $\mathcal{A}$ 的半稳定对象，满足
\[
\phi(F_1)>\phi(F_2)>\cdots>\phi(F_n).
\]

\SourceParagraph{ch05-5.7.1-004}{409}
注意，若 $f\colon E\to F$ 是半稳定对象之间的非零映射，通过通常地考虑 $\operatorname{im}f\cong\operatorname{coim}f$，就得到 $\phi(E)\le\phi(F)$。由此，容易说明，Harder--Narasimhan 滤过若存在，就唯一。
\BookStatementNumber{36}
\begin{definition}\label{def:5.36}
若 $\mathcal{A}$ 中每个非零对象，都有 Harder--Narasimhan 滤过，就称稳定性函数 $Z$ 具有 Harder--Narasimhan 性质。
\end{definition}

\SourceParagraph{ch05-5.7.1-005}{409}
遵循 Rudakov\BookCite{ref407}的思想，可以证明\BookCite{ref065}，命题2.4，这来自对 $\mathcal{A}$ 的一个相当弱的条件，禁止相关类型的无穷升链或降链。

\SourceParagraph{ch05-5.7.1-006}{409}
现在可以证明\BookRef{prop:5.35}{命题5.35}。
\begin{proof}
给定 $D$ 上的稳定性条件，中心荷在心 $\mathcal{A}=\mathcal{P}((0,1])\subset D$ 上，定义稳定性函数，公理(d)的分解给出 Harder--Narasimhan 滤过。

\SourceParagraph{ch05-5.7.1-007}{410}
反过来，给定 $D$ 上的有界 $t$ 结构，以及其心 $\mathcal{A}\subset D$ 上具有 Harder--Narasimhan 性质的稳定性函数 $Z$，对每个 $\phi\in(0,1]$，把 $\mathcal{P}(\phi)$ 定义为 $\mathcal{A}$ 中相位 $\phi$ 的半稳定对象，加上零对象。公理(b)，随后确定所有 $\phi\in\mathbb{R}$ 的 $\mathcal{P}(\phi)$。稳定性函数给出(a)，半稳定对象间的相位约束与 $t$ 结构给出(c)，有界上同调塔及其有限 Harder--Narasimhan 滤过给出(d)。
\end{proof}

\SourceParagraph{ch05-5.7.1-008}{410}
同样的论证说明，任意 $\mathcal{A}=\mathcal{P}((\phi-1,\phi])$，都是 Abel 范畴。此外，各个 $\mathcal{P}(\phi)$ 本身，也是 Abel 范畴\BookCite{ref065}，命题5.2。为说明这一点，只须证明，$\mathcal{P}(\phi)$ 中态射 $f\colon E\to F$，作为 $\mathcal{A}=\mathcal{P}((\phi-1,\phi])$ 的态射时，其核与余核仍属于 $\mathcal{P}(\phi)$。半稳定性首先约束像与商的相位，再结合中心荷的可加性，迫使这些非零核与余核，也具有相位 $\phi$；这个心中不含相位高于 $\phi$ 的对象。于是，短正合列中两项属于 $\mathcal{P}(\phi)$ 时，第三项也属于它。另一方面，长度小于1的区间 $I$，其 $\mathcal{P}(I)$，只需是准 Abel 范畴，定义见文献\BookCite{ref411,ref065}。

\SourceParagraph{ch05-5.7.1-009}{410}
\begin{definition*}[局部有限性]
若存在 $\epsilon>0$，使每个准 Abel 范畴 $\mathcal{P}((\phi-\epsilon,\phi+\epsilon))$ 都具有有限长度，就称稳定性条件局部有限。这里的子对象与商，取准 Abel 范畴中的严格子对象与严格商；有限长度意味着，相关无穷链最终稳定。
\end{definition*}
下面的一些结果，假定这一条件。

\SourceParagraph{ch05-5.7.1-010}{410}
虽然可以构造不满足局部有限性的稳定性条件\BookCite{ref065}，例题5.6，它们比较病态。物理上，预期局部有限性，与\BookSectionRef{sec:5.4.4}{5.4.4}的主张相联系，即，小相位区间 $I$ 内的稳定对象，可以由有效场论描述。它也与质量隙的思想有关：在当前考虑的 $g_s\to0$ 极限中，非奇异弦紧化的稳定对象中心荷，应当与零保持正距离。

\SourceParagraph{ch05-5.7.1-011}{410}
与\BookSectionRef{sec:5.3.3}{5.3.3}及\BookSectionRef{sec:5.5}{5.5}比较。那里看到，自然的物理定义，从 SCFT 及相应物理膜集合出发，用它构造同伦范畴，最终构造导出范畴。若知道最初的物理膜构成 Abel 范畴，这个过程没有歧义。反过来，刚才的结果说明，在导出范畴上赋予稳定性结构，会得到一组优先选择的 Abel 范畴，即每个 $\phi$ 对应的 $\mathcal{P}((\phi-1,\phi])$。此外，在稳定性结构空间中移动时，根据公理(c)，这些 Abel 范畴按 $\Pi$ 稳定性的规则改变。现在可以回答\BookSectionRef{sec:5.5}{5.5}的问题：这一条件，在多大程度上确定稳定膜的集合？

\subsection{稳定性条件空间}\label{sec:5.7.2}

\SourceParagraph{ch05-5.7.2-intro-001}{410}
为了定义这个空间，先在 $D$ 上局部有限稳定性条件的集合 $\operatorname{Stab}(D)$ 上，赋予拓扑。它由扩充距离
\[
\begin{aligned}
d(\sigma_1,\sigma_2)=\sup_{0\ne E\in D}\max\biggl\{
&|\phi_{\sigma_2}^-(E)-\phi_{\sigma_1}^-(E)|,\\
&|\phi_{\sigma_2}^+(E)-\phi_{\sigma_1}^+(E)|,
\left|\log\frac{m_{\sigma_2}(E)}{m_{\sigma_1}(E)}\right|\biggr\}
\in[0,\infty]
\end{aligned}
\]
诱导。

\SourceParagraph{ch05-5.7.2-intro-002}{411}
关键在于，在适当的局部线性拓扑下，中心荷的任意充分小形变，都能提升为稳定性条件的唯一小形变。Bridgeland 的下面这个基本结果，精确表述这一点。
\BookStatementNumber{37}
\begin{theorem}[\BookCite{ref065}]\label{thm:5.37}
对于每个连通分支 $\Sigma\subset\operatorname{Stab}(D)$，存在复线性子空间 $V(\Sigma)\subset\operatorname{Hom}_{\mathbb{Z}}(K(D),\mathbb{C})$，带有良好定义的线性拓扑，以及局部同胚
\[
\mathcal{Z}\colon\Sigma\longrightarrow V(\Sigma),
\]
把稳定性条件 $(Z,\mathcal{P})$ 映到其中心荷 $Z$。
\end{theorem}

\SourceParagraph{ch05-5.7.2-intro-003}{411}
若 $X$ 是光滑射影簇，用 $\operatorname{Stab}(X)$，表示 $D(X)$ 上局部有限稳定性条件的集合，其中，中心荷 $Z$ 通过 Chern 特征标映射 $\operatorname{ch}\colon K(X)\to H^*(X,\mathbb{Q})$ 分解。由\BookRef{thm:5.37}{定理5.37}，立刻可知，$\operatorname{Stab}(X)$ 是有限维复流形。根据\BookRef{prop:5.35}{命题5.35}，它的点，参数化 $D(X)$ 上的有界 $t$ 结构，以及满足相应条件的额外中心荷数据 $Z$。

\SourceParagraph{ch05-5.7.2-intro-004}{411}
\BookStatementNumber{38}
\begin{example}\label{ex:5.38}
令 $\mathcal{A}$ 为复数域上非奇异射影曲线 $X$ 的凝聚 $\mathcal{O}_X$ 模范畴，取 $Z(E)=-\deg E+\mathrm{i}\operatorname{rk}E$。应用\BookRef{prop:5.35}{命题5.35}，就在有界导出范畴 $D(\mathcal{A})$ 上，得到稳定性条件。
\end{example}

\subsubsection{自然群作用}\label{sec:5.7.2.1}

\SourceParagraph{ch05-5.7.2.1-001}{411}
任意三角范畴的稳定性条件空间，都有自然群作用。第一个推广\BookSectionRef{sec:5.3.4.2}{5.3.4.2}的规范不变性，第二个表达\BookSectionRef{sec:5.6.2}{5.6.2}的单值化作用。
\BookStatementNumber{39}
\begin{lemma}\label{lem:5.39}
空间 $\operatorname{Stab}(D)$，具有群 $\widetilde{\operatorname{GL}}^{+}(2,\mathbb{R})$ 的右作用，该群是 $\operatorname{GL}^{+}(2,\mathbb{R})$ 的泛覆盖；还有 $D$ 的正合自等价群 $\operatorname{Aut}(D)$ 的左作用。这两个作用对易。
\end{lemma}

\SourceParagraph{ch05-5.7.2.1-002}{411}
\begin{proof}
首先注意，$\widetilde{\operatorname{GL}}^{+}(2,\mathbb{R})$，可以看作二元组 $(T,f)$ 的集合。其中，$f\colon\mathbb{R}\to\mathbb{R}$ 是递增同胚，满足 $f(\phi+1)=f(\phi)+1$；$T\colon\mathbb{R}^2\to\mathbb{R}^2$ 是保持定向的线性同构；并且，二者在 $S^1=\mathbb{R}/2\mathbb{Z}=(\mathbb{R}^2\setminus\{0\})/\mathbb{R}_{>0}$ 上诱导的映射相同。

\SourceParagraph{ch05-5.7.2.1-003}{411}
给定 $\sigma=(Z,\mathcal{P})\in\operatorname{Stab}(D)$，以及 $(T,f)\in\widetilde{\operatorname{GL}}^{+}(2,\mathbb{R})$，定义新稳定性条件 $\sigma'=(Z',\mathcal{P}')$，令 $Z'=T^{-1}\circ Z$、$\mathcal{P}'(\phi)=\mathcal{P}(f(\phi))$。注意，$\sigma,\sigma'$ 的半稳定对象相同，只是相位重新标记。

\SourceParagraph{ch05-5.7.2.1-004}{411}
对于第二个作用，$\Phi\in\operatorname{Aut}(D)$，在 $K(D)$ 上诱导自同构 $\varphi$。若 $\sigma=(Z,\mathcal{P})$ 是 $D$ 上的稳定性条件，则定义 $\Phi(\sigma)=(Z\circ\varphi^{-1},\mathcal{P}')$，其中，$\mathcal{P}'(t)=\Phi(\mathcal{P}(t))$。这也直接表明，两个作用对易。
\end{proof}

\SourceParagraph{ch05-5.7.2.1-005}{411}
\BookRef{lem:5.39}{引理5.39}的两个群作用，一般都未必自由。特别地，若 $\sigma=(Z,\mathcal{P})$ 的中心荷 $Z\colon K(D)\to\mathbb{C}$ 的像，落在 $\mathbb{C}$ 中一条实直线上，某个 $\widetilde{\operatorname{GL}}^{+}(2,\mathbb{R})$ 的子群，就会固定 $\sigma$。不过，存在自由作用的子群 $\mathbb{C}\subset\widetilde{\operatorname{GL}}^{+}(2,\mathbb{R})$。$\lambda\in\mathbb{C}$，把 $\sigma$ 映到 $\lambda(\sigma)=(Z',\mathcal{P}')$，其中，
\[
Z'(E)=\mathrm{e}^{-\mathrm{i}\pi\lambda}Z(E),\qquad
\mathcal{P}'(\phi)=\mathcal{P}(\phi+\operatorname{Re}\lambda).
\]
注意，任意整数 $n$ 的移位函子 $[n]$，在 $\operatorname{Stab}(D)$ 上的作用，与 $n\in\mathbb{C}$ 的作用一致。

\section{稳定性条件空间的例子}\label{sec:5.8}

\subsection{椭圆曲线}\label{sec:5.8.1}

\SourceParagraph{ch05-5.8.1-001}{412}
令 $X$ 为非奇异亏格1射影曲线，$D$ 为其凝聚层有界导出范畴。这里，中心荷通过数值 Grothendieck 群 $N(D)\cong\mathbb{Z}\oplus\mathbb{Z}$ 分解；对层而言，这两个坐标是秩与次数。

\SourceParagraph{ch05-5.8.1-002}{412}
由\BookRef{lem:5.39}{引理5.39}与\BookRef{thm:5.37}{定理5.37}，有局部同胚
\[
\mathcal{Z}\colon\operatorname{Stab}(X)
\longrightarrow\operatorname{Hom}_{\mathbb{Z}}(N(D),\mathbb{C}),
\]
其像，是二维复向量空间 $\operatorname{Hom}_{\mathbb{Z}}(N(D),\mathbb{C})$ 的某个开子集。事实上，
\BookStatementNumber{40}
\begin{theorem}[\BookCite{ref065}]\label{thm:5.40}
群 $\widetilde{\operatorname{GL}}^{+}(2,\mathbb{R})$ 在 $\operatorname{Stab}(X)$ 上的作用，自由且传递，所以，
\[
\operatorname{Stab}(X)\cong\widetilde{\operatorname{GL}}^{+}(2,\mathbb{R})
\cong\mathbb{C}\times\mathcal{H},
\]
其中，$\mathcal{H}$ 是上半平面。
\end{theorem}

\SourceParagraph{ch05-5.8.1-003}{412}
\begin{proof}
首先，若 $E$ 是 $X$ 上的不可分解层，或者更一般地，是导出范畴中的不可分解对象，那么，它在任意 $\sigma\in\operatorname{Stab}(X)$ 中，都必须半稳定。否则，其 Harder--Narasimhan 滤过，给出非平凡三角形 $A\to E\to B$，满足 $\operatorname{Hom}_D(A,B)=0$。Serre 对偶随后给出
\[
\operatorname{Hom}_D^1(B,A)=\operatorname{Hom}_D(A,B)^*=0,
\]
从而 $E$ 是直和 $A\oplus B$，矛盾。

\SourceParagraph{ch05-5.8.1-004}{412}
取 $\sigma=(Z,\mathcal{P})\in\operatorname{Stab}(X)$。反证假定，中心荷 $Z$ 的像，包含在 $\mathbb{C}$ 的某条实直线中。由于 $\sigma$ 局部有限，其心 $\mathcal{A}$ 必须具有有限长度。若 $A,B$ 是 $\mathcal{A}$ 中不同构的单对象，则
\[
\operatorname{Hom}_D(A,B)=\operatorname{Hom}_D(B,A)=0,
\]
结合 Serre 对偶，得到 $\chi(A,B)=0$；自配对也满足 $\chi(A,A)=0$。这意味着，$\mathcal{A}$ 的所有单对象的数值类，都位于 $N(D)$ 的同一条直线上，因而 $D$ 中所有对象的数值类也如此，与 $N(D)$ 的秩为2矛盾。因此，$Z$ 作为从 $N(D)\otimes\mathbb{R}=\mathbb{R}^2$ 到 $\mathbb{C}\cong\mathbb{R}^2$ 的映射，是同构；从而 $\widetilde{\operatorname{GL}}^{+}(2,\mathbb{R})$ 的作用自由。

\SourceParagraph{ch05-5.8.1-005}{412}
设 $A,B$ 是 $X$ 上的线丛，$\deg A<\deg B$。二者不可分解，所以，在 $\sigma$ 中半稳定，设其相位分别为 $\phi,\psi$。存在映射 $A\to B$ 与 $B\to A[1]$，给出 $\phi\le\psi\le\phi+1$，从而 $Z$ 保持定向。于是，通过 $\widetilde{\operatorname{GL}}^{+}(2,\mathbb{R})$ 的作用，可以假定 $Z(E)=-\deg E+\mathrm{i}\operatorname{rk}E$，并且，对某个 $x\in X$，摩天层 $\mathcal{O}_x$ 的相位为1。随后，$X$ 上所有通常意义的半稳定向量丛，在 $\sigma$ 中也半稳定，相位位于 $(0,1)$。由此很快可知，$\sigma$ 就是\BookRef{ex:5.38}{例题5.38}中的标准稳定性条件。
\end{proof}

\SourceParagraph{ch05-5.8.1-006}{413}
容易说明，$D$ 的自等价，由移位、$X$ 的自同构、与线丛作张量积，以及 Fourier--Mukai 变换生成。$X$ 的自同构，以及次数零线丛的张量积，在 $\operatorname{Stab}(X)$ 上作用平凡，得到
\begin{equation}\label{eq:5.200}
\operatorname{Stab}(X)/\operatorname{Aut}D
\cong\operatorname{GL}^{+}(2,\mathbb{R})/\operatorname{SL}(2,\mathbb{Z}).
\end{equation}
容易看出，在模堆或轨形的意义下，这是椭圆曲线模空间上的 $\mathbb{C}^*$ 束，参数化由 $X$ 的复结构与非零全纯1形式组成的数据的等价类。

\SourceParagraph{ch05-5.8.1-007}{413}
现在可以检验，它是否符合弦理论的预期。根据\BookRef{property:5.1}{性质5.1}，稳定性条件空间，应当是镜像 $Y$ 的复结构模空间，再加上全纯1形式 $\Omega$ 的选择。由于 $Y$ 再次是二维环面，二者完全一致。

\SourceParagraph{ch05-5.8.1-008}{413}
其他与弦理论无关的一维例子，也已有研究。文献\BookCite{ref078}处理了算术亏格为1的不可约奇异曲线，再次得到\BookEqRef{eq:5.200}{5.200}。S. Okada\BookCite{ref379}证明，
\[
\operatorname{Stab}(\mathbb{P}^1)\cong\mathbb{C}^2,
\]
E. Macr\`i\BookCite{ref342}证明，对任意亏格 $g\ge2$ 的光滑射影曲线 $C$，有
\[
\operatorname{Stab}(C)\cong\widetilde{\operatorname{GL}}^{+}(2,\mathbb{R})
\cong\mathbb{C}\times\mathcal{H}.
\]

\subsection{K3 曲面}\label{sec:5.8.2}

\SourceParagraph{ch05-5.8.2-001}{413}
令 $X$ 为 K3 曲面，即单连通、典范丛平凡的光滑射影曲面。杯积在 Picard 群 $\operatorname{Pic}(X)$ 上定义对称配对，它是第二上同调的代数部分。Mukai 把它扩展到
\[
\mathcal{N}(X)=\mathbb{Z}\oplus\operatorname{Pic}(X)\oplus\mathbb{Z},
\]
定义
\[
\langle(r_1,D_1,s_1),(r_2,D_2,s_2)\rangle
=D_1\cdot D_2-r_1s_2-r_2s_1.
\]
代数 Mukai 格 $\mathcal{N}(X)$ 是偶幺模格，符号型为 $(2,\rho)$，其中，$\rho$ 是 $X$ 的 Picard 群秩。完整的 Mukai 上同调格 $H^0\oplus H^2\oplus H^4$，则是符号型 $(4,20)$ 的偶幺模格；代数子格一般不幺模。

\SourceParagraph{ch05-5.8.2-002}{413}
$X$ 上凝聚层有界导出范畴 $D(X)$ 的对象 $E$，具有对应的 Mukai 向量，
\[
v(E)=\operatorname{ch}(E)\sqrt{\operatorname{td}_X}\in\mathcal{N}(X),
\]
其中，$\operatorname{td}_X$ 是 $X$ 的 Todd 类。$\mathcal{N}(X)$ 上的配对、对象的 Mukai 向量，以及 $D(X)$ 的结构，由下面这个漂亮的公式联系起来；它是 Riemann--Roch 定理的简单结果：
\[
-\langle v(E),v(F)\rangle
=\sum_{i\in\mathbb{Z}}(-1)^i\dim\operatorname{Hom}_{D(X)}(E,F[i]),
\qquad E,F\in D(X).
\]

\SourceParagraph{ch05-5.8.2-003}{414}
回顾\BookRef{def:5.33}{定义5.33}的扭转函子。由非零对象 $E\in D(X)$ 定义的 $\Phi_E$，是 $D(X)$ 的自等价，当且仅当 $E$ 为球面对象\BookCite{ref422}。二维中，球面条件为
\[
\dim\operatorname{Hom}_{D(X)}(E,E[i])
=\begin{cases}
1,&i=0,2,\\
0,&\text{其他情形}.
\end{cases}
\]
代入上面的公式，就得到球面对象的 Mukai 向量满足
\[
\langle v(E),v(E)\rangle=-2.
\]

\SourceParagraph{ch05-5.8.2-004}{414}
换句话说，球面对象的 Mukai 向量，是格 $\mathcal{N}(X)$ 的根。因此，这个格的根系，对描述 $D(X)$ 的自等价群起重要作用。由此，根系也对描述 $\operatorname{Stab}(X)$ 起基本作用，并不令人意外。

\SourceParagraph{ch05-5.8.2-005}{414}
由于 Mukai 配对为幺模的，并且 Mukai 向量实现了 $D(X)$ 的数值 K 格，\BookRef{thm:5.37}{定理5.37}的中心荷映射，可以通过这一配对，表示为
\[
\mathcal{Z}\colon\operatorname{Stab}(X)
\longrightarrow\mathcal{N}(X)\otimes\mathbb{C}.
\]

\SourceParagraph{ch05-5.8.2-006}{414}
在 $\mathcal{N}(X)\otimes\mathbb{C}$ 内，考虑如下开集的两个连通分支之一 $\mathcal{P}$：其中向量的实部与虚部，在 $\mathcal{N}(X)\otimes\mathbb{R}$ 中张成正定二维平面。对每个根 $\delta\in\mathcal{N}(X)$，用 $\delta^\perp$ 表示它在 $\mathcal{N}(X)\otimes\mathbb{C}$ 内，对 Mukai 配对的正交补。于是，有 Bridgeland 的下面这个结果。
\BookStatementNumber{41}
\begin{theorem}[\BookCite{ref066}]\label{thm:5.41}
$\operatorname{Stab}(X)$ 有一个连通分支，其中心荷映射 $\mathcal{Z}$，是到开集
\[
\mathcal{P}^0=\mathcal{P}\setminus\bigcup_\delta\delta^\perp
\subset\mathcal{N}(X)\otimes\mathbb{C}
\]
的正规覆盖。
\end{theorem}

\SourceParagraph{ch05-5.8.2-007}{414}
猜想还有更强的结论：预期，这个稳定性条件分支单连通，并由相应的自等价保持。若如此，就得到自等价群的几何描述；其中，在上同调上作用平凡、且保持该分支的部分，可以通过区域 $\mathcal{P}^0$ 的基本群描述，而其余对称通过 Mukai 格上的作用体现。\BookRef{thm:5.41}{定理5.41}的证明及进一步讨论，见文献\BookCite{ref066}。

\subsection{一个三维例子：\texorpdfstring{$\mathcal{O}_{\mathbb{P}^2}(-3)$}{Oₚ²(−3)}}\label{sec:5.8.3}

\SourceParagraph{ch05-5.8.3-intro-001}{415}
取 $D$ 为 $X=\operatorname{Tot}(\mathcal{O}_{\mathbb{P}^2}(-3))$ 上凝聚层有界导出范畴的子范畴，其对象的上同调层，支撑在零截面 $\mathbb{P}^2\subset X$ 上。注意，由导出 McKay 对应，它等价于 $\mathbb{C}^3/\mathbb{Z}_3$ 上支撑在原点的轨形层导出范畴，\BookSectionRef{sec:5.6}{5.6}通过镜像对称研究过它。本小节的稳定性空间，均针对这个支撑受限的范畴 $D$。

\SourceParagraph{ch05-5.8.3-intro-002}{415}
Grothendieck 群 $K(D)$，是秩3自由 Abel 群。遵循文献\BookCite{ref064}，具体描述 $\operatorname{Stab}(D)$ 中一个较大的连通开区域；猜想，其所在连通分支，是唯一的连通分支。它由区域 $D(g)$ 的并组成，每个区域由 $g\in G$ 标记，其中，$G$ 是具有表现
\[
G=\langle\tau_0,\tau_1,\tau_2\mid
\tau_i\tau_j\tau_i=\tau_j\tau_i\tau_j\text{，对所有 }i\ne j\rangle
\]
的仿射辫群。每个区域中，相位在 $(0,1]$ 的稳定对象，位于某个箭图表示范畴 $\mathcal{A}(g)$ 中，其他稳定对象由移位得到。这些区域，如下拼接为连通区域。
\BookStatementNumber{42}
\begin{theorem}\label{thm:5.42}
存在连通开子集 $\operatorname{Stab}^0(D)\subset\operatorname{Stab}(D)$，可写为区域的不交并，
\[
\operatorname{Stab}^0(D)=\bigsqcup_{g\in G}D(g).
\]
每个区域 $D(g)$，由中心荷映射 $\mathcal{Z}$，同胚映到三维复向量空间 $\operatorname{Hom}_{\mathbb{Z}}(K(D),\mathbb{C})$ 的某个局部闭子集。对于不同区域 $D(g_1),D(g_2)$，它们的闭包在 $\operatorname{Stab}^0(D)$ 中相交，当且仅当，对某个 $i\in\{0,1,2\}$，有 $g_1g_2^{-1}=\tau_i^{\pm1}$。
\end{theorem}

\subsubsection{箭图子范畴}\label{sec:5.8.3.1}

\SourceParagraph{ch05-5.8.3.1-001}{415}
令 $(E_0,E_1,E_2)$ 为 $\mathbb{P}^2$ 上的完备例外向量丛集合。$D^b(\mathbb{P}^2)$ 中的完备例外集合，除去移位，都可以取这种形式。文献\BookCite{ref062}证明，对所用集合，存在范畴等价，
\[
\mathrm{R}\operatorname{Hom}_X\!\left(\bigoplus_{i=0}^2\pi^*E_i,-\right)
\colon D^b(\operatorname{Coh}X)\longrightarrow D^b(\operatorname{mod}_{\mathrm{fg}}B),
\]
其中，$\operatorname{mod}_{\mathrm{fg}}B$ 是代数
\[
B=\operatorname{End}_X\!\left(\bigoplus_{i=0}^2\pi^*E_i\right)
\]
的有限生成右模范畴。限制到 $D$ 时，对应于有限维幂零模的导出范畴。代数 $B$，可以描述为如下带关系箭图的路径代数：
\begin{figure}[H]
\centering
\BookImage{02_Unnumbered_Diagrams/Chapter_05/Ch05_Diagram_PDFp415_01_Bookp405.pdf}
\end{figure}

\SourceParagraph{ch05-5.8.3.1-002}{416}
拉回右模导出范畴的标准 $t$ 结构，就在 $D$ 上得到有界 $t$ 结构，其心，等价于 $B$ 的幂零模范畴。以此方式得到的 Abel 子范畴 $\mathcal{A}\subset D$，称为例外子范畴。若 $D$ 的 Abel 子范畴，形如 $\Phi(\mathcal{A})$，其中，$\mathcal{A}$ 为例外子范畴，$\Phi\in\operatorname{Aut}(D)$ 为自等价，就称它为箭图子范畴。

\SourceParagraph{ch05-5.8.3.1-003}{416}
任意箭图子范畴 $\mathcal{A}\subset D$，都等价于上述形式的某个代数的幂零模范畴。因此，它具有三个单对象 $\{S_0,S_1,S_2\}$，对应于箭图的三个一维表示。这些 $S_i$，是\BookEqRef{eq:5.179}{5.179}意义下的球面对象，因而给出\BookEqRef{eq:5.176}{5.176}中的自等价，
\[
\Phi_{S_i}\in\operatorname{Aut}(D).
\]

\SourceParagraph{ch05-5.8.3.1-004}{416}
注意，三个单对象 $S_i$，完全确定相应箭图子范畴 $\mathcal{A}\subset D$。它们之间的 Ext 群，可以从箭图读出。这里采用右 $B$ 模，单对象的 Ext 箭方向，与 $B$ 的箭图方向相反，故
\[
\operatorname{Hom}_D^1(S_1,S_0)=\mathbb{C}^a,\qquad
\operatorname{Hom}_D^1(S_2,S_1)=\mathbb{C}^b,\qquad
\operatorname{Hom}_D^1(S_0,S_2)=\mathbb{C}^c,
\]
其他 $\operatorname{Hom}^1$ 群为零。Serre 对偶，随后确定其他群。

\SourceParagraph{ch05-5.8.3.1-005}{416}
\BookStatementNumber{43}
\begin{example}\label{ex:5.43}
取 $\mathcal{A}$ 为 $D$ 的例外子范畴，对应于 $\mathbb{P}^2$ 上的例外集合 $(\mathcal{O},\mathcal{O}(1),\mathcal{O}(2))$。其单对象为
\[
S_0=i_*\mathcal{O},\qquad
S_1=i_*\Omega^1(1)[1],\qquad
S_2=i_*\mathcal{O}(-1)[2],
\]
其中，$i\colon\mathbb{P}^2\hookrightarrow X$ 是零截面包含映射，$\Omega$ 是 $\mathbb{P}^2$ 的余切丛。这时，$(a,b,c)=(3,3,3)$。
\end{example}

\SourceParagraph{ch05-5.8.3.1-006}{416}
计算自等价 $\Phi_{S_i}$，在 $K(D)$ 上诱导的自同构 $\varphi_{S_i}$。扭转函子 $\Phi_S$，由三角形
\[
\mathrm{R}\operatorname{Hom}_D(S,E)\otimes^{\mathrm L}S
\longrightarrow E\longrightarrow\Phi_S(E)
\]
定义，因此，在 K 理论层面，
\[
\varphi_S([E])=[E]-\chi(S,E)[S].
\]
用列向量坐标，把 $\varphi_{S_i}$ 相对于基 $([S_0],[S_1],[S_2])$ 的矩阵记为 $P_i$，则
\[
\begin{aligned}
P_0&=\begin{pmatrix}1&-3&3\\0&1&0\\0&0&1\end{pmatrix},\qquad
P_1=\begin{pmatrix}1&0&0\\3&1&-3\\0&0&1\end{pmatrix},\\
P_2&=\begin{pmatrix}1&0&0\\0&1&0\\-3&3&1\end{pmatrix}.
\end{aligned}
\]

\subsubsection{辫群作用}\label{sec:5.8.3.2}

\SourceParagraph{ch05-5.8.3.2-001}{417}
文献\BookCite{ref062}说明，在某个单对象处倾斜箭图子范畴 $\mathcal{A}\subset D$，得到另一个箭图子范畴。为了更详细地描述这一过程，需要定义一个作用于球面对象三元组的辫群。

\SourceParagraph{ch05-5.8.3.2-002}{417}
三股环形辫群 $CB_3$，是 $\mathbb{C}^*$ 中三个无序不同点的配置空间的基本群。它由按循环群 $i\in\mathbb{Z}_3$ 编号的三个元素 $\tau_i$，以及一个元素 $r$ 生成，满足关系
\[
\begin{aligned}
r\tau_i r^{-1}&=\tau_{i+1},\qquad i\in\mathbb{Z}_3,\\
\tau_i\tau_j\tau_i&=\tau_j\tau_i\tau_j,\qquad i\ne j.
\end{aligned}
\]
这一表现的证明，见文献\BookCite{ref303}。若以三个单位根作为基点，$\tau_1,r$，对应于按下面方式移动点得到的环路：
\begin{figure}[H]
\centering
\BookImage{02_Unnumbered_Diagrams/Chapter_05/Ch05_Diagram_PDFp417_01_Bookp407.pdf}
\end{figure}

\SourceParagraph{ch05-5.8.3.2-003}{417}
用 $G\subset CB_3$，表示由 $\tau_0,\tau_1,\tau_2$ 生成的子群，就是前面引入的群。

\SourceParagraph{ch05-5.8.3.2-004}{417}
$D$ 中的球面三元组，是由球面对象组成的三元组 $(S_0,S_1,S_2)$。群 $CB_3$，按公式
\[
\begin{aligned}
\tau_1(S_0,S_1,S_2)&=(S_1[-1],\Phi_{S_1}(S_0),S_2),\\
r(S_0,S_1,S_2)&=(S_2,S_0,S_1),
\end{aligned}
\]
作用于相应的球面三元组集合。下面的结果，允许完整理解 $D$ 的箭图子范畴的倾斜过程。
\BookStatementNumber{44}
\begin{proposition}\label{prop:5.44}
令 $\mathcal{A}\subset D$ 为箭图子范畴，单对象为 $(S_0,S_1,S_2)$。对于每个 $i=0,1,2$，倾斜后的箭图子范畴 $L_{S_i}(\mathcal{A})$ 的三个单对象，由球面三元组 $\tau_i(S_0,S_1,S_2)$ 给出。
\end{proposition}

\SourceParagraph{ch05-5.8.3.2-005}{417}
因此，对每个 $g\in G$，有从 $\mathcal{A}$ 出发，反复倾斜得到的箭图子范畴 $\mathcal{A}(g)\subset D$，其三个单对象，由球面三元组
\[
(S_0(g),S_1(g),S_2(g))=g(S_0,S_1,S_2)
\]
给出。

\SourceParagraph{ch05-5.8.3.2-006}{418}
注意，任意箭图子范畴的三个单对象，本来没有确定的顺序；上述定义，则为 $\mathcal{A}(g)$ 的单对象，指定了一个顺序。

\SourceParagraph{ch05-5.8.3.2-007}{418}
令 $P_i(g)\in\operatorname{SL}(3,\mathbb{Z})$，为扭转函子 $\Phi_{S_i(g)}$ 在 $K(D)$ 上诱导的自同构，相对于固定基 $([S_0],[S_1],[S_2])$ 的矩阵。定义辫群在球面三元组上作用的公式，给出这组矩阵的变换律，
\[
\begin{aligned}
P_0(\tau_1g)&=P_1(g),\\
P_1(\tau_1g)&=P_1(g)P_0(g)P_1(g)^{-1},\\
P_2(\tau_1g)&=P_2(g),\\[3pt]
P_0(rg)&=P_2(g),\qquad
P_1(rg)=P_0(g),\qquad P_2(rg)=P_1(g).
\end{aligned}
\]

\SourceParagraph{ch05-5.8.3.2-008}{418}
引入图 $\Gamma(D)$，顶点为 $D$ 的箭图子范畴；两个子范畴，若相差某个单对象处的倾斜，就由一条边相连。文献\BookCite{ref062}证明，不同 $g\in G$，给出不同子范畴 $\mathcal{A}(g)\subset D$。因此，$\Gamma(D)$ 的每个连通分支，就是 $G$ 相对于生成元 $\tau_0,\tau_1,\tau_2$ 的 Cayley 图。

\subsubsection{\texorpdfstring{$X$}{X} 上的稳定性条件}\label{sec:5.8.3.3}

\SourceParagraph{ch05-5.8.3.3-001}{418}
给定 $g\in G$，令 $\mathcal{A}(g)\subset D$ 为对应箭图子范畴。任意非零对象 $E\in\mathcal{A}(g)$，其类是非零的非负整系数组合，
\[
[E]=\sum_{i=0}^2n_i[S_i(g)],\qquad
n_0,n_1,n_2\ge0\text{，且不全为零}.
\]
因此，为在 $D$ 上定义稳定性条件，只须选择三个复数 $z_i$，位于上半平面加负实轴的集合
\[
H=\{z\in\mathbb{C}:z=r\exp(\mathrm{i}\pi\phi),\ r>0,\ 0<\phi\le1\},
\]
再令 $Z(S_i(g))=z_i$。由于 $\mathcal{A}(g)$ 具有有限长度，Harder--Narasimhan 性质自动满足。把这个稳定性条件，记为 $\sigma(g,z_0,z_1,z_2)$。
\BookStatementNumber{45}
\begin{lemma}\label{lem:5.45}
若 $\sigma=\sigma(g,z_0,z_1,z_2)$，是上述类型的稳定性条件，且 $E\in D$ 在 $\sigma$ 中稳定，则存在包含 $\sigma$ 的开子集 $U\subset\operatorname{Stab}(D)$，使 $E$ 在 $U$ 的每个稳定性条件中，都稳定。
\end{lemma}

\SourceParagraph{ch05-5.8.3.3-002}{418}
\begin{proof}
这来自文献\BookCite{ref065}第8节的论证。只须检查如下类的集合有限：$\gamma\in K(D)$，存在 $F\in D$，满足 $[F]=\gamma$，以及 $m_\sigma(F)\le m_\sigma(E)$。这里的心，具有有限长度及仅三个单对象。由于三个 $z_i$ 位于同一个适当半平面，可以选择实线性函数，在它们上都取正值；质量界，于是约束相应 Jordan--H\"older 重数的总和。对任意 $F$，把 Harder--Narasimhan 因子移入心后，得到同样的约束，所以，可能的类只有有限个。
\end{proof}

\SourceParagraph{ch05-5.8.3.3-003}{418}
每个 $g\in G$，对应稳定性条件集合，
\[
\begin{aligned}
D(g)=\{&\sigma(g,z_0,z_1,z_2):(z_0,z_1,z_2)\in H^3,\\
&\text{至多一个 }z_i\in\mathbb{R}\}\subset\operatorname{Stab}(D).
\end{aligned}
\]
根据定义，这些子集互不相交，因为它们对应于不同的心。
\BookStatementNumber{46}
\begin{proposition}\label{prop:5.46}
有开子集
\[
\operatorname{Stab}^0(D)=\bigsqcup_{g\in G}D(g)\subset\operatorname{Stab}(D).
\]
若 $g_1\ne g_2$，区域 $D(g_1),D(g_2)$ 的闭包，在 $\operatorname{Stab}^0(D)$ 中相交，当且仅当，对某个 $i\in\{0,1,2\}$，有 $g_1=\tau_i^{\pm1}g_2$。
\end{proposition}

\SourceParagraph{ch05-5.8.3.3-004}{419}
\begin{proof}
设 $\sigma=\sigma(g,z_0,z_1,z_2)\in D(g)$。需要说明，$\sigma$ 有开邻域，包含在 $\operatorname{Stab}^0(D)$ 内。单对象 $S_i=S_i(g)\in\mathcal{A}(g)$，在 $\sigma$ 中稳定；在 $\sigma$ 的充分小开邻域 $U$ 中，它们仍稳定。反复使用一个容易证明的事实：若 $\mathcal{A},\mathcal{A}'\subset D$ 是有界 $t$ 结构的心，且 $\mathcal{A}\subset\mathcal{A}'$，则二者相等。

\SourceParagraph{ch05-5.8.3.3-005}{419}
首先假定，对每个 $i$，都有 $\operatorname{Im}z_i>0$。缩小 $U$，可以假定，在 $U$ 中所有稳定性条件 $(Z,\mathcal{P})$ 下，每个 $S_i$ 的相位都在 $(0,1)$ 中。$\mathcal{A}(g)$，是包含 $S_i$ 的最小扩张封闭子范畴，因此，包含在所有这些稳定性条件的心 $\mathcal{P}((0,1])$ 中。这说明，$\mathcal{P}((0,1])=\mathcal{A}(g)$，所以，$U\subset D(g)$。

\SourceParagraph{ch05-5.8.3.3-006}{419}
现在假定，一个 $z_i$ 在实轴上；不失一般性，取 $z_0$，所以，$\sigma$ 位于 $D(g)$ 的边界。这时，$z_0\in\mathbb{R}_{<0}$，而 $i=1,2$ 的 $\operatorname{Im}z_i>0$。缩小 $U$，可以假定，对 $U$ 中所有 $(Z,\mathcal{P})$，都有 $\operatorname{Re}Z(S_0)<0$，以及 $i=1,2$ 的 $\operatorname{Im}Z(S_i)>0$。

\SourceParagraph{ch05-5.8.3.3-007}{419}
对象 $S'=\Phi_{S_0}(S_2)$，属于 $\mathcal{A}(g)$，实际上是泛扩张，
\[
0\longrightarrow S_2\longrightarrow S'
\longrightarrow S_0^{\oplus m}\longrightarrow0,
\qquad m=\dim\operatorname{Hom}_D^1(S_0,S_2).
\]
由于 $\operatorname{Hom}_D(S_0,S')=0$，对象 $S'$ 属于 $\mathcal{P}((0,1))$；利用相位的连续性，再缩小 $U$，可以假定，$U$ 中所有稳定性条件都如此。

\SourceParagraph{ch05-5.8.3.3-008}{419}
把 $U$ 分成 $U_+$ 与 $U_-$，分别由 $\operatorname{Im}Z(S_0)\ge0$ 与 $\operatorname{Im}Z(S_0)<0$ 给出。前面的论证说明，$U_+\subset D(g)$。另一方面，$U_-$ 中任意稳定性条件，都使 $S_0$ 稳定，且相位位于 $(1,3/2)$。因此，心 $\mathcal{P}((0,1])$，包含 $S_0[-1],S',S_1$。它们是有限长度范畴 $\mathcal{A}(\tau_0g)$ 的单对象，所以，$U_-\subset D(\tau_0g)$。由于这些区域只允许一个中心荷位于实轴，交界处对应于一个单对象的倾斜，得到所述邻接条件。
\end{proof}

\SourceParagraph{ch05-5.8.3.3-009}{419}
综合起来，就证明了\BookRef{thm:5.42}{定理5.42}。

\section{进一步方向与开放问题}\label{sec:5.9}

\subsection{弦理论与数学的关系}\label{sec:5.9.1}

\SourceParagraph{ch05-5.9.1-intro-001}{419}
这里汇集一些尚未解决的细节，以及更实质的问题；解决它们，会显著澄清整个图景。

\subsubsection{范畴结构}\label{sec:5.9.1.1}

\SourceParagraph{ch05-5.9.1.1-001}{420}
同调代数的许多构造，在物理讨论中，几乎没有起作用，甚至完全没有出现，例如，张量积、推出与拉回的导出版本。此外，虽然实质地使用了对应与导出自等价，它们主要作为单值化出现；物理上，这些操作相当复杂。

\SourceParagraph{ch05-5.9.1.1-002}{420}
为这些构造给出更直接的物理定义，也许有用，既补全图景，也指导相关背景下的物理工作，例如，具有更少超对称，或没有超对称的 Dirichlet 膜理论。

\SourceParagraph{ch05-5.9.1.1-003}{420}
已知物理构造的一个例子，是\BookSectionRef{sec:4.6.1}{4.6.1}讨论的对应。那里核 $\mathcal{F}$ 的共形场论对应物，是乘积 SCFT $X\times Y$ 中的边界态 $\mathcal{F}$。按第4章中 $X\to Y$ 的变换约定，可以把它解释为“缺陷线”或“界面”：世界面右半由 SCFT $X$ 描述，左半由 SCFT $Y$ 描述。把 $X$ 的边界态 $B$ 放在右侧，再取它接近 $\mathcal{F}$ 的极限，就得到 SCFT $Y$ 中对应的边界态 $\Phi_{\mathcal{F}}(B)$。这一构造，已用于表述有理共形场论中的对偶\BookCite{ref161,ref162}。

\SourceParagraph{ch05-5.9.1.1-004}{420}
拓扑扭曲 SCFT 以后，可以用 Hilbert 空间上的显式迹，替换该极限操作。把 $\mathcal{F}$ 看作以 $Y$ 为参数的一族 $X$ 边界态；为匹配协变 Fourier--Mukai 约定，在传入边界采用对偶 $B^\vee$，得到族指标
\[
[\Phi_{\mathcal{F}}(B)]
=\operatorname{Ind}_X\mathcal{H}_{B^\vee,\mathcal{F}}
=\operatorname{Tr}_{\mathcal{H}_{B^\vee,\mathcal{F}}}(-1)^{\mathsf F}
\in K(Y).
\]
这里，迹作为族指标理解。相应的 BRST 复形，由沿 $X$ 的导出映射复形给出；在 $y\in Y$ 的纤维上，是 $\mathrm{R}\operatorname{Hom}_X(B^\vee,\mathcal{F}_y)$。保留整个复形，得到导出变换；取指标，则得到其 K 理论类。因此，导出对应，可以看作缺陷线构造的拓扑扭曲版本。随着对这一物理构造的理解发展，这也许会带来有趣后果。

\subsubsection{稳定性条件与弦性 K\"ahler 模空间}\label{sec:5.9.1.2}

\SourceParagraph{ch05-5.9.1.2-001}{420}
令 $X$ 为单连通 Calabi--Yau 三维流形，$D=D^b(X)$。根据\BookRef{property:5.1}{性质5.1}，也许预期
\[
\operatorname{Aut}(D)\backslash\operatorname{Stab}(X)/\mathbb{C}
\]
与镜像 $Y$ 的复结构模空间识别。未除去 $\mathbb{C}$ 的整体作用以前，还包括非零全纯3形式的尺度与相位选择；这里的右商，除去这一自由度。

\SourceParagraph{ch05-5.9.1.2-002}{420}
实际上，利用\BookRef{thm:5.37}{定理5.37}，容易看到，这种对整个稳定性空间的识别，不可能成立。简单地说，$\operatorname{Stab}(X)$ 太大，也太平坦。

\SourceParagraph{ch05-5.9.1.2-003}{420}
具体取 $X$ 为五次三维流形。弦性 K\"ahler 模空间 $\mathcal{M}_K(X)$，基本上就是镜像三维流形 $Y$ 的复模空间。众所周知，它是具有两个穿孔及一个特殊点的二维球面。穿孔是大体积极限点与锥形点，特殊点称为 Gepner 点。镜像 $Y$ 的周期，在 $\mathcal{M}_C(Y)$ 上定义全纯函数，满足四阶 Picard--Fuchs 方程，在这些特殊点具有正则奇点。

\SourceParagraph{ch05-5.9.1.2-004}{421}
镜像对称下，$Y$ 的 Lagrange 子流形周期，对应于 $D$ 中对象的中心荷。因此，来自 $\mathcal{M}_K(X)$ 的物理稳定性条件，其中心荷 $Z\colon K(D)\to\mathbb{C}$，满足 $Y$ 的 Picard--Fuchs 方程。若预期的物理稳定性条件存在，并按这些周期连续变化，由于这些周期不满足非零常系数线性关系，与\BookRef{thm:5.37}{定理5.37}比较，可知，包含它们的稳定性条件分支，是四维的。上述双商，是三维空间，其中，$\mathcal{M}_K(X)$ 占一维子流形。这一周期轨迹，在稳定性空间中的嵌入，高度超越。

\SourceParagraph{ch05-5.9.1.2-005}{421}
更一般地，对于单连通 Calabi--Yau 三维流形 $X$，可以猜想，$\operatorname{Stab}(X)$ 是弦性 K\"ahler 模空间的某种扩展版本。通常 K\"ahler 模空间的切空间，可与 $H^{1,1}(X)$ 识别；扩展版本的候选切空间，则为
\[
\bigoplus_p H^{p,p}(X).
\]
若要从 $\operatorname{Aut}(D)\backslash\operatorname{Stab}(X)/\mathbb{C}$ 中，选出 K\"ahler 模空间这一子流形，就需要在稳定性条件空间上，定义额外结构。研究镜像对称中，其他扩展模空间出现的情形，可以为这一结构提供线索。

\subsubsection{扩展模空间}\label{sec:5.9.1.3}

\SourceParagraph{ch05-5.9.1.3-001}{421}
代数几何中，至少有三种情形，会出现扩展模空间：普适展开空间、大量子上同调，以及 Barannikov--Kontsevich 的扩展模空间。这些空间，都带有与 Frobenius 结构密切相关的丰富几何结构，也都与 SCFT 模空间密切相关。每一种，都可以与稳定性条件空间联系，不过，这些联系都远未精确建立。这里只简要概述，并给出一些文献。

\SourceParagraph{ch05-5.9.1.3-002}{421}
如 Takahashi\BookCite{ref439,ref280}所解释，$n$ 维孤立超曲面奇点 $X_0$ 的展开空间 $T$，应当与该奇点的 Milnor 纤维 $X_t$ 上 Fukaya 范畴的稳定性条件空间有关。注意，
\[
\mu=\dim_{\mathbb C}H_n(X_t,\mathbb{C})=\dim_{\mathbb C}T.
\]
给定 $H_n(X_t,\mathbb{C})$ 的基 $L_1,\ldots,L_\mu$，K. Saito 的本原形式理论，为纤维 $X_t$ 上合适的一族全纯 $n$ 形式 $\Omega_t$，提供相应的平坦结构。取适当非退化区域与分支，周期
\[
Z(L_i)=\int_{L_i}\Omega_t
\]
作为相关平坦结构的坐标。由于这些周期，类似于中心荷，这正是\BookRef{thm:5.37}{定理5.37}所提示的图景。

\SourceParagraph{ch05-5.9.1.3-003}{422}
Fano 簇 $Z$ 的大量子上同调，似乎与 $Z$ 及对应局部 Calabi--Yau 空间 $\operatorname{Tot}(\omega_Z)$ 的导出范畴的稳定性条件空间有关。若 $Z$ 的量子上同调，在一般点半单，Dubrovin 说明，如何把预势，从 $H^*(Z,\mathbb{C})$ 的一个开子集解析延拓，在配置空间
\[
M\subset\operatorname{Con}_n(\mathbb{C})
=\{(u_1,\ldots,u_n)\in\mathbb{C}^n:i\ne j\Rightarrow u_i\ne u_j\}
\]
的稠密开子集上，给出 Frobenius 结构。根据 Dubrovin 的一个猜想\BookCite{ref143}，$Z$ 的量子上同调在一般点半单，当且仅当 $D^b(Z)$ 具有完备强例外集合 $(E_0,\ldots,E_{n-1})$。前者允许定义上述扩展模空间 $M$；后者允许像\BookRef{thm:5.42}{定理5.42}那样，通过倾斜理解稳定性条件空间。此外，在适当坐标及选择下，量子上同调的 Stokes 矩阵 $S_{ij}$，等于 Gram 矩阵 $\chi(E_i,E_j)$：前者控制 $M$ 上 Frobenius 结构的解析延拓，后者控制倾斜或突变过程。更多内容，见文献\BookCite{ref064,ref059}。

\SourceParagraph{ch05-5.9.1.3-004}{422}
最后，Barannikov 与 Kontsevich\BookCite{ref030}研究了扩展的形式形变空间。若 $X$ 是复射影簇，则 $X$ 的形变形式胚，其切空间为 $H^1(X,TX)$，包含在更大的分次形式胚中，后者的切空间由
\[
\bigoplus_{p,q}H^p(X,\Lambda^qTX)
\]
给出，并描述范畴 $D^b(X)$ 的 $A_\infty$ 形变。设 $X_1,X_2$ 是镜像 Calabi--Yau 三维流形。$X_1$ 的复形变，对应于 $X_2$ 的 K\"ahler 形变。进入扩展模空间以后，可以设想，参数化 $D^b(X_1)$ 形变的 Barannikov--Kontsevich 空间的某种整体版本，与 $X_2$ 的稳定性条件空间互为镜像。非常示意地，可以写为
\[
\operatorname{Def}(D^b(X_1))\cong\operatorname{Stab}(D^b(X_2)),
\]
不过，要使维数匹配，还应扩展 $\operatorname{Stab}(X_2)$，使其分次切空间，包含 $X_2$ 的全部上同调。这样的同构，将是完全留在代数几何范围内的镜像对称陈述。

\subsection{稳定对象的计数与穿墙}\label{sec:5.9.2}

\SourceParagraph{ch05-5.9.2-001}{423}
给定三维 Calabi--Yau 范畴中的对象 $E\in D$，例如 Calabi--Yau 三维流形上的向量丛，Ext 群满足对偶，
\[
\operatorname{Ext}_D^i(E,E)\cong
\operatorname{Ext}_D^{3-i}(E,E)^*.
\]
特别地，形变空间 $\operatorname{Ext}_D^1(E,E)$，与阻碍空间 $\operatorname{Ext}_D^2(E,E)$ 对偶。这意味着，适当处理自同构以后，$E$ 局部所在的对象模问题，预期由与变量数同样多的方程切出，因而具有虚维数零。若方程横截，模空间局部就是孤立点 $[E]$。因此，再固定保证模空间紧的一些条件，就希望直接计数这些点，即满足条件的对象 $E$。

\SourceParagraph{ch05-5.9.2-002}{423}
当然，一般不能期望，定义 $E$ 的方程始终横截，给出真正零维模空间。不过，虚循环技术\BookCite{ref333,ref035,ref443}，仍然允许给紧对象模空间赋予整数，把它解释为固定类型对象的“个数”。

\SourceParagraph{ch05-5.9.2-003}{423}
$E$ 的拓扑类型，由数值 K 群中的类 $[E]$ 给出。与层理论中一样，紧性需要更多条件；这正是稳定性发挥作用的地方。给定 $D$ 上的稳定性条件 $\sigma=(Z,\mathcal{P})$，预期可以构造给定类 $\alpha$ 的 $\sigma$ 半稳定对象模空间。在这一类型没有严格半稳定对象，并满足所需紧性等条件时，可以使用虚循环理论，得到计数该类 $\sigma$ 稳定对象的整数 $n_\sigma(\alpha)$。

\SourceParagraph{ch05-5.9.2-004}{423}
使上一段精确的技术构造，很困难。其中一个棘手问题，是一般范畴 $D$ 的对象模“空间”，究竟是什么。簇上的向量丛，在适当稳定性等条件下，具有准射影模空间；不过，构造涉及不变量理论，以及把向量丛嵌入 Grassmann 簇，这种方法不易推广到任意对象。文献采用不同办法，最一般的包括代数模堆\BookCite{ref269}与形式模空间\BookCite{ref313}。

\SourceParagraph{ch05-5.9.2-005}{423}
一个非常有趣的问题，是改变 $\sigma$ 时，不变量 $n_\sigma(\alpha)$ 如何变化。两个非零对象的相位，在稳定性空间的一道墙上交换顺序时，稳定对象集合发生剧烈变化，如前面多处讨论，因此，不变量也变化。不过，所有不变量的集合 $\{n_\sigma(\alpha)\}_\alpha$，按普遍规律改变。Joyce 的一系列论文\BookCite{ref269}，以及 Kontsevich 与 Soibelman 很新的工作\BookCite{ref313}，研究了它。Denef 与 Moore\BookCite{ref111}从物理观点研究穿墙。Joyce\BookCite{ref272}用这些思想，在满足 Calabi--Yau 条件的 Abel 范畴的稳定性条件空间上，定义平坦联络，不过，推广到导出范畴似乎有困难。Bridgeland 与 Toledano Laredo\BookCite{ref059}，把 Joyce 构造，解释为 Abel 范畴的稳定性空间上，一个不规则联络的等单值化族。

\SourceParagraph{ch05-5.9.2-006}{424}
这是一个迅速发展的领域。不再进一步展开技术细节，只提一下，它还联系到拓扑弦\BookCite{ref381,ref354,ref386,ref437}、模形式\BookCite{ref111}、Ringel--Hall 代数\BookCite{ref269,ref270,ref271,ref273,ref059}，以及箭图突变的组合学\BookCite{ref313}。

\subsection{边界态的直接构造}\label{sec:5.9.3}

\SourceParagraph{ch05-5.9.3-001}{424}
确定物理 D 膜的集合，并检验它是否由稳定性条件决定，最令人满意的方法，是表述并证明 Donaldson--Uhlenbeck--Yau\BookRef{thm:5.8}{定理5.8}的物理类似物；该定理控制大体积 B 膜的问题。

\SourceParagraph{ch05-5.9.3-002}{424}
也许，最好的方法，是使用边界重整化群。给定拓扑理论中的对象 $E$，在 SCFT 中表述一个简单边界条件 $E'$，拓扑扭曲后对应于 $E$，但未必满足边界共形不变性的约束。

\SourceParagraph{ch05-5.9.3-003}{424}
然后，应用边界 RG，流向不动点。物理上，这相当于取世界面关联函数的长距离极限。由于拓扑扭曲理论不依赖世界面度规，这一操作给出边界态 $E'$：它的拓扑扭曲仍为 $E$，但现在已共形不变。

\SourceParagraph{ch05-5.9.3-004}{424}
若这个极限仍为单个稳定膜，即具有唯一的恒等算符，或世界体语言中的 $U(1)$ 规范理论，那么，结果应当是与 $E$ 对应的物理膜。一般而言，却可能得到一组物理膜 $E_i$。预期，它们正是\BookRef{conj:5.32}{猜想5.32}所要求的因子。若如此，就说明，物理膜，是\BookSectionRef{sec:5.7}{5.7}的稳定对象。

\SourceParagraph{ch05-5.9.3-005}{424}
这一一般方案，包含 Donaldson 用于证明 Donaldson--Uhlenbeck--Yau 定理的 Yang--Mills 流。以非线性 $\sigma$ 模型定义的 SCFT 与一个丛 $E$ 为起点，选择一般联络，用它像\BookSectionRef{sec:3.5.2.2}{3.5.2.2}那样定义边界条件。已证明，在首阶（当 $\operatorname{rk}E=1$ 时则对所有阶），RG 流为\BookCite{ref451}
\[
\Lambda\frac{\partial A_i}{\partial\Lambda}
=\frac{\delta S_{\mathrm{BI}}}{\delta A_i},
\]
其中，右边按联络空间相应度规的梯度理解，$S_{\mathrm{BI}}$ 是\BookEqRef{eq:3.142}{3.142}的 Born--Infeld 作用量。取 $\Lambda$ 为能标，向红外流动时，这是下降流；在大体积极限中，化为 Yang--Mills 流。

\SourceParagraph{ch05-5.9.3-006}{424}
要使这种论证在更一般情形中可信，还需要发展许多要素。虽然已在物理上理解某些具体的快子凝聚例子，它们先验地涉及耦合的大幅变化，可能使算符维数，因而流本身，也发生很大变化。

\SourceParagraph{ch05-5.9.3-007}{425}
一种简化，是承认\BookSectionRef{sec:5.7}{5.7}的结论：对于固定稳定性结构，适当相位区间内的稳定对象，位于某个 Abel 子范畴中。于是，只以对应于该子范畴对象的边界条件，作为流的出发点。可以预期，这避开了对快子凝聚一般结果的需求。最具体的版本，再次是非线性 $\sigma$ 模型与丛，并说明，$\alpha'$ 修正和世界面瞬子引起的稳定性形变，确实如预言。即使这一点，目前看起来也很困难。

\SourceParagraph{ch05-5.9.3-008}{425}
这一论证提出的主要问题之一，是 RG 不动点条件，是否也有辛商的解释。这在 Donaldson 的证明中相当重要，因为它说明，相应泛函沿适当轨道凸，从而使流的终点，不依赖于轨道内的出发点。对于所提出的推广，这并不清楚。若不成立，某些流，可能终止于非超对称、但共形的边界条件，使整个问题显著复杂化。

\subsection{几何 Langlands 对偶的应用}\label{sec:5.9.4}

\SourceParagraph{ch05-5.9.4-001}{425}
拓扑开弦理论，对几何 Langlands 对偶，有重要应用。在最初级的层面，可以如下表述其范畴性主张。令 $C$ 为紧、无边界的 Riemann 曲面，$G$ 为紧约化 Lie 群，$G_{\mathbb{C}}$ 为其复化，$\mathcal{M}_{\mathrm{flat}}(G,C)$ 为 $C$ 上稳定平坦 $G_{\mathbb{C}}$ 联络的模空间。$G$ 的 Langlands 对偶，是另一个紧约化 Lie 群 ${}^LG$，由交换根数据给出，特别地，交换权格与余权格，以及根与余根。令 $\operatorname{Bun}({}^LG,C)$，为 $C$ 上全纯 ${}^LG_{\mathbb C}$ 主丛的模堆。一项几何 Langlands 范畴性主张，是 $\mathcal{M}_{\mathrm{flat}}(G,C)$ 上凝聚层的导出范畴，应当与 $\operatorname{Bun}({}^LG,C)$ 上 D 模的导出范畴等价；这里先以相应稳定轨迹及所需扭曲为背景。

\SourceParagraph{ch05-5.9.4-002}{425}
文献\BookCite{ref294}说明，可以通过规范理论，给出如下物理论证。首先，把 $\mathcal{M}_{\mathrm{flat}}(G,C)$，与另一个模空间作镜像识别，后者粗略地可以描述为 $\operatorname{Bun}({}^LG,C)$ 的余切丛。其次，说明 $T^*\operatorname{Bun}({}^LG,C)$ 上、取典范辛形式的 A 膜范畴，等价于其非交换形变上的 B 膜范畴。后者，又对应于 $\operatorname{Bun}({}^LG,C)$ 上的解析 D 模范畴。

\SourceParagraph{ch05-5.9.4-003}{425}
解释第一部分，需要熟悉超对称规范理论，会远远超出本书范围。这里的目标，是解释第二部分，即 A 膜与非交换 B 膜的关系。只要靶空间 $X$，是某个复流形 $Y$ 的余切丛总空间，这种关系就会出现。假定辛形式 $\omega$，与 $T^*Y$ 上的典范辛形式成比例，并且 B 场为零。由于 $Y$ 复，可以把 $\omega$，看作全纯辛形式 $\Omega$ 的实部。若 $q^i$ 为 $Y$ 的全纯坐标，$p_i$ 为余切纤维上的对偶坐标，则
\[
\Omega=\frac1\hbar\,\mathrm{d}p_i\wedge\mathrm{d}q^i
=\mathrm{d}\Theta.
\]
由于 $\omega$ 与 $\Omega$ 都恰当，$X$ 的闭 A 模型相当简单：没有非平凡的闭弦瞬子，量子上同调环，与经典环同构。

\SourceParagraph{ch05-5.9.4-004}{426}
希望理解 $X=T^*Y$ 上的 A 膜范畴。关键观察，是 $X$ 上存在自然的余各向同性 A 膜，至多允许与 $X$ 上的平坦线丛作张量积。取与所用复结构及 $B-F$ 约定相容的符号，其曲率二形式恰当，并为
\[
F=-\operatorname{Im}\Omega.
\]
用 $I$ 表示由 $Y$ 的复结构，诱导的 $X$ 上自然复结构，则 $F=\omega I$，因而 $\omega^{-1}F=I$，平方为 $-1$。因此，$X$ 上平凡线丛的任意曲率为 $F$ 的幺正联络，都定义余各向同性 A 膜。例如，可取联络1形式为
\[
A=-\operatorname{Im}\Theta.
\]
称这一 A 膜，为典范余各向同性膜。

\SourceParagraph{ch05-5.9.4-005}{426}
接下来，希望理解典范余各向同性 A 膜的自同态，即 BRST 闭开弦顶点算符的代数。$Y$ 为仿射空间时，这很容易。为理解一般情形，希望用每个同构于 $\mathbb{C}^n$ 开子集的坐标片覆盖 $Y$，再论证，可以逐片局部计算，并把结果“粘合”。更准确地说，希望相关代数，是某个代数层的上同调，其局部结构，与 $Y=\mathbb{C}^n$ 时相同。

\SourceParagraph{ch05-5.9.4-006}{426}
一般而言，定义顶点算符关联函数的路径积分，并不在靶空间中具有局域性。例如，A 模型的圆盘瞬子贡献，显然依赖于整个 $X$。不过，可以说明，没有满足典范余各向同性膜边界条件的非平凡圆盘瞬子\footnote{事实上，对于任意辛流形 $X$ 上、余维为零的余各向同性 A 膜，都是如此。}。此外，$\sigma$ 模型微扰理论，把路径积分的被积量，在到 $X$ 的常值映射附近作 Taylor 展开。容易说明，每一微扰项，只依赖于 $X$ 中某点的无穷小邻域。结合这些观察，可以说明，作为 $\hbar$ 的形式幂级数，开弦顶点算符代数，是一个代数层的上同调；它局部同构于 $X=\mathbb{C}^n\times\mathbb{C}^n$ 时的类似层。

\SourceParagraph{ch05-5.9.4-007}{427}
把这些观察，用于 $X=T^*Y$ 上的典范余各向同性 A 膜。局部地，用全纯坐标 $q^1,\ldots,q^n$，把 $Y$ 识别为 $\mathbb{C}^n$ 的一个区域。除去 BRST 恰当项，A 模型在圆盘 $\Sigma$ 上的作用量，为
\[
S=\frac1\hbar\int_{\partial\Sigma}\phi^*(p_i\,\mathrm{d}q^i),
\]
其中，$\phi\colon\Sigma\to X$。这个作用量，与 $Y$ 上 Hamilton 函数为零的粒子的作用量相同，只是 $q^i$ 为全纯坐标，而非普通实坐标。此外，可以说明，BRST 不变开弦顶点算符，可取为 $p,q$ 的全纯函数。因此，局部量子化直接得到 $T^*Y$ 上全纯函数代数的非交换形变，对应于全纯 Poisson 双向量
\[
P=\hbar\frac{\partial}{\partial p_i}\wedge\frac{\partial}{\partial q^i}.
\]
形变乘积有显式公式，
\[
\begin{aligned}
(f\star g)(p,q)
=\left.\exp\!\left[\frac\hbar2\left(
\frac{\partial^2}{\partial p_i\partial\tilde q^i}
-\frac{\partial^2}{\partial q^i\partial\tilde p_i}\right)\right]
 f(p,q)g(\tilde p,\tilde q)\right|_{\tilde p=p,\,\tilde q=q}.
\end{aligned}
\]
这个乘积，称为 Moyal--Wigner 乘积。

\SourceParagraph{ch05-5.9.4-008}{427}
Moyal--Wigner 公式，是 $\hbar$ 的形式幂级数，收敛半径可能为零。为处理这一点，可以把函数限制为关于纤维坐标 $p_i$ 的多项式。这样的局部函数，可以看作微分算符的符号；此时，Moyal--Wigner 乘积，化为相应 Weyl 符号的乘积，并是 $\hbar$ 的多项式。因此，只要限制到 $p_i$ 的多项式算符，开弦顶点算符层，局部由 $Y$ 上全纯微分算符层，或其 $\hbar$ 版本，描述。

\SourceParagraph{ch05-5.9.4-009}{427}
局部地，$Y$ 的全纯微分算符层 $\mathcal{D}(Y)$，与作用在线丛 $\mathcal{L}$ 上的全纯微分算符层 $\mathcal{D}(Y,\mathcal{L})$，没有区别。因此，开弦顶点算符层，可能带有这种全局扭曲。文献\BookCite{ref294}说明，在所讨论的量子化方案中，线丛，或相应扭曲类，的选择，描述剩余的全局歧义。此外，经典问题，在 $p_i\mapsto-p_i$，同时反转 $\Sigma$ 定向时对称。若要求量子化保持这一对称，开弦顶点算符代数，就必须与其反代数同构。众所周知\BookCite{ref052}，
\[
\mathcal{D}(Y,\mathcal{L})^{\mathrm{op}}
\cong\mathcal{D}(Y,\mathcal{L}^{-1}\otimes K_Y).
\]
因此，这个对称，要求扭曲为半典范扭曲；若 $K_Y$ 有平方根，可以取 $\mathcal{L}=K_Y^{1/2}$。取哪个平方根并不重要，因为它们只相差平坦线丛，与平坦线丛作张量积，不改变这一微分算符层。若平方根不是实际线丛，$K_Y^{1/2}$ 的记号，按局部平方根及其扭曲微分算符层理解。

\SourceParagraph{ch05-5.9.4-010}{428}
结论是，$X=T^*Y$ 上典范余各向同性 A 膜 $\alpha$ 的开弦顶点算符层，同构于非交换代数层 $\mathcal{D}(Y,K_Y^{1/2})$。由此，可以给任意 A 膜 $\beta$，赋予扭曲 D 模，即这个代数层的模层。考虑条带 $\Sigma=I\times\mathbb{R}$ 上、靶空间为 $X$ 的 A 模型，其中，$I$ 是单位区间；在两条边界，分别施加膜 $\alpha,\beta$ 的边界条件。量子化以后，得到 $Y$ 上的向量空间层，它是插入在 $\alpha$ 边界上的开弦顶点算符层的模。
\begin{example*}
取 $\beta$ 为 $T^*Y$ 的零截面，携带平凡的物理线丛，并作相应半典范解释。所得层，就是 $K_Y^{1/2}$ 的局部截面层，带有 $\mathcal{D}(Y,K_Y^{1/2})$ 的重言作用。
\end{example*}

\SourceParagraph{ch05-5.9.4-011}{428}
可以非严格地论证，从 A 膜到 D 模，或 D 模复形，的映射，延伸为 $X$ 上 A 膜范畴，与 $Y$ 上相应半典范约定下 D 模导出范畴的等价；平方根实际存在时，可通过张量积，与通常的 D 模表述识别。这个论证，依赖于猜想：任意广义 Calabi--Yau，都存在广义复膜范畴\BookCite{ref284,ref390}。这里不再展开。

\SourceParagraph{ch05-5.9.4-012}{428}
这些思想的一个简单应用，是物理上解释 G. Laumon 与 M. Rothstein 的 D 模 Fourier 变换定理\BookCite{ref320,ref405}。令 $Y$ 为 $n$ 维 Abel 簇，看作向量空间 $V$ 对格 $\Gamma$ 的商。令 $Y^\vee$ 为对偶环面，$X=T^*Y$。$X$ 同构于 $V^\vee\times(V/\Gamma)$。对 $Y$ 作 T 对偶，得到 $X^\vee=V^\vee\times(V^\vee/\Gamma^\vee)$。注意，$Y$ 是 $X$ 的 Lagrange 子流形，所以，这个 T 对偶，是镜像对称，即 $X^\vee$ 与 $X$ 互为镜像。作为实流形，$X^\vee$ 同构于 $Y^\vee=V^\vee/\Gamma^\vee$ 的切丛。不过，容易说明，与 $X$ 的典范辛形式 T 对偶的复结构，并非乘积复结构。一种描述方法，是选择 $V^\vee$ 上的实坐标 $v^j$；若 $y^j$ 为环面上的对应实坐标，那么，$X^\vee$ 的复坐标为 $y^j+\mathrm{i}v^j$。文献\BookCite{ref320,ref405}把这个具有相应全纯仿射结构的空间，称为 $Y^\vee$ 的扭曲切丛。

\SourceParagraph{ch05-5.9.4-013}{428}
于是，$X$ 上的 A 膜范畴，等价于 $Y^\vee$ 的这个扭曲切丛上的 B 膜范畴。承认 A 膜与 D 模之间的关系，就推知，$Y$ 上的 T 对偶，诱导 $Y$ 上 D 模导出范畴，与这个扭曲丛上的凝聚层导出范畴的等价。这就是文献\BookCite{ref320,ref405}的主要定理。

\SourceParagraph{ch05-5.9.4-014}{428}
几何 Langlands 对偶，可以看作这一例子的非 Abel 推广。不使用 $T^*Y$，而考虑 Riemann 曲面 $C$ 上 Higgs ${}^LG$ 丛的模空间 $\mathcal{M}_{\mathrm{Higgs}}({}^LG,C)$。
\begin{definition*}[Higgs 丛]
Higgs ${}^LG$ 丛，是二元组 $(E,\varphi)$，其中，$E$ 是 $C$ 上的全纯 ${}^LG_{\mathbb C}$ 主丛，$\varphi$ 是 $\operatorname{ad}(E)\otimes K_C$ 的全纯截面。
\end{definition*}
$Y^\vee$ 的扭曲切丛的对应物，是 $C$ 上平坦 $G_{\mathbb C}$ 联络的模空间 $\mathcal{M}_{\mathrm{flat}}(G,C)$。如文献\BookCite{ref294}所解释，Montonen--Olive 猜想，意味着，对 $\mathcal{M}_{\mathrm{Higgs}}({}^LG,C)$ 的某个典范辛结构选择，这两个模空间互为镜像。这个辛结构，是某个全纯辛形式 $\Omega$ 的实部。因此，该 Higgs 模空间上，存在典范余各向同性 A 膜。

\SourceParagraph{ch05-5.9.4-015}{429}
虽然 $\mathcal{M}_{\mathrm{Higgs}}({}^LG,C)$，不完全是余切丛，却与余切丛双有理：在底层主丛 $E$ 稳定的开轨迹上，有把 $(E,\varphi)$ 映到 $E$ 的遗忘映射。这个开轨迹，可以识别为稳定 ${}^LG$ 丛模空间 $\mathcal{M}({}^LG,C)$ 的余切丛 $T^*\mathcal{M}({}^LG,C)$；该识别，把 $\Omega$ 映到余切丛上的典范全纯辛形式。因此，前面的论证适用，推知，镜像对称，把 $\mathcal{M}_{\mathrm{flat}}(G,C)$ 上的 B 膜，与 $\mathcal{M}({}^LG,C)$ 上的扭曲 D 模联系起来。严格地说，这个物理论证，只适用于镜像位于开子集 $T^*\mathcal{M}({}^LG,C)\subset\mathcal{M}_{\mathrm{Higgs}}({}^LG,C)$ 内的 B 膜；几何 Langlands 的数学表述，通过把稳定丛模空间，替换为 $C$ 上全纯 ${}^LG$ 丛的模堆，避开这一限制。
